%%%%%%%%%%%%%%%%%%%%%%%%%%%%%%%%%
%               %
%   Les choses à vérifier   %
%    dans verifications.txt
%               %
%%%%%%%%%%%%%%%%%%%%%%%%%%%%%%%%%


% This is part of Mes notes de mathématique
% Copyright (c) 2011-2016
%   Laurent Claessens
% See the file fdl-1.3.txt for copying conditions.

% TODO : mettre à jour pour suivre les recommandations
%http://openclassrooms.com/courses/guide-des-bonnes-pratiques-en-latex

\documentclass[a4paper,twoside,11pt]{book}

\usepackage[utf8]{inputenc}
\usepackage[T1]{fontenc}

\usepackage{etex}
\usepackage{ifthen}
\usepackage{etoolbox}           % Ceci devrait remplacer ifthen. 

\usepackage{latexsym}
\usepackage{amsfonts}
\usepackage{amsmath}
\usepackage{amsthm}
\usepackage{amssymb}
\usepackage{mathrsfs}           
\usepackage{mathabx}           % For \divides et \widehat.
\usepackage{bbm}

\usepackage{wrapfig}
\usepackage{framed}
\let\Sun\undefined      
\let\Moon\undefined      
\let\Venus\undefined
\let\Mars\undefined
\let\Jupiter\undefined
\let\Saturn\undefined
\let\Uranus\undefined
\let\Mercury\undefined
\let\Venus\undefined
\let\Mars\undefined
\let\Jupiter\undefined
\let\Saturn\undefined
\let\Uranus\undefined
\let\Neptune\undefined
\let\Pluto\undefined
\let\Earth\undefined
\let\Aries\undefined
\let\Taurus\undefined
\let\Gemini\undefined
\let\Leo\undefined
\let\Libra\undefined
\let\Scorpio\undefined
\usepackage{marvosym}       % marvosym redefines the previous symbols
\usepackage{tikz}           % Configuration at 1829426939
\usepackage{calc}
\usetikzlibrary{calc}
\usetikzlibrary{patterns}
\usepackage{color}
\usepackage{graphicx}                   % Pour l'inclusion d'image en pfd.
 
\usepackage{subfigure}
\usepackage{fancyvrb}
\usepackage{stmaryrd}       % Pour le \obslash
\usepackage{xstring}        % Utilisé pour les références vers wikipédia
\usepackage{cases}
\usepackage{lscape}         % pour l'environnement landscape, utilisé dans la correction corr0076.tex
\usepackage{multicol}
\usepackage{xspace}
\usepackage[normalem]{ulem}		% Pour le barré, commande \sout
\usepackage[all]{xy}
\let\second\undefined      % le paquet mathabx définit \second
\let\degree\undefined       % le paquet mathabx définit \degree

\usepackage[cdot,thinqspace,amssymb]{SIunits}    %   1410612643

\usepackage{textcomp}
\usepackage{lmodern}
\usepackage[a4paper,margin=2cm,left=2.6cm]{geometry} 

\usepackage{hyperref}     
\usepackage{makeidx}
\usepackage[nottoc]{tocbibind}      % Biblio inclue  dans la table des matières.
\usepackage[numbers]{natbib}        % le champ URL dans le fichier bibtex
\usepackage[refpage]{nomencl}       % Some configuration at   1338719836
\usepackage{array}

\usepackage[fr]{exocorr}
\usepackage[english,frenchb]{babel}
\usepackage{listingsutf8}   % Has to be called after babel


% Some  'configuration' for tikz       1829426939
\newcounter{defHatch}
\newcounter{defPattern}
\setcounter{defHatch}{0}
\setcounter{defPattern}{0}
\newcommand{\utilde}[1]{\underline{#1}}

% This is part of Mes notes de mathématique
% Copyright (c) 2011-2016
%   Laurent Claessens
% See the file fdl-1.3.txt for copying conditions.

% This file contains the ``configuration'' of some packages.

\hypersetup{
colorlinks=true,
citecolor=blue,
linkcolor=violet,
urlcolor=blue,     % couleur des url
filecolor=Violet   % couleur des textes qui sont des liens
}

\definecolor{dkgreen}{rgb}{0,0.4,0}
\definecolor{gray}{rgb}{0.5,0.5,0.5}
\definecolor{mauve}{rgb}{0.58,0,0.82}

% This 'lstset' is from Lilian Besson 
\lstset{ %
  inputencoding=utf8/latin1,
  backgroundcolor=\color{white},  % choose the background color; you must add \usepackage{color} or \usepackage{xcolor}
  basicstyle=\ttfamily, % \texttt\small,              % the size of the fonts that are used for the code, FIXME \ttfamily
  breakatwhitespace=false,        % sets if automatic breaks should only happen at whitespace
  breaklines=true,                % sets automatic line breaking
  captionpos=b,                   % sets the caption-position to bottom
  commentstyle=\small\color{dkgreen},   % comment style
%  deletekeywords={...},          % if you want to delete keywords from the given language
%  escapeinside={\%*}{*)},        % if you want to add LaTeX within your code
  frame=single,                   % adds a frame around the code
  keywordstyle=\small\color{blue},      % keyword style
  language=python,                % the language of the code
  fontadjust=false,
  % if you want to add more keywords to the set
%  morekeywords={define,domain,objects,init,goal,problem,action,parameters,precondition,effect,types,requirements,strips,typing},
  numbers=left,                   % where to put the line-numbers; possible values are (none, left, right)
  numbersep=5pt,                  % how far the line-numbers are from the code
  numberstyle=\tiny\color{gray},  % the style that is used for the line-numbers
  rulecolor=\color{black},        % if not set, the frame-color may be changed on line-breaks within not-black text (e.g. comments (green here))
  showspaces=false,               % show spaces everywhere adding particular underscores; it overrides 'showstringspaces'
  showstringspaces=false,         % underline spaces within strings only
  showtabs=false,                 % show tabs within strings adding particular underscores
  stepnumber=1,                   % the step between two line-numbers. If it's 1, each line will be numbered
  stringstyle=\small\color{mauve},      % string literal style
  tabsize=2,                      % sets default tabsize to 2 spaces
  prebreak = \raisebox{0ex}[0ex][0ex]{\ensuremath{\hookleftarrow}}, % pour la fin des lignes.
  aboveskip={1.5\baselineskip},
  title=\lstname                  % show the filename of files included with \lstinputlisting; also try caption instead of title
%  title=\tiny{File \textcolor{blue}{\url{\lstname}}}          % show the filename of files included with \lstinputlisting; also try caption instead of title
  %% FIXME title !
}


% Directories

% g@addto@macro\input@path gives the path in which LaTeX has to 
%   search for its \input

% \exoDirectory provides the directory in which \Exo will search for the files
%    exo*.tex and corr*.tex

% The paths for 'phystricks' are given in 'src_pictures/Directories.py'


\renewcommand{\exoDirectory}{src_exocorr/}


% For each part of the document, we begin with
% \emptyInputPath
% \addInputPath{ path-to-the-files-for-the-part  }
% In such a way, we are sure to cause an error if files are badly sorted.

\makeatletter
\providecommand{\input@path}{}
\newcommand{\addInputPath}[1]{ \g@addto@macro\input@path{ {#1/} } }
\newcommand{\emptyInputPath}{  \renewcommand{\input@path}{}  }
\makeatother


%%%%%%%%%%%%%%%%%%%%%%%%%%
%
%   Trucs mathématiques
%
%%%%%%%%%%%%%%%%%%%%%%%%

% ENSEMBLES DE NOMBRES
\newcommand{\eA}{\mathbbm{A}}
\newcommand{\eB}{\mathbbm{B}}
\newcommand{\eC}{\mathbbm{C}}
\newcommand{\eD}{\mathbbm{D}}
\newcommand{\eE}{\mathbbm{E}}
\newcommand{\eF}{\mathbbm{F}}
\newcommand{\eG}{\mathbbm{G}}
\newcommand{\eH}{\mathbbm{H}}
\newcommand{\eK}{\mathbbm{K}}
\newcommand{\eL}{\mathbbm{L}}
\newcommand{\eM}{\mathbbm{M}}
\newcommand{\eN}{\mathbbm{N}}
\newcommand{\eP}{\mathbbm{P}}
\newcommand{\eQ}{\mathbbm{Q}}
\newcommand{\eR}{\mathbbm{R}}
\newcommand{\eS}{\mathbbm{S}}
\newcommand{\eT}{\mathbbm{S}}
\newcommand{\eZ}{\mathbbm{Z}}


% ENSEMBLES de fonctions
\newcommand{\aL}{\mathcal{L}}       % Les applications linéaires
\newcommand{\cL}{L}       % Les applications linéaires continues
\newcommand{\aC}{\mathcal{C}}       % Les fonctions C^1, C^2 etc
\newcommand{\swS}{\mathscr{S}}          % L'ensemble des fonctions Schwartz
\newcommand{\swD}{\mathscr{D}}          % L'ensemble des fonctions Cinfinie à support compact.
\newcommand{\swE}{\mathscr{E}}          % L'espace des fonctions qu'on peut déformer (le grand epsilon)
                                  % Les espaces de distributions correspondants sont les mêmes avec un prime.
\newcommand{\comC}{\mathcal{C}}       % Le commutant d'un endomorphisme.

\DeclareMathOperator{\Pol}{Pol}        
\DeclareMathOperator{\Poly}{\mathcal{P}}        % Space of polynomials
\newcommand{\sdS}{\mathcal{S}}      % L'ensemble des subdivisions d'un intervalle.
\newcommand{\TF}{\mathcal{F}} %% Transformée de Fourier.
\newcommand{\mtu}{\mathbbm{1}}              % La matrice unité
\newcommand{\caract}{\mathbbm{1}}    % Characteristic function of a set
\newcommand{\catC}{\mathscr{C}}     % \catX is for the categories
\newcommand{\catD}{\mathscr{D}}     
\newcommand{\catM}{\mathscr{M}}     
\newcommand{\oB}{\mathfrak{B}}          % The space of bounded operators
\newcommand{\oK}{\mathcal{K}}           % L'espace des opérateurs compacts
\newcommand{\oL}{\mathscr{L}}           % Le L pour l'idéal de Schatten-von Neumann. Cela est aussi l'ensemble des opérateurs linéaires sur des espaces vectoriels.
\newcommand{\oP}{\mathscr{P}}           % Le P est pour l'ensemble des projections dans une algèbre de VN.

\newcommand{\euler}{\mbox{\rm e}}
\newcommand{\dist}{\operatorname {dist}}
\newcommand{\sii}{\mbox{\rm \scriptsize i}}

% LES NEWCOMMAND UN PEU ACTIFS

\newcommand*{\conclusion}{\emph{Conclusion~:~}}

% DECLARE MATH OPERATORS
\DeclareMathOperator{\fl}{fl}
\DeclareMathOperator{\SimplePrec}{sp}
\DeclareMathOperator{\NaN}{NaN}
\DeclareMathOperator{\supp}{supp}
\DeclareMathOperator{\Iso}{Iso}
\DeclareMathOperator{\Isom}{Isom}       % The group of isometries
\DeclareMathOperator{\Aut}{Aut}
\DeclareMathOperator{\Ob}{Ob}           % The ``set'' of object of a category
\DeclareMathOperator{\val}{val}     % valuation d'un polynôme
\DeclareMathOperator{\res}{res}     % Le résultant de deux polynômes
\DeclareMathOperator{\Inv}{Inv}     % L'application inverse
\DeclareMathOperator{\SP}{SP}           
\DeclareMathOperator{\Conf}{Conf}   
\DeclareMathOperator{\gsl}{\mathfrak{sl}}
\DeclareMathOperator{\go}{\mathfrak{o}}
\DeclareMathOperator{\gsu}{\mathfrak{su}}
\DeclareMathOperator{\gsp}{\mathfrak{sp}}   
\DeclareMathOperator{\so}{\mathfrak{so}}    
\DeclareMathOperator{\Spin}{Spin}
\DeclareMathOperator{\mSpin}{Spin}      % La commande \mSpin dénote l'application Spin qui va de SL(2,C) vers L^+ flèche.
\DeclareMathOperator{\spin}{\mathfrak{spin}}
\DeclareMathOperator{\Cl}{Cl}
\DeclareMathOperator{\Cliff}{Cl}
\DeclareMathOperator{\CCliff}{\Cliff^{\eC}}     % Changement de notation par rapport à avant.
\DeclareMathOperator{\volume}{vol}
\DeclareMathOperator{\esssup}{ess-\sup}
\DeclareMathOperator{\gpAff}{Aff}
\DeclareMathOperator{\Vect}{Vect}
\DeclareMathOperator{\eae}{eae}         % espace affine engendré
\DeclareMathOperator{\gpSymp}{Symp}
\DeclareMathOperator{\horsp}{hor}
\DeclareMathOperator{\Dim}{Dim}
\DeclareMathOperator{\Harm}{Harm}       % The space of harmonic forms
\DeclareMathOperator{\Sign}{Sign}       % The sign function.
\DeclareMathOperator{\Rank}{Rank}
\DeclareMathOperator{\Res}{Res}
\DeclareMathOperator{\ResW}{\Res_W}
\DeclareMathOperator{\sgrad}{sgrad}     % symplectic gradient
\DeclareMathOperator{\stab}{\mathfrak{Stab}}
\DeclareMathOperator{\ad}{ad}
\DeclareMathOperator{\Ad}{Ad}
\DeclareMathOperator{\AD}{\textbf{Ad}}
\DeclareMathOperator{\Der}{\texttt{Der}}
\DeclareMathOperator{\Inn}{Inn}
\DeclareMathOperator{\Out}{Out}
\DeclareMathOperator{\Diff}{Diff}
\DeclareMathOperator{\biDiff}{bi-Diff}
\DeclareMathOperator{\Hol}{Hol}
\DeclareMathOperator{\Ray}{Ray}
\DeclareMathOperator{\mfsp}{\mathfrak{sp}}
\DeclareMathOperator{\Fr}{Fr}
\DeclareMathOperator{\Rad}{Rad}
\DeclareMathOperator{\niv}{Level}
\DeclareMathOperator{\Supp}{Supp}
\DeclareMathOperator{\sech}{sech}
\DeclareMathOperator{\Prim}{Prim}
\DeclareMathOperator{\Trans}{Trans}
\DeclareMathOperator{\Verm}{Verm}       % Pour le module de Verma
\DeclareMathOperator{\Irr}{Irr}
\DeclareMathOperator{\vol}{Vol}
\DeclareMathOperator{\Op}{Op}           % Le truc de la quantification de Weyl
\DeclareMathOperator{\rDi}{Di}
\DeclareMathOperator{\rRac}{Rac}
\DeclareMathOperator{\Maxpprod}{\texttt{pprod}}
\DeclareMathOperator{\Maxproj}{\texttt{proj}}
\DeclareMathOperator{\Maxcom}{\texttt{com}}
\DeclareMathOperator{\Maxcombis}{\texttt{combi6}}
\DeclareMathOperator{\Maxtables}{\texttt{table6}}
\DeclareMathOperator{\Maxtablesc}{\texttt{table6c}}
\DeclareMathOperator{\Maxdecomps}{\texttt{decomp6}} % Ces commandes sont pour Maxima.
\DeclareMathOperator{\Maxdecompsc}{\texttt{decomp6c}}
\DeclareMathOperator{\Maxtableqc}{\texttt{table4c}}
\DeclareMathOperator{\Maxtableq}{\texttt{table4}}
\DeclareMathOperator{\Maxdecompq}{\texttt{decomp4}}
\DeclareMathOperator{\Maxdecompqc}{\texttt{decomp4c}}
\DeclareMathOperator{\Maxomega}{\texttt{omega}}
\DeclareMathOperator{\Maxsymple}{\texttt{symple}}
\DeclareMathOperator{\Maxcycle}{\texttt{cycle}}
\DeclareMathOperator{\Maxsolve}{\texttt{solve}}
\DeclareMathOperator{\Maxdelxistar}{\texttt{delxistar}}
\DeclareMathOperator{\Maxxistar}{\texttt{xistar}}


\DeclareMathOperator{\signe}{sgn}
\DeclareMathOperator{\Vol}{Vol}
\DeclareMathOperator{\Int}{Int}     % Intérieur d'un ensemble.
\DeclareMathOperator{\Adh}{Adh}     % Adhérence d'un ensemble.
\DeclareMathOperator{\Ind}{Ind}     % l'indice d'un chemin en analyse complexe
\DeclareMathOperator{\Turn}{Turn} % Le nombre de tours d'une courbe fermée
\DeclareMathOperator{\IR}{Ind}       % indice de rotation
\DeclareMathOperator{\Cond}{Cond}   % conditionnement d'une matrice
\DeclareMathOperator{\Diam}{Diam}
\DeclareMathOperator{\id}{Id}
\DeclareMathOperator{\Graph}{Graph}
\DeclareMathOperator{\Conv}{\mathcal{C}onv}
\DeclareMathOperator{\pr}{\texttt{proj}}
\DeclareMathOperator{\dom}{dom}
\DeclareMathOperator{\Graphe}{Gr}
\DeclareMathOperator{\Spec}{Spec}   % spectre d'un opérateur
\DeclareMathOperator{\arctg}{arctg}
\DeclareMathOperator{\cotg}{cotg}
\DeclareMathOperator{\cosec}{cosec}
\DeclareMathOperator{\arcsinh}{arcsinh}
\DeclareMathOperator{\sinc}{sinc}   % Le sinus cardinal
\DeclareMathOperator{\PGL}{PGL}   % le groupe projectif
\DeclareMathOperator{\SO}{SO}
\DeclareMathOperator{\SL}{SL}
\DeclareMathOperator{\PSL}{PSL}   % Le groupe modulaire SL(2,Z)/Z2
\DeclareMathOperator{\gS}{S}        % le groupe des matrices symétriques et aussi les nombres complexes de norme 1.
\DeclareMathOperator{\SU}{SU}
\DeclareMathOperator{\gU}{U}
\DeclareMathOperator{\gO}{O}            % On mets un g devant les groupes dont le nom est juste une lettre, ou est ambigu.
\DeclareMathOperator{\diag}{diag}   
\DeclareMathOperator{\Hom}{Hom}
\DeclareMathOperator{\Domain}{Dom}      % Domaine of an operator
\DeclareMathOperator{\Domaine}{Dom}
\DeclareMathOperator{\Dom}{Domaine}
%\DeclareMathOperator{\Reel}{Re}        % La partie réelle d'un nombre complexe
%\DeclareMathOperator{\Imag}{Im}        % La partie imaginaire d'un nombre complexe
                                       % Ne plus utilisaer

\DeclareMathOperator{\real}{Re}        % Real and imaginary part of a complex number.
\DeclareMathOperator{\imag}{Im}        % These names are from Sage.

\DeclareMathOperator{\Image}{Image}        % ... avec \Image qui donne l'image d'une fonction ou d'un opérateur.
\DeclareMathOperator{\rang}{rang}
\DeclareMathOperator{\Kernel}{Ker}
\DeclareMathOperator{\Span}{Span}
\DeclareMathOperator{\End}{End}     % L'ensemble des endomorphismes
\DeclareMathOperator{\Cyl}{Cyl}
\DeclareMathOperator{\tr}{Tr}       % la trace
\DeclareMathOperator{\Tr}{Tr}       % la trace
\renewcommand{\det}{\mathop{\mathrm{det}}\nolimits}           % le déterminant
\DeclareMathOperator{\Majorant}{Maj}
\DeclareMathOperator{\codim}{codim} % pour la codimension.
\DeclareMathOperator{\diam}{diam} % le diamètre d'un ensemble.
\DeclareMathOperator{\Var}{Var}     % Variance d'une variable aléatoire.
\DeclareMathOperator{\Fun}{\texttt{Fun}}     % Ensemble des applications d'un ensemble vers l'autre.
\DeclareMathOperator{\Cov}{Cov}     % la covariance.
\DeclareMathOperator{\gr}{gr}     % le groupe engendré
\DeclareMathOperator{\pgcd}{pgcd}     
\DeclareMathOperator{\ppcm}{ppcm}     
\DeclareMathOperator{\Frob}{Frob}     
\DeclareMathOperator{\Card}{Card}       % Le cardinal d'un ensemble.
\DeclareMathOperator{\Stab}{Stab}       % Le stabilisateur d'un point sous l'action d'un groupe.
\DeclareMathOperator{\Frac}{Frac}       % le corps des fractions d'un anneau
\DeclareMathOperator{\Aff}{Aff}         %  l'espace affine engendré


\newcommand{\dpt}[3]{#1\colon #2\to #3}
\newcommand{\subdem}[1]{\par\noindent {\it #1.} }       % TODO : use subproof instead

\newcommand{\slim}{\mathrm{s\lim}}
\newcommand{\uwlim}{\mathrm{uw\lim}}

\newcommand{\sod}{\mathfrak{so}(2)}
\newcommand{\SOdn}{\SO(2,n)}
\newcommand{\sodn}{  {\mathfrak{so}}(2,n)   }
\newcommand{\soun}{\mathfrak{so}(1,n)}
\newcommand{\SLdc}{\SL(2,\eC)}
\newcommand{\sldr}{\mathfrak{sl}(2,\eR)}
\newcommand{\SOun}{\SO(1,n)}
\newcommand{\gud}{\Gamma_{(2)}}
\newcommand{\Sput}{\Spin(1,3)}
\newcommand{\Sppq}{\Spin(p,q)}
\newcommand{\sppq}{\mathfrak{spin}(p,q)}
\newcommand{\Sopq}{\SO(p,q)}
\newcommand{\sopq}{\mathfrak{so}(p,q)}
\newcommand{\gl}{\mathfrak{gl}}
\newcommand{\Fix}{\operatorname{Fix}}
\newcommand{\cat}{\operatorname{cat}}

\newcommand{\ecarts}{ecarts}
\newcommand{\angl}{\quext{Anglais ?}}
\newcommand{\intr}[1]{\mathaccent 23 {#1}}
\newcommand{\UU}{\intr{U}}
%TODO : supprimer ces commandes.
\newcommand{\dref}[1]{}
\newcommand{\dixref}[1]{}
\newcommand{\lref}[1]{} 
\newcommand{\leref}[1]{[Landsman {(#1)}]}
\newcommand{\lsref}[1]{[Landsman {\bf #1}]}

\newcommand{\ecart}{ecart} % quand tu trouveras une meilleure traduction...

%--------- Alphabets math

\newcommand{\mfa}{\mathfrak{a}}
\newcommand{\mfb}{\mathfrak{b}}
\newcommand{\mfg}{\mathfrak{g}}
\newcommand{\mfs}{\mathfrak{s}}
\newcommand{\mfM}{\mathfrak{M}}

\newcommand{\scrC}{\mathscr{C}}
\newcommand{\scrD}{\mathscr{D}}         %Demande le paquetage mathrsfs.
\newcommand{\scrE}{\mathscr{E}}
\newcommand{\scrM}{\mathscr{M}}
\newcommand{\scrS}{\mathscr{S}}
\newcommand{\hS}{\mathscr{S}}           % C'est lui qui donne la singularité
\newcommand{\hF}{\mathscr{F}}           % \hF donne la partie libre de l'espace

\newcommand{\cA}{\mathfrak{A}}          % Pour les C^* algebres; comme ça je peux choisir.
\newcommand{\cB}{\mathcal{B}}           % Le mathfrak{B} est pour l'ensemble des operateurs bornés.
\newcommand{\cun}{\mtu}             % L'unite dans les $C^*$-algèbres.
\newcommand{\cI}{\mathfrak{I}}

\newcommand{\Cinf}{C^{\infty}}
\newcommand{\vnM}{\mathfrak{M}}         % Le M des algèbres de von Neumann
\newcommand{\hodge}{\star}         % the Hodge dual


%\newcommand{\bmodE}{\mathcal{\bar E}}      % Pour le module conjugué
\newcommand{\modM}{\mathfrak{M}}    
\newcommand{\modN}{\mathfrak{N}}    


\newcommand{\pH}{\mathscr{H}}           % L'espace de Hilbert pour la physique
\newcommand{\rR}{\mathcal{R}}           % Les \r? sont les lettres pour les rayons de l'espace de Hilbert.
\newcommand{\rC}{\mathcal{C}}

\newcommand{\cdA}{\mathscr{A}}
\newcommand{\cdE}{\mathscr{E}}          % Les ensembles de fonctions continuement d\'erivables
\newcommand{\cdD}{\mathscr{D}}          % L'ensemble des fonctions à support compact.

%-------Overline,underline, hat,tilde
\newcommand{\uw}{\underline{w}}
\newcommand{\uv}{\underline{v}}
\newcommand{\uW}{\underline{W}}
\newcommand{\uvH}{\underline{H}}
\newcommand{\ovH}{\overline{H}}
\newcommand{\ovN}{\overline{N}}
\newcommand{\ovx}{\overline{x}}
\newcommand{\ovy}{\overline{y}}
\newcommand{\ovj}{\overline{j}}
\newcommand{\ova}{\overline{a}}
\newcommand{\os}{\overline{s}}
\newcommand{\oJ}{\overline{J}}
\newcommand{\oX}{\overline{X}}
\newcommand{\oY}{\overline{Y}}
\newcommand{\ovR}{\overline{R}}
\newcommand{\olG}{\overline{\lG}}
\newcommand{\olR}{\overline{\lR}}
\newcommand{\olS}{\overline{\lS}}

\newcommand{\uG}{\underline{G}}
\newcommand{\uX}{\underline{X}}
\newcommand{\uAN}{\underline{AN}}

\newcommand{\oui}{\overline{1}_i}
\newcommand{\oalpha}{\overline{\alpha}}
\newcommand{\obeta}{\overline{\beta}}
\newcommand{\oeta}{\overline{\eta}}
\newcommand{\oxi}{\overline{\xi}}
\newcommand{\ogamma}{\overline{\gamma}}
\newcommand{\odelta}{\overline{\delta}}
\newcommand{\tA}{\widetilde{A}}
\newcommand{\tK}{\widetilde{K}}
\newcommand{\tN}{\widetilde{N}}
\newcommand{\tR}{\widetilde{R}}
\newcommand{\tilr}{\widetilde{r}}
\newcommand{\tx}{\tilde{x}}
\newcommand{\ty}{\tilde{y}}
\newcommand{\tE}{\tilde{E}}
\newcommand{\tF}{\tilde{F}}
\newcommand{\tH}{\tilde{H}}
\newcommand{\tX}{\tilde{X}}
\newcommand{\tY}{\tilde{Y}}
\newcommand{\utX}{\underline{X}}                %Il faut encore trouver comment sousligner avec un tilde.
\newcommand{\utE}{\underline{E}}
\newcommand{\utH}{\underline{H}}
\newcommand{\talpha}{\tilde{\alpha}}
\newcommand{\tomega}{\tilde{\omega}}
\newcommand{\tgamma}{\tilde{\gamma}}
\newcommand{\tnab}{\widetilde{\nabla}}
\newcommand{\expotilde}{\widetilde{\hphantom{X}}}       % Waiting to know how to create a suitable tilde in twists_general.tex

\newcommand{\hpsi}{\hat{\psi}}
\newcommand{\ha}{\hat{a}}
\newcommand{\hg}{\hat{g}}
\newcommand{\hs}{\hat{s}}
\newcommand{\hu}{\hat{u}}
\newcommand{\hv}{\hat{v}}
\newcommand{\hw}{\hat{w}}
\newcommand{\hx}{\hat{x}}
\newcommand{\hy}{\hat{y}}
\newcommand{\hB}{\hat{B}}
\newcommand{\hX}{\hat{X}}
\newcommand{\hY}{\hat{Y}}
%\Overline,underline, hat,tilde----------

%------Produits star
\newcommand{\stG}{\star^{G}}
\newcommand{\stW}{\star^{W}}
\newcommand{\stWt}{\star^{W}_{\theta}}
\newcommand{\stWh}{\star^{W}_{\hbar}}
\newcommand{\stX}{\star^{X}}
\newcommand{\stM}{\ast_M}
\newcommand{\stt}{\star_{\theta}}
%\newcommand{\st}{\ast}
%\Produits star-------------
\newcommand{\us}[1]{\frac{1}{#1}}
\newcommand{\dsd}[2]{\frac{\partial #1}{\partial #2}}
\newcommand{\me}[1]{(-1)^{#1}}
\newcommand{\xdp}[2]{#1\to #2}
\newcommand{\brak}[2]{\langle #1,#2\rangle}
\newcommand{\dsdd}[3]{\left.\frac{d}{d#2}#1\right|_{#2=#3}}
\newcommand{\Dsddb}[4]{\frac{d}{d#2}\Big[#1\Big]_{#3=#4}}
\newcommand{\Dsdd}[3]{ \Dsddb{#1}{#2}{#2}{#3}   }
\newcommand{\Dsddc}[3]{\frac{d}{d#2}\Big(#1\Big)_{#2=#3}}
\newcommand{\Dsddp}[3]{\frac{d}{d#2}\Big(#1\Big)_{#2=#3}}
\newcommand{\dDsdd}[5]{\frac{d}{d#2}\frac{d}{d#4}
           \Big[#1\Big]_{ \begin{subarray}{l}#4=#5\\#2=#3\end{subarray} }}

\newcommand{\DDsdd}[5]{\frac{d}{d#2}\frac{d}{d#4}
           \Big[#1\Big]_{ \begin{subarray}{l}#4=#5\\#2=#3\end{subarray} }}

\newcommand{\mfo}{\vartheta}

\newcommand{\hbeta}{^{\beta}}
\newcommand{\hkappa}{^{\kappa}}
\newcommand{\bgamma}{_{\gamma}}
\newcommand{\bdelta}{_{\delta}}
\newcommand{\hmu}{^{\mu}}
\newcommand{\hnu}{^{\nu}}
\newcommand{\bab}{_{\alpha\beta}}

\newcommand{\heta}{^{\eta}}
\newcommand{\bxi}{_{\xi}}
\newcommand{\hsigma}{^{\sigma}}

%Extension de l'alphabet grec------------


\newcommand{\lA}{\mathfrak{a}}
\newcommand{\lB}{\mathfrak{b}}
\newcommand{\lF}{\mathfrak{f}}
\newcommand{\lG}{\mathfrak{g}}      % Pour les algèbes de Lie en geneal; comme ça je peux choisir,  
\newcommand{\lH}{\mathfrak{h}}      % mais le mal du \mG est déjà loin !
\newcommand{\lI}{\mathfrak{i}}
\newcommand{\lJ}{\mathfrak{j}}
\newcommand{\lK}{\mathfrak{k}}
\newcommand{\lL}{\mathfrak{L}}
\newcommand{\lM}{\mathfrak{m}}
\newcommand{\lN}{\mathfrak{n}}
\newcommand{\lP}{\mathfrak{p}}
\newcommand{\lQ}{\mathfrak{q}}
\newcommand{\lR}{\mathfrak{r}}
\newcommand{\lS}{\mathfrak{s}}
\newcommand{\lU}{\mathfrak{u}}
\newcommand{\lX}{\mathfrak{x}}
\newcommand{\lZ}{\mathfrak{z}}

\newcommand{\lW}{\mathcal{W}}

\newcommand{\iA}{\mathcal{A}}
\newcommand{\iK}{\mathcal{K}}
\newcommand{\iN}{\mathcal{N}}           % Pour les éléments de décomposition d'Iwasawa
\newcommand{\iP}{\mathcal{P}}
\newcommand{\iR}{\mathcal{R}}
\newcommand{\iAH}{\mathcal{A_H}}
\newcommand{\iKH}{\mathcal{K_H}}
\newcommand{\iKQ}{\iK_{\sQ}}
\newcommand{\iNH}{\mathcal{N_H}}
\newcommand{\iPH}{\mathcal{P_H}}
\newcommand{\iRH}{\mathcal{R_H}}

\newcommand{\curR}{\mathrm{R}}      % La courbure scalaire d'un triple spectral

\newcommand{\SUR}{\mathrm{R}}
\newcommand{\SUA}{\mathrm{A}}       % les SUx sont pour les parties de SU(1,n).
\newcommand{\SUN}{\mathrm{N}}

\newcommand{\suqA}{\mathcal{A}}   % The algebra of SU_q(n)

\newcommand{\sA}{\mathcal{A}}
\newcommand{\sG}{\mathcal{G}}
\newcommand{\sH}{\mathcal{H}}           % Pour les morceaux de SO(2,n) et SO(1,n)
\newcommand{\sK}{\mathcal{K}}           % En fait, les morceaux de AdS_l par Iwasawa vont aussi êres notes avec des \sX
\newcommand{\sN}{\mathcal{N}}
\newcommand{\sP}{\mathcal{P}}
\newcommand{\sQ}{\mathcal{Q}}
\newcommand{\sR}{\mathcal{R}}
\newcommand{\sS}{\mathcal{S}}
\newcommand{\sZ}{\mathcal{Z}}

\newcommand{\etS}{\mathcal{S}}      % L'ensemble des etats sur une $C^*$-algebre. 
\newcommand{\etP}{\mathcal{P}}      % Les états purs

\newcommand{\tsA}{\widetilde{\mathcal{A}}}
\newcommand{\tsN}{\widetilde{\mathcal{N}}}
\newcommand{\tsR}{\widetilde{\mathcal{R}}}

\newcommand{\ovf}{\overline{ f }}
\newcommand{\cvec}{\mathfrak{X}}
\newcommand{\Wedge}{\bigwedge}  
\newcommand{\LogOu}{\vee}
\newcommand{\LogEt}{\wedge}
\newcommand{\cuppr}{\sharp}     % En attendant de trouver mieux.
\newcommand{\svec}{\mathcal{B}}     % Les vecteurs C^{\infty} d'une action
\newcommand{\Dir}{\mathcal{D}}

\newcommand{\yG}{\mathcal{G}}  % L'algèbre de Lie dans le truc sur YM

%---------Constructions d'un besoin passager

\newcommand{\nomscript}[1]{\emph{#1}}
\newcommand{\dtau}{\partial_{\tau}}
\newcommand{\du}{\partial_{u}}
\newcommand{\dphi}{\partial_{\phi}}
\newcommand{\frZ}[2]{   \frac{2(#1,#2)}{(#1,#1)}     }
\newcommand{\heC}{^{\eC}}
\newcommand{\beC}{_{\eC}}
\newcommand{\beR}{_{\eR}}
\newcommand{\heR}{^{\eR}}
\newcommand{\blF}{_{\lF}}
\newcommand{\etalH}{\eta_{\lH}}
\newcommand{\lHeR}{\lH_{\eR}}
\newcommand{\lGeR}{\lG_{\eR}}
\newcommand{\lFeC}{\lF^{\eC}}
\newcommand{\lGeC}{\lG^{\eC}}
\newcommand{\lHeC}{\lH^{\eC}}
\newcommand{\lbha}{\beta^{\alpha}}
\newcommand{\lbba}{\beta_{\alpha}}
\newcommand{\aba}{a_{\alpha}}
\newcommand{\abb}{a_{\beta}}
\newcommand{\abg}{a_{\gamma}}
\newcommand{\abd}{a_{\delta}}
\newcommand{\abmb}{a_{-\beta}}
\newcommand{\abmg}{a_{-\gamma}}
\newcommand{\abmd}{a_{-\delta}}
\newcommand{\abab}{a_{\alpha+\beta}}
\newcommand{\abma}{a_{-\alpha}}
\newcommand{\abmr}{a_{-\rho}}
\newcommand{\abr}{a_{\rho}}
\newcommand{\abbp}{a_{\beta'}}
\newcommand{\xbg}{x_{\gamma}}
\newcommand{\xbd}{x_{\delta}}
\newcommand{\hbb}{h_{\beta}}
\newcommand{\xbma}{x_{-\alpha}}
\newcommand{\xbb}{x_{\beta}}
\newcommand{\xbmb}{x_{-\beta}}
\newcommand{\xbmab}{x_{\alpha-\beta}}
\newcommand{\xbmamb}{x_{-\alpha-\beta}}
\newcommand{\rmg}[1]{ \big( \rho(#1)-\gamma(#1) \big) }
\newcommand{\lRlR}{[\lR,\lR]}
\newcommand{\cloi}{\overline{i}}
\newcommand{\cloj}{\overline{j}}
\newcommand{\cloip}{\overline{i'}}
\newcommand{\dD}{\scrD}

\newcommand{\AutA}{\Aut(\lA)}
\newcommand{\IntA}{\Int(\lA)}
\newcommand{\AutB}{\Aut(\lB)}
\newcommand{\IntB}{\Int(\lB)}
\newcommand{\RM}{\pr_{\sQ}\sR}
\newcommand{\Oexp}[3]{ \Omega_2\Big(  e^{\ad#1}#2,e^{\ad#1}#3  \Big)  }
\newcommand{\Xrnz}{X\times(\eR^N\setminus\{o\})}
\newcommand{\DxaDxb}{D_x^{\alpha} D\bxi\hbeta}
\newcommand{\abxi}{|\xi|}
\newcommand{\baz}[2]{\{#1e_i\}_{#2}}
\newcommand{\decompss}[3]{%
\begin{equation}
\begin{split}  
\mfs_1&=\{#1\}\\
\mfs_2&=\{#2\}#3
\end{split}  
\end{equation}
}
\newcommand{\delE}[2]{\delta_{#1}E_{#2}}
\newcommand{\dcr}[1]{[[#1]]}
\newcommand{\dga}[2]{\gamma_{#1}\gamma_{#2}}
\newcommand{\tga}[3]{\gamma_{#1}\gamma_{#2}\gamma_{#3}}
\newcommand{\qga}[4]{\gamma_{#1}\gamma_{#2}\gamma_{#3}\gamma_{#4}}
\newcommand{\rhoM}{\rho^M}
\newcommand{\hperp}{^{\perp}}
   % Mes produits scalaires                 % Je crois que je vais unifier sous \braket pour le produit < x | y > et sous \scal pour < x , y >.
\newcommand{\braket}[2]{ \langle #1|#2\rangle }
\newcommand{\ket}[1]{ | #1\rangle }
\newcommand{\bra}[1]{ \langle #1| }
\newcommand{\scalp}[2]{  (#1|#2) }
\newcommand{\scald}[2]{ \scal{#1}{#2} }
\newcommand{\scalh}[2]{ \braket{#1}{#2} }
\newcommand{\ketbra}[2]{|#1\rangle\,\langle #2|}


\newcommand{\dptvb}[3]{#1\stackrel{#2}{\longrightarrow}#3}
\newcommand{\ovv}{\overline{v}}
\newcommand{\ovX}{\overline{X}}
\newcommand{\ovS}{\overline{S}}

\newcommand{\bghd}[3]{#1_{#2}^{\phantom{#2}#3}}

% These commands were \mathbf instead of overline but I have the ``Too many math alphabets used in version normal`` error.
\newcommand{\bE}{\overline{ E }}       
\newcommand{\bA}{\overline{ A }}
\newcommand{\bB}{\overline{ B }}
\newcommand{\BX}{\overline{ X }}

\newcommand{\gab}{g_{\alpha\beta}}
\newcommand{\sbeta}{\sigma_{\beta}}
\newcommand{\salpha}{\sigma_{\alpha}}

\newcommand{\quextproj}{\quext{In project\ldots}}
%\newcommand{\tb}{\tilde{b}}
\newcommand{\gamsai}{\gamma_{\alpha j}}
\newcommand{\bsa}{{}_{(\alpha)}{}}
\newcommand{\gamaj}{\gamma_{\alpha j}}
\newcommand{\gamai}{\gamma_{\alpha i}}
\newcommand{\psisa}{\psi\bsa}  
\newcommand{\opK}{\mathfrak{K}}     % Compact operators
\newcommand{\opB}{\mathfrak{B}}     % Bounded operators
\newcommand{\invtible}{^{\times}}   % Cette commande est en attendant de trouver un symbole plus spécifique à mettre sur les ensembles pour désigner leur partir inversible.
\newcommand{\osint}{\widetilde{\int}}
\newcommand{\Lie}[1]{\mathfrak{Lie}(#1)}
\newcommand{\quexto}[1]{ \footnote{\textsf{#1}}  }
\newcommand{\qvect}[4]{(#1,#2,#3,#4)}
\newcommand{\osiint}{\widetilde{\iint}}
\newcommand{\osiiint}{\widetilde{\iiint}}
% La commande suivante est tirée de symbols-letter.pdf pour écrire une intégrale avec une barre dedans
\def\Xint#1{\mathchoice
   {\XXint\displaystyle\textstyle{#1}}%
   {\XXint\textstyle\scriptstyle{#1}}%
   {\XXint\scriptstyle\scriptscriptstyle{#1}}%
   {\XXint\scriptscriptstyle\scriptscriptstyle{#1}}%
   \!\int}
\def\XXint#1#2#3{{\setbox0=\hbox{$#1{#2#3}{\int}$}
     \vcenter{\hbox{$#2#3$}}\kern-.5\wd0}}
\newcommand{\ddashint}{\Xint=}
\newcommand{\dashint}{\Xint-}


% Le premier argument est optionnel, c'est pour ajouter un [math.QA] par exemple pour la nouvelle numérotation de arXiv. Comme tu le vois, la valeur par défaut est vide.
\newcommand{\arxiv}[2][]{%
\newline
\ifthenelse{\equal{#1}{}}{%                     Tester si un argument optionnel est passé ou non.
    \href{http://arxiv.org/abs/#2}{{\tt arXiv:#2}}%     Si tu ne mets pas ce %, il y a un problème d'espace.
            }
            {%
    \href{http://arxiv.org/abs/#2}{{\tt arXiv:#2}[#1]}%     Si tu ne mets pas ce %, il y a un problème d'espace.
}%                                  Ce %-ci aussi est indispensable pour un espace à éviter avant le . ajouté par bibtex.
}               % Fin de la commande \arxiv


% -- L'environement suivant est taxé de la classe article.cls, sauf que j'ai enlevé la possibilité que ce soit sur une page de titre.
\newcommand\abstractname{Abstract}
\makeatletter
  \newenvironment{abstract}{%
      \if@twocolumn
        \section*{\abstractname}%
      \else
        \small
        \begin{center}%
          {\bfseries \abstractname\vspace{-.5em}\vspace{\z@}}%
        \end{center}%
        \quotation
      \fi}
      {\if@twocolumn\else\endquotation\fi}
\makeatother

\newcommand{\PB}[2]{\left\{#1,#2\right\}}
                    % Les opérateurs définis pour Maxima
\newcommand{\LoL}{\mathscr{L}}      % Lorentz group
\newcommand{\LoP}{\mathscr{P}}      % Poincaré group
\newcommand{\f}{\frac}


%%%%%%%%%%%%%%%%%%%%%%%%%%
%
%   Les théorèmes et choses attenantes
%
%%%%%%%%%%%%%%%%%%%%%%%%



\setcounter{tocdepth}{3}        % Profondeur de la table des matières
\setcounter{secnumdepth}{3}     % Profondeur dans le texte

\newcounter{numtho}
\newcounter{numprob}
\newcounter{numTheme}           % For the thematic index, see InternalLinks

\makeatletter
\@addtoreset{numtho}{chapter}
\@addtoreset{equation}{chapter}
\@addtoreset{CountExercice}{chapter}
\makeatother

\newlength{\EnvSpace}
\setlength{\EnvSpace}{9pt}      % C'est la distance que je veux mettre avant et après les théorèmes, remarques, \ldots

\newtheoremstyle{MyTheorems}%
        {\EnvSpace}{\EnvSpace}%
        {\itshape}%
        {}%
        {\bfseries}{.}%
        {\newline}%
        {}%
\newtheoremstyle{MyExamples}%
        {\EnvSpace}{\EnvSpace}%
        {}%
        {}%
        {\bfseries}{.}%
        {\newline}%
        {}%
\newtheoremstyle{MyRemarks}%
        {\EnvSpace}{\EnvSpace}%
        {}%
        {}%
        {\bfseries}{.}%
        {\newline}%
        {}%


\newcounter{numloiphyz}

\newcounter{CounterExample}
\renewcommand{\theCounterExample}{\thechapter.\arabic{CounterExample}}
\renewcommand{\thenumtho}{\thechapter.\arabic{numtho}}

% The ifthenelse line here serves to write something in parenthesis if the optional argument is given.
\newenvironment{example}[1][]%
{\vspace{\EnvSpace}\refstepcounter{numtho}\noindent{\bf Exemple \thenumtho}%
\ifthenelse{\equal{#1}{}}{}{(#1)}%
\newline}%
{\phantom{a}\hfill $\triangle$\vspace{\EnvSpace}}
%TODO : example est un environnement et non un théorème seulement pour cette histoire de triangle. Il faudrait trouver quelque chose.

\newenvironment{Aretenir}{\refstepcounter{numtho}\begin{oframed}\noindent{\bf À retenir \thenumtho}\newline}{\end{oframed}\vspace{\EnvSpace}}

\theoremstyle{MyRemarks}    \newtheorem{remark}[numtho]{Remarque}

                \newtheorem{amusement}[numtho]{Amusement}
                \newtheorem{erreur}[numtho]{Error}
                \newtheorem{normaltext}[numtho]{}      

\theoremstyle{MyTheorems}
            \newtheorem{lemma}[numtho]{Lemme}
            \newtheorem{corollary}[numtho]{Corollaire}
            \newtheorem{theorem}[numtho]{Théorème}      
            \newtheorem{definition}[numtho]{Définition}      
            \newtheorem{proposition}[numtho]{Proposition}      
            \newtheorem{theoremDef}[numtho]{Théorème-définition}      
            \newtheorem{corollaryDef}[numtho]{Corollaire-définition}      
            \newtheorem{propositionDef}[numtho]{Proposition-définition}      
            \newtheorem{lemmaDef}[numtho]{Lemme-définition}
			\newtheorem{loiphyz}[numloiphyz]{Loi numéro}

            \newtheorem{exo}[CountExercice]{Exercice}       % C'est provisoire, pour Chafaï

% La numérotation des équations change dans les corrigés
\renewcommand{\theequation}{\thechapter.\arabic{equation}}
\renewcommand{\theCountExercice}{\arabic{CountExercice}}       % Ce compteur est défini dans SystemeCorr.sty
\newcommand{\defe}[2]{\textbf{#1}\index{#2}}

\renewcommand{\labelenumi}{\theenumi}
\renewcommand{\theenumi}{(\arabic{enumi})}



%%%%%%%%%%%%%%%%%%%%%%%%%%
%
%   Les macros qui font des choses
%
%%%%%%%%%%%%%%%%%%%%%%%%

\newcommand{\mA}{\mathcal{A}}
\newcommand{\mB}{\mathcal{B}}
\newcommand{\mC}{\mathcal{C}}
\newcommand{\mCC}{\mathcal{CC}}
\newcommand{\mD}{\mathcal{D}}
\newcommand{\mE}{\mathcal{E}}
\newcommand{\mF}{\mathcal{F}}
\newcommand{\mG}{\mathcal{G}}
\newcommand{\mH}{\mathcal{H}}
\newcommand{\mI}{\mathcal{I}}
\newcommand{\mJ}{\mathcal{J}}
\newcommand{\mK}{\mathcal{K}}
\newcommand{\mL}{\mathcal{L}}
\newcommand{\mM}{\mathcal{M}}
\newcommand{\mN}{\mathcal{N}}
\newcommand{\mO}{\mathcal{O}}
\newcommand{\mP}{\mathcal{P}}
\newcommand{\mQ}{\mathcal{Q}}
\newcommand{\mR}{\mathcal{R}}
\newcommand{\mS}{\mathcal{S}}
\newcommand{\mT}{\mathcal{T}}
\newcommand{\mU}{\mathcal{U}}
\newcommand{\mV}{\mathcal{V}}
\newcommand{\mW}{\mathcal{W}}
\newcommand{\mZ}{\mathcal{Z}}


\newcommand{\scal}[2]{ \langle #1,#2\rangle }
\newcommand{\vect}[1]{\overrightarrow{#1}}

\newcommand{\tq}{\text{ tel que }}          %TODO : il y a un paquet qui doit être changé en \st
\newcommand{\tqs}{\text{ tels que }}
\newcommand{\quext}[1]{ \footnote{\textsf{#1}}  }
\newcommand{\info}[1]{\texttt{#1}}

\newcommand{\normal}{\lhd}  % Cette notation n'est plus censée être utilisée.

\newcommand{\Borelien}{\mathcal{B}\text{or}}       % Les boréliens
\newcommand{\Lebesgue}{\mathcal{L}\text{eb}}       % La tribu de Lebesgue
\newcommand{\Baire}{\mathcal{B}\text{a}}       % La tribu de Baire
\newcommand{\tribA}{\mathcal{A}}            % Une tribu A
\newcommand{\tribB}{\mathcal{B}}            
\newcommand{\tribC}{\mathcal{C}}            
\newcommand{\tribD}{\mathcal{D}}            
\newcommand{\tribE}{\mathcal{E}}            % Une tribu E
\newcommand{\tribF}{\mathcal{F}}            % Une tribu F
\newcommand{\tribM}{\mathcal{M}}            % Une tribu M
\newcommand{\tribN}{\mathcal{N}}       
\newcommand{\tribT}{\mathcal{T}}

\newcommand{\affE}{\mathcal{E}}            % Un espace affine E
\newcommand{\affF}{\mathcal{F}}            
\newcommand{\affG}{\mathcal{G}}            

\newcommand{\statS}{\mathcal{S}}            % Un modèle statistique
\newcommand{\partP}{\mathcal{P}}            % L'ensemble des parties d'un ensemble
\newcommand{\pP}{\mathcal{P}}            % L'ensemble des nombres premiers

\newcommand{\polyP}{\mathcal{P}}            % L'ensemble des polynômes

\newcommand{\dB}{\mathscr{B}}       % la distribution de Bernoulli
\newcommand{\dE}{\mathscr{E}}       % la distribution exponentielle
\newcommand{\dirE}{\mathcal{E}}           % Dirichlet form
\newcommand{\dG}{\mathscr{G}}       % la distribution géométrique.
\newcommand{\dM}{\mathscr{M}}       % la distribution multinomiale
\newcommand{\dN}{\mathscr{N}}       % la distribution normale.
\newcommand{\dP}{\mathscr{P}}       % la distribution de Poisson.
\newcommand{\dT}{\mathscr{T}}       % la distribution de Student
\newcommand{\dU}{\mathscr{U}}       % la distribution uniforme

\newcommand{\hL}{\mathscr{L}}       
%\newcommand{\cL}{\hL}           % Pour la partie Chafai

\newcommand{\modE}{\mathcal{E}}         % Le E des modules
\newcommand{\modF}{\mathcal{F}}         % Le F des modules
\newcommand{\hH}{\mathscr{H}}           % Le H des espaces de Hilbert

\newcommand{\ellE}{\mathcal{E}}         % Le E des ellipsoïde




% Configuration for nomencl        1338719836

\renewcommand{\nomname}{Liste des notations}
%
% La syntaxe est facile, par exemple 
%       $\SL(2,\eR)$\nomenclature[G]{$\SL(2,\eR)$}{Le groupe de matrices deux par deux de déterminant 1.}
\renewcommand{\nomgroup}[1]{%
    \ifthenelse{\equal{#1}{A}}{\item[\textbf{Algèbre}]}{}%
    \ifthenelse{\equal{#1}{B}}{\item[\textbf{Ensembles de matrices}]}{}%
    \ifthenelse{\equal{#1}{G}}{\item[\textbf{Géométrie}]}{}%
    \ifthenelse{\equal{#1}{M}}{\item[\textbf{Chaînes de Markov}]}{}%
    \ifthenelse{\equal{#1}{R}}{\item[\textbf{Théorie des groupes}]}{}%
    \ifthenelse{\equal{#1}{P}}{\item[\textbf{Probabilités et statistique}]}{}%
    \ifthenelse{\equal{#1}{Y}}{\item[\textbf{Analyse}]}{}%
    \ifthenelse{\equal{#1}{T}}{\item[\textbf{Topologie et théorie des ensembles}]}{}%
}
% Note pour moi-même : si cette liste est changée, il faut changer mon raccourcis dans Vim.
\newcommand*{\Sp}{\textup{Sp}}
\newcommand*{\GL}{\textup{GL}}      % Le groupe linéaire; je crois qu'ontologiquement c'est mieux que le DeclareMathOperator


\newcommand{\subprooflabel}[1]{\underline{\bf #1}}

\newenvironment{subproof}{\let\Olddescriptionlabel\descriptionlabel\let\descriptionlabel\subprooflabel \begin{description}}{\end{description}\let\descriptionlabel\Olddescriptionlabel}

\newcommand{\InternalLinks}[1]
{
    \refstepcounter{numTheme} \paragraph{Thème \arabic{numTheme} : #1} 
}

\newenvironment{probleme}{\refstepcounter{numprob}\tiny\fbox{\bf Problèmes et choses à faire}\\}{\normalsize}

\newcommand{\TextePourISBN}{
    \begin{center}
ISBN : Y-Y-YYYYYYY-Y-Y
    \end{center}
}

\newcommand{\LicenceFDL}{
\begin{center}

            \includegraphics[width=1cm]{gfdl-logo-small.png}

Copyright 2011-2016  Laurent Claessens, Carlotta Donadello

Permission is granted to copy, distribute and/or modify this document under the terms of the \href{http://www.gnu.org/licenses/fdl-1.3.html}{GNU Free Documentation License}, Version 1.3 or any later version published by the Free Software Foundation; with no Invariant Sections, no Front-Cover Texts, and no Back-Cover Texts. A copy of the license is included in the chapter entitled ``GNU Free Documentation~License''.

\end{center}
}

\newcommand{\LogoEtLicence}{
\notbool{isFrido}{
\LicenceFDL
} 
{
    %\LicenceCC
    \LicenceFDL
    \TextePourISBN
}
}


% FIN DE MES CHOSES %%%%%%%

\newcounter{exoNico}
\setcounter{exoNico}{1}
\newcommand{\exerNico}{\stepcounter{exoNico}{\bf Exercice }\arabic{exoNico}. }



\newtheorem{theo}{Th{\'e}or{\`e}me}[section]
\newtheorem{defn}{D{\'e}finition}
\newtheorem{prop}{Proposition}     % redef encore dans Chafaï
%\newtheorem{lem}{Lemme}[section]


%\newtheorem{rem}{Remarque}[section]
%\newcommand{\R}{\mathbb{R}}
%\newcommand{\Rn}{\eR} 
%\newcommand{\Nn}{\eN}
\newcommand{\dem}{\textbf{D{\'e}monstration.}}
\newcommand{\vc}[1]{\boldsymbol{#1}}
\newcommand{\p}{\textrm{P}}
%\newcommand{\e}{\textrm{E}}
\newcommand{\mbt}{arbre binaire markovien}
\newcommand{\mbts}{arbres binaires markoviens}

%\newcommand{\ea}{\end{array}}


%%%%%%%%%%%%%%%%%%%%%%%%%%%%%%%%%%%%%%
%
% les petis yeux 
%
%%%%%%%%%%%%%%%%%%%%%%%%%%%%%%%%%%%%%%%%%%%%%

\newcommand{\coolexo}{$\circledast\circledast$}
\newcommand{\boringexo}{$\circleddash\circleddash$}
\newcommand{\minsyndical}{$\odot\odot$}
\newcommand{\mortelexo}{$\obslash\oslash$}



%%%%%%%%%%%%%% TRUCS DE PIERRE %%%%%%%%%%%%%%%%%%%%%%


% Le paquet array est là pour faire fonctionner l'environement arrowcases dans les trucs de Pierre.

% À régler par l'utilisateur
\newlength{\arrowsep}\setlength{\arrowsep}{3pt}
\newlength{\arrowlength}\setlength{\arrowlength}{1cm}

\newenvironment{arrowcases}%
{\begin{cases}}
{\end{cases}}



\makeatletter %% \limite[condition]x x_0
\newcommand*{\limite}[3][\@empty]{\lim_{\substack{#2\rightarrow#3\\#1}}}
\makeatother

\newcommand*\sev{<} % 

\let\ssi\iff
\newcommand*{\ideal}[1]{\{#1\}}
\newcommand*{\fleche}[1]{\stackrel{#1}\longrightarrow}

\setcounter{CountExercice}{0}

\newcommand{\Acplx}{A_\cdot}
\newcommand{\Bcplx}{B_\cdot}
\newcommand{\toisom}{\fleche\simeq}
\newcommand{\D}{\partial}
%\newcommand{\cat}[1]{{\bf #1}}
%\newcommand{\donc}{\Rightarrow}
%\newcommand{\im}{\text{im}}
%\newcommand{\coker}{\text{coker}}
\newcommand{\lied}{\mathcal L}
\newcommand*{\nom}[1]{\textsc{#1}}
%\makeatletter
%\newcommand*{\attention}[1]{\@latex@warning{#1}{!\small\bf #1!}\marginpar{Warning}}
%\makeatother
\newcommand*{\inner}{\imath}
\newcommand*{\newexo}{}
\newcommand*{\principe}{}
\newcommand*{\etape}{}
\newcommand*{\preuve}{}
\newcommand*{\exr}{\item}

\newcommand*{\crochets}[1]{\Bigl[ #1 \Bigr]}
\newcommand*{\llbrack}[1]{\left\lbrack #1 \right\lbrack}
\newcommand*{\rlbrack}[1]{\left\rbrack #1 \right\lbrack}
\newcommand*{\lrbrack}[1]{\left\lbrack #1 \right\rbrack}
\newcommand*{\rrbrack}[1]{\left\rbrack #1 \right\rbrack}


\newcommand*{\alg}[1]{\mathcal{#1}} % Algèbre
\newcommand*{\TT}{\ens T}% Tore !
\newcommand*{\topologie}{\mathscr{T}}
\newcommand*{\Topologie}{\textcursive{T}}
\newcommand{\LL}{\text{\textup{L}}} %% Espace de Lebesgue droit
\newcommand{\Ll}{\mathcal{L}} %% Lebesgue ronde
\newcommand{\sigmaalgebre}[1]{\mathcal{#1}} %% Une sigma algèbre...
\DeclareMathOperator{\SymMatrix}{Sym}
\DeclareMathOperator{\ASymMatrix}{ASym}
\newcommand{\Sym}{\SymMatrix}
\newcommand{\ASym}{\ASymMatrix}
%\newcommand{\transpose}[1]{{\vphantom{#1}}^{\mathit t}{\/#1}}
\newcommand*{\dprime}{{\prime\prime}}

%% Maths : Symboles divers
\newcommand{\surj}{\vers}
\newcommand{\isom}{\simeq}
\newcommand*{\Tau}{\alg T}
\newcommand{\cdv}{\mathfrak{X}} % Champs de vecteurs



\newcommand*{\abs}[1]{\left\vert#1\right\vert} % Valeur absolue.
\newcommand*{\module}[1]{\left\vert#1\right\vert} % Valeur absolue.
\newcommand*{\norme}[1]{\left\Vert#1\right\Vert} % norme
\newcommand*{\ordre}[1]{\left\vert#1\right\vert} % L'ordre d'un élément.
\newcommand*{\scalprod}[2]{\left\langle #1,#2\right\rangle}
\let\dual\ast

\newcommand*{\pardef}{\stackrel{\text{def}}{=}} % Par définition.
\newcommand*{\iffdefn}{\stackrel{\text{def}}{\iff}} % Par définition.
\newcommand*{\Defn}[1]{\emph{#1}} %
\newcommand*{\tensor}{\otimes}
\newcommand*{\pder}[2]{\frac{\partial #1}{\partial #2}}

\DeclareRobustCommand{\sfrac}[3][\mathrm]{\hspace{0.1em}%
  \raisebox{0.4ex}{$#1{\scriptstyle
#2}$}\hspace{-0.1em}/\hspace{-0.07em}%
  \mbox{$#1{\scriptstyle #3}$}}


\newcommand{\telque}{\vert\,}
\newcommand{\donc}{\Rightarrow}

%
\newcommand{\OC}{Octave}
\newcommand{\ML}{Matlab}
\newcommand{\SB}{Stixbox}
%
\newcommand{\FIG}[4]{%  fname scale floatparams caption
 \begin{figure}[#3]%
  \begin{center}%
  \includegraphics[scale=#2]{#1}%
  \caption{#4}%
  \label{fi:#1}%
  \end{center}%
 \end{figure}
}

% 1410612643  :   L'option amssymb sert à éviter un conflit avec la commande \square de amssymb. Note que cette dernière n'est plus accessible. Si tu en as besoin, faudra RTFM
%                    The change from SIunits to siunitx causes a too many math alphabet error.
%                    ! LaTeX Error: Too many math alphabets used in version normal.
%                    Thus we stick on the old SIunits


\makeindex
\makenomenclature

\begin{document}


% This is a hack in order to be warned of references to the future.
% the package hyperref make it fail. 

\ifthenelse{\value{VerifForwardRef}=1}{
    \let\Oldref\ref
    \renewcommand{\ref}[1]{%
    \ifthenelse{\value{page}<\pageref{#1}}{ {\huge ATTENTION} }{}%
    \Oldref{#1}%
    }
    \newcommand{\url}[1]{}
    \newcommand{\href}[2]{}
    \renewcommand{\caption}[1]{}
    \newcommand{\texorpdfstring}[2]{#1}
}{}



\thispagestyle{empty}
%\vfill
\begin{center}
  \begin{minipage}{15cm}
    \hrule\par
    \vspace{2mm}
    \begin{center}
    \Huge \bfseries Field theory from a bundle point of view \par
    \end{center}
    \hrule\par
  \end{minipage}
\end{center}

\vspace{2cm}

\begin{center}
    Laurent \textsc{Claessens}\\
    \today\\
    \url{http://student.ulb.ac.be/~lclaesse/lectures.pdf}
\end{center}

\vfill

\begin{center}

        \ifpdf
            %\includegraphics[width=9cm]{CreativeCommons.png}
            \includegraphics{CreativeCommons.png}
        \else
            \includegraphics{CreativeCommons.eps}
        \fi

        Field theory from a bundle point of view by \href{http://student.ulb.ac.be/~lclaesse/research_page.html}{Laurent Claessens} is licensed under a \href{http://creativecommons.org/licenses/by-sa/3.0/}{Creative Commons Attribution-ShareAlike 3.0 Unported License}.
\end{center}


    \tableofcontents

    \newpage

%\pagestyle{empty}

\pagestyle{headings}

%+++++++++++++++++++++++++++++++++++++++++++++++++++++++++++++++++++++++++++++++++++++++++++++++++++++++++++++++++++++++++++
\section{Warning}
%+++++++++++++++++++++++++++++++++++++++++++++++++++++++++++++++++++++++++++++++++++++++++++++++++++++++++++++++++++++++++++

During the firsts days of my thesis, I decided to write down everything I was learning. Some parts of this text were written in 2003 wile others were written yesterday; don't expect a high quality everywhere. This document thus takes the point of view of the learner with some consequences. As far as I can judge my own work:
\begin{enumerate}
    \item
        There are \emph{much} more details in the proofs in this text that what you can find in other textbooks.
    \item
        This is not a text in which you can get a deep understanding of what you are reading.
    \item
        There are (many) holes in the text.
\end{enumerate}

There are still open questions in the sense that there are points I didn't understand when I wrote. I think that these points are clearly indicated with footnotes or special environment ``Problem and misunderstanding''. Let me know if you know some answers.

By the way, let me know any comment or suggestions.


\chapter{General differential geometry} \label{Chapitre_FB}
\section{Differentiable manifolds}
%+++++++++++++++++++++++++++++++++

More precisions in \cite{kobayashi}, chapter II (\S 1 and 2) and III. (\S 1, 2, 3 and 7); \cite{madore,Helgason} can also be useful. An other source for this chapter is \cite{ms_book}. A systematic exposition of manifolds and such can be found in \cite{dgbook}.

\subsection{Definition and examples}
%-----------------------------------

A $n$-dimensional \defe{differentiable manifold}{differentiable!manifold}\index{manifold} is a set $M$ and a system of charts $\{(\mU_{\alpha},\varphi_{\alpha})\}_{\alpha\in I}$ where each set $\mU_{\alpha}$ is open in $\eR^n$ and the maps $\dpt{\varphi_{\alpha}}{\mU_{\alpha}}{M}$ are injective and satisfy the three following conditions:

\begin{itemize}
\item every $x\in M$ is contained in at least one set $\varphi_{\alpha}(\mU_{\alpha})$,
\item for any two charts $\dpt{\varphi_{\alpha}}{\mU_{\alpha}}{M}$ and $\dpt{\varphi_{\beta}}{\mU_{\beta}}{M}$, the set
\[
   \varphi_{\alpha}^{-1}( \varphi_{\alpha}(\mU_{\alpha})\cap\varphi_{\beta}(\mU_{\beta}) )
\]
is an open subset of $\mU_{\alpha}$, 
\item the map
\[
  \dpt{  (\varphi_{\beta}^{-1}\circ\varphi_{\alpha})  }{   \varphi_{\alpha}^{-1}( \varphi_{\alpha}(\mU_{\alpha})\cap\varphi_{\beta}(\mU_{\beta})  )   }{\mU_{\beta}}
\]
is differentiable\footnote{In the sequel, by ``differentiable'' we always mean smooth. If this map is differentiable, $C^k$, analytic,\ldots then the manifold is said to be differentiable, $C^k$, analytic,\ldots} as map from $\eR^n$ to $\eR^n$.

\end{itemize}

Each time we say ``manifold``, we mean ``differentiable manifold``. We will only consider manifolds with Hausdorff topology (see later for the definition of a topology on a manifold). Any open set of $\eR^n$ is a differentiable manifold if we choose the identity map as chart system. Most of surfaces $z=f(x,y)$ in $\eR^3$ are manifolds, depending on certain regularity conditions on~$f$. 

If $M_1$ and $M_2$ are two differentiable manifolds, a map $\dpt{f}{M_1}{M_2}$ is \defe{differentiable}{differentiable!map} if $f$ is continuous and for each two coordinate systems $\dpt{\varphi_1}{\mU_1}{M_1}$ and $\dpt{\varphi_2}{\mU_2}{M_2}$, the map $\varphi_2^{-1}\circ f\circ\varphi_1$ is differentiable on its domain. One can show that if $\dpt{f}{M_1}{M_2}$ and $\dpt{g}{M_2}{M_3}$ are differentiable, then $\dpt{g\circ f}{M_1}{M_3}$ is differentiable.

\subsubsection{Example: the sphere}\index{sphere}

The sphere $S^n$ is the set
\[
  S^n=\{  (x_1,\ldots, x_{n+1})\in\eR^{n+1}\tq \|x\|=1  \}
\]
for which we consider the following open set in $\eR^n$:
\[
   \mU=\{  (u_1,\ldots,u_n)\in\eR^n\tq\|u\|<1  \}
\]
and the charts $\dpt{\varphi_i}{\mU}{S}$, and $\dpt{\tilde{\varphi}_i}{\mU}{S}$
\begin{subequations}
\begin{align}
   \varphi_i(u_1,\ldots,u_n)&=(u_1,\ldots,u_{i-1}, \sqrt{  1-\|u\|^2  },u_i,\ldots,u_n )\\
   \tilde{\varphi}_i(u_1,\ldots,u_n)&=(u_1,\ldots,u_{i-1}, -\sqrt{  1-\|u\|^2  },u_i,\ldots,u_n ).
\end{align}
\end{subequations}
These map are clearly injective. To see that $\varphi(\mU)\cup\tilde{\varphi}(\mU)=S$, consider $(x_1,\ldots,x_{n+1})\in S$. Then at least one of the $x_i$ is non zero. Let us suppose $x_1\neq 0$, thus $x_1^2=1-(x_2^2+\ldots+x_{n+1}^2)$ and
\begin{equation}\label{eq:xupm}
   x_1=\pm\sqrt{1-(\ldots)}.
\end{equation}
If we put $u_i=x_{i+1}$, we have $x=\varphi(u)$ or $x=\tilde{\varphi}(u)$ following the sign in relation \eqref{eq:xupm}. The fact that $\varphi^{-1}\circ\tilde{\varphi}$ and $\tilde{\varphi}^{-1}\circ\varphi$ are differentiable is a ``first year in analysis exercise``.

\subsubsection{Example: projective space}

On $\eR^{n+1}\setminus\{o\}$, we consider the equivalence relation $v\sim\lambda w$ for all non zero $\lambda\in\eR$, and we put
\[
  \eR P^n=\left(\eR^{n+1}\setminus\{o\}\right)/\sim.
\]
This is the set of all the one dimensional subspaces of $\eR^{n+1}$. This is the real \defe{projective space}{projective!real space} of dimension $n$. We set $\mU=\eR^n$ and
\[
  \varphi_i(u_1,\ldots,u_n)=\Span\{ (u_1,\ldots,u_{i-1},1,u_i,\ldots,u_n) \}.
\]
One can see that this gives a manifold structure to $\eR P^n$. Moreover, the map
		\begin{equation}
		\begin{aligned}
			A \colon S^n &\to \eR P^n\
			v&\mapsto \Span v
		\end{aligned}
	\end{equation}	
is differentiable.

Let us show how to identify $\eR\cup\{ \infty \}$ to $\eR P^1$, the set of directions in the plane $\eR^2$. Indeed consider any vertical line $l$ (which does contain the origin). A non vertical vector subspace of $\eR^2$ intersects $l$ in one and only one point, while the vertical vector subspace is associated with the infinite point.


\subsection{Topology on manifold and submanifold}
%------------------------------------------------

A subset $V\subset M$ is \defe{open}{topology!on manifold} if for every chart $\dpt{\varphi}{\mU}{M}$, the set $\varphi^{-1}(V\cap\varphi(\mU))$ is open in $\mU$.

\begin{theorem}
This definition gives a topology on $M$ which has the following properties:

\begin{enumerate}
\item the charts maps are continuous,
\item the sets $\varphi_{\alpha}(\mU_{\alpha})$ are open.
\end{enumerate}

\end{theorem}

\begin{proof}
First we prove that the open system defines a topology. For this, remark that $\varphi_{\alpha}^{-1}$ is injective (if not, there should be some multivalued points). Then $\varphi^{-1}(A\cap B)=\varphi^{-1}(A)\cap\varphi^{-1}(B)$. If $V_1$ and $V_2$ are open in $M$, then
\[
  \varphi^{-1}(V_1\cap V_2\cap\varphi(\mU))=\varphi^{-1}(V_1\cap\varphi(\mU))\cap\varphi^{-1}(V_2\cap\varphi(\mU))
\]
which is open in $\eR^n$. The same property works for the unions.

Now we turn our attention to the continuity of $\dpt{\varphi}{\mU}{M}$; for an open set $V$ in $M$, we have to show that $\varphi^{-1}(V)$ is open in $\mU\subset\eR^n$. But the definition of the topology on $M$, is precisely the fact that $\varphi^{-1}(V\cap\varphi(\mU))$ is open.
\end{proof}

If $M$ is a differentiable manifold and $N$, a subset of $M$, we say that $N$ is a \defe{submanifold}{submanifold} of dimension $k$ if $\forall\,p\in N$, there exists a chart $\dpt{\varphi}{\mU}{M}$ around $p$ such that
\[
   \varphi^{-1}(\varphi(\mU)\cap N)=\eR^k\cap\mU:=\{(x_1,\ldots,x_k,0\ldots,0)\in\mU\}.
\]

In this case, $N$ is itself a manifold of dimension $k$ for which one can choose the $\varphi$ of the definition as charts.

Let us consider $M$ and $N$, two differentiable manifolds, $\dpt{f}{M}{N}$ a $\Cinf$ map and $x\in M$. We say that $f$ is an \defe{immersion}{immersion} at $x$ if $\dpt{df_x}{T_xM}{T_{f(x)}N}$ is injective and that $f$ is a \defe{submersion}{submersion} if $df_x$ is surjective.

If $M$ and $M$ are two analytic manifolds, a map $\dpt{\phi}{M}{N}$ is \defe{regular}{regular}\label{PgDefRegular} at $p\in M$ if it is analytic at $p$ and $\dpt{d\phi_p}{T_pM}{T_{\phi(p)}N}$ is injective.

\begin{proposition}
Let $M$ be a submanifold of the manifold $N$. If $p\in M$, then there  exists a coordinate system $\{x_1,\ldots,x_n\}$ on a neighbourhood of $p$ in $N$ such that $x_1(p)=\ldots=x_n(p)=0$ and such that the set
\[
   U=\{q\in V\tq x_j(q)=0\,\forall\, m+1\leq j\leq n\}
\]
gives a local chart of $M$ containing $p$.
\label{prop:var_coord}
\end{proposition}

\begin{proof}
No proof.
\end{proof}

The sense of this proposition is that one can put $p$ at the center of a coordinate system on $N$ such that $M$ is just a submanifold of $N$ parametrised by the fact that its last $m-n$ components are zero.

Now we can give a characterization for a submanifold: $N$ is a submanifold of $M$ when $N\subset M$ (as set) and the identity $\dpt{\iota}{N}{M}$ is regular.\label{pg:caract_subvar}

\begin{proposition}
The own topology of a submanifold is finer than the induced one from the manifold.
\label{prop:topo_sub_manif}
\index{topology!on submanifold}
\end{proposition}

\begin{proof}
Let $M$ be a manifold of dimension $n$ and $N$ a submanifold\footnote{In the whole proof, we should say ``there exists a sub-neighbourhood such that\ldots``} of dimension $k<n$. We consider $V$, an open subset of $N$ for the induced topology, so $V=N\cap\mO$ for a certain open subset $\mO$ of $M$. The aim is to show that $V$ is an open subset in the topology of $N$. 

Let us define $\mP=\varphi^{-1}(\mO)$.  The charts of $N$ are the projection to $\eR^k$ of the ones of $M$. We have to consider $W=\varphi^{-1}(V)$, since $N$ is a submanifold, $\varphi^{-1}(\mO\cap N)=\eR^k\cap\mP$. It is clear that $W=\eR^k\cap\mP$ is an open subset of $\eR^k$ because it is the projection on the $k$ first coordinates of an open subset of $\eR^n$.

The subset $V$ of $N$ will be open in the sense of the own topology of $N$ if $\varphi'{}^{-1}(V\cap\varphi'(\mU'))$ is open in $\eR^k$ where $\varphi'$ is the restriction of $\varphi$ to his $k$ first coordinates: $\varphi'(a)=\varphi(a,0)$ and $\mU'$ is the projection of $\mU$.
\end{proof}


\begin{lemma}
Let $V,M$ be two manifolds and $\varphi\colon V\to M$, a differentiable map. We suppose that $\varphi(V)$ is contained in a submanifold $S$ of $M$. If $\dpt{\varphi}{V}{S}$ is continuous, then it is differentiable.
\label{lem:var_cont_diff}
\end{lemma}

\begin{remark}
The map $\varphi$ is certainly continuous as map from $V$ to $M$ (this is in the assumptions). But this don't imply that it is continuous for the topology on $S$ (which is the induced one from $M$). So the continuity of $\dpt{\varphi}{V}{S}$ is a true assumption.
\end{remark}

\begin{proof}
Let $p\in V$. By proposition \ref{prop:var_coord}, we have  a coordinate system $\{x_1,\ldots,x_m\}$ valid on a neighbourhood $N$ of $\varphi(p)$ in $M$ such that the set 
\[
  \{r\in N\tq x_j(r)=0\, \forall s<j\leq m  \}
\]
with the restriction of $(x_1,\ldots x_s)\in N_S$ form a local chart which contains $\varphi(p)$. From the continuity of $\varphi$, there exists a chart $(W,\psi)$ around $p$ such that $\varphi(W)\subset N_S$. The coordinates $x_j(\varphi(q))$ are differentiable functions of  the coordinates of $q$ in $W$. In particular, the coordinates $x_j(\varphi(q))$ for $1\leq j\leq s$ are differentiable and $\dpt{\varphi}{V}{S}$ is differentiable because its expression in a chart is differentiable.
\end{proof}

A consequence of this lemma: if $V$ and $S$ are submanifolds of $M$ with $V\subset S$, and if $S$ has the induced topology from $M$, then $V$ is a submanifold of $S$. Indeed, we can consider the inclusion $\dpt{\iota}{V}{S}$: it is differentiable from $V$ to $M$ and continuous from $V$ to $S$ then it is differentiable from $V$ to $S$ by the lemma. Thus $V=\iota^{-1}(S)$ is a submanifold of $S$ (this is a classical result of differential geometry).

\begin{proposition}
A submanifold is open if and only if it has the same dimension as the main manifold.
\label{prop:subvar_ouvert}
\end{proposition}

\begin{proof}
\subdem{Necessary condition}
We consider some charts $\dpt{\varphi_i}{U_i}{M}$ on some open subsets $U_i$ of $\eR^n$. If $N$ is open in $M$, then this can be written as 
\[
  N=\bigcup_iU_i.
\]
If we choose the charts on $M$ in such a manner that $\dpt{\varphi_i}{U_i\cap \eR^k}{N}$ are charts of $N$, we must have $\varphi_i(U_i\cap\eR^k)=\varphi_i(U_i)$. Then it is clear that $k=n$ is necessary.
\subdem{Sufficient condition}
If $N$ has same dimension as $M$, the charts $\dpt{\varphi_i}{U_i}{M}$ are trivially restricted to $N$.
\end{proof}


%+++++++++++++++++++++++++++++++++++++++++++++++++++++++++++++++++++++++++++++++++++++++++++++++++++++++++++++++++++++++++++
\section{Tangent and cotangent bundle}
%+++++++++++++++++++++++++++++++++++++++++++++++++++++++++++++++++++++++++++++++++++++++++++++++++++++++++++++++++++++++++++

\subsection{Tangent vector}
%--------------------------

As first attempt, we define a tangent vector of $M$ at the point $x\in M$ as the ``derivative'' of a path $\dpt{\gamma}{(-\epsilon,\epsilon)}{M}$ such that $\gamma(0)=x$. It is denoted by
\[
\gamma'(0)= \dsdd{\gamma(t)}{t}{0}.
\]
The question is to correctly define de derivative in the right hand side. Such a definition is achieved as follows. A \defe{tangent vector}{} to the manifold $M$ is a linear map $X\colon  C^{\infty}(M)\to \eR$ which can be written under the form
\begin{equation}  \label{eq_deftgpath}
   Xf=(f\circ X)'(0)=\Dsdd{f(X(t))}{t}{0}
\end{equation}
for a certain path $X\colon \eR\to M$. Notice the abuse of notation between the tangent vector and the path which defines it.

A more formal way to define a tangent vector is to say that it is an equivalent class of path in the sense that two path are equivalent if and only if they induced maps by \eqref{eq_deftgpath} are equals. 

Using the chain rule $d(g\circ f)(a)=dg(f(a))\circ df(a)$ for the differentiation in $\eR^n$, one sees that this equivalence notion doesn't depend on the choice of $\varphi$. In other words, if $\varphi$ and $\tilde{\varphi}$ are two charts for a neighbourhood of $x$, then $(\varphi^{-1} \circ\gamma)'(0)=(\varphi^{-1} \circ\sigma)'(0)$ if and only if $(\tilde{\varphi}^{-1} \circ\gamma)'(0)=(\tilde{\varphi}^{-1} \circ\sigma)'(0)$. The space of all tangent vectors at $x$ is denoted by $T_xM$. There exists a bijection $[\gamma]\leftrightarrow (\varphi^{-1}\circ\gamma)'(0)$ between $T_xM$ and $\eR^n$, so $T_xM$ is endowed with a vector space structure.

If $(\mU,\varphi)$ is a chart around $X(0)$, we can express $Xf$ using only well know objects by defining the function $\tilde f =f\circ\varphi$ and $\tX=\varphi^{-1}\circ X$
\[
  Xf=\Dsdd{ (\tilde f \circ\tX)(t) }{t}{0}=\left.\dsd{\tilde f }{x^{\alpha}}\right|_{x=\tX(0)}\left.\frac{d\tX^{\alpha}}{dt}\right|_{t=0}.
\]
In this sense, we write
\begin{equation}
  X=\frac{d\tX^{\alpha}}{dt} \dsd{}{x^{\alpha}}
\end{equation}
and we say that $\{\partial_1,\ldots,\partial_n\}$ is a basis of $T_xM$. As far as notations are concerned, from now a tangent vector is written as $X=X^{\alpha}\partial_{\alpha}$ where $X^{\alpha}$ is related to the path $\dpt{X}{\eR}{M}$ by $X^{\alpha}=d\tX^{\alpha}/dt$. We will no more mention the chart $\varphi$ and write
\[
  Xf=\Dsdd{f(X(t))}{t}{0}.
\]
Correctness of this short notation is because the equivalence relation is independent of the choice of chart. When we speak about a tangent vector to a given path $X(t)$ without specification, we think about $X'(0)$.

All this construction gives back the notion of tangent vector when $M\subset \eR^m$. In order to see it, think to a surface in $\eR^3$. A tangent vector is precisely given by a derivative of a path: if $\dpt{c}{\eR}{\eR^n}$ is a path in the surface, a tangent vector to this curve is given by
\[
   \lim_{t\to 0}\frac{c(t_0)-c(t_0+t)}{t}
\]
which is a well know limit of a difference in $\eR^n$.

\label{pg:vecto_vecto}Let us precise how does a tangent vector acts on maps others than $\eR$-valued functions. If $V$ is a vector space and $\dpt{f}{M}{V}$, we define
\[
   Xf=(Xf^i)e_i
\]
where $\{e_i\}$ is a basis of $V$ and the functions $\dpt{f^i}{M}{\eR}$, the decomposition of $f$ with respect to this basis. If we consider a map $\dpt{\varphi}{M}{N}$ between two manifolds, the natural definition is $Xf:=dfX$. More precisely, if we consider local coordinates $x^{\alpha}$ and a function $\dpt{f}{M}{\eR}$,
\begin{equation}\label{eq:dvp_phi}
   (d\varphi X)f=\Dsdd{  (f\circ\varphi\circ X)(t) }{t}{0}=\dsd{f}{x^{\alpha}}\dsd{\varphi^{\alpha}}{x\hbeta}\frac{dX\hbeta}{dt}.
\end{equation}
Now we are in a notational trouble: when we write $X=X^{\alpha}\partial_{\alpha}$, the ``$X^{\alpha}$``{} is the derivative of the ``$X^{\alpha}$``{} which appears in the path $X(t)=(X^1(t),\ldots,X^n(t))$ which gives $X$ by $X=X'(0)$. So equation \eqref{eq:dvp_phi} gives
\begin{equation}
   X(\varphi):=d\varphi X=X\hbeta(\partial_{\beta}\varphi^{\alpha})\partial_{\alpha}.
\end{equation}

\subsection{Differential of a map}
%------------------------------------------

Let $\dpt{f}{M_1}{M_2}$ be a differentiable map, $x\in M_1$ and $X\in T_xM_1$, i.e. $\dpt{X}{\eR}{M_1}$ with $X(0)=x$ and $X'(0)=X$. We can consider the path $Y=f\circ X$ in $M_2$. The tangent vector to this path is written $df_x X$.

\begin{proposition}
If $\dpt{f}{M_1}{M_2}$ is a differentiable map between two differentiable manifolds, the map
		\begin{equation}
		\begin{aligned}
			df_x \colon T_xM_1 &\to T_{f(x)}M_2\
			X'(0)&\mapsto (f\circ X)'(0)
		\end{aligned}
	\end{equation}	
is linear.
\end{proposition}

\begin{proof}
We consider local coordinates $\dpt{x}{\eR^n}{M_1}$ and $\dpt{y}{\eR^m}{M_2}$. The maps $\dpt{f}{M_1}{M_2}$ and $\dpt{y^{-1}\circ f\circ x}{\eR^n}{\eR^m}$ will sometimes be denoted by the same symbol $f$. We have $(x^{-1}\circ X)(t)=(x_1(t),\ldots,x_n(t))$ and $(y^{-1}\circ Y)(t)=\big( y_1(x_1(t),\ldots,x_n(t),\ldots, y_m(x_1(t),\ldots,x_n(t)  \big)$, so that
\[
  Y'(0)=\left(   \sum_{i=1}^n \dsd{y_1}{x_i}x_i'(0),\ldots,\sum_{i=1}^n \dsd{y_m}{x_i}x_i'(0)   \right)\in\eR^m
\]
which can be written in a more matricial way under the form
\[
   Y'(0)=\left( \dsd{y_i}{x_j}x'_j(0) \right).
\]
So in the parametrisations $x$ and $y$, the map $df_x$ is given by the matrix $\partial y^i/\partial x_j$ which is well defined from the only given of $f$.
\end{proof}


Let $\dpt{x}{\mU}{M}$ and $\dpt{y}{\mV}{M}$ be two charts systems around $p\in M$. Consider the path $c(t)=x(0,\ldots,t,\ldots 0)$ where the $t$ is at the position $k$. Then, with respect to these coordinates,
\[
  c'(0)f=\Dsdd{ f(c(t))  }{t}{0}=\dsd{f}{x^i}\frac{dc^i}{dt}=\dsd{f}{x^k},
\]
so $c'(0)=\partial/\partial x^k$. Here, implicitly, we wrote $c^i=(x^i)^{-1}\circ c$ where $(x^i)^{-1}$ is the $i$th component of $x^{-1}$ seen as element of $\eR^n$. We can make the same computation with the system $y$. With these abuse of notation,
\begin{equation}
   \dsd{}{x^i}=\sum_j\dsd{y^j}{x^i}\dsd{}{y^j}
\end{equation}
as it can be seen by applying it on any function $\dpt{f}{M}{\eR}$. More precisely if $\dpt{x}{\mU}{M}$ and $\dpt{y}{\mU}{M}$ are two charts (let $\mU$ be the intersection of the domains of $x$ and $y$), let $\dpt{f}{M}{\eR}$ and $\ovf=f\circ x$, $\tilde f =f\circ y$. The action of the vector $\partial_{x^i}$ of the function $f$ is given by
\[
  \partial_{x^i}f=\dsd{\ovf}{x^i}
\]
where the right hand side is a real number that can be computed with usual analysis on $\eR^n$. This real \emph{defines} the left hand side. Now, $\ovf=\tilde f \circ y^{-1}\circ x$, so that
\[
   \dsd{\ovf}{x^i}=\dsd{ (\tilde f \circ y^{-1}\circ x) }{x^i}=\dsd{\tilde f }{y^j}\dsd{y^j}{x^i}
\]
where $\dsd{\tilde f }{y^j}$ is precisely what we write now by $\partial_{y^j}f$ and $\dsd{y^j}{x^i}$ must be understood as the derivative with respect to $x^i$ of the function $\dpt{(y^{-1}\circ x)}{\eR^n}{\eR^n}$.

Let $\dpt{f}{M}{N}$ and $\dpt{g}{N}{\eR}$; the definitions gives
\[
  (df_xX)g=\Dsdd{(g\circ f)(X(t))}{t}{0}
          =\dsd{g}{y^i}\dsd{f^i}{x^{\alpha}}\frac{dX^{\alpha}}{dt}.
\]
This shows that $\dsd{f^i}{x^{\alpha}}\frac{dX^{\alpha}}{dt}$ is $(df_xX)^i$.  But $dX^{\alpha}/dt$ is what we should call $X^{\alpha}$ in the decomposition $X=X^{\alpha}\partial_{\alpha}$ then the matrix of $df$ is given by $\dsd{f^i}{x^{\alpha}}$. So we find back the old notion of differential.

\begin{remark}
If $X\in T_xM$ and $f$ is a \emph{vector valued} function on $M$, then one can define $Xf$ by exactly the same expression. In this case,
\[
  df X=\Dsdd{f(v(t))}{t}{0}=Xf.
\]
\end{remark}

A map $\dpt{f}{M_1}{M_2}$ is an \defe{immersion}{immersion} at $p\in M_1$ if $\dpt{df_p}{T_pM_1}{T_{f(p)}M_2}$ is injective. It is a \defe{submersion}{submersion} if $df_p$ is surjective.

\subsection{Tangent and cotangent bundle}
%--------------------------------------

If $M$ is a $n$ dimensional manifold, as set the tangent bundle\index{tangent!space} is the \emph{disjoint} union of tangent spaces
\[
  TM=\bigcup_{x\in M}T_xM.
\]

\begin{theorem}
	The tangent bundle admits a $2n$ dimensional manifold structure for which the projection 
	\begin{equation}
		\begin{aligned}
			\pi \colon TM &\to M\
			T_pM&\mapsto p
		\end{aligned}
	\end{equation}	
	is a submersion.
\end{theorem}

The structure is easy to guess. If $\dpt{\varphi_{\alpha}}{\mU_{\alpha}}{M}$ is a coordinate system on $M$ (with $\mU_{\alpha}\subset\eR^n$), we define $\dpt{\psi_{\alpha}}{\mU_{\alpha}\times \eR^n}{TM}$ by
\[
  \psi( \underbrace{x_1,\ldots x_n}_{\in\mU_{\alpha}},\underbrace{a_1,\ldots a_n}_{\in\eR^n}  )
          =\sum_i a_i\left.\dsd{}{x_i}\right|_{\varphi(x_1,\ldots,x_n)}.
\]
The map $\psi_{\beta}^{-1}\circ\psi_{\beta}$ is differentiable because
\[
(\psi_{\beta}^{-1}\circ\psi_{\beta})(x,a)=( y(x),\sum_i a_i\left.\dsd{y_j}{x_i}\right|_{y(x)}  )
\]
which is a composition of differentiable maps. The set $TM$ endowed with this structure is called the \defe{tangent bundle}{tangent!bundle}.

%---------------------------------------------------------------------------------------------------------------------------
\subsection{Vector space structure on the tangent space}
%---------------------------------------------------------------------------------------------------------------------------

If \( X,Y\in T_pM\) are tangent vectors, one can define \( X+Y\) and \( \lambda X\) for every \( \lambda\in\eR\). The second one is easy:
\begin{equation}
    \lambda X=\Dsdd{ X(\lambda t) }{t}{0}.
\end{equation}
In order to define the sum of two vectors one has to consider a neighbourhood \( \mU\) of \( p\) in \( M\) and a chart \( \varphi\colon \mU\to \mO\) where \( \mO\) is an open set in \( \eR^n\). Then one consider a basis \( \{ e_i \}_{1\leq i\leq n}\) of \( \eR^n\) at the point \( \varphi(p)\). With these choices we define the ``basis'' path
\begin{equation}
    \gamma_i(t)=\varphi^{-1}(te_i)
\end{equation}
and we write
\begin{equation}
    \partial_i=\frac{ \partial  }{ \partial x_i }=\Dsdd{ \varphi^{-1}(te_i) }{t}{0}.
\end{equation}
The vectors \( \partial_i\) form a basis of \( T_pM\) in the sense of the following lemma.

\begin{lemma}
    The action of a vector \( X\in T_pM\) on a function \( f\colon M\to \eR\) can be decomposed into
    \begin{equation}
        Xf=\sum_{i=1}^n X_i(\partial_if)
    \end{equation}
    with \( X_i\in\eR\)
\end{lemma}

\begin{proof}
    Let \( \varphi\colon M\to \eR^n\) be a chart of a neighbourhood of \( p\) with \( \varphi(p)=0\). We determine the value of \( X_i\) using the function
    \begin{equation}
        f_i(x)=\varphi(x)_i,
    \end{equation}
    that is the \( i\)th component of the point \( \varphi(x)\in\eR^n\). Then if we write \( \varphi\big( X(t) \big)=\sum_j a_j(t)e_j\) we have
    \begin{subequations}
        \begin{align}
            X(f_i)=\Dsdd{ f_i\big( X(t) \big) }{t}{0}=\Dsdd{ \big[ \sum_ja_j(t)e_j \big]_i }{t}{0}=\Dsdd{ ai_(t) }{t}{0}=a_i'(0).
        \end{align}
    \end{subequations}
    Notice that \( a_i(0)=0\) since \( X(0)=p\) and \( \varphi(p)=0\). The combination \( f\circ\varphi^{-1}\) is an usual function from \( \eR^n\) to \( \eR\), so that we can use the chain rule on it. The following computation thus make sense:
    \begin{subequations}
        \begin{align}
            Xf&=\Dsdd{ f\big( X(t) \big) }{t}{0}\\
            &=\Dsdd{ f\Big( \varphi^{-1}\varphi\big( X(t) \big) \Big) }{t}{0}\\
            &=\Dsdd{ (f\circ\varphi^{-1})\big( \sum_ja_j(t)e_j \big) }{t}{0}\\
            &=\sum_k \frac{ \partial (f\circ\varphi^{-1}) }{ \partial x_k }\big( \underbrace{\sum_ja_j(0)e_j}_{=\varphi(p)=0} \big)\underbrace{\frac{ d\big[ \sum_ja_j(t)e_j \big]_k  }{ dt }}_{=a'_k(0)}\\
            &=\sum_k a'_k(0)\frac{ \partial (f\circ\varphi^{-1}) }{ \partial x_k }(0).
        \end{align}
    \end{subequations}
    Now using the definition of a derivative of a function \( \eR^n\to \eR\) and of the ``basis'' tangent vector \( \partial_k\),
    \begin{subequations}
        \begin{align}
            \frac{ \partial (f\circ\varphi^{-1}) }{ \partial x_k }(0)&=\Dsdd{ (f\circ\varphi^{-1})(te_k) }{t}{0}\\
            &=\partial_k f
        \end{align}
    \end{subequations}
   At the end of the day we have
   \begin{equation}
       Xf=\sum_k a'_k(0)\partial_kf.
   \end{equation}
\end{proof}

This lemma allows us to define the sum in \( T_pM\) as
\begin{equation}
    \left( \sum_kX_k\partial_k \right)+\left( \sum_kY_k\partial_k \right)=\sum_k (X_k+Y_k)\partial_k
\end{equation}
when \( X_k\) and \( Y_k\) are reals.

The tangent space \( T_pM\) is thus a vector space.


\subsection{Commutator of vector fields}

If $X$, $Y\in\cvec(M)$, one can define the \defe{commutator}{commutator of vector fields} $[X,Y]$ in the following way. First remark that, if $\dpt{f}{M}{\eR}$, the object $X(f)$ is also a function from $M$ to $\eR$ by $X(f)(x)=X_x(f)$, so we can apply $Y$ on $X(f)$. The definition of $[X,Y]_x$ is
\begin{equation}
  [X,Y]_xf=X_x(Yf)-Y_x(Xf).
\end{equation}
If $X=X^i\partial_i$ and $Y=Y^j\partial_j$, then 
$XY(f)=X^i\partial_i(Y^j\partial_jf)
     =X^i\partial_i Y^j\partial_j f+X^iY^j\partial^2_{ij}f$.
From symmetry $\partial^2_{ij}f=\partial^2_{ji}f$, the difference $XYf-YXf$ is only $X^i\partial_iY^j-Y^i\partial_iX^j$, so that
\begin{equation}
  [X,Y]^i=XY^i-YX^i
\end{equation}
where $X^i$ and $Y^i$ are seen as functions from $M$ to $\eR$.


\subsection{Some Leibnitz formulas}

See \cite{kobayashi}, chapter I, proposition 1.4.

\begin{lemma}
If $M$ and $N$ are two manifolds, we have a canonical isomorphism
\[
     T_{(p,q)}(M\times N)\simeq T_pM+T_qN.
\]
\label{lemLeibnitz}
\end{lemma}

\begin{proof}
A $Z\in T_{(p,q)}(M\times N)$ is the tangent vector to a curve $(x(t),y(y))$ in $M\times N$. We can consider $X\in T_pM$ given by $X=x'(0)$ and $Y\in T_qN$ given by $Y=y'(0)$. The isomorphism is the identification $(X,Y)\simeq Z$. Indeed, let us define $\oX\in T_{(p,q)}(M\times N)$, the tangent vector to the curve $(x(t),q)$, and $\oY\in T_{(p,q)}(M\times N)$, the tangent vector to the curve $(p,y(t))$. Then $Z=\oX+\oY$ because for any $\dpt{f}{M\times N}{\eR}$,
\begin{equation}
 Zf=\dsdd{f(x(t),y(t))}{t}{0}
   =\dsdd{f(x(t),y(0))}{t}{0}+\dsdd{f(x(0),y(t))}{t}{0}
   =\oX f+\oY f.
\end{equation}
\end{proof}

\begin{proposition}[Leibnitz formula]
Let us consider $M,N,V$, three manifold; a map $\dpt{\varphi}{M\times N}{V}$ and a vector $Z\in T_{(p,q)}(M\times N)$ which corresponds (lemma \ref{lemLeibnitz}) to $(X,Y)\in T_pM+T_qN$.

If we define $\dpt{\varphi_1}{M}{V}$ and  $\dpt{\varphi_2}{N}{V}$ by $\varphi_1(p')=\varphi(p',q)$ and $\varphi_2(q')=\varphi(p,q')$, we have the \defe{Leibnitz formula}{leibnitz formula}:
\begin{equation}
    d\varphi(Z)=d\varphi_1(X)+d\varphi_2(Y).
\end{equation}
\label{Leibnitz}
\end{proposition}
\begin{proof}
 Since $Z=\oX+\oY$, we just have to remark that
\[
                  d\varphi(\oX)=\dsdd{\varphi(x(t),q)}{t}{0}=d\varphi_1(X),
\]
so $d\varphi(Z)=d\varphi(\oX+\oY)=d\varphi_1(X)+d\varphi_2(Y)$.
\end{proof}
One of the most important application of the Leibnitz rule is the corollary \ref{cor_PrincLeib} on principal bundles. 


\subsection{Cotangent bundle}

A form on a vector space $V$ is a linear map $\dpt{\alpha}{V}{\eR}$. The set of all forms on $V$ is denoted by $V^*$ and is called the \defe{dual space}{dual!of a vector space} of $V$. On each point of a manifold, one can consider the tangent bundle which is a vector space. Then one can consider, for each $x\in M$ the dual space $T^*_xM:=(T_xM)^*$ which is called the \defe{cotangent bundle}{cotangent bundle}. A $1$-\defe{differential form}{differential!form} on $M$ is a smooth map $\dpt{\omega}{M}{T^*M}$ such that $\omega_x:=\omega(x)\in T^*_xM$. So, for each $x\in M$, we have a $1$-form $\dpt{\omega_x}{T_xM}{\eR}$.

Here, the smoothness is the fact that for any smooth vector field $X\in\cvec(M)$, the map $x\to\omega_x(X_x)$ is smooth as function on $M$. One often considers vector-valued forms. This is exactly the same, but $\omega_xX_x$ belongs to a certain vector space instead of $\eR$. The set of $V$-valued $1$-forms on $M$ is denoted by $\Omega(M,V)$ \nomenclature{$\Omega(M,V)$}{$V$ valued $1$-forms} and simply $\Omega(M)$ if $V=\eR$
The cotangent space $T^*_pM$ of $M$ at $p$ is the dual space of $T_pM$, i.e. the vector space of all the (real valued) linear\footnote{When we say \emph{a form}, we will always mean \emph{a linear form}.} $1$-forms on $T_pM$. In the coordinate system $\dpt{x}{\mU}{M}$, we naturally use, on $T^*_pM$, the dual basis of the basis $\{\partial/\partial_{x^i},\ldots\partial/\partial_{x^i}\}$ of $T_pM$. This dual basis is denoted by $\{dx_1,\ldots,dx_n\}$, the definition being as usual:
\begin{equation}\label{eq:dx_v}
  dx_i(\partial^j)=\delta^j_i.
\end{equation}
The notation comes from the fact that equation \eqref{eq:dx_v} describes the action of the differential of the projection $\dpt{x_i}{\mU}{\eR}$ on the vector $\partial^j$.

If $(\mU_{\alpha},\varphi_{\alpha})$ is a chart of $M$, then the maps
		\begin{equation}
		\begin{aligned}
			\phi_{\alpha} \colon \mU_{\alpha}\times\eR^n &\to T^*M\
			(x,a)&\mapsto a^idx_i|_x
		\end{aligned}
	\end{equation}	
give to $T^*M$ a $2n$ dimensional manifold structure such that the canonical projection $\dpt{\pi}{T^*M}{M}$ is an immersion.

When $V$ is a finite-dimensional vector space, we denote by $V^*$ its dual\footnote{The vector space of all the linear map $V\to \eR$.} and we often use the identifications $V\simeq V^*\simeq T_vV\simeq T_wV\simeq T^*_vV$ where $v$ and $w$ are any elements of $V$. Note however that there are no \emph{canonical} isomorphism between these spaces, unless we consider some basis.

\subsection{Exterior algebra}

Here are some recall without proof about forms on vector space. If $V$ is a vector space, we denote by $\Lambda^kV^*$ the space of all the $k$-form on $V$. We define $\dpt{\wedge}{\Lambda^kV^*\times\Lambda^lV^*}{\Lambda^{k+l}V^*}$ by
\begin{equation}
  (\omega^k\wedge\eta^l)(v_1,\ldots,v_{k+l})
 =\us{k!l!}\sum_{\sigma\in S_{k+l}}sgn(\sigma)\omega(v_{\sigma(1)},\ldots,v_{\sigma(k)})\eta(v_{\sigma(k+1)},v_{\sigma(k+1)})
\end{equation}
If $\{e_1,\ldots,e_n\}$ is a basis of $V$, the dual basis $\{\sigma^1,\ldots,\sigma^n\}$ of $V^*$ is defined by $\sigma^i(e_j)=\delta^i_j$.

If $I=\{1\leq i_1\leq\ldots i_k\leq n\}$, we write $\sigma^I=\sigma^{i_1}\wedge\ldots\sigma^{i_k}$ any $k$-form can be decomposed as
\[
  \omega=\sum_{I}\omega_I\sigma^I.
\]
The exterior algebra is provided with the \defe{interior product}{interior!product} denoted by $\iota$. It is defined by\label{pg_DefProdExt}
\begin{equation}
\begin{aligned}
 \iota(v_0)\colon\Lambda^kW&\to \Lambda^{k-1}W \\ 
(\iota(v_0)\omega)(v_1,\ldots,v_{k-1})& =\omega(v_0,v_1,\ldots,v_{k-1}).
\end{aligned}
\end{equation}

\subsection{Pull-back and push-forward}

Let $\dpt{\varphi}{M}{N}$ be a smooth map, $\alpha$ a $k$-form on $N$,
and $Y$ a vector field on $N$. Consider the map $\dpt{d\varphi}{T_xM}{T_{\varphi(x)}M}$. The aim is to extend it to a map from the tensor algebra of ${T_xM}$ to the one of $T_{\varphi(x)}M$. See \cite{kobayashi} for precise definition of the tensor algebra.

The \defe{pull-back}{pull-back!of a $k$-form} of $\varphi$ on a $k$-form $\alpha$ is the map
\[
 \dpt{\varphi^*}{\Omega^k(N)}{\Omega^k(M)}
\]
 defined by
\begin{equation}\label{306e1}
 (\varphi^*\alpha)_m(v_1,\ldots,v_k)
 =\alpha_{\varphi(m)}(d\varphi_mv_1,\ldots,d\varphi_mv_k)
\end{equation}
for all $m\in M$ and $v_i\in\cvec(M)$.

Note the particular case $k=0$. In this case, we take --instead of $\alpha$-- a function $\dpt{f}{N}{\eR}$ and the definition \eqref{306e1} gives $\dpt{\varphi^*f}{M}{\eR}$ by
\[
     \varphi^*f=f\circ\varphi.
\]

The \defe{push-forward}{push-forward!of a $k$-form} of $\varphi$ on a $k$-form is the map
\[
 \dpt{\varphi_*}{\Omega^k(M)}{\Omega^k(N)}
\]
defined by $\varphi_*=(\varphi^{-1})^*$. For $v\in T_nN$, we explicitly have:
\[
                   (\varphi_*\alpha)_n(v)=\alpha_{\varphi^{-1}(n)}(d\varphi_n^{-1} v).
\]

Let now $\dpt{\varphi}{M}{N}$ be a diffeomorphism. The \defe{pull-back}{pull-back!of a vector field} of $\varphi$ on a vector field is the map
\[
           \dpt{\varphi^*}{\cvec(N)}{\cvec(M)}
\]
defined by
\[
              (\varphi^*Y)(m)=[(d\varphi^{-1})_m\circ Y\circ\varphi](m),
\]
or
\[
 (\varphi^*Y)_{\varphi^{-1}(n)}=(d\varphi^{-1})_nY_n,
\]
for all $n\in N$ and $m\in M$. Notice that \[\dpt{(d\varphi^{-1})_n}{T_nN}{T_{\varphi^{-1}(n)}M},\] and that  $\varphi^{-1}(n)$ is well defined because $\varphi$ is an homeomorphism.

The \defe{push-forward}{push-forward!of a vector field} is, as before, defined by $\varphi_*=(\varphi^{-1})^*$. In order to show how to manipulate these notations, let us prove the following equation:
\[ 
   f_{*\xi}=(df)_{\xi}.
\]
For $\dpt{\varphi}{M}{N}$ and $Y$ in $\cvec(N)$, we just defined $\dpt{\varphi^*}{\cvec(N)}{\cvec(M)}$, by
\begin{eqnarray}
 \label{2112r1}(\varphi^*Y)_{\varphi^{-1}(n)}=(d\varphi^{-1})_nY_n.
\end{eqnarray}
Take $\dpt{f}{M}{N}$; we want to compute $f_*=(f^{-1})^*$ with $\dpt{(f^{-1})^*}{\cvec(M)}{\cvec(N)}$. Replacing the ``$^{-1}$``\ on the right places, the definition \eqref{2112r1} gives us
\[
 \Big[(f^{-1})^*X\Big]_{f(m)}=(df)_mX_m,
\]
if $X\in\cvec(M)$, and $m\in M$.

We can rewrite it without  any indices: the coherence of the spaces automatically impose the indices: $(f^{-1})^*X=(df)X$. It can also be rewritten as $(f^{-1})^*=df$, and thus $f_*=df$. From there to $f_{* \xi}=(df)_{\xi}$, it is straightforward.

\subsection{Differential of \texorpdfstring{$k$}{k}-forms}

The differential of a $k$-form is defined by the following theorem.

\begin{theorem}
Let $M$ be a differentiable manifold. Then for each $k\in \eN$, there exists an unique map
\[
  \dpt{d}{\Omega^k(M)}{\Omega^{k+1}(M)}
\]
such that

\begin{enumerate}
\item $d$ is linear,
\item for $k=0$, we find back the $\dpt{d}{\Cinf(M)}{\Omega^(M)}$ previously defined,
\item if $f$ is a function and $\omega^k$ a $k$-form, then 
\begin{equation}
d(f\omega^k)=df\wedge\omega^k+fd\omega^k,
\end{equation}


\item $d(\omega^k\wedge\eta^l)=d\omega^k\wedge\eta^l+(-1)^k\omega^k\wedge d\eta^l$,
\item $d\circ d=0$.
\end{enumerate}
\end{theorem}

An explicit expression for $d\omega^k$ is actually given by
\begin{equation}
   d\omega^k=\sum d\omega_I\wedge dx^I
\end{equation}
if $\omega^k=\sum\omega_I dx^I$.
An useful other way to write it is the following. If $\omega$ is a $k$-form and $X_1,\ldots,X_{p+1}$ some vector fields,
\begin{equation}\label{eq:formule_domega}
\begin{split}
  (k+1)d\omega(X_1,\ldots,X_{p+1})&=\sum_{i=1}^{p+1}(-1)^{i+1}X_i\omega(X_1,\ldots\hat{X}_i,\ldots,X_{p+1})\\
                                  &\quad+\sum_{i<j}(-1)^{i+j}\omega([X_i,X_j],X_1,\ldots,\hX_i,\hX_j,\ldots,X_{p+1}).
\end{split}
\end{equation}
Let us show it with $p=1$. Let $\omega=\omega_i dx^i$ and compute $d\omega(X,Y)=\partial_i\omega_j(dx^i\wedge dx^j)(X,Y)$. For this, we have to keep in mind that the $\partial_i$ acts only on $\omega_j$ while, in equation \eqref{eq:formule_domega}, a term $X\omega(Y)$ means --pointwise-- the action of $X$ on the function $\dpt{\omega(Y)}{M}{\eR}$. So we have to use Leibnitz formula:
\[
  (\partial_i\omega_j)X^iY^j=(X\omega_j)Y^j
                            =X(\omega_j Y^j)-\omega_j XY^j.
\]
On the other hand, we know that $[X,Y]^i=XY^i-YX^i$, so
\begin{equation}
   d\omega(X,Y)=X\omega(Y)-Y\omega(X)-\omega([X,Y]).
\end{equation}


\subsubsection{Hodge operator}
%/////////////////////////////
Let us take a manifold $M$ endowed with a metric $g$.  We can define a map $\dpt{r}{T^*_xM}{T_xM}$ by, for $\alpha\in T^*_xM$,
\[
   \scal{r(\alpha)}{v}=\alpha(v).
\]
for all $v\in T_xM$, where $\scal{\cdot}{\cdot}$ stands for the product given by the metric $g$. If we have $\alpha,\beta\in T^*_xM$, we can define
\[
   \scal{\alpha}{\beta}=\scal{r(\alpha)}{r(\beta)}.
\]
With this, we define an inner product on $\Lambda^p(T^*_xM)$:
\[
   \scal{\alpha_1\wedge\ldots\alpha_p}{\beta_1\wedge\ldots\beta_p}=\det_{ij}\scal{\alpha_i}{\beta_j}.
\]

The \defe{Hodge operator}{Hodge operator} is $\dpt{\star}{\Lambda^p(T^*_xM)}{\Lambda^{n-p}(T^*_xM)}$ such that for any $\phi\in\Lambda^p(T^*_xM)$,
\begin{equation}
   \phi\wedge(\star\psi)=\scal{\phi}{\psi}\Omega=\sqrt{|\det(g)|}dx^1\wedge\ldots\wedge dx^n.
\end{equation}

\subsubsection{Volume form and orientation}
%//////////////////////////////////////////

Let $M$ be a $n$ dimensional smooth manifold. A \defe{volume form}{volume!form} on $M$ is a nowhere vanishing $n$-form and the manifold itself is said to be \defe{orientable}{orientable manifold} if such a volume form exists. Two volume forms $\mu_1$ and $\mu_2$ are describe the same orientation if there exists a function $f>0$ such that\footnote{Recall that the space of $n$-forms is one-dimensional.} $\mu_1=f\mu_2$. 

\begin{proposition}
There exists only two orientations on a connected orientable manifold.
\end{proposition}
\begin{probleme}
    Check if the statement of that proposition is correct. Find a reference.
\end{probleme}

One says that the \emph{ordered} basis $(v_1,\cdots,v_n)$ of $T_xM$ is \defe{positively oriented}{positive!orientation} with respect to the volume form $\mu$ is $\mu_x(v_1,\cdots,v_n)>0$.

\subsection{Musical isomorphism}\label{subsec_musique}\index{musical isomorphism}
%---------------------------------

In some literature, we find the symbols $v^{\flat}$ and $\alpha^{\sharp}$. What does it mean ? For $X\in\cvec(M)$ and $\omega\in\Omega^2(M)$, the \defe{flat}{flat} operation $v^{\flat}\in\Omega^1(M)$ is simply defined by the inner product:
\begin{equation}
  v^{\flat}=i(v)\omega
\end{equation}
 In the same way, we define the \defe{sharp}{sharp} operation by taking a $1$-form $\alpha$ and defining $\alpha^{\sharp}$ by
\begin{equation}
   i(\alpha^{\sharp})\omega=\alpha.
\end{equation}
An immediate property is, for all $v\in\cvec(M)$, $v^{\flat\sharp}=v$, and for all $\alpha\in\Omega^1(M)$, $\alpha^{\sharp\flat}=\omega$.

\subsection{Lie derivative}
%------------------------

Consider $X\in\cvec(M)$ and $\alpha\in\Omega^p(M)$. Let $\dpt{\varphi_t}{M}{M}$ be the flow of $X$. The \defe{Lie derivative}{Lie!derivative!of a $p$-form} of $\alpha$ is
\begin{equation}\label{liesurforme}
         \mL_X\alpha=\lim_{t\to 0}\us{t}[(\varphi^*_t\alpha)-\alpha]=\dsdd{\varphi^*_t\alpha}{t}{0}.
\end{equation}
More explicitly, for $x\in M$ and $v\in T_xM$,
\[
             (\mL_X\alpha)_x(v)=\lim_{t\to 0}\us{t}\left[(\varphi_t^*\alpha)_x(v)-\alpha_x(v)\right]
\]
In the definition of the \defe{Lie derivative}{Lie!derivative!of a vector field} for a vector field, we need an extra minus sign:
\begin{equation}		\label{EqDefLieDerivativeVect}
            (\mL_XY)_x=\dsdd{\varphi_{-t*}Y_{\varphi_t(x)}}{t}{0}.
\end{equation}
Why a minus sign ? Because $Y_{\varphi_t(x)}\in T_{\varphi_t(x)}M$, but $\dpt{(d\varphi_{-t})_a}{T_aM}{T_{\varphi_{-t}(a)}M}$ so that, if we want, $\varphi_{-t*}Y_{\varphi_t(x)}$ to be a vector at $x$, we can't use $\varphi_{t*}$.

These two definitions can be embedded in only one. Let $X\in\cvec(M)$ and $\varphi_t$ its integral curve\footnote{\textit{i.e.} for all $x\in M$, $\varphi_0(x)=x$ and $\dsdd{\varphi_{u+t}(x)}{t}{0}=X_{\varphi_u(x)}$.}\index{integral!curve}. We know that $\varphi_{t*}$ is an isomorphism $\dpt{\varphi_{t*}}{T_{\varphi^{-1}(x)}M}{T_xM}$. It can be extended to an isomorphism of the tensor algebras at $\varphi^{-1}(x)$ and $x$. We note it $\tilde{\varphi}_t$. For all tensor field $K$ on $M$, we define
\[
            (\mL_XK)_x=\lim_{t\to 0}[K_x-(\tilde{\varphi_t}K)_x].
\]

On a Riemannian manifold $(M,g)$, a vector field $X$ is a \defe{\href{http://en.wikipedia.org/wiki/Killing_vector_fields}{Killing vector field}}{killing!vector field} if $\mL_Xg=0$.



\begin{lemma}
Let $\dpt{f}{(-\epsilon,\epsilon)\times M}{\eR}$ be a differentiable map with $f(0,p)=0$ for all $p\in U$. Then there exists $\dpt{g}{(-\epsilon,\epsilon)\times M}{\eR}$, a differentiable map such that $f(t,p)=tg(t,p)$ and
\[
                g(0,q)=\left.\dsd{f(t,q)}{t}\right|_{t=0}.
\]
\end{lemma}
\begin{proof}
Take
\[
                g(t,q)=\int_0^1\dsd{f(ts,p)}{(ts)}ds,
\]
and use the change of variable $s\to ts$.
\end{proof}

\begin{lemma}
If $\varphi_t$ is the integral curve of $X$, for all function $\dpt{f}{M}{\eR}$, there exists a map $g$, $g_t(p)=g(t,p)$ such that
$f\circ\varphi_t=f+tg_t$ and $g_0=Xf$.
\end{lemma}

\begin{proof}
Consider $\overline{f}(t,p)=f(\varphi_t(p))-f(p)$, and apply the lemma:
\[
          f\circ\varphi_t=tg_t(p)+f(p).
\]
Thus we have
\[
      Xf=\lim_{t\to 0}\us{t}[f(\varphi_t(p))-f(p)]=\lim_{t\to 0}g_t(p)=g_0(p).
\]
\end{proof}

One of the main properties of the Lie derivative is the following:
\begin{theorem}		\label{ThoLieDerrComm}
Let $X$, $Y\in\cvec(M)$ and $\varphi_t$ be the integral curve of $X$. Then
\[
         [X,Y]_p=\lim_{t\to 0}\us{t}[Y-d\varphi_tY](\varphi_t(p)),
\]
or
\[
          \mL_XY=[X,Y].
\]

\end{theorem}
\begin{proof}
Take $\dpt{f}{M}{\eR}$ and the function given by the lemma: $\dpt{g_t}{M}{\eR}$ such that $f\circ \varphi_t=f+tg_t$ and $g_0=Xf$. Then put $p(t)=\varphi_t^{-1}(p)$. The rest of the proof is a computation:
\[
            (\varphi_{t*}Y)_pf=Y(f\circ\varphi_t)_{p(t)}=(Yf)_{p(t)}+t(Yg_t)_{p(t)},
\]
so
\begin{equation}
\begin{split}
  \lim_{t\to 0}\us{t}[Y_p-(\varphi_{t*}Y)_p]f&=\lim_{t\to 0}\us{t}[(Yf)_p-(Yf)_{p(t)}]-\lim_{t\to 0}(Yg_t)_{p(t)}\\
                                         &=X_p(Yf)-Y_pg_0\\
                                         &=[X,Y]_pf.
\end{split}
\end{equation}

\end{proof}

A second important property is
\begin{theorem}
For any function $f\colon M\to V$,
\[
           \mL_Xf=Xf.
\]
\end{theorem}

\begin{proof}
If $X(t)$ is the path which defines the vector $X$, it is obvious that at $t=0$, $X(t)$ is an integral curve to $X$, so that we can take $X(t)$ instead of $\varphi_t$ in \eqref{liesurforme}. Therefore we have:
\begin{equation}
    \mL_Xf=\dsdd{\varphi_t^*f}{t}{0}
          =Xf
\end{equation}
by definition of the action of a vector on a function.

\end{proof}


\chapter{Lie groups and subgroups}
% This is part of (almost) Everything I know in mathematics
% Copyright (c) 2013-2015
%   Laurent Claessens
% See the file fdl-1.3.txt for copying conditions.

Most of this chapter come (often very directly) from \cite{Helgason}. Other sources are \cite{Knapp_reprez,DirkEnvFiniteDimNilLieAlg,SamelsonNotesLieAlg,SternLieAlgebra}.

Do you know what is violet and commutative ? Answer in the footnote\footnote{An abelian grape !}.

\section{Lie groups}
%+++++++++++++++++++++++++++++

A \defe{Lie group}{Lie!group}\index{group!Lie} is a manifold $G$ endowed with a group structure such that the inversion map $\dpt{i}{G}{G}$, $i(x)=x^{-1}$ and the multiplication $\dpt{m}{G\times G}{G}$, $m(x,y)=xy$ are differentiable. The \defe{Lie algebra}{Lie!algebra}\index{algebra!Lie} of the Lie group $G$ is the tangent space of $G$ at the identity: $\yG=T_eG$. 

 It is immediate to see that $g\mapsto g^{-1}$ is a smooth homeomorphism and that, for any fixed $g_0, g_1$, the maps
\[ 
\begin{split}
   g&\mapsto g_0g,\\
g&\mapsto gg_0,\\
g&\mapsto g_0gg_1
\end{split}  
\]
are smooth homeomorphisms. When $A\subset G$, we define $A^{-1}=\{ g^{-1}\tq g\in G \}$.

\subsection{Connected component of Lie groups}
%---------------------------------------------

\begin{proposition}		\label{PropUssGpGenere} 
If $G$ is a connected Lie group and $\mU$, a neighbourhood of the identity $e$, then $G$ is generated by $\mU$ in the sense that $\forall g\in G$, there exists a \emph{finite} number of $g_{i}\in \mU$ such that
\[ 
  g=g_1\ldots g_n.
\]
Notice that the number $n$ is function of $g$ in general.
\end{proposition}

\begin{proof}
Eventually passing to a subset, we can suppose that $\mU$ is open. In this case, $\mU^{-1}$ is open because it is the image of $\mU$ under the homeomorphism $g\mapsto g^{-1}$. Now we consider $V=\mU\cap\mU^{-1}$. The main property of this set is that $V=V^{-1}$. Let
\[ 
  [V]=\{ g_1\ldots g_n\tq g_{i}\in V \};
\]
we will prove that $[V]=G$ by proving that it is closed and open in $G$ (the fact that $G$ is connected then concludes).

We begin by openness of $[V]$. Let $g_0=g_1\cdots g_n\in[V]$. We know that $g_0V$ is open because the multiplication by $g_0$ is an homeomorphism. It is clear that $g_0V\subset [V]$ and that $g_0=g_0e\in g_0V$. Hence $g_0\in g_0V\subseteq[V]$. It proves that $[V]$ is open because $g_0V$ is a neighbourhood of $g_0$ in $[V]$.

We now turn our attention to the closeness of $[V]$. Let $h\in\overline{ [V] }$. The set $hV$ is an open set which contains $h$ and $hV\cap [V]\neq \emptyset$ because an open which contains an element of the closure of a set intersects the set (it is almost the definition of the closure). Let $g_0\in hV\cap[V]$. There exists a $h_{1}\in V$ such that $g_0=hh_1$. For this $h_1$, we have $hh_1=g_0=g_1\cdots g_n$, and therefore
\[ 
  h=g_1\cdots g_n h_1^{-1}\in[V].
\]
This proves that $h\in[V]$ because $h_1^{-1}\in V$ from the fact that $V=V^{-1}$.
\end{proof}
 Remark that this proof emphasises the topological aspect of a Lie group: the differential structure was only used to prove thinks like that $A^{-1}$ is open when $A$ is open.

\begin{proposition}
Let $G$ be a Lie group and $G_0$, the identity component of $G$. We have the following:
\begin{enumerate}
\item $G_0$ is an open invariant subgroup of $G$,
\item $G_0$ is a Lie group,
\item the connected components of $G$ are lateral classes of $G_0$. More specifically, if $x$ belongs to the connected component $G_1$, then $G_1=xG_0=G_0x$.
\end{enumerate}

\end{proposition}

\begin{proof}
We know that when $M_{1}$ is open in the manifold $M$, one can put on $M_{1}$ a differential structure of manifold of same dimension as $M$ with the induced topology. Since $G_0$ is open, it is a smooth manifold. In order for $G_0$ to be a Lie group, we have to prove that it is stable under the inversion and that $gh\in G_0$ whenever $g$, $h\in G_0$.

First, $G_0^{-1}$ is connected because  it is homeomorphic to $G_0$ in $G$. The element $e$ belongs to the intersection of $G_0$ and $G_0^{-1}$, so $G_0\cup G_0^{-1}$ is connected as non-disjoint union of connected sets. Hence $G_0\cup G_0^{-1}=G_0$ and we conclude that $G_0^{-1}\subseteq G_0$. The set $G_0G_0$ is connected because it is the image of $G_0\times G_0$ under the multiplication map, but $e\in G_0G_0$, so $G_0G_0\subseteq G_0$ and  $G_0$ is thus closed for the multiplication. Hence $G_0$ is a Lie group.

For all $x\in G$, we have $e=xex^{-1}\in xG_0x^{-1}$, but $xG_0x^{-1}$ is connected. Hence $xG_0x^{-1}\subseteq G_0$, which proves that $G_0$ is an invariant subset of $G$.

Lateral classes $xG_0$ are connected because the left multiplication is an homeomorphism. They are moreover \emph{maximal} connected subsets because, if $xG_0\subset H$ (proper inclusion) with a connected $H$, then $G_0\subset x^{-1}H$ (still proper inclusion). But the definition of $G_0$ is that this proper inclusion is impossible. Therefore, the sets of the form $xG_0$ are maximally connected sets. It is clear that $\cup_{g\in G}gG_0=G$.

Notice that the last point works with $G_0x$ too.
\end{proof}

%+++++++++++++++++++++++++++++++++++++++++++++++++++++++++++++++++++++++++++++++++++++++++++++++++++++++++++++++++++++++++++
\section{Two words about Lie algebra}
%+++++++++++++++++++++++++++++++++++++++++++++++++++++++++++++++++++++++++++++++++++++++++++++++++++++++++++++++++++++++++++

\subsection{The Lie algebra of \texorpdfstring{$\SU(n)$}{SU2}}
%---------------------------------------

Let consider $G=\SU(n)$; the elements are complexes $n\times n$ matrices $U$ such that $UU^{\dag}=\mtu$ and $\det U=1$. An element of the Lie algebra is given by a path $\dpt{u}{\eR}{G}$ in the group with $u(0)=\mtu$. Since for all $t$, $u(t)u(t)^{\dag}=\mtu$,
\begin{equation}
\begin{split}
  0&=\Dsdd{u(t)u(t)^{\dag}}{t}{0}\\
   &=u(0)\Dsdd{u(t)^{\dag}}{t}{0}+\Dsdd{u(t)}{t}{0}u(0)^{\dag}\\
   &=[d_tu(t)]^{\dag}+[d_tu(t)].
\end{split}
\end{equation}
So a general element of the Lie algebra $\mathfrak{su}(n)$ is an anti-hermitian matrix. 

An element of \( \SU(n)\) has also a determinant equal to \( 1\). What condition does it implies on the elements of the Lie algebra ? If \( g(t)\) is a path in \( \SU(n)\) with \( g(0)=\mtu\) we have
\begin{equation}
    \det\begin{pmatrix}
        g_{11}(t)    &   g_{12}(t)    &   \ldots    \\
        f_{21}(t)    &   g_{22}(t)    &   \ldots    \\
        \vdots    &   \vdots    &   \ddots
    \end{pmatrix}=g_{11}(t)M_{11}(t)+g_{12}(t)M_{12}(t)+\ldots=1
\end{equation}
where \( M_{ij}\) is the minor of \( g\). If we derive the left hand side we get
\begin{equation}
    g'_{11}(0)M_{11}(0)+g_{11}(0)M'_{11}(0)+g'_{12}(0)M_{12}(0)+g_{12}(0)M'_{12}(0)+\ldots
\end{equation}
where the numbers \( g'_{ij}(0)\) are the matrix entries of the tangent matrix, that is the matrix elements of a general element in \( \gsu(n)\). Since \( g(0)=\mtu\) we have \( M_{11}(0)=1\), \( g_{11}(0)=1\), \( M_{12}(0)=0\) and \( g_{12}(0)=0\). Thus we have
\begin{equation}
    (\det g)'(0)=X_{11}+M'_{11}(0)
\end{equation}
where \( X=g'(0)\). By induction we found that the trace of \( X\) appears. Thus the elements of \( \gsu(n)\) have vanishing trace.


\subsection{What is  \texorpdfstring{$g^{-1} dg$}{g-1dg} ?}\label{SubSecgmudg}
%--------------------------------

The expression $g^{-1} dg$ is often written in the physical literature. In our framework, the way to gives a sense to this expression is to consider it pointwise acting on a tangent vector. More precisely, the framework is the data of a manifold $M$, a Lie group $G$ and a map $\dpt{g}{M}{G}$. Pointwise, we have to apply $g(x)^{-1} dg_x$ to a tangent vector $v\in T_xM$.

Note that $\dpt{dg_x}{T_xM}{T_{g(0)}G\neq T_eG}$, so $dg_x\notin \yG$. But the product $g(x)^{-1} dg_x v$ is defined by
\begin{equation}
	g(x)^{-1} dg_x v=\Dsdd{ g(x)^{-1} g(v(t)) }{t}{0}\in\yG.
\end{equation}


\subsection{Invariant vector fields}\index{invariant!vector field}
%-----------------------------------

If $G$ is a Lie group, a vector field $\tilde X\in\Gamma^{\infty}(TG)$ is \defe{left invariant}{left invariant!vector field} if 
\begin{equation}
	(dL_g)\tilde X=\tilde X.
\end{equation}
To each element $X$ in the Lie algebra $\mG$, we have an associated left invariant vector field given on the point $g\in G$ by the path
\begin{equation}
	\tilde X_g(t)=ge^{tX}.
\end{equation}
In the same way, we associate a \defe{right invariant}{right!invariant!vector field} given by
\begin{equation}
	\utilde X_g(t)=e^{tX}g
\end{equation}
The invariance means that $(dL_h)_g\tilde X_g=\tilde X_{hg}$ and $(dR_h)_g\utilde X_g=\utilde X_{gh}$. The invariant vector fields are important because they carry the structure of the tangent space at identity (the Lie algebra). More precisely we have the following result:

\begin{theorem}
	The map $X\to X_e$ is a bijection between the left invariant vector fields on a Lie group and its Lie algebra $T_eG$.
\end{theorem}
Invariant vector fields are also often used in order to transport a structure from the identity of a Lie group to the whole group by $A_g(X_g)=A_e(dL_{g^{-1}}X_g)$ where $A_e$ is some structure and $X_g$, a vector at $g$.


\begin{proposition}
	Let $G$ be a Lie group and $\mG$ the vector space of its left invariant vector fields.
	\begin{enumerate}

		\item
			The map
			\begin{equation}
				\begin{aligned}
					\mG\colon &\to T_eG \\
					\tilde X&\mapsto \tilde X_e 
				\end{aligned}
			\end{equation}
			is a vector space isomorphism.
		\item
			We have $[\mG,\mG]\subset \mG$ and $\mG$ is a Lie algebra. Here, the commutator is the bracket of vector fields.

	\end{enumerate}
\end{proposition}
\begin{proof}
	No proof.
\end{proof}
%---------------------------------------------------------------------------------------------------------------------------
\subsection{Invariant vector field}
%---------------------------------------------------------------------------------------------------------------------------

A vector field \( X\) on a Lie group \( G\) is \defe{left invariant}{left invariant!vector field} if \( dL_g(X)=X\) for every \( g\in G\). Here \( L_g\colon G\to G\) is the left translation defined by \( L_g(h)=gh\). More explicitly, the left invariance is expressed by
\begin{equation}
    \Dsdd{ gX_h(t) }{t}{0}=X_{gh}
\end{equation}
where \( X_h(t)\) is the path defining the tangent vector \( X_h\in T_h G\).




We want to prove that the vector space of left invariant vector fields is isomorphic to the tangent vector space \( T_eG\) to \( G\) at identity. If \( X\in T_eG\), we introduce the left invariant vector field \( X^L=dLX\), more explicitly:
\begin{equation}
    X^L_g=\Dsdd{ gX(t) }{t}{0}.
\end{equation}
Then we consider \( \alpha_X\colon I\to G\) the integral curve of maximal length to \( X^L\) trough \( X_e\). Here, \( I\) is the interval on which \( \alpha_X\) is defined. This is the solution of
\begin{subequations}
    \begin{numcases}{}
        \Dsdd{ \alpha_X(t_0+t) }{t}{0}=X_{\alpha_X(t_0)}\\
        \alpha_X(0)=e.
    \end{numcases}
\end{subequations}

\begin{proposition}     \label{PROPooWEYCooCvyHNr}
    Let \( X\in T_eG\). The integral curve has \( \eR\) as domain and for every \( s,t\in\eR\),
    \begin{equation}
        \alpha_X(s+t)=\alpha_X(s)\alpha_X(t).
    \end{equation}
\end{proposition}

\begin{proof}
    Let \( \alpha\) be any integral curve for \( X^L\) and \( y\in G\). If we put \( \alpha_1(t)=y\alpha(t)\), we have
    \begin{equation}
        \Dsdd{ \alpha_1(t) }{t}{0}=X^L_y,
    \end{equation}
    so that \( \alpha_1\) is an integral curve for \( X^L\) trough the point \( y\). 

    Let now \( I\) be the maximal domain of \( \alpha_X\), and \( t_1\in I\). If we set \( x_1=\alpha_X(t_1)\), the path
    \begin{equation}
         \alpha_1(t)=x_1\alpha_X(t)
    \end{equation}
    is an integral curve of \( X^L\) trough \( x_1\) and has the same maximal definition domain \( I\). On the other hand, the maximal integral curve starting at \( e\) being \( \alpha_X\), the maximal integral curve starting at \( \alpha_X(t_1)\) is
    \begin{equation}
        \alpha_2\colon t\mapsto \alpha_X(t+t_1).
    \end{equation}
    Its domain is \( I-t_1\), but since it starts at \( x_1\), it has to be the same as \( \alpha_1\), then \( I\subset I-t_1\) which proves that \( I=\eR\).

    For each \( s\) and \( t\) in \( \eR\), the maximal integral curve starting at \( \alpha_X(s)\) can be written as
    \begin{equation}
        c(t)=\alpha_X(s)\alpha_X(t)
    \end{equation}
    as well as
    \begin{equation}
        d(t)=\alpha_X(s+t),
    \end{equation}
    so again by unicity, \( \alpha_X(s+t)=\alpha_X(s)\alpha_X(t)\).
\end{proof}


%---------------------------------------------------------------------------------------------------------------------------
\subsection{Integral curve and exponential}
%---------------------------------------------------------------------------------------------------------------------------

\begin{definition}
    If \( \alpha_X\) is the integral curve to \( X^L\), we define the \defe{exponential}{exponential!on Lie group}
    \begin{equation}        \label{EqdefExpoLieTgFGp}
        \begin{aligned}
            \exp\colon T_eG&\to G \\
            X&\mapsto \alpha_X(1). 
        \end{aligned}
    \end{equation}
\end{definition}

This definition works on Lie groups thanks to the group structure that allows to build a natural vector field \( X^L\) from the data of a single vector \( X\). On general manifolds, one has not a notion of exponential. However, if one has a Riemannian manifold, one consider the geodesic.

In the case of groups for which the Killing form defines a scalar product, the notion of exponential associated with the Riemannian structure propagated from the Killing form coincides with the definition \eqref{EqdefExpoLieTgFGp}.

A very important point\cite{ooOLNIooDLmxkR} is that when \( G\) is acting on $M$, one can reconstruct the action of \( G\) only knowing the action of \( \lG\). Let \( X\in \lG\) and \( x\in M\). We consider the path
\begin{equation}
    \begin{aligned}
        \gamma\colon \eR&\to M \\
        t&\mapsto \exp(tX)\cdot x. 
    \end{aligned}
\end{equation}
This map satisfies \( \gamma(0)=x\). We also have, using proposition \ref{PROPooWEYCooCvyHNr},
\begin{equation}
    \gamma(t_0+u)=\exp(uX)\cdot \underbrace{\exp(t_0X)\cdot x}_{=\gamma(t_0)},
\end{equation}
then
\begin{equation}
    \Dsdd{ \gamma(t_0+u) }{u}{0}=X\big( \gamma(t_0) \big)
\end{equation}
We conclude that \( \gamma\) satisfies the differential equation
\begin{equation}        \label{EQooFGSIooUplbmN}
    \gamma'(t)=-X\big( \gamma(t) \big).
\end{equation}
When \( M=\eR^n\), the Cauchy-Lipschitz theorem \ref{ThokUUlgU} provides unicity of the solution on a maximal domain providing the map \( X\colon \eR^n\to \eR^n\) has nice properties.

If we have to determine the transformations of \( \eR^n\) that satisfies some properties, the strategy is then the following :
\begin{itemize}
    \item Suppose the searched group to be a connected Lie group.
    \item Write the condition with the groupe element \( \exp(tX)\) and differentiate with respect to \( t\).
    \item This provides an equation for \( X\). Typically a differential equation for the map \( X\colon \eR^n\to \eR^n\). Solve it.
    \item The group action is then retrieved solving the differential equation \eqref{EQooFGSIooUplbmN}.
\end{itemize}
An example of using this is when we will determine the conformal group around definition \ref{DEFooVKNBooFBWQQM}.  % position 10906-29466 : provides a more precise reference to the result instead of the definition.

%--------------------------------------------------------------------------------------------------------------------------- 
\subsection{Adjoint map}
%---------------------------------------------------------------------------------------------------------------------------

The ideas of this short note comes from \cite{Lie}. A more traumatic definition of the adjoint group can be found in \cite{Helgason}, chapter II, \S 5. Let $G$ be a Lie group, and $\mG$, its Lie algebra. We define the \defe{adjoint map}{adjoint!map} at the point $x\in G$ by 
\begin{equation}
\begin{aligned}
 \AD_x\colon G&\to G \\ 
 \AD_xy&=xyx^{-1} 
\end{aligned}
\end{equation}
Then we define 
\[
\dpt{Ad_x:=(d\AD_x)_e}{\mG}{\mG};
\]
the chain rule applied on $\AD_{xy}=\AD_x\circ\AD_y$ leads to $Ad_{xy}=Ad_x\circ Ad_y$, and thus we can see $Ad$ as a group homomorphism $\dpt{Ad}{G}{GL(\mG)}$, $Ad(x)=Ad_x$.

\begin{definition}
This homomorphism is the \defe{adjoint representation}{adjoint!representation!Lie group on its Lie algebra}\index{representation!adjoint} of the group $G$ in the vector space $\mG$.
\end{definition}


Finally, we define
\[
 \dpt{ad:=d(Ad)_1}{\mG}{L(\mG,\mG)}
\]
 where we identify $T_1GL(\mG)$ with $L(\mG,\mG)$.

\begin{lemma}
If $\dpt{f}{G}{G}$ is an automorphism of $G$ (i.e.: $f(xy)=f(x)f(y)$), then $df_e$ is an automorphism of $\mG$: $df[X,Y]=[df X,dfY]$
\label{lem:auto_1}
\end{lemma}

\begin{proof}
First, remark that $f(\AD_xy)=\AD_{f(x)}f(y)$. Now, $\Ad_x X=(d\AD_x)_eX$, so that one can compute:
\begin{equation}
\begin{split}
   df(\Ad_xX)&=\Dsdd{f(\AD_xX(t))}{t}{0}\\
             &=\Dsdd{   \AD_{f(x)}f(X(t))  }{t}{0}\\
	     &=(d\AD_{f(x)})_{f(e)}df X\\
	     &=\Ad_{f(x)}df X.
\end{split}
\end{equation}
On the other hand, we need to understand how does the $\ad$ work.
\[
  \ad XY=\Dsdd{\Ad_{X(t)}}{t}{0}Y=\Dsdd{\Ad_{X(t)}Y}{t}{0}
\]
because $\dpt{\Ad_{X(t)}}{\mG}{\mG}$ is linear, so that $Y$ can enter the derivation (for this, we identify $\mG$ and $T_X\mG$). Since $\Ad_{X(t)}Y$ is a path in $\mG$ the \emph{true space} is
\[   
(\ad X)Y=\Dsdd{ \Ad_{X(t)}Y }{t}{0}\in T_{[X,Y]}\mG\simeq\mG.
\]
For the same reason of linearity, $df$ can get in the derivative in the expression $df\Dsdd{  \Ad_{X(t)}Y  }{t}{0}$. Thus
\begin{equation}
\begin{split}
(\ad X)Y&=\Dsdd{  df\big(\Ad_{X(t)}Y\big)  }{t}{0}\\
        &=\Dsdd{  \Ad_{ f(X(t)) }df Y  }{t}{0}\\
	&=\Dsdd{ \Ad_{f(X(t))} }{t}{0}df Y\\
	&=\ad(dfX)df Y\\
	&=[dfX,df Y]
\end{split}
\end{equation}
because $f(X(t))$ is a path which gives $df X$.

\end{proof}

One can show that $[X,Y]$ is tangent to the curve
\begin{equation}
  c(t)=e^{-\sqrt{s}X}e^{-\sqrt{s}Y}e^{\sqrt{s}X}e^{\sqrt{s}Y}.
\end{equation}

\begin{lemma}
	In the case of Lie algebra, the bracket is given by the derivative of the adjoint action:
	\begin{equation}
		\Dsdd{ \Ad( e^{tX})Y }{t}{0}=[X,Y]
	\end{equation}
\end{lemma}

\begin{proof}
	Let us make $[\tilde X,\tilde Y]_e$ act on a function $f$. Using the definition \eqref{EqDefLieDerivativeVect} and the property of theorem \ref{ThoLieDerrComm}, we have
	\begin{equation}
		\begin{aligned}[]
			[\tilde X,\tilde Y]_ef&=\Dsdd{ (d\varphi_{-t}^X)\tilde Y }{t}{0}f\\
			&=\Dsdd{ (d\varphi_{-t}^X)_{\varphi_t^X(e)}\big( \tilde Y_{\varphi_t^X(e)} \big) }{t}{0}f\\
			&=\Dsdd{ \tilde Y_{ e^{tX}}\cdot(f\circ\varphi_{-t}^X) }{t}{0}
		\end{aligned}
	\end{equation}
	Now, we use the fact that, by definition, $\varphi_t^X(x)=x e^{tX}$, so that $\varphi_s^Y( e^{tX})= e^{tX} e^{sY}$ and we get
	\begin{equation}
		\begin{aligned}[]
			[\tilde X,\tilde Y]_ef&=\Dsdd{ \Dsdd{ f\big( \varphi_{-t}^X( e^{tX} e^{sY}) \big) }{s}{0} }{t}{0}\\
			&=\Dsdd { \Dsdd{ f( e^{tX} e^{sY} e^{-tX}) }{s}{0} } {t}{0}\\
			&=\Dsdd{ \Dsdd{ f\big(  e^{s\Ad( e^{tX})Y} \big) }{s}{0} }{t}{0}\\
			&=\Dsdd{ \big( \Ad( e^{tX})Y \big)_e\cdot f }{t}{0}
		\end{aligned}
	\end{equation}
	
\end{proof}


\section{Fundamental vector field}\label{sec:fond_vec}
%++++++++++++++++++++++++++++++++++++

If $\yG$ is the Lie algebra of a Lie group $G$ acting on a manifold $M$ (the action of $g$ on $x$ being denoted by $x\cdot g$), the \defe{fundamental vector field}{fundamental!vector field} associated with $A\in\yG$ is given by
\begin{equation}			\label{EqDefChmpFond}
   A^*_x=\Dsdd{ x\cdot e^{-tA} }{t}{0}.
\end{equation}
We always suppose that the action is effective. If the action of $G$ is transitive, the fundamental vectors at point $x\in M$ form a basis of $T_xM$. More precisely, we have the

\begin{lemma}
For any $v\in T_xM$, there exists a $A\in\yG$ such that $v=A^*_x$, in other terms
\[ 
  \Span\{ A^*_{x}\tq A\in\yG \}=T_{x}M.
\]
\label{LemFundSpansTan}
\end{lemma}

\begin{proof}
The vector $v$ is given by a path $v(t)$ in $M$. Since the action is transitive, one can write $v(t)=x\cdot c(t)$ for a certain path $c$ in $G$ which fulfills $c(0)=e$. We have to show that $v$ depends only on $c'(0)\in\yG$. We consider 
\begin{equation}  \label{eq_def_RGM}
\begin{aligned}
 R\colon G\times M&\to M \\ 
R(g,x)&= x\cdot g,
\end{aligned}
\end{equation}
so
\begin{equation}\label{eq:v_Rc}
   v=\Dsdd{ R(c(t),x) }{t}{0}=dR_{(e,x)}\big[  (d_tc(t),x)+(c(0),x)   \big].
\end{equation}

\end{proof}



\begin{lemma}
If $A$, $B\in\yG$ are such that $A^*=B^*$, and if the action is effective, then $A=B$.
\label{lem:As_Bs_A_B}
\end{lemma}

\begin{proof}
 We consider once again the map \eqref{eq_def_RGM} and we look at 
\[
  v=\Dsdd{ R(c(t),x) }{t}{0}
   =(dR)_{(e,x)}\Dsdd{ (c(t),x) }{t}{0},
\]
keeping in mind that $c(t)=e^{-tA}$. In order to treat this expression, we define
\begin{subequations}
\begin{align}
  \dpt{R_1}{G}{M},\quad  R_1(h)&=R(h,x),\\
  \dpt{R_2}{M}{M},\quad  R_2(y)&=R(g,y).
\end{align}
\end{subequations}
So
\[
  v=dR_1(X)+dR_2(0)=dR_1c'(0)
\]
and the assumption $A^*_x=B^*_x$ becomes $dR_1 A=dR_1 B$. This makes, for small enough $t$, $R_1(e^{tA}e^{-tB})=x\cdot e^{tA}e^{-tB}=x$; if the action is effective, it imposes $A=B$.
\end{proof}

\begin{lemma}
If we consider the action of a matrix group, $R_g$ acts on the fundamental field by
\[
  dR_g(A^*_{\xi})=\big( \Ad(g^{-1})A \big)^*_{\xi\cdot g}.
\]
\label{lem:dRgAstar}
\end{lemma}

\begin{proof}
Just notice that $e^{-t\Ad(g^{-1})A}=\AD_{g^{-1}}(e^{-tA})=g^{-1} e^{-tA}g$, thus
\begin{equation}
  \big( \Ad(g^{-1})A \big)^*_{\xi\cdot g}=\Dsdd{ \xi\cdot ge^{-t\Ad(g^{-1})A} }{t}{0}=dR_g(A^*_{\xi}).
\end{equation}
\end{proof}

%+++++++++++++++++++++++++++++++++++++++++++++++++++++++++++++++++++++++++++++++++++++++++++++++++++++++++++++++++++++++++++
\section{Exponential map}
%+++++++++++++++++++++++++++++++++++++++++++++++++++++++++++++++++++++++++++++++++++++++++++++++++++++++++++++++++++++++++++

A \defe{topological group}{topological!group} is a group $G$ equipped with a topological structure such that the maps $(x,y)\in G^2\to xy\in G$ and $x\in G\to x^{-1}\in G$ are continuous.

\begin{remark}\label{rem:ouvert}
From the existence of an unique inverse for any element of $G$, the multiplication and the inversion are also open maps.
\end{remark}

A \defe{Lie group}{Lie!group} is a group $G$ which is in the same times an analytic manifold such that the group operations (multiplication and inverse) are analytic. In particular, we \emph{do not} suppose that they are diffeomorphism.

The concept of normal neighbourhood will be widely used for the study of the relations between a Lie group and its algebra. Let $M$ be a differentiable manifold. If $V$ is a neighbourhood of zero in $T_pM$ on which the exponential $\dpt{\exp_p}{T_pM}{M}$ is a diffeomorphism, then $\exp_pV$ is  \defe{normal neighbourhood}{normal!neighbourhood} of $p$.

\begin{lemma}[\cite{Lie}]		\label{lemsur5d}
	Let $G$, $H$ be two Lie groups with algebras $\mG$ and $\mH$. Let $\dpt{\phi}{G}{H}$ be a homomorphism differentiable at $e$, the unit in $G$. Then for all $X\in\mG$, the following formula holds:
	\[
		\phi(\exp X)=\exp(d\phi_eX).
	\]
\end{lemma}

\begin{corollary}\label{Ad_e}
An useful formula:
\[
   \Ad(e^X)=e^{\ad X}.
\]
\end{corollary}

\begin{corollary}
Another useful corollary of lemma \ref{lemsur5d} is the particular case $\phi=\AD(e^X)$:
\[
   e^Xe^Ye^{-X}=e^{Ad(e^Y)X}.
\]
\label{cor:eXeYe-X}
\end{corollary}

\begin{proposition}
	Let $G$ be a connected Lie group.
	\begin{enumerate}

		\item
			All the left invariant vector fields are complete. That means that the map $X\mapsto  e^{X}$ is defined for every $X\in \mG$.
		\item
			The map $\exp\colon \mG\to G$ is a local diffeomorphism in a neighbourhood of $0$ in $\mG$.
	\end{enumerate}
\end{proposition}

\begin{proof}
	\begin{enumerate}

		\item
			The flow is a one parameter subgroup. Thus if $ e^{tX}$ is defined for $t\in[0,a]$, by composition, $ e^{2a}$ is defined. So $ e^{tX}$ is defined for every value of $t$ in $\eR$.
		\item
			Let us consider the manifold $G\times \mG$ and the vector field $\Xi$ defined by
			\begin{equation}
				\Xi_{(g,X)}=\tilde X_g\oplus 0\in T_g(G)\oplus T_X\mG\simeq T_{(g,X)}(G\times \mG).
			\end{equation}
			The flow of that vector field is given by
			\begin{equation}
				\Phi_t(g,X)=\big( g\exp(tX),X \big).
			\end{equation}
			In particular, $\Xi$ is a complete vector field, and we consider the global diffeomorphism
			\begin{equation}
				\begin{aligned}
					\Phi_1\colon G\times \mG&\to G\times \mG \\
					(g,X)&\mapsto \big( g\exp(X),X \big). 
				\end{aligned}
			\end{equation}
			On the point $(e,X)$ we have $\Phi_1(e,X)=(\exp(X),X)$. Thus the exponential is the projection on the first component of $\Phi_1(e,X)$ and we can write
			\begin{equation}
				\exp(X)=\pr_1\circ\Phi_1(e,X).
			\end{equation}
			It is a smooth function since both the projection and $\Phi_1$ are smooth.

			Now, the differential $(d\exp)_0$ is the identity on $\mG$, so that the theorem of inverse function makes $\exp$ a local diffeomorphism.
	\end{enumerate}	
\end{proof}


\begin{theorem}
For any $p\in M$, there exist a $\delta>0$ and a neighbourhood $W$ of $p$ in $M$ such that for every $q\in W$, we have

\begin{itemize}
\item $\exp_q$ is a diffeomorphism on $B\bdelta(0)\subset T_qM$,
\item $\exp_q B\bdelta(0)$ contains $W$
\end{itemize}
\end{theorem}
This theorem says that everywhere on a differentiable manifold, one can find a neighbourhood which is a normal neighbourhood of each of its points. Such a neighbourhood is said a \emph{totally} normal neighbourhood.

\begin{lemma}
In a Lie group, $e$ is an isolated fixed point for the inversion.
\end{lemma}

\begin{proof}
One can use an exponential map in a neighbourhood of $e$. In this neighbourhood, an element $g$ can be written as $g=e^X$ for a certain $X\in\lG$. The equality $g=g^{-1}$ gives (because the exponential is a diffeomorphism) $X=-X$, so that $X=0$ and $g=e$.
\end{proof}

\begin{lemma}
Let $\lG$ be a Lie algebra and $A$, a linear operator on $\lG$ (see as a common vector space) such that $\forall t\in\eR$, the map $e^{tA}$ is an automorphism of $\lG$. Then $A$ is a derivation of $\lG$.
\label{lem:autom_derr}
\end{lemma}

\begin{proof}
Let us consider $X$, $Y\in\lG$;  the assumption is 
\[
  e^{tA}[X,Y]=[e^{tA}X,e^{tA}Y]. 
\]
Since $e^{tA}$ is a linear map, it has a ``good behavior''\ with the derivations: 
\[
\Dsddc{e^{tA}[X,Y]}{t}{0}=\Dsddc{e^{tA}}{t}{0}[X,Y]=A[X,Y].
\]
Using on the other hand the linearity of $\ad$, we can see
\[
  (\ad(e^{tA}X))(e^{tA}Y)
\]
as a product ``matrix times vector''. Then
\begin{equation}
\begin{split}
  \Dsddc{[e^{tA}X,e^{tA}Y]}{t}{0}&=\Dsddc{(\ad e^{tA}X)Y}{t}{0}+\Dsddc{(\ad X)(e^{tA}Y)}{t}{0}\\
                                 &=(\ad AX)Y+(\ad X)(AY).
\end{split}
\end{equation}
Finally, $A[X,Y]=[AX,Y]+[X,AY]$.

\end{proof}

As notational convention, if $G$ and $H$ are Lie groups, their Lie algebra are denoted by $\lG$ and $\lH$.

\begin{lemma}		\label{LemAlgEtGroupesGenere}
	Let $\lG$ be a Lie algebra ans $\lS$ be a subset of $\lG$. The algebra of the group generated by $ e^{\lS}$ is the algebra generated by $\lS$.
\end{lemma}


\section{Universal enveloping algebra}  \label{subsec:env_alg}
%----------------------------------------

Let $\mA$ be a Lie algebra. One knows that the composition law $(X,Y)\to[X,Y]$ is often non associative. In order to build an associative Lie algebra which ``looks like''\ $\mA$, one considers $T(\mA)$, the tensor algebra of $\mA$ (as vector space) and $\mJ$ the two-sided ideal in $T(\mA)$ generated by elements of the form 
\[
   X\otimes Y-Y\otimes X -[X,Y]
\]
for $X$, $Y\in\mA$. The \defe{universal enveloping algebra}{universal!enveloping algebra} of $\mA$ is the quotient \nomenclature{$U(\mA)$}{Universal enveloping algebra} 
\begin{equation}
     U(\mA)=T(\mA)/\mJ.
\end{equation}
For $X\in\mA$, we denote by $X^*$\nomenclature{$X^*$}{Image of a tensor in the universal enveloping algebra} the image of $X$ by canonical projection $\dpt{\pi}{T(\mA)}{U(\mA)}$ and by $1$ the unit in $U(\mA)$. One has $1\neq 0$ if and only if $\mA\neq\{0\}$.

\begin{proposition}[\cite{Helgason}]
Le $V$ be a vector space on $K$. Then there is a natural bijection between the representations of $\mA$\index{representation! of $U(\mA)$} on $V$ and the ones of $U(\mA)$ on $V$. If $\rho$ is a representation of $\mA$ on $V$, the corresponding $\rho^*$ of $U(\mA)$ is given by
\[
   \rho(X)=\rho^*(X^*)
\]
($X\in\mA$).
\end{proposition}

Let $\{X_1,\ldots,X_n\}$ be a basis of $\mA$. For a $n$-uple of complex numbers $(t_i)$, one defines
\begin{equation}
X^*(t)=\sum_{i=1}^nt_iX^*_i.
\end{equation}
On the other hand, we consider a $n$-uple of positive integers $M=(m_1+\ldots m_n)$, and the notation 
\begin{equation}
\begin{split}
   |M|&=m_1+\ldots+m_n\\
   t^M&=t_1^{m_1}\cdots t_n^{m_n}.
\end{split}
\end{equation}

When $|M|>0$, we denote by $X^*(M)\in U(\mA)$ the coefficient of $t^M$ in the expansion of $(|M|!)^{-1} (X^*(t))^{|M|}$. If $|M|=0$, the definition is $X^*(0)=1$. Once again a proposition without proof.

\begin{proposition}
The smallest vector subspace of $U(\mA)$ which contains all the elements of the form $X^*(M)$ is $U(\mA)$ itself:
\[
   U(\mA)=\Span\{ X^*(M):M\in \eN^n \}.
\]
\end{proposition}

\begin{corollary} \label{cor:/24}
    Let $\mA$ be a Banach algebra of dimension $n$, $\mB$ a Banach subalgebra of dimension $n-r$ and a basis $\{X_1,\ldots,X_n\}$ of $\mA$ such that the $n-r$ last basis vectors are in $\mB$. We denotes by $B$ the vector subspace of $U(\mA)$ spanned by the elements of the form $X^*(M)$ with $m=(0,\ldots,0,m_{r+1},\ldots,m_n)$. Then $B$ is a subalgebra of $U(\mA)$.
\end{corollary}

\begin{definition}
Two Lie groups $G$ and $G'$ are \defe{isomorphic}{isomorphism!of Lie groups} when there exists a differentiable group isomorphism between $G$ and $G'$.

They are \defe{locally isomorphic}{locally!isomorphic!Lie groups} when there exists neighbourhoods $\mU$ and $\mU'$of $e$ and $e'$ and a differentiable diffeomorphism $\dpt{f}{\mU}{\mU'}$ such that

$\forall x,y,xy\in\mU$, $f(xy)=f(x)f(y)$, \\and

$\forall x',y',x'y'\in\mU'$, $f^{-1}(x'y')=f^{-1}(x')f^{-1}(y')$.
\end{definition}

The following universal property of the \emph{universal} enveloping algebra explains the denomination:
\begin{proposition}
Let $\dpt{\sigma}{\mG}{\mU(\mG)}$ the canonical inclusion and $A$ an unital complex associative algebra. A linear map $\dpt{\varphi}{\mG}{A}$ such that
\begin{equation}
\varphi[X,Y]=\varphi(X)\varphi(Y)-\varphi(Y)\varphi(X)
\end{equation}
can be extended in only one way to an algebra homomorphism $\dpt{\varphi_0}{\mU(\mG)}{A}$ such that $\varphi_0\circ\sigma=\varphi$ and $\varphi(1)=1$
\label{prop:extunifmap}
\end{proposition}
For a proof, see \cite{Knapp_reprez}.

\subsection{Adjoint map in \texorpdfstring{$\mU(\mG)$}{U(G)}}   \label{ssadjunif}
%/////////////////////////////////////////////

We know that $\dpt{ \Ad(g) }{ \mG }{ \mG }$ fulfils
\[ 
  \Ad(g)[X,y]=[  \Ad(g)X,\Ad(g)Y  ],
\]
and we can define $\dpt{ \Ad(g) }{ \mG }{ \mU(\mG) }$ by $\Ad(g)X=X$ where in the right hand side, $X$ denotes the class of $X$ for the quotients of the tensor algebra which defines the universal enveloping algebra.

When $[A,B]$ is seen in $\mU(\mG)$, we have $[A,B]=A\otimes B-B\otimes A$. Then $\dpt{ \Ad(g) }{ \mG }{ \mU(\mG) }$ fulfils proposition \ref{prop:extunifmap} and is extended in an unique way to $\dpt{ \Ad(g) }{ \mU(\mG) }{ \mU(\mG) }$ with $\Ad(g)1=1$.
 
\begin{lemma} 
    If $D\in\mU(\mG)$, the following properties are equivalent:
    \begin{itemize}
        \item $D\in\mZ(\mG)$
        \item $D\otimes X=X\otimes D$ for all $X\in \mG$
        \item $e^{\ad X}D=D$ for all $X\in\mG$
        \item $\Ad(g)D=D$ for all $g\in G$.
    \end{itemize}
     \label{lem:equivDAd}
\end{lemma}

%---------------------------------------------------------------------------------------------------------------------------
\subsection{Invariant fields}
%---------------------------------------------------------------------------------------------------------------------------

If $X\in\lG$, we have the associated left invariant vector field on $G$ given by $\tilde X_x=dL_xX$. That field is left invariant as operator on the functions because
\begin{equation}
    \tilde X_x(u)=\tilde X_e(L^*_xu)
\end{equation}
as the following computation shows
\begin{equation}
        \tilde X_e(L^*u)=\Dsdd{ (L_x^*u)\big(  e^{tX} \big) }{t}{0}
        =\Dsdd{ u\big( x e^{tX} \big) }{t}{0}
        =\Dsdd{ u\big( \tilde X_x(t) \big) }{t}{0}
        =\tilde X_x(u)
\end{equation}
because the path defining $\tilde X_x$ is $x e^{tX}$.

We can perform the same construction in order to build left invariant fields based on $\mU(\lG)$. If $X$ and $Y$ are elements of $\lG$, the  differential operator on $ C^{\infty}(G)$ associated to $XY\in\mU(\lG)$ is given by
\begin{equation}
    (XY)(f)=\DDsdd{ f\big( X(s)Y(t) \big) }{t}{0}{s}{0}
\end{equation}
The path defining the field $\widetilde{XY}$ is
\begin{equation}
    \widetilde{XY}_x=xX(s)Y(t).
\end{equation}
Thus we have
\begin{equation}        \label{EqInvarUgField}
    \widetilde{(XY)}_e(L^*u)=\widetilde{(XY)}_xu
\end{equation}

\begin{lemma}       \label{LemAdesthioo}
    If \( X,Y\in\lG\) we have
    \begin{equation}
        [\ad(X),\ad(Y)]=\ad([X,Y]).
    \end{equation}
\end{lemma}

\begin{proof}
    Let \( f\in\lG\) and compute the action of \( [\ad(X),\ad(Y)]\):
    \begin{subequations}
        \begin{align}
            [\ad(X),\ad(Y)]f&=\ad(X)[Yf,fY]-\ad(Y)(Xf-fX)\\
            &=(XY-YX)f+f(YX-XY)\\
            &=\ad([X,Y])f.
        \end{align}
    \end{subequations}
\end{proof}

%---------------------------------------------------------------------------------------------------------------------------
                    \subsection{Representation of Lie groups}
%---------------------------------------------------------------------------------------------------------------------------

\begin{proposition}
    Let $G$ be a Lie group and $\mG$ its Lie algebra. A representation $\varphi\colon G\to \End(V)$ of the group induces a representation $\phi\colon \mU(\mG)\to \End(V)$ of the universal enveloping algebra with the definitions
    \begin{subequations}
        \begin{align}
            \phi(X)     &=d\varphi_e(X),\\
            \phi(XY)    &=\phi(X)\circ\phi(Y)
        \end{align}
    \end{subequations}
    where $e$ is the unit in $G$ and $X$, $Y$ are any elements of $\mG$.
\end{proposition}

\begin{proof}
    We have
    \begin{equation}
        \phi(X)=\Dsdd{ \varphi( e^{tX})v }{t}{0}=d\varphi_e(X)v.
    \end{equation}
    Notice that, by linearity of the action of $\varphi( e^{tX})$ on $v$, one can leave $v$ outside the derivation. Now, neglecting the second order terms in $t$ in the derivative, and using the Leibnitz formula, we have
    \begin{equation}
        \begin{aligned}[]
            \phi([X,Y])v    &=  \Dsdd{ \varphi( e^{tXY} e^{-tXY}) }{t}{0}v\\
                    &=  \Dsdd{ \varphi( e^{tXY})\varphi(\mtu) }{t}{0}v+\Dsdd{ \varphi(\mtu)\varphi( e^{-tXY}) }{t}{0}v\\
                    &=  \phi(XY)v-\phi(YX)v\\
                    &=  \big( \phi(X)\phi(Y)-\phi(Y)\phi(X) \big)v\\
                    &=  [\phi(X),\phi(Y)]v,
        \end{aligned}
    \end{equation}
    which is the claim.
\end{proof}

\section{Lie subgroup}
%------------------------

Now a great theorem without proof:
\begin{theorem} \label{tho:loc_isom}
Two Lie groups are locally isomorphic if and only if their Lie algebras are isomorphic.
\end{theorem}
%TODO : a proof

\begin{definition}
Let $G$ be a Lie group. A submanifold $H$ of $G$ is a \defe{Lie subgroup}{Lie!subgroup} of $G$ when
\begin{enumerate}
\item as group, $H$ is a subgroup of $G$,
\item $H$ is a topological group.
\end{enumerate}
\end{definition}

\begin{remark}\label{rem:sub_Lie}
The definition doesn't include that $H$ has the same topology as $G$ (or the induced one). In some literature, the definition of a Lie subgroup include the fact for $H$ to be a topological subgroup. This choice make some proofs much easier and others more difficult; be careful when you try to compare different texts.
\end{remark}

\begin{lemma}
A Lie subgroup is a Lie group.
\end{lemma}
(without proof)

\begin{theorem}		\label{ThoSubGpSubAlg}		\label{tho:gp_alg}
If $G$ is a Lie group, then
\begin{enumerate}
\item\label{ThoSubGpSubAlgi} if $\lH$ is the Lie algebra of a Lie subgroup $H$ of $G$, then it is a subalgebra of $\lG$,
\item Any subalgebra of $\lG$ is the Lie algebra of one and only one connected Lie subgroup of $G$.
\end{enumerate}

\begin{probleme}
À mon avis, il faut dire ``connexe et simplement connexe'', et non juste ``connexe''.
\end{probleme}

\end{theorem}
\begin{proof}

\subdem{First item}
Let $\dpt{i}{H}{G}$ be the identity map; it is a homomorphism from $H$ to $G$, thus $di_e$ is a homomorphism from $\lH$ to $\lG$. Conclusion: $\lH$ is a subalgebra of $\lG$. 

\subdem{Characterization for $\lH$}
Before to go on with the second point, we derive an important characterization of $\lH$ :
\begin{equation}\label{eq:path_alg}
\lH=\{X\in\lG:\text{the map } t\to\exp tX\text{ is a path in $H$}\}.
\end{equation}
For that, consider $\dpt{\exp_H}{\lH}{H}$ and $\dpt{\exp_G}{\lG}{G}$; from unicity of the exponential, for any $X\in\lH$, $\exp_HX=\exp_GX$, so that one can simply write ``$\exp$''\ instead of ``$\exp_h$''\ or ``$\exp_G$''.

Now, if $X\in\lH$, the map $t\to\exp tX$ is a curve in $H$. But it is not immediately clear that such a curve in $H$ is automatically build from a vector in $\lH$ rather than in $\lG$.  More precisely, consider a $X\in\lG$ such that $t\to\exp tX$ is a path (continuous curve) in H. By lemma \ref{lem:var_cont_diff}, the map $t\to\exp tX$ is differentiable and thus by derivation, $X\in\lH$.
The characterisation \eqref{eq:path_alg} is proved.

Thus $\lH$ is a Lie subalgebra of $\lG$.

\subdem{Second item}
For the second part, we consider $\lH$ any subalgebra of $\lG$ and $H$, the smallest subgroup of $G$ which contains $\exp\lH$. We also consider a basis $\{X_1,\ldots,X_n\}$ of $\lG$ such that $\{X_{r+1},\ldots,X_n\}$ is a basis of $\lH$.

By corollary \ref{cor:/24}, the set of linear combinations of elements of the form $X(M)$ with $M=(0,\ldots,0,m_{r+1},\ldots,m_r)$ form a subalgebra of $U(\lG)$. If $X=x_1X_1+\ldots+x_nX_n$, we define $|X|=(x_1^2+\ldots+x_n^2)^{1/2}$ ($x_i\in\eR$).

Let us consider a $\delta>0$ such that $\exp$ is a diffeomorphism (normal neighbourhood) from $B_{\delta}=\{X\in\lG:|X|<\delta\}$ to a neighbourhood $N_e$ of $e\in G$ and such that $\forall x,y,xy\in N_e$,
\begin{equation}\label{eq:coord_xy}
   (xy)_k=\sum_{M,N}C^{[k]}_{MN}x^My^N
\end{equation}
holds\footnote{The validity of this second condition is assured during the proof of theorem \ref{tho:loc_isom} which is not given here.}. We note $V=\exp(\lH\cap B_{\delta})\subset N_e$. The map
\[
   \exp(x_{r+1}X_{r+1}+\ldots+x_nX_n)\to(x_{r+1},\ldots,x_n)
\]
is a coordinate system on $V$ for which $V$ is a connected manifold. But $\lH\cap B_{\delta}$ is a submanifold of $B_{\delta}$, then $V$ is a submanifold of $N_e$ and consequently of~$G$.

Let $x$, $y\in V$ such that $xy\in N_e$ (this exist: $x=y=e$); the canonical coordinates of $xy$ are given by \eqref{eq:coord_xy}. Since $x_k=y_k=0$ for $1\leq k\leq r$, $(xy)_k=0$ for the same $k$ because for $(xy)_k$ to be non zero, one need $m_1=\ldots=m_r=n_1=\ldots=n_r=0$ -- otherwise, $x^M$ or $y^N$ is zero. Now we looks at $C^{[k]}_{MN}$ for such a $k$ (say $k=1$ to fix ideas) : $[k]=(\delta_{11},\ldots,\delta_{1k})=(1,0,\ldots,0)$ and by definition of the $C$'s,
\[
   X(M)X(N)=\sum_PC_{MN}^PX(P).
\]
But we had seen that the set of the $X(A)$ with $A=(0,\ldots,0,a_{r+1},\ldots,a_n)$ form a subalgebra of $U(\lG)$. Then, only terms with $P=(0,\ldots,0,p_{r+1},\ldots,p_n)$ are present in the sum; in particular, $C_{MN}^{[k]}=0$ for $k=1,\ldots,r$. Thus $VV\cap N_e\subset V$.

The next step is to consider $\mV$, the set of all the subset of $H$ whose contains a neighbourhood of $e$ in $V$. We can check that this fulfils the six axioms of a topological group\index{topological!group} :

\begin{enumerate}
\item The intersection of two elements of $\mV$ is in $\mV$;
\item the intersection of all the elements of $\mV$ is $\{e\}$;
\item any subset of $H$ which contains a set of $\mV$ is in $\mV$;
\item If $\mU\in\mV$, there exists a $\mU_1\in\mV$ such that $\mU_1\mU_1\subset\mU$ because $VV\cap N_e\subset V$;
\item if $\mU\in\mV$, then $\mU^{-1}\in\mV$ because the inverse map is differentiable and transforms a neighbourhood of $e$ into a neighbourhood of $e$;
\item if $\mU\in\mV$ and $h\in H$, then $h\mU h^{-1}\in\mV$.
\end{enumerate}

To see this last item, we denote by $\log$ the inverse map of $\dpt{\exp}{B_{\delta}}{N_e}$. By definition of $V$, it sends $V$ on $\lH\cap B_{\delta}$. If $X\in\lG$, there exists one and only one $X'\in\lG$ such that $he^{tX}h^{-1}=e^{tX'}$ for any $t\in\eR$. Indeed we know that $he^{X}h^{-1}=e^{\Ad_hX}$, then $X'$ must satisfy $e^{tX'}=e^{\Ad_htX}$. If it is true for any $t$, then, by derivation, $X'=\Ad_hX$.

The map $X\to X'$ is an automorphism of $\lG$ which sent $\lH$ on itself. So one can find a $\delta_1$ with $0<\delta_1<\delta$ such that 
\[
   h\exp({B_{\delta_1}\cap\lH})h^{-1}\subset V.
\]
Indeed, $he^{\lH} h^{-1}\subset\lH$, so that taking $\delta_1<\delta$, we get the strict inclusion. We can choose $\delta_1$ even smaller to satisfy $he^{B_{\delta_1}}h^{-1}\subset N_e$. Since the map $X\to\log(he^{X}h^{-1})$ from $B_{\delta_1\cap\lH}$ to $B_{\delta}\cap\lH$ is regular, the image of $B_{\delta_1}\cap\lH$ is a neighbourhood of $0$ in $\lH$. Thus $he^{B_{\delta_1}\cap\lH}h^{-1}$ is a neighbourhood of $e$ in $V$. Finally, $h\mU h^{-1}\in\mV$ and the last axiom of a topological group is checked.

This is important because there exists a topology on $H$ such that $H$ becomes a topological group and $\mV$ is a family of neighbourhood of $e$ in $H$. In particular, $V$ is a neighbourhood of $e$ in $H$.

For any $z\in G$, we define the map $\dpt{\phi_z}{zN_e}{B_{\delta}}$ by
\begin{equation}
  \phi_z(ze^{x_1X_1+\ldots+x_nX_n})=(x_1,\ldots,x_n),
\end{equation}
and we denote by $\varphi_z$ the restriction of $\phi_z$ to $zV$. If $z\in H$, then $\varphi_z$ sends the neighbourhood $zV$ of $z$ in $H$ to the open set $B_{\delta}\cap\lH$ in $\eR^{n-r}$. Indeed, an element of $zV$ is a $ze^Z$ with $Z\in\lH\cap B_{\delta}$ which is sent by $\varphi_z$ to an element of $\lH\cap B_{\delta}$. (we just have to identify $x_1X_1+\ldots+x_nX_n$ with $(x_1,\ldots,x_n)$).

Moreover, if $z_1,z_2\in H$, the map $\varphi_{z_1}\circ\varphi_{z_2}^{-1}$ is the restriction to an open subset of $\lH$ of $\phi_{z_1}\circ\phi_{z_2}$. Then $\varphi_{z_1}\circ\varphi_{z_2}^{-1}$ is differentiable. Conclusion: $(H,\varphi_z: z\in H)$ is a differentiable manifold.

Recall that the definition of $\lH$ was to be a subalgebra of $\lG$; therefore $V=e^{\lH\cap B_{\delta}}$ is a submanifold of $G$. But the left translations are diffeomorphism of $H$ and $H$ is the smallest subgroup of $G$ containing $e^{\lH}$. Thus $H$ is a manifold on which the multiplication is diffeomorphic and consequently, $H$ is a Lie subgroup of $G$.

Rest to prove that the Lie algebra of $H$ is $\lH$ and the unicity part of the theorem.

We know that $\dim H=\dim\lH$ and moreover for $i>r$, the map $t\to\exp tX_i$ is a curve in $H$. Now, the fact that $\lH$ is the set of $X\in\lG$ such that $t\to\exp tX$ is a path in $H$ show that $X_i\in\lH$. Then the Lie algebra of $H$ is $\lH$ and $H$ is a connected group because it is generated by $\exp\lH$ which is a connected neighbourhood of $e$ in $H$.

We turn our attention to the unicity part. Let $H_1$ be a connected Lie subgroup of $G$ such that $T_eH_1=\lH$. Since $\exp_{\lH}X=\exp_{\lH_1}X$, $H=H_1$ as set. But $\exp$ is a differentiable diffeomorphism from a neighbourhood of $0$ in $\lH$ to a neighbourhood of $e$ in $H$ and $H_1$, so as Lie groups, $H$ and $H_1$ are the same.


\end{proof}

We state a corollary without proof :

\begin{corollary} 
If $H_1$ and $H_2$ are two Lie subgroups of the Lie group $G$  such that $H_1=H_2$ as topological groups, then $H_1=H_2$ as Lie groups.
 \label{cor:top_subgroup}
\end{corollary}

\begin{proposition}
Let $G_1$ and $G_2$ be two Lie groups with same Lie algebra such that $\pi_0(G_1)=\pi_0(G_2)$ and $\pi_1(G_1)=\pi_1(G_2)$, then $G_1$ and $G_2$ are isomorphic.
\end{proposition}

\begin{proof}
The assumptions of equality of Lie algebras and of the $\pi_0$ make that the universal covering $\tilde G_1$ and $\tilde G_2$ of $G_1$ and $G_2$ are the same. But we know that $G_i=\tilde G_i/\pi_1(G_i)$. Now equality $\pi_1(G_1)=\pi_1(G_2)$ concludes that $G_1=G_2$.
\end{proof}

\begin{lemma}
Let $\lG$ admit a direct sum decomposition (as vector space) $\lG=\lM\oplus\lN$. Then there exists open and bounded neighbourhoods $\mU_m$ and $\mU_n$ of $0$ in $\lM$ and $\lN$ such that the map 
		\begin{equation}
		\begin{aligned}
			\phi \colon \mU_m\times\mU_n &\to G\
			(A,B)&\mapsto e^Ae^B
		\end{aligned}
	\end{equation}	
is a diffeomorphism between $\mU_m\times\mU_n$ and an open neighbourhood of $e$ in $G$.
 \label{lem:decomp}
\end{lemma}


\begin{proof}
Let $\{X_1,\ldots,X_n\}$ be a basis of $\lG$ such that $X_i\in\lM$ for $1\leq i\leq r$ and $X_j\in\lN$ for $r<j\leq n$. We consider $\{t_1,\ldots,t_n\}$, the canonical coordinates of $\exp(x_1X_1+\ldots+x_rX_r)\exp(x_{r+1}X_{r+1}+\ldots+x_nX_n)$ in this coordinate system. By properties of the exponential, the function $\varphi_j$ defined by $t_j=\varphi_j(x_1,\ldots,x_n)$ is differentiable at $(0,\ldots,0)$. If $x_i=\delta_{ij}s$, then $t_i=\delta_{ij}s$ and the Jacobian of 
\[
   \dsd{(\varphi_1,\ldots,\varphi_n)}{(x_1,\ldots,x_n)}
\]
is $1$ for $x_1=\ldots=x_n=0$. Thus $d\varphi_e$ is a diffeomorphism and so $\varphi$ is a locally diffeomorphic.
\end{proof}

\begin{theorem}  
Let $G$ be a Lie group whose Lie algebra is $\lG$ and $H$, a closed subgroup (not specially a \emph{Lie} subgroup) of $G$. Then there exists one and only one analytic structure on $H$ for which $H$ is a topological Lie subgroup of $G$.
\label{tho:diff_sur_ferme}
\end{theorem}

\begin{remark}
A \textit{topological} Lie subgroup\index{topological!Lie subgroup} is stronger that a common Lie subgroup because it needs to be a topological subgroup: it must carry \emph{exactly} the induced topology. In our definition of a Lie group, this feature doesn't appears. 
\end{remark}

\begin{proof}
   Let $\lH$ be the subspace of $\lG$ defined by 
\begin{equation}\label{eq:lH_de_G}
  \lH=\{X\in\lG\tq \forall t\in\eR,\, e^{tX}\in H\}.
\end{equation}
We begin to show that $\lH$ is a subalgebra of $\lG$; i.e. to show that $t(X+Y)\in\lH$ and $t^2[X,Y]\in\lH$ if $X$, $Y\in\lH$. Remark that $X\in\lH$ and $s\in\eR$ implies $sX\in\lH$. Consider now $X$, $Y\in\lH$ and the classical formula :
\begin{subequations}
\begin{align}
\left(  \exp(\frac{t}{n}X)\exp(\frac{t}{n}Y)  \right )^n
                       =\exp( t(X+Y)+\frac{t^2}{2n}[X,Y]+o(\frac{1}{n^2}) ),\\
\left(  \exp(-\frac{t}{n}X)\exp(-\frac{t}{n}Y)\exp(\frac{t}{n}X)\exp(\frac{t}{n}Y)   \right)^{n^2}
                       =\exp\left( t^2[X,Y]+o(\frac{1}{n})\right).
\end{align}
\end{subequations}
The left hand side of these equations are in $H$ for any $n$; but, since $H$ is closed, it keeps in $H$ when $n\to\infty$. The right hand side, at the limit, is just $\exp(t(X+Y))$ and $\exp(t^2[X,Y])$, which keeps in $H$ for any $t$. Thus $X+Y$ and $[X,Y]$ belong to $\lH$. The space $\lH$ is thus a Lie subalgebra of $\lG$.

Let $H^*$ be the connected Lie subgroup of $G$ whose Lie algebra is $\lH$ (existence and unicity from \ref{tho:gp_alg}). From the proof of theorem \ref{tho:gp_alg}, we know that $H^*$ is the smallest subgroup of $G$ containing $\exp\lH$, then it is made up from products and inverses of elements of the type $e^X$ with $X\in\lH$, and thus is is included in $H$ by definition of $\lH$. So, $H^*\subset H$.
  
We will show that if we put on $H^*$ the induced topology from $G$ and if $H_0$ denotes the identity component of $H$, then $H^*=H_0$ as topological groups. For this, we first have to show the equality as set and then prove that if $N$ is a neighbourhood of $e$ in $H^*$, then it is a neighbourhood of $e$ in $H_0$. In facts, the equality as set can be derives from this second fact. Indeed, since $H_0$ is a connected topological group, it is generated by any neighbourhood of $e$, so if one can show that any neighbourhood $N$ of $e$ in $H^*$ is a neighbourhood of $e$ in $H$, then $H^*$ is a neighbourhood of $e$ in $H_0$ and then $H_0$ should be generated by $H^*$, so that $H_0\subset H^*$ (as set). Moreover, the most general element of $H^*$ is product and inverse of $e^X$ with $X\in\lH$ and $e^X$ is connected to $e$ by the path $e^{tX}$ ($\dpt{t}{1}{0}$). Then $H^*\subset H_0$, and $H^*=H_0$ as set. Immediately, $H^*=H_0$ as topological groups from our assertion about neighbourhoods of $e$. Let us now prove it.
  
We consider a neighbourhood $N$ of $e$ in $H^*$ and suppose that this is not a neighbourhood of $e$ in $H$. Thus there exists a sequence $c_k\in H\setminus N$ such that $c_k\to e$ in the sense of the topology on $G$. Indeed, a neighbourhood of $e$ in the sense of $H$ must contains at least a point which is not in $N$ because if we have an open set of $H$ around $e$ included in $N$, then $N$ is a neighbourhood of $e$ for $H$. So we consider a suitable sequence of such open set around $e$ and one element not in $N$ in each of them. There is the $c_k$'s\quext{Je crois qu'on utilise l'axiome du choix.}.

Using lemma \ref{lem:decomp} with a decomposition $\lG=\lH\oplus\lM$ (i.e. $\lM$ : a complementary for $\lH$ for $\lG$), one can find sequences $A_k\in\mU_m$ and $B_k\in\mU_n$ such that
\[
   c_k=e^{A_k}e^{B_k}.
\]
Here, $\mU_m$ is an open neighbourhood of $0$ in $\lM$ and $\mU_h$, an open neighbourhood of $0$ in $\lH$.

As $e^{B_k}\in N$ and $c_k\in H\setminus N$, $A_k\neq 0$ and $\lim A_k=\lim B_k=0$ (because $(A,B)\to e^Ae^B$ is a diffeomorphism and $e^0e^0=e$ -- and also because all is continuous and thus has a good behaviour with respect to the limit). The set $\mU_m$ is open and bounded --this is a part of the lemma. Then there exist a sequence of positive reals numbers $r_k\in$ such that $r_kA_k\in\mU_m$ and $(r_k+1)A_k\notin\mU_m$. We know that $\mU_m$ is a bounded open subset of the vector space $\lM$, then the whole sequences $r_kA_k$ and $(r_k+1)A_k$ are in a compact domain of $\lM$. Then --by eventually considering subsequences-- there are no problems to consider limits of these sequences in $\lM$ : $r_kA_k\to A\in\lM$ (not necessary in $\mU_m$). Since $A_k\to 0$, the point $A$ is the common limit of $r_kA_k\in\mU_m$ and of $(r_k+1)A_k\notin\mU_m$. Thus $A$ is in the boundary of $\mU_m$; in particular, $A\neq 0$.

On the other hand, consider two integers $p,q$ with $q>0$. One can find sequences $s_k,t_k\in\eN$ and $0\leq t_k<q$ such that $pr_k=qs_k+t_k$. It is clear that 
\begin{equation}
  \lim_{k\to\infty}\frac{t_k}{q}A_k=0,
\end{equation}
thus
\[
   \exp \frac{p}{q}A=\lim \exp\frac{pr_k}{a}A_k=\lim (\exp A_k)^{s_k},
\]
which belongs to $H$. By continuity, $\exp tA\in H$ for any $t\in\eR$ and finally $A\in\lH$; this contradict $A\neq 0$ so that $A\in\lM$ (because by definition, $A\in\lM$ and the sum $\lG=\lH\oplus\lM$ is direct).

By its definition, $H^*$ has an analytic structure of Lie subgroup of $G$; but we had just proved that the induced topology from $G$ is the one of $H_0$ which by definition is a submanifold of $G$. So the set $H_0=H^*$ becomes a submanifold of $G$ whose topology is compatible with the analytic structure: thus it is a Lie subgroup of $G$. From analyticity, this structure is extended to the whole $H$.

\begin{probleme}
Est-ce bien vrai, tout \c ca ? En particulier, je n'utilise pas que $H_0$ est ouvert dans $H$ (ce qui est un tho de topo classique : je ne vois pas pourquoi Helgason fait tout un cin\'ema --que je ne comprends pas-- dessus). En prenant $N=H^*$, on a juste d\'emontr\'e que $H_0$ est un voisinage de $e$ dans $H$, mais ça, on le savait bien avant.
\end{probleme}

The unicity part comes from the corollary \ref{cor:top_subgroup}.
\end{proof}


With the notations and the structure of theorem \ref{tho:diff_sur_ferme}, the subgroup $H$ is discrete if and only if $\lH=\{0\}$. Indeed, recall the definition \eqref{eq:lH_de_G} :
\[
  \lH=\{X\in\lG: \forall t\in\eR, e^{tX}\in H\},
\]
and the fact that there exists a neighbourhood of $e$ in $H$ on which the exponential map is a diffeomorphism.

\begin{remark}
This fact should not be placed after the following lemma. In fact, we use here just the existence of normal neighbourhood (which is a common result) while the following lemma gives much more than normal neighbourhood.
\end{remark}

The lemma (without proof) :

\begin{lemma} 
 Let $G$ be a Lie group and $H$, a Lie subgroup of $G$ ($\lG$ and $\lH$ are the corresponding Lie algebras). If $H$ is a topological subspace of $G$ (cf remark \ref{rem:sub_Lie}), then there exists an open neighbourhood $V$ of $0$ in $\mG$ such that
 \begin{enumerate}
 \item $\exp$ is a diffeomorphism between $V$ and an open neighbourhood of $e$ in~$G$,
 \item $\exp(V\cap\lH)=(\exp V)\cap H$.
 \end{enumerate}
\label{lem:sugroup_normal}
\end{lemma}

Now a theorem with proof.

\begin{theorem}
Let $G$ and $H$ be two Lie groups and $\dpt{\varphi}{G}{H}$ a continuous homomorphism. Then $\varphi$ is analytic.
\end{theorem}

\begin{proof}
The Lie algebra of the product manifold $G\times H$ as $\lG\times\lH$ is given in \ref{lemLeibnitz}. We define
\begin{equation}
  K=\{(g,\varphi(g)):g\in G\}\subset G\times H.
\end{equation}
It is clear that $K$ is closed in $G\times H$ because $G$ is closed and $\varphi$ is continuous.
By theorem \ref{tho:diff_sur_ferme}, there exists an unique differentiable structure on $G\times H$ such that $K$ is a topological Lie subgroup of $G\times H$ (i.e. : Lie subgroup + induced topology). The Lie algebra of $K$ is
\begin{equation}
  \lK=\{(X,Y)\in\lG\times\lH:\forall t\in\eR, (e^{tX},e^{tY})\in K\}.
\end{equation}
Let $N_0$ be an open neighbourhood of $0$ in $\lH$ such that $\exp$ is diffeomorphic between $N_0$ and an open neighbourhood $N_e$ of $e$ in $H$. We define $M_0$ and $M_e$ in the same way, for $G$ instead of $H$. We can suppose $\varphi(M_e)\subset N_e$ : if it is not, we consider a smaller $M_e$ : the openness of $N_e$ and the continuity of $\varphi$ make it coherent.

The lemma \ref{lem:sugroup_normal} allow us to consider $M_0$ and $N_0$ small enough to say that 
\[
   \dpt{\exp}{(M_0\times N_0)\cap\lK}{(M_e\times N_e)\cap K}
\]
is diffeomorphic. Now, we are going to show that for any $X\in\lG$, there exists an unique $Y\in\lH$ such that $(X,Y)\in\lK$. The unicity is easy: consider $(X,Y_1),(X,Y_2)\in\lK$; then $(0,Y_1-Y_2)\in\lK$ (because a Lie algebra is a vector space). Then the definition of $\lK$ makes for any $t\in\eR$, $(e,\exp{t(Y_1-Y_2)})\in K$. Consequently, $\exp t(Y_1-Y_2)=\varphi(e)=e$ and then $Y_1-Y_2=0$.

In order to show the existence, let us consider a $r>0$ such that $X_r=(1/r)X$ keeps in $M_0$. This exists because the sequence $X_r\to 0$ (then it comes $M_0$ from a certain $r$). From the definitions, $\exp$ is diffeomorphic between $M_0$ and $M_e$, then $\exp X_r\in M_e$ and $\varphi(\exp X_r)\in N_e$ because $\varphi(M_e)\subset N_e$.

From this, there exists an unique $Y_r\in N_0$ such that $\exp Y_r=\varphi(\exp X_r)$ and an unique $Z_r\in(M_0\times N_0)\cap\lK$ satisfying  $\exp Z_r=(\exp X_r,\exp Y_r)$. But $\exp$ is bijective from $M_0\times N_0$, so that $Z_r=(X_r,Y_r)$ and we can choose $Y=rY_r$ as a $Y\in\lH$ such that $(X,Y)\in\lK$ (it is not really a choice: the unicity was previously shown). We denotes by $\dpt{\psi}{\lG}{\lH}$ the map which gives the unique $Y\in\lH$ associated with $X\in\lG$ such that $(X,Y)\in\lK$. This is a homomorphism between $\lG$ and $\lH$.

By definition, $(X,\psi(X))\in\lK$, i.e. $(\exp tX,\exp t\psi(X))\in K$ or 
\begin{equation}
  \varphi(\exp tX)=\exp t\psi(X).
\end{equation}
Let us now consider a basis $\{X_1,\ldots,X_n\}$ of $\lG$. Since $\varphi$ is a homomorphism,
\begin{equation}\label{eq:coord_vp_exp}
   \varphi\big((\exp t_1X_1)(\exp t_2X_2)\ldots(\exp t_nX_n)\big)
     =\big(\exp t_1\psi(X_1)\big)\ldots\big( \exp t_n\psi(X_n) \big) 
\end{equation}
Now, we apply lemma \ref{lem:decomp} on the decomposition of $\lG$ into the $n$ subspace spanned by the $n$ vector basis (this is $n$ applications of the lemma), the map
\[
  (\exp t_1X_1)\ldots(\exp t_nX_n)\to (t_1,\ldots,t_n)
\]
is a coordinate system around $e$ in $G$. In this case, the relation \eqref{eq:coord_vp_exp} shows that $\varphi$ is differentiable at $e$. Then it is differentiable anywhere in $G$.
\end{proof}


\begin{proposition}
Let $G$ be a Lie group and $H$, a Lie subgroup of $G$ ($\lG$ and $\lH$ are the corresponding Lie algebras). We suppose that $H$ has at most a countable number of connected components. Then 
\begin{equation}
  \lH=\{ X\in\lG:\forall t\in\eR,e^{tX}\in H \}
\end{equation}
\end{proposition}

\begin{proof}
We will once again use the lemma  \ref{lem:decomp} with $\lN=\lH$ and $\lM$, a complementary vector space of $\lH$ in $\lG$. We define 
\[
   V=\exp\mU_m\exp\mU_h
\]
where $\mU_m$ and $\mU_h$ are the sets given by the lemma. We consider on $V$ the induced topology from $G$. If we define
\[
   \mA=\{A\in\mU_m:e^{A}\in H\},
\]
we have 
\begin{equation}\label{eq:union_A}
   H\cap V=\bigcup_{A\in\mA}e^{A}e^{\mU_h}.
\end{equation}
First, the definition of $V$ makes clear that the elements of the form $\exp A\exp\mU_h$ are in $V$. They are also in $H$ because $\exp A\in H$ (definition of $\mA$) and $\exp\mU_h$ still by definition. In order to see the inverse inclusion, let us consider a $h\in H\cap V$. We know that 
\begin{equation}\label{eq:AB_to_exp}
(A,B)\to\exp A\exp B 
\end{equation}
is a diffeomorphism between $\mU_m\times\mU_h$ and a neighbourhood of $e$ in $G$ which we called $V$. Thus any element of $V$ (\emph{a fortiori} in $V\cap H$) can be written as $\exp A\exp B$ with $A\in\mU_m$ and $B\in\mU_h$. Then $h=e^Ae^B$ for some $A\in\mU_m$, $B\in\mU_h$. Since $H$ is a group and $e^B\in H$, in order the product to belongs to $H$, $e^A$ must lies in $H$ : $A\in\mA$.

\begin{remark}\label{rem:union_disj}
Note that since \eqref{eq:AB_to_exp} is diffeomorphic, the union in right hand side of \eqref{eq:union_A} is disjoint. Each member of this union is a neighbourhood in $H$ because it is a set $h\exp\mU_h$ where $\exp\mU_h$ is a neighbourhood of $e$ in $H$.
\end{remark}

Now we consider the map $\dpt{\pi}{V}{\mU_m}$, 
\[
  \pi(e^{X}e^Y)=X
\]
if $X\in\mU_m$ and $Y\in\mU_h$. This is a continuous map which sends $H\cap V$ into $\mA$. The identity component of $H\cap V$ (in the sense of topology of $V$) is sent to a countable subset of $\mU_m$. Indeed by remark \ref{rem:union_disj}, identity component of $H\cap V$ is only one of the terms in the union \eqref{eq:union_A}, namely $A=0$. But we know that $\pi^{-1}(o)=\exp\mU_h$, thus $\exp\mU_h$ is the identity component of $H\cap V$ for the topology of $V$.

Let us consider a $X\in\lG$ such that $\exp tX\in H$ for any $t\in\eR$, and the map $\dpt{\varphi}{\eR}{G}$, $\varphi(t)=\exp tX$. This is continuous, then there exists a connected neighbourhood $\mU$ of $0$ in $\eR$ such that $\varphi(\mU)\subset V$. Then $\varphi(\mU)\subset H\cap V$ and the connectedness of $\varphi(\mU)$ makes $\varphi(\mU)\subset\exp\mU_h$. But $\exp\mU_h$ is an arbitrary small neighbourhood of $e$ in $H$; the conclusion is that $\varphi$ is a continuous map from $\eR$ into $H$. Indeed, we had chosen $X$ such that $\exp tX\in H$.

Moreover, we know that 
\[
  e^{(t_0+\epsilon)X}=e^{t_0X}e^{\epsilon X},
\]
but $\exp \epsilon X$ can be as close to $e$ as we want (this proves the continuity at $t_0$). Then $\varphi$ is a path in $H$. 

In definitive, we had shown that $\exp tX\in H$ implies that $t\to\exp tX$ is a path. Now equation \eqref{eq:path_alg} gives the thesis.

\end{proof}

\begin{corollary}
Let $G$ be a Lie group and $H_1$, $H_2$, two subgroups both having a finite number of connected components (each for his own topology). If $H_1=H_2$ as sets, then $H_1=H_2$ as Lie groups.
\end{corollary}

\begin{proof}
The proposition shows that $H_1$ and $H_2$ have same Lie algebra. But any Lie subalgebra of $\lG$ is the Lie algebra of exactly one connected subgroup of $G$ (theorem \ref{tho:gp_alg}). Then as Lie groups, ${H_1}_0={H_2}_0$. Since $H_1$ and $H_2$ are topological groups, the equality of they topology on one connected component gives the equality everywhere (because translations are differentiable).
\end{proof}

\begin{definition}
A \defe{differentiable subgroup}{differentiable!subgroup} is a connected Lie subgroup.
\end{definition}

\begin{corollary}
Let $G$ be a Lie group, and $K$, $H$ two differentiable subgroups of $G$. We suppose $K\subset H$. Then $K$ is a differentiable subgroup of the Lie group $H$.
\end{corollary}

\begin{proof}
The Lie algebras of $K$ and $H$ are respectively denoted by $\lK$ and $\lH$. We denote by $K^*$ the differentiable subgroup of $H$ which has $\lK$ as Lie algebra. The differentiable subgroups $K$ and $K^*$ have same Lie algebra, and then coincide as Lie groups.
\end{proof}

\label{pg:ex_topo_Lie}
Consider the group $T=S^1\times S^1$ and the continuous map $\dpt{\gamma}{\eR}{T}$ given by
\[
  \gamma(t)=(e^{it},e^{i\alpha t})
\]
with a certain irrational $\alpha$ in such a manner that $\gamma$ is injective and $\Gamma=\gamma(\eR)$ is dense in $T$.

The subset $\Gamma$ is not closed because his complementary in $T$ is not open: any neighbourhood of element $p\in T$ which don't lie in $\Gamma$ contains some elements of $\Gamma$. We will show that the inclusion map $\dpt{\iota}{\Gamma}{T}$ is continuous. An open subset of $T$ is somethings like
\[
  \mO=(e^{iU},e^{iV})
\]
where $U,V$ are open subsets of $\eR$. It is clear that 
\[
   \iota^{-1}(\mO)=\{ \gamma(t)\tq t\in U+2k\pi,\alpha t\in V+2m\pi \},
\]
but the set of elements $t$ of $\eR$ which satisfies it is clearly open. Then $\Gamma$ has at least the induced topology from $T$ (as shown in proposition \ref{prop:topo_sub_manif}). In fact, the own topology of $\Gamma$ is \emph{more} than the induced: the open subsets of $\Gamma$ whose are just some small segments clearly doesn't appear in the induced topology. Thus the present case is an example (and not a counter-example) of the next theorem at page \pageref{tho:H_ferme}.

This example show the importance of the condition for a topological subspace to have \emph{exactly} the induced topology. If not, any Lie subgroup were a topological Lie subgroup because a submanifold has at least the induced topology. We will go further with this example after the proof.


\begin{theorem} 
Let $G$ be a Lie group and $H$, a Lie subgroup of $G$.
\begin{enumerate}
\item If $H$ is a topological Lie subgroup of $G$, then it is closed in $G$,
\item If $H$ has at most a countable number of connected components and is closed in $G$, then $H$ is a topological subgroup of $G$.
\end{enumerate}
\label{tho:H_ferme}
\end{theorem}

\begin{proof}
\subdem{First point} It is sufficient to prove that if a sequence $h_n\in H$ converges (in $G$) to $g\in G$, then $g\in H$ (this is almost the definition of a closed subset). We consider $V$, a neighbourhood of $0$ in $\lG$ such that 

\begin{itemize}
\item $\exp$ is diffeomorphic between $V$ and an open neighbourhood  of $e$ in $G$,
\item $\exp(V\cap \lH)=(\exp V)\cap H$.
\end{itemize}
This exists by the lemma \ref{lem:sugroup_normal}; we can suppose that $V$ is bounded. Consider $\mU$, an open neighbourhood of $0$ in $\lG$ contained in $V$ such that $\exp-\mU\exp\mU\subset\exp V$.

Since $h_n\to g$, there exists a $N\in\eN$ such that $n\geq N$ implies $h_n\in g\exp\mU$ (i.e $h_n$ is the product of $g$ by an element rather close to $e$; since the multiplication is differentiable, the notion of ``not so far''\ is good to express the convergence notion). From now we only consider such elements in the sequence. So, $h_N^{-1} h_n\in(\exp V)\cap H$ ($n\geq N$) because
\[
   h_N^{-1} h_n\in\exp-\mU g^{-1} g\exp\mU\subset\exp V.
\]
(note that $H$ is a group, then $h_i^{-1}\in H$) From the second point of the definition of $V$, there exists a $X_n\in V\cap\lH$ such that $h^{-1}_N h_n=\exp X_n$ for any $n\geq N$.

Since $V$ in bounded, there exists a subsequence out of $(X_i)$ (which is also called $X_i$) converging to a certain $Z\in\lG$. But $\lH$ is closed in $\lG$ because it is a vector subspace (we are in a finite dimensional case), then $Z\in\lH$ and thus the sequence $(h_i)$ converges to $h_N\exp Z$; therefore $g\in H$.

\subdem{Second point} The subgroup $H$ is closed in $G$ and has a countable number of connected component. Since $H$ is closed, theorem \ref{tho:diff_sur_ferme} it has an analytics structure for which it is a topological Lie subgroup of $G$. We denotes by $H'$ this Lie group.

The identity map $\dpt{I}{H}{H'}$ is continuous\quext{pourquoi ?} (see error \ref{err:gross}). Thus any connected component of $H$ is contained in a connected component of $H'$, the it has only a countable number of connected components. By corollary \ref{cor:top_subgroup}, $H=H'$ as Lie group.

\end{proof}

Now we take back our example with $G=S^1\times S^1$, $H=\gamma(\eR)$. In this case, the theorem doesn't works. Let us see why as deep as possible. We have $\lG=\eR\oplus\eR=\eR^2$ and $\lH=\eR$, a one-dimensional vector subspace of $\lG$. ($\lH$ is a ``direction ''\ in $\lG$) First, we build the neighbourhood $V$ of $0$ in $\lG$. It is standard to require that $\exp$ is diffeomorphic between $V$ and an open around $(1,1)\in S^1\times S^1$. It also must satisfy $e^{V\cap\lH}=e^V\cap H$. This second requirement is impossible.

Intuitively. We can see $V\subset\lG$ as a little disk tangent to  the torus. The exponential map deposits it on the torus, as well that $e^V$ covers a little area on $G$. Then $e^V\cap H$ is one of these amazing open subset of $\Gamma$ which are dense in a certain domain of $G$.

On the other hand, $V\cap\lH$ is just a little vector in $\lH$; the exponential deposits it on a small line in $G$. This is not the same at all. Then lemma \ref{lem:sugroup_normal} fails in our case. Let us review the proof of this lemma until we find a problem. 

Let $W_0\subset\lG$  be a neighbourhood of $0$ which is in bijection with an open around $e$ in $G$. We consider $N_0$, an open subset of $H$ such that $N_0\subset W_0$ and $N_0$ is in bijection with $N_e$, a neighbourhood of $e$ in $G$. Until here, no problems. But now the proof says that there exists an open $U_e$ in $G$ such that $N_e=U_e\cap H$. This is false in our case. Indeed, $N_e=e^{N_0}$ is just a segment in $G$ while any subset of $G$ of the form $U_e\cap H$ is an ``amazing''\ open.

So we see that deeply, the obstruction for a Lie subgroup to be a topological Lie subgroup resides in the fact that the topology of a submanifold is \emph{more} than the induced topology, so that we can't automatically find the open $U_e$ in $G$.


\chapter{Lie algebras}
\input{helgaLie}

\chapter{Fibre bundle}
\section{Vector bundle}
%++++++++++++++++++++++

Let $M$ be a smooth manifold. A \defe{$V$-vector bundle}{vector!bundle}\index{bundle!vector} of rank $r$ on $M$ is a smooth manifold $F$ and a smooth projection $\dpt{p}{F}{M}$ such that

\begin{itemize}
\item for any $x\in M$, the fiber $F_x:=p^{-1}(x)$ is a vector space of dimension $r$ on the same field that $V$ (let's say $\eK=\eR$ or $\eC$).
\item for any $x\in M$, there exists an open neighbourhood $\mU$ of $x$ and a ``chart diffeomorphism``{} $\dpt{\phi}{p^{-1}(\mU)}{\mU\times V}$ such that for any $l\in p^{-1}(y)$,
   \begin{itemize}
      \item $\phi(l)=(y,\phi_y(l))$
      \item $\dpt{\phi_y}{E_y}{V}$ is a vector space isomorphism.
   \end{itemize}
\end{itemize}

The pair $(\mU,\phi)$ is a \emph{local trivialization}; $M$ is the \emph{base space}; $F$, the \emph{total space}, $p$ the \emph{projection} and $r$, the \emph{rank} of the bundle. The denominations of total and base spaces will also be used in the same way for principal bundles.

We will sometimes use charts diffeomorphism $\dpt{\phi}{\mU\times V}{p^{-1}(\mU)}$ instead of $\dpt{\phi}{p^{-1}(\mU)}{\mU\times V}$. Since they are diffeomorphism, this difference don't affect anything.

\subsection{Transition functions}
%--------------------------------

The trivializations will be denoted by Greek indices: $\mU_{\alpha}$, $\phi_{\alpha}$,\ldots The symbol $\mU_{\alpha}{}_{\beta}$ naturally denotes $\mU_{\alpha}\cap\mU_{\beta}$. If we consider two local trivializations $(\mU_{\alpha},\phi_{\alpha})$ and $(\mU_{\beta},\phi_{\beta})$, we have to look at $\dpt{\phi_{\alpha}\circ\phi_{\beta}^{-1}}{\mU_{\alpha}{}_{\beta}\times\eK^r}{\mU_{\alpha}{}_{\beta}\times\eK^r}$. We define the \defe{transition functions}{transition function} $\dpt{g_{\alpha}{}_{\beta}}{\mU_{\alpha}{}_{\beta}}{GL(r,\eK)}$ by
\begin{equation}
\phi_{\alpha}\circ\phi_{\beta}^{-1}(x,v)=(x,g\bab(x)v). 
\end{equation}
These functions take their values in $GL(r,\eK)$ because $\dpt{\phi_y}{E_y}{V}$ is a vector space isomorphism. Since $(\phi_{\alpha}\circ\phi_{\beta})^{-1}=\phi_{\beta}\circ\phi_{\alpha}^{-1}$, it is clear that $g\bab(x)=g_{\alpha\beta}(x)^{-1}$.

If $x\in\mU_{\alpha\beta\gamma}=\mU_{\alpha}\cap\mU_{\beta}\cap\mU\bgamma$, we have $\phi_{\alpha}\circ\phi\bgamma^{-1}(x,v)=(x,g_{\alpha\gamma}(x)v)$, but also $\phi_{\alpha}\circ\phi\bgamma^{-1}=\phi_{\alpha}\circ\phi_{\beta}^{-1}\phi_{\beta}\circ\phi\bgamma^{-1}$, then
\begin{equation}
  (x,g_{\alpha\gamma}(x)v)=(\phi_{\alpha}\circ\phi_{\beta}^{-1})(x,g_{\beta\gamma}(x)v)
               =(x,g\bab(x)_{\beta\gamma}(x)v).
\end{equation}
Thus $g_{\alpha\gamma}(x)=g\bab(x)g_{\beta\gamma}(x)$. So, as linear maps, we have
\begin{equation}\label{eq:g_compat}
  g\bab\circ g_{\alpha\gamma}\circ g_{\gamma\alpha}=\mtu.
\end{equation}

\subsection{Inverse construction}\label{subsec:inv_g}
%-------------------------------------------

Let us consider a manifold $M$, an open covering $\{\mU_{\alpha}:\alpha\in I\}$ and some functions $\dpt{g\bab}{\mU\bab}{GL(r,\eK)}$ which fulfill relations \eqref{eq:g_compat}. We will build a vector bundle $E\stackrel{p}{\longrightarrow}M$ whose transition functions are the $g_{\alpha\beta}$'s. Let $\tilde{E}$ be the disjoint union
\[
  \tilde{E}=\bigsqcup_{\alpha\in I}\mU_{\alpha}\times\eK^r,
\]
i.e. triples of the form $(x,v,\alpha)\in M\times\eK^r\times I$ with the condition that $x\in\mU_{\alpha}$. We define an equivalence relation on $\tilde{E}$ by $(x,v,\alpha)\sim(y,w,\beta)$ if and only if $x=y$ and $w=g\bab(x)v$. Next, we define $E=\tilde{E}/\sim$ and $\dpt{\omega}{\tilde{E}}{E}$, the canonical projection. The projection $\dpt{p}{E}{M}$ is naturally defined by $p([x,v,\alpha])=x$. The chart diffeomorphism is $\dpt{\varphi_{\alpha}}{\mU_{\alpha}\times\eK^r}{p^{-1}(\mU_{\alpha})}$, 
\[
  \varphi_{\alpha}(x,v)=\omega(x,v,\alpha).
\]
Now we have to prove that $E$ endowed with the $\varphi_{\alpha}$'s is a vector bundle.

First we prove that $\varphi_{\alpha}$ is surjective. For this we remark that a general element in $p^{-1}(\mU_{\alpha})$ can be written under the form $\omega(x,v,\alpha)$ with $x\in\mU\bab$. But
\begin{equation}
\begin{split}
  \varphi_{\alpha}(x,g\bab(x)w)&=\omega(x,g\bab(x)w,\alpha)\\
                         &=\omega(x,g_{\alpha\beta}(x)g\bab(x)w,\beta)\\
			 &=\omega(x,w\beta),
\end{split}
\end{equation}
then $\varphi_{\alpha}$ is surjective. Now we suppose $\varphi_{\alpha}(x,v)=\varphi_{\alpha}(y,w)$. Then $\omega(x,v,\alpha)=\omega(y,w,\alpha)$ and $x=y$, $w=g_{\alpha\alpha}v$ which immediately gives $v=w$. Then $\varphi_{\alpha}$ is injective.

Finally, we have 
\begin{equation}
  (\varphi\alpha\circ\varphi_{\beta}^{-1})(\omega(x,v,\alpha))=\varphi_{\alpha}(x,g_{\alpha\beta}(x)v)
                                                   =\omega(x,g_{\alpha\beta}(x)v,\alpha),
\end{equation}
which proves that the maps $g$ are the transition functions of the vector bundle $E$.

\subsection{Equivalence of vector bundle}
%----------------------------------------

Let $E\stackrel{p}{\longrightarrow}M$ and $F\stackrel{p'}{\longrightarrow}M$ be two vector bundles on $M$. They are \defe{equivalent}{equivalence!of vector bundle} if there exists a smooth diffeomorphism $\dpt{f}{E}{F}$ such that 

\begin{itemize}
\item $p'\circ f=p$,
\item $\dpt{f|_{E_x}}{E_x}{F_x}$ is a vector space isomorphism.
\end{itemize}

Let $E$ and $F$ be two equivalent vector bundles, $\{\mU_{\alpha}\tq \alpha\in I\}$, an open covering which trivialize $E$ and $F$ in the same time and $\phi^E_{\alpha}$, $\phi^F_{\alpha}$ the corresponding trivializations. A map $\dpt{f}{E}{F}$ reads ``in the trivialization''\ as $\dpt{\phi^F_{\alpha}\circ f|_{p^{-1}(\mu_{\alpha})}\circ\phi^E{}^{-1}_{\alpha}}{\mU_{\alpha}\times\eK^r}{\mU_{\alpha}\times\eK^r}$ and defines a map $\dpt{\lambda_{\alpha}}{\mU_{\alpha}}{GL(r,\eK)}$ by
\begin{equation}
(\phi^F_{\alpha}\circ f|_{p^{-1}(\mu_{\alpha})}\circ\phi^E{}^{-1}_{\alpha})(x,v)=(x,\lambda_{\alpha}(x)v).
\end{equation}
If we denote by $g^E$ the transition functions for $E$ (and $g^F$ for $F$),
\[
 \phi^F_{\alpha}\circ\phi^F_{\beta}{}^{-1}= (\phi^F_{\alpha}\circ f\circ\phi_{\alpha}^E{}^{-1})\circ
                                        (\phi^E_{\alpha}\circ\phi^E_{\beta}{}^{-1})\circ
					(\phi^E_{\beta}\circ f^{-1}\circ\phi_{\beta}^E{}^{-1}),
\]
so that 
\begin{equation}\label{eq:g_l_g_l}
  g\bab^F(x)=\lambda_{\alpha}(x) g^E\bab(x)\lambda(x)^{-1}.
\end{equation}

Once again we have an inverse construction. We consider a vector bundle $E$ on $M$ with transition functions $g^E$ and some maps $\dpt{\lambda_{\alpha}}{\mU_{\alpha}}{GL(r,\eK)}$; then we define $g^F\bab(x)$ by equation \eqref{eq:g_l_g_l}. 

From subsection \ref{subsec:inv_g}, one can construct a vector bundle $F$ on $M$ whose transition functions are these $g^F$. With the trivializations $\phi^F$ of $F$, one can define $\dpt{f}{E}{F}$ by
\[
(\phi^F_{\alpha}\circ f\circ\phi^E_{\alpha}{}^{-1})(x,v)=(x,\lambda_{\alpha}(x)v).
\]

When a basis space $B$ is given, we denote by $\Vect(B)$ the set of isomorphism classes of vector bundles over $B$. In the complex case, we denote it by $\Vect_{\eC}(B)$.

\begin{proposition}
Any vector bundle over $\eR^n$ is trivial.
\end{proposition}

\begin{proof}
Let $\dpt{p}{F}{M}$ be a vector bundle on $M=\eR^n$ and $\{\mU_{\alpha}\}$ be covering of $\eR^n$ by local trivializations. Now consider a partition of unity\index{partition of unity} related to the covering $\mU_{\alpha}$: a set of functions $\dpt{f_{\alpha}}{M}{\eR}$ such that
\begin{itemize}
\item $f_{\alpha}>0$,
\item $\forall x\in M$, one can find a neighbourhood of $x$ in which only a \emph{finite} number of $f_{\alpha}$ is non zero,
\item $\forall x\in M$, $\sum_{\alpha} f_{\alpha}(x)=1$.
\item $f_{\alpha}=0$ outside of $\mU_{\alpha}$.
\end{itemize}
Using that partition of unity, we build the trivialization function $\dpt{f}{F}{\eR^n\times V}$ by $f(l)=(x,\sum_{\alpha} f_{\alpha}(x)\phi_{\alpha x}(l))$.
\end{proof}

The following two propositions have some importance in K-theory.
\begin{proposition}		\label{PropEoplusEprimetriv}
Let $\pi\colon E\to B$ be a complex vector bundle over a basis compact, Hausdorff, connected basis $B$. Then there exists a vector bundle $E'$ such that $E\oplus E'$ is trivial.
\end{proposition}

\begin{proposition}		\label{PropmapfEEsun}
Let $f\colon A\to B$ be a map between the topological spaces $A$ and $B$, and consider a vector bundle $\pi\colon E\to B$. Then there exists one and only one vector bundle $\pi'\colon E'\to A$ and a map $f'\colon E'\to E$ such that $f'|_{E'_x}\colon E'_x\to E_{f(x)}$ is an isomorphism. The vector bundle $E'$ is unique up to isomorphism.
\end{proposition}
Proofs can be found in \cite{VB_and_K}. Let us denote by $f^*(E)$ the function given by proposition \ref{PropmapfEEsun}. It satisfies the following properties
\begin{equation}		\label{EqPropfstarEVect}
\begin{split}
	(fg)^*(E)		&=g^*\big( f^*(E) \big)\\
	\id^*(E)		&=E\\
	f^*(E_1\oplus E_2)	&=f^*(E_1)\oplus f^*(E_2)\\
	f^*(E_1\otimes E_2)	&=f^*(E_1)\otimes f^*(E_2).
\end{split}
\end{equation}



\subsection{Sections of vector bundle}
%-------------------------------------

A \defe{section}{section!of vector bundle} of the vector bundle $p\colon E\to M$ is a smooth map $\dpt{s}{M}{E}$ such that $p\circ s=\id|_M$. The set of all the sections is denoted by $\Gamma^{\infty}(M)$ or simply $\Gamma(E)$.\nomenclature{$\Gamma(E)$}{Space of sections of the vector bundle $E$}

If $(\mU_{\alpha},\phi_{\alpha})$ is a local trivialization, one can describe the section $s$ by a function $\dpt{s_{\alpha}}{\mU_{\alpha}}{V}$ defined by $\phi_{\alpha}(s(x))=(x,s_{\alpha}(x))$, or equivalently by
\[
s(x)=\phi_{\alpha}^{-1}(x,s_{\alpha}(x)).
\]
As usual when we define such a local quantity, we have to ask ourself how are related $s_{\alpha}$ and $s_{\beta}$ on $\mU_{\alpha}\cap\mU_{\beta}$. The best is $s_{\alpha}=s_{\beta}$, but most of the time it is not. Here, we compute
\[
  \phi_{\beta}\circ\phi_{\alpha}^{-1}\circ\phi_{\alpha}(s(s))=(x,g_{\alpha\beta}(x)s_{\alpha}(x)),
\]
which is obviously also equal to $(x,s_{\beta}(x))$. Then
\begin{equation}\label{eq:tr_sec}
s_{\beta}(x)=g_{\alpha\beta}(x)s_{\alpha}(x)
\end{equation}
without summation.

\section{Vector valued differential forms}	\label{SecVectValFiffFor}
%+++++++++++++++++++++++++++++++++++++++++++

Let $E$ be a vector bundle over $M$. A \defe{$E$-valued $p$-form}{vector-valued differential form}\index{differential!form!vector-valued} is a section
\[ 
  e\in\Gamma\big( E\otimes\Wedge^pT^*M \big).
\]
We denote by $\Omega(M,E)=\Gamma\big( E\otimes\Wedge^pT^*M \big)$\nomenclature[D]{$\Omega(M,E)$}{the set of $E$-valued differential forms} the set of $E$-valued differential forms. An element of $\Omega^1(M,E)=\Gamma\big( E\otimes\Wedge T^*M\big)$ always reads  $\sum_is_i\otimes\omega_i$ for some sections $s_i$ and usual differential forms $\omega_i$.

A form of $\Omega^p(M,E)$ can be seen as a fiber morphism $\underbrace{TM\otimes\cdots\otimes TM}_{\text{$p$ times}}\to E$ by associating 
\[ 
  s\otimes\omega(X_1,\cdots,X_p)=s(x)\omega(X_1,\cdots,X_p)\in E_x
\]
to the element $(s\otimes \omega)\in\Omega^p(M,E)$. There exists a wedge product between vector-valued forms. If $e\in\Omega^p(M,E_1)$ and $f\in\Omega^q(M,E_2)$, then we define $e\wedge f\in\Omega^{p+q}(M,E_1\otimes E_2)$ by
\begin{equation}	\label{EqDefwedgevecor}
(e\wedge f)(v_1,\cdots,v_{p+q})=\frac{1}{ p!q! }\sum_{\pi\in S_{p+q}}(-1)^{\pi} e(v_{\pi(1)},\cdots v_{\pi(p)})\otimes f(v_{\pi(p+1)},\cdots,v_{\pi(p+q)})\in E_1\otimes E_2.
\end{equation}
where $(-1)^{\pi}$ stands for the sign of the permutation $\pi$. For example when $e$, $f\in \Omega^1(M,E)$, we have
\[ 
  (e\wedge f)(X,Y)=e(X)\otimes f(Y)-e(Y)\otimes f(X)\in E\otimes E.
\]

When $M$ is a differentiable manifold, the \defe{fundamental $1$-form}{fundamental!$1$-form} is the element $\theta\in\Omega(M,TM)$ such that
\[ 
  \iota(X)\theta=X
\]
for every $X\in \Gamma(TM)$. 

\section{Lie algebra valued differential forms}	\label{SecLiaAlgformval}
%+++++++++++++++++++++++++++++++++++++++++++++++++

An important particular case of vector valued forms is given by Lie algebra valued forms. That case appears for example in the connection theory over principal bundle\footnote{So in Maxwell and other gauge field theories.}. If $\omega$ and $\eta$ are elements of $\Omega^1(M,\mG)$ for some Lie algebra $\mG$, we define
\[ 
		(\omega\wedge\eta)(X,Y)=\omega(X)\otimes\eta(Y)-\omega(Y)\otimes\eta(X).
\]
Combining with the Lie bracket, we define\nomenclature[D]{$[\omega\wedge\eta]$}{Combination of the wedge and the bracket in the case of Lie algebra-valued forms}
\begin{equation}	\label{EqDefomegawedgebeta}
[\omega\wedge\eta](X,Y):=[\omega(X),\eta(Y)]-[\omega(Y),\eta(X)].
\end{equation}
Using the proposition \ref{prop:XY_YX}, we often implicitly transforms the tensor product into a product \eqref{eq:yGyGb} and put
\begin{equation}	\label{EqAbuswesgeomom}
  (\omega\wedge\omega)(X,Y)=[\omega(X),\omega(Y)].
\end{equation}
Let us point out the fact that that kind of formula only holds for a ``wedge square'', but not for a general product $\omega\wedge\eta$. Remark that for $\omega\in\Omega^1(M,\mG)$ and $\beta\in\Omega^2(M,\mG)$, a simple computation of definition \eqref{EqDefwedgevecor} yields
\begin{equation}	\label{EqomwedgebetaXYZ}
(\omega\wedge\beta)(X,Y,Z)=\omega(X)\otimes\beta(Y,Z)-\omega(Y)\otimes\beta(X,Z)+\omega(Z)\otimes\beta(X,Y),
\end{equation}
so that, using the same trick as for equation \eqref{EqAbuswesgeomom}, we find
\[ 
  (\omega\wedge\beta-\beta\wedge\omega)(X,Y,Z)=[\omega(X),\beta(Y,Z)]-[\omega(Y),\beta(X,Z)]+[\omega(Z),\beta(X,Y)].
\]
But that expression is exactly what we find by exchanging the tensor product by Lie bracket in expression \eqref{EqomwedgebetaXYZ}. So we define
\begin{equation}	\label{EqDefCrochwedgedeux}
[\omega\wedge\beta]=\omega\wedge\beta-\beta\wedge\omega
\end{equation}
when $\omega\in\Omega^1(M,\mG)$ and $\beta\in\Omega^2(M,\mG)$. The reader should remark that this is what one would expect from generalisation of definition \eqref{EqDefomegawedgebeta}. 
\section{Principal bundle}
%+++++++++++++++++++++++++

Let $M$ be a manifold and $G$, a Lie group whose unit is denoted by $e$. A $G$-\defe{principal bundle}{principal!bundle}\index{bundle!principal} on $M$ is a smooth manifold $P$, a smooth map $\dpt{\pi}{P}{M}$ and a right action of $G$ on $P$ denoted by $\xi\cdot g$ with $g\in G$ and $\xi\in P$ such that

\begin{itemize}
\item $\pi(\xi\cdot g)=\pi(\xi)$,
\item $\forall \xi\in\pi^{-1}(x)$, $\pi^{-1}(x)=\{\xi\cdot g\tq g\in G\}\simeq G$,
\item $\forall x\in M$, there exists a neighbourhood $\mU_{\alpha}$ of $x$ in $M$, a diffeomorphism $\dpt{\phi_{\alpha}}{\pi^{-1}(\mU_{\alpha})}{\mU_{\alpha}\times G}$ and a diffeomorphism $\dpt{\phi_{\alpha x}}{P}{G}$ such that

\begin{itemize}
\item $\phi_{\alpha}(\xi)=(x,\phi_{\alpha x}(\xi))$,
\item $\phi_{\alpha x}(\xi\cdot g)=\phi_{\alpha x}(\xi)\cdot g$.
\end{itemize}
\end{itemize}
The group $G$ is often called the \defe{structure group}{structure!group}. We suppose that the action is effective. We will sometimes use the notation $P(G,M)$ to precise that $P$ is a principal bundle over $M$ with structure group $G$.

\begin{lemma}\label{lem:phixh}
The map $\phi_{\alpha}^{-1}$ fulfills
\[
  \phi\alpha^{-1}(x,h)\cdot g=\phi_{\alpha}^{-1}(x,hg).
\]
\end{lemma}

\begin{proof}

From the definition of a principal bundle, any $\xi\in P$ can be written under the form $\xi=\phi_{\alpha}^{-1}(x,\phi_{\alpha x}(\xi))$ with $\phi_x$ satisfying $\phi_x(\xi\cdot h)=\phi_x(\xi)h$ for a certain function $\dpt{\phi_x}{P}{G}$.  We consider in particular $\xi=\phi_{\alpha}^{-1}(x,h)\cdot g$. Then $\xi\cdot g^{-1}=\phi_{\alpha}^{-1}(x,h)$. But $\xi\cdot g^{-1}=\phi_{\alpha}^{-1}(x,\phi_{\alpha x}(\xi)g^{-1})$, then $h=\phi{\alpha_x}(\xi)g^{-1}$ and $\phi_{\alpha x}(\xi)=hg$. So we have
\[
\xi=\phi_{\alpha}^{-1}(x,h)\cdot g=\phi_{\alpha}^{-1}(x,\phi_{\alpha x}(\xi))=\phi_{\alpha}^{-1}(x,hg).
\]

\end{proof}

Let 
\[
   R=\{ (x,y)\in P\times P\tq x=y\cdot g\,\textrm{ for a certain $g\in G$}  \}.
\]

\begin{proposition}
The function $\dpt{u}{R}{G}$ defined by the condition
\[
  p\cdot u(p,q)=q.
\]
is differentiable.
\end{proposition}

\begin{proof}
Let $\mU$ be an open subset of $M$ and $\dpt{\sigma}{\mU}{P}$, a section. We consider a differentiable map $\dpt{\rho}{\pi^{-1}(\mU)}{G}$ such that $\rho(\xi\cdot g)=\rho(\xi)\cdot g$ and $\rho(\sigma(x))=e$. Such a map is given by
\[
   \rho(\xi)=\phi_x(\sigma(x))^{-1}\phi_x(\xi)
\]
where $x=\pi(\xi)$. We naturally define $R_{\mU}=R\cap( \pi^{-1}(\mU)\times\pi^{-1}(\mU) )$ and we pick $(\xi,\eta)\in R_{\mU}$. Let $s\in G$ be the one such that $\xi\cdot s=\eta$, so that $\rho(\xi)\cdot s=\rho(\eta)$. Then the restriction of $u$ to $R_{\mU}$ is given by $u(\xi,\eta)=\rho(\xi)^{-1}\rho(\eta)$ which makes $u|_{\mU}$ differentiable. Since this reasoning can be made on every chart open $\mU$, $u$ is differentiable everywhere on $P$.
\end{proof}

The following is a corollary of Leibnitz rule.
\begin{corollary}  \label{cor_PrincLeib}
If $P$ is a $G$-principal bundle and $v$, $a$ are curves in $P$ and $G$ respectively, we can consider the curve $u(t)=v(t)a(t)$. We have:
\[
         \dsdd{u(t)}{t}{0}=\dsdd{v(t)a(0)}{t}{0}+\dsdd{v(0)a(t)}{t}{0}.
\]
\end{corollary}
 The proof is direct. This result is often written as
\begin{equation}
               \dot{u}_t=\dot{v}_ta_t+v_t\dot{a}_t.
\end{equation}
A main application is 
\begin{equation}\label{eq:rdotht}
  \Dsdd{ r\cdot h(t) }{t}{0}=\Dsdd{r\cdot e^{th'(0)}}{t}{0}.
\end{equation}

\subsection{Transition functions}
%--------------------------------


Let $(\mU_{\alpha},\phi_{\alpha})$ be a local trivialization of $P$. This induces transition functions $\dpt{g\bab}{\mU_{\alpha}\cap\mU_{\beta}}{G}$ defined by
\begin{equation}
	\begin{aligned}
		\phi_{\alpha}\circ\phi_{\beta}^{-1}\colon \mU_{\alpha}\cap\mU_{\beta}\times G&\to \mU_{\alpha}\cap\mU_{\beta}\times G \\
		(x,a)&\mapsto (x,g\bab(x)a).\label{eq:transi_princ} 
	\end{aligned}
\end{equation}
Clearly, $g_{\alpha\alpha}=e$ and $g\bab g_{\alpha\beta}=e$ on $\mU_{\alpha}\cap\mU_{\beta}$. Then the triviality
\[
  \phi_{\alpha}\circ\phi_{\beta}^{-1}\circ\phi_{\beta}\circ\phi\bgamma^{-1}\circ\phi\bgamma\circ\phi_{\alpha}^{-1}=\id
\]
implies the compatibility conditions
\begin{equation}
g\bab g_{\beta\gamma} g_{\gamma\alpha}=e
\end{equation}
on $\mU_{\alpha}\cap\mU_{\beta}\cap\mU\bgamma$.

There is an inverse construction. Let $\{\mU_{\alpha}\tq\alpha\in I\}$ be an open covering of $M$ and $\dpt{g\bab}{\mU_{\alpha}\cap\mU_{\beta}}{G}$ a family of functions such that $g_{\alpha\alpha}=e$, $g\bab g_{\alpha\beta}=e$ on $\mU_{\alpha}\cap\mU_{\beta}$ and $g\bab g_{\beta\gamma}g_{\gamma\alpha}=e$ on $\mU_{\alpha}\cap\mU_{\beta}\cap\mU\bgamma$. Then the following construction gives a $G$-principal bundle whose transition functions are the $g\bab$'s.

\begin{itemize}
\item $\tilde{P}=\bigsqcup_{\alpha\in I}\mU_{\alpha}\times G$  (disjoint union),
\item if $(x,a)\in\mU_{\alpha}\times G$ and $(y,b)\in\mU_{\beta}\times G$, then $(x,a)\sim(y,b)$ if and only if $x=y$ and $b=g\bab(x)a$, 
\item $\dpt{\pi}{\tilde{P}}{M}$ is defined by $\pi[(x,a)]=x$ where $[(x,a)]$ is the class of $(x,a)$ for~$\sim$,
\item the action is defined by $[(x,a)]\cdot g=[(x,ag)]$.
\end{itemize}

\begin{theorem}
Let $G$ be a Lie group; $M$, a differentiable manifold; $\{\mU_{\alpha}\}_{\alpha\in I}$, an open covering of $M$ and some functions $\dpt{\varphi\bab}{\mU_{\alpha}\cap\mU_{\beta}}{G}$ such that $\varphi\bab(x)=\varphi_{\alpha\gamma}(x)\varphi_{\gamma\beta}(x)$. Then there exists a principal bundle $P$ whose transition functions are the $\varphi_{\alpha}$'s for the covering $\{\mU_{\alpha}\}_{\alpha\in I}$.
\end{theorem}

\begin{proof}
We consider the topological space 
\begin{equation}
    E=\bigcup_{\alpha\in I}(G\times\mU_{\alpha}\times I)
\end{equation}
where we put the discrete topology on $I$. Each $G\times\mU_{\alpha}\times\{\alpha\}$ is a manifold. Thus $E$ has a structure of differentiable manifold induced from the one of $G\times M$. We consider on $E$ the equivalence relation given by the following subset of $E\times E$:
\[
   R=\left\{\big(  (g,x,\alpha),(h,y,\beta)\big)\in E\times E\tq y=x\text{ and }  h=\varphi_{\alpha\beta}(x)g \right\}.
\]
We will show that $P=E/R$ has a structure of principal bundle. We begin by defining an action of $G$ on $P$ by
\[
  [ (g,x,\alpha)\cdot h ]=[ (gh,x,\alpha) ].
\]
In order to see that this definition is correct, let us consider $[g',x,\beta]=[g,x,\alpha]$. From the definition of the equivalence class, $g'=\varphi_{\alpha\beta}(x)g$. Then $[(g',x,\beta)]\cdot h=[(\varphi_{\alpha\beta}(g)gh,x,\beta)]$, and the form of $R$ shows that this is well $[(gh,x,\alpha)]$. Since the map $(g,h)\to gh$ is differentiable on $G$, the so defined action is a differentiable action of $G$ on $P$ and $G$ is a transformation group on $P$\quext{Faut voir comment ça correspond à la définition de l'autre texte.}.

If $[(g,x,\alpha)]=[(gh,x,\alpha)]$, then $gh=\varphi_{\alpha\alpha}g=g$ and $h=e$. So the action is effective. 

Now we consider the quotient $P/G$. A typical element is
\[
   \overline{ (s,x,i) }=\{ [s,x,i]\cdot g\tq g\in G \}.
\]
The projection $\dpt{\pi}{P}{M}$, $[(s,x,\alpha)]\to x$ is well defined and we can consider $\dpt{\varphi}{P/G}{M}$, $\varphi\overline{(s,x,\alpha)}=x$. It provides a bijection between $P/G$ and $M$. So we can identify $P/G$ and $M$. Now we are going to show that $P$ endowed with the projection $\dpt{\pi}{P}{X}$ is a principal bundle. 

We consider the map 
		\begin{equation}
		\begin{aligned}
			h_{\alpha} \colon G\times\mU_{\alpha} &\to P\
			(g,x)&\mapsto \omega(g,x,\alpha)
		\end{aligned}
	\end{equation}	
%
where $\dpt{\omega}{E}{P=E/R}$ is the canonical projection. Since 
\[
  (\pi\circ h_{\alpha})(g,x)=(\pi\circ\omega)(g,x,\alpha)=\pi[(g,x,\alpha)]=x,
\]
the map $h_{\alpha}$ actually is $\dpt{h_{\alpha}}{G\times\mU_{\alpha}}{\pi^{-1}(\mU_{\alpha})}$. In order to see that $h_{\alpha}$ is surjective on $\pi^{-1}(\mU_{\alpha})$, let us take a general element of $\pi^{-1}(\mU_{\alpha})$ under the form $\omega(g,x,\beta)$ with $x\in\mU_{\alpha}\cap\mU_{\beta}$. Then $(g,x,\beta)\in[ (\varphi\bab(x)g,x,\alpha) ]$ and therefore $\omega(g,x,\beta)=h_{\alpha}(\varphi\bab(x)g,x)$. For the injectivity, remark that $\omega(g,x,\beta)=\omega(h,y,\alpha)$ implies $x=y$ and $h=\varphi_{\beta\beta}(x)g=g$. In particular, $h_{\alpha}(g,x)=h_{\alpha}(h,y)$ implies $x=y$ and $g=h$.

Now we will prove that the inverse of $h_{\alpha}$ is continuous. For this we consider an open set $\Omega\subset G\times\mU_{\alpha}$ and we have to show that $h_{\alpha}(\Omega)$ is open in $\pi^{-1}(\mU_{\alpha})$. 

We recall the \defe{quotient topology}{topology!quotient}: if $A$ is a topological space with an equivalence relation $\sim$ and the canonical projection $\dpt{\varphi}{A}{A/\sim}$, then $V\subset A/\sim$ is open if and only if $\varphi^{-1}(V)\subset A$ is open. So in our case, we have to check the openness of $V=\omega^{-1}( h_{\alpha}(\Omega) )$ in $E$. We consider the open covering 
\[
  \{ G\times\mU_{\alpha}\times\{\alpha\} \}_{\alpha\in I}
\]
of $E$ and we will show that the intersection of $V$ with any of these open set is open. We have to show that 
$\omega^{-1}\big(  h_{\alpha}(\Omega)\cap (G\times\mU_{\alpha}\times\{\beta\})  \big)$ is open for any $\beta\in I$. For this, we define a map $\dpt{\alpha}{G\times(\mU_{\alpha}\cap\mU_{\beta})\times\{\beta\}}{G\times\mU_{\alpha}}$ by
\begin{equation}
  \alpha_{\beta}(g,x,\beta)=(\varphi\bab(x)g,x)
\end{equation}
which is continuous. The set $(h_{\alpha}\circ\alpha_{\beta})^{-1}(h_{\alpha}(\Omega))=\alpha^{-1}_{\beta}(\Omega)$ is open because $h_{\alpha}\circ\alpha_{\beta}$ is the restriction of $\omega$ to $G\times (\mU_{\alpha}\cap\mU_{\beta})\times\{\beta\}$. Then $h_{\alpha}$ is an homeomorphism from $G\times\mU_{\alpha}$ tp $\pi^{-1}(\mU_{\alpha})$. Since it is build from differentiable functions, it is moreover a diffeomorphism.

So we have a chart system $\{ (h_{\alpha},\mU_{\alpha}) \}_{\alpha\in I}$ where $h_{\alpha}$ fulfils the ``good'' properties with respect to $\pi$. It remains to be proved that the $\varphi\bab$'s are the transition functions and that $\pi^{-1}(\pi(\xi))=\xi\cdot G$ for every $\xi\in P$. We begin by the latter. For $\xi=[(g,x,\alpha)]$, $\pi(\xi)=x$ and we have to study the set
\[
  \pi^{-1}(x)=\{ [(h,x,\beta)]\tq h\in G,\beta\in I \}.
\]
Clearly, $[(h,x,\beta)]\cdot G\subset\pi^{-1}(x)$. The fact that there is nothing else than $[(h,x,\beta)]\cdot G$ in $\pi^{-1}(x)$ is seen by
\[
  [h,x,\beta]=[\varphi\bab(x)g,x,\alpha]\in[(h,x,\alpha)]\cdot G.
\]

In order to check the change of charts, let us consider $g'=h_{\beta,x}^{-1}\circ h_{\alpha,x}(g)$ where
\begin{equation}
  h_{\alpha,x}(g)=h_{\alpha}(g,x)=\omega(g,x,\alpha).
\end{equation}
The fact that $h_{\beta}(g',x)=g_{\alpha}(g,x)$ concludes the proof. To see this fact, remark that $h_{\beta,x}(h^{-1}_{\beta,x}\circ h_{\alpha,x}(g))=h_{\alpha,x}(g)$, so that $h_{\alpha}(g',x)=h_{\alpha}(g,x)$ implies $\omega(g',x,\beta)=\omega(g,x,\alpha)$ which proves that $g'=\varphi\bab(g)$.
\end{proof}

The \defe{trivial bundle}{trivial!principal bundle} is simply $P=M\times G$ and $\pi(x,g)=x$ with the action $(x,a)\cdot g=(x,ag)$.

\subsection{Morphisms and such\texorpdfstring{\ldots}{...}}
%----------------------------------

An \defe{homomorphism}{homomorphism!of principal bundle} between $P(G,M)$ and $P'(G',M')$ is a differentiable map $\dpt{h}{P}{P'}$ such that $\forall \xi\in P,g\in G$, 
\begin{equation}\label{eq:def_princ_homo}
   h(\xi\cdot g)=h(\xi)\cdot h_G(g)
\end{equation}
where $\dpt{h_G}{G}{G'}$ is a Lie group homomorphism. From the definition, $h$ maps a fiber to only one fiber, but it is not specially surjective on any fiber. So $h$ induces a homomorphism $\dpt{h_M}{M}{M'}$ such that $\pi'\circ h=h_M\circ\pi$.

An \defe{isomorphism}{isomorphism!of principal bundle} is a homomorphism $\dpt{g}{P(G,M)}{P'(G',M')}$ such that 

\begin{itemize}
\item $h_P$ is a diffeomorphism $P\to P'$,
\item $h_G$ is a Lie group homomorphism $G\to G'$, and
\item $h_M$ is a diffeomorphism $M\to M'$.
\end{itemize}

A principal bundle is \defe{trivial}{trivial!principal bundle} if one can find an isomorphism $\dpt{h}{G\times M}{P}$ such that $\pi\circ h=\id\circ\pr_2$, i.e. the following diagram commutes:
\begin{equation}\label{diag:princ_triv}
\xymatrix{ G\times M \ar[d]^{\pr_2} \ar[r]^h & P\ar[d]^{\pi} \\M \ar[r] _{\id}&M}
\end{equation}
We say that $P$ is \defe{locally trivial}{locally!trivial!principal bundle} if for every $x\in M$, there exists an open neighbourhood $\mU$ in $M$ such that $\pi^{-1}(\mU)$ endowed with the induced structure of principal bundle is trivial.

\subsection{Frame bundle: first}\label{pg:frame_bundle}
%--------------------------------

In the ideas, the building of a vector bundle is just to put a vector space on each point of the base manifold. A principal bundle is to put something on which a group acts on each point. If you have a vector bundle on a manifold, you can consider, on each point $x\in M$, the set of all the basis of the fiber $E_x$ over $x$. The group $GL(r,\eK)$ naturally acts on this set which becomes a candidate to be a $GL(r,\eK)$-principal bundle.

More formally, we consider a vector bundle $\dptvb{F}{p}{M}$, and for each $x$, the set of the basis of the vector space $F_x=p^{-1}(x)$. We define
\[
  P=\bigcup_{x\in M}(\textrm{basis of $F_x$}).
\]
We naturally consider the projection $\dpt{\pi}{P}{M}$, $\pi(b_x)=x$ if $b_x$ is a basis of $F_x$.

Let $\dpt{\phi^F_{\alpha}}{p^{-1}(\mU_{\alpha})}{\mU_{\alpha}\times\eK^r}$ be a local trivialization of $F$, and $\{\overline{e}_1,\ldots,\overline{e}_r\}$, the canonical basis of $\eK^r$. We naturally define
\[
  \ovS_{\alpha i}(x)=\phi^F_{\alpha}{}^{-1}(x,\overline{e}_i).
\]
The set $\{\ovS_{\alpha 1}(x),\ldots,\ovS_{\alpha r}(x)\}$ is a ``reference''{} basis of $F_x$ with respect to the trivialization $\phi_{\alpha}$. If we choose another basis $\{\ovv_1,\ldots\ovv_r\}$ of $F_x$, we can find a matrix $A\in GL(r,\eK)$ such that $\ovv_k=A^l_k\ovS_{\alpha l}(x)$. This gives a bijection
\begin{equation}
	\begin{aligned}
		\phi_{\alpha}^P\colon \pi^{-1}(\mU_{\alpha})&\to \mU_{\alpha}\times GL(r,\eK) \\
		(\ovv_1,\ldots,\ovv_r)&\mapsto (x,A). 
	\end{aligned}
\end{equation}
One can give to $P$ a $GL(r,\eK)$-principal bundle structure such that the $\phi_{\alpha}^P$ are diffeomorphism.

Let $(\mU_{\alpha},\phi_{\alpha}^F)$ be a local trivialization of  $F$ and $\dpt{g\bab^F}{\mU_{\alpha} \cap\mU_{\beta}}{GL(r,\eK)}$. In this case, $(\mU_{\alpha},\phi_{\alpha}^P)$ is a trivialization of $P$ whose transition function is $g\bab^P=g\bab^F$. Indeed
\[
  \phi_{\alpha}^P\circ\phi_{\beta}^P{}^{-1}(x,A)=\phi_{\alpha}^P(\{\ovv_1,\ldots,\ovv_r\})
\]
where $\ovv_s=(\phi^F_{\beta})^{-1}(x,A^l_s\overline{e}_l)$. In order to see it, recall that $\ovv_s=A^l_s\ovS_{\alpha l}(x)$ and that $\phi_{\alpha}^F{}^{-1}(x,\overline{e}_s)=\ovS_{\alpha s}(x)$. Then
\[
   \ovv_s=(\phi_{\beta}^F)^{-1}(x,A^l_s\overline{e}_l)=A^l_s\ovS_{\alpha s}(x).
\]
On the other hand, from the definition of $\phi_{\beta}^P$, the basis $(\phi_{\beta}^P)^{-1}(x,A)$ is the one obtained by applying $A$ on $S$. With all this,
\begin{equation}
\begin{split}
  \phi_{\alpha}^P\circ(\phi_{\beta}^P)^{-1}(x,A)&=\phi_{\alpha}^P\{  (\phi_{\beta}^F)^{-1}(x,A^l_s\overline{e}_l)\}_{s=1,\ldots r}\\
                 &=\phi_{\alpha}^P\{(\phi_{\alpha}^F)^{-1}\circ(\phi_{\alpha}^E\circ\phi_{\beta}^F{}^{-1})(x,A^l_s\overline{e}_l)  \}_{s=1,\ldots r}\\
		 &=\phi_{\alpha}^P\{  (\phi^E_{\alpha})^{-1}(x,g\bab^F(x)^s_iA^l_s\overline{e}_l   ) \}_{i=1,\ldots r}\\
		 &=(x,g^F\bab(x)A).
\end{split}
\end{equation}
The last product $g^F\bab(x)A$ is a matricial product.

\subsection{Frame bundle: second} \label{subsec_frbundle}
%------------------------------

\subsubsection{Basis}\label{subsubsecframebundle}
%////////////////////////////////////////////////

If $M$ is a $m$-dimensional manifold, a \defe{frame}{basis!of $T_xM$} of $T_xM$ is an isomorphism $\dpt{b}{\eR^m}{T_xM}$. In our purpose, we will always deal with (pseudo)Riemannian manifold. So, the tangents spaces $T_xM$ comes with a metric, and we ask a frame to be isometric. In other words, we ask $b$ to be an isometry from $(\eR^m,\cdot)$ to $(T_xM,g_x)$, where the dot denotes the (pseudo)euclidian product on $\eR^m$. Such a frame is given by a base point $x$ of $M$ and a matrix $S$ in $\SO(g_x)$:
\begin{equation}
                 \label{r1504e1}b(v)=(Sv)^i(\partial_i)_x,
\end{equation}
if the vector $v$ is written as $v=v^i\oui$ in the canonical orthogonal frame $\{\oui\}$ of $\eR^m$ and $\SO(g_x)$ is the set of the $m\times m$ matrix $A$ such that $A^tg_xA=g_x$.

This frame intuitively corresponds to the basis of $T_xM$ (see as a ``true''\ vector space) that we would have written by $\{Se_i\}_x$ if $e_i=\dsd{}{x^i}$. In order to follow this idea, we will effectively denote by $\baz{S}{x}$ the map $\dpt{b}{\eR^m}{T_xM}$ given by \eqref{r1504e1}.

We will often write the frame $b$ as $\baz{b}{x}$, making no differences in notation between the $b$ of $\SO(M)$ and the $b$ of $\SO(g_x)$ which implement it.

\begin{remark}
One has to distinguish a \emph{frame} and a \emph{basis}: a basis is only a free and generator set while a frame can be interpreted as an ordered basis.
\end{remark}


\subsubsection{Construction}
%////////////////////////////

We just saw how to build a frame bundle over a manifold. One can get another expression of the frame bundle when we express a basis of $T_xM$ by means of an isomorphism between $\eR^n$ and $T_xM$. If $M$ is a $n$-dimensional manifold, a \defe{frame}{frame} at $x$ is an ordered basis  
\[
   b=({\bf b}_1,\ldots,{\bf b}_n)
\]
of $T_xM$. It is clear that any frame defines an isomorphism (linear bijective map)
\begin{equation}
\begin{aligned}
	\tilde{b}\colon\eR^n&\to T_xM \\ 
e_i&\mapsto {\bf e_i} 
\end{aligned}
\end{equation}
where $\{e_i\}$ is the canonical basis of $\eR^n$. It is also clear that any isomorphism gives rise to a frame. Then we see a frame of $M$ at $x$ as an isomorphism $\dpt{\tilde{b}}{\eR^n}{T_xM}$. Let $B(M)_x$ be the of all the frames of $M$ at $x$; we define
\[
   B(M)=\bigcup_{x\in M}B(M)_x.
\]
For all $b\in B(M)_x$, we define $p_B(b)=x$ and the action $B(M)\times GL(n,\eR)\to B(M)$ by $b\cdot g=({\bf b}'_1,\ldots,{\bf b}'_n)$ where
\begin{equation}
  {\bf b}'_j={\bf b}_i\bghd{g}{i}{j}.
\end{equation}
It is easy to see that $\dpt{\widetilde{b\cdot g}=\tilde{b}\circ g}{\eR^n}{T_xM}$. So we can give to
\begin{equation}
\xymatrix{
    GL(n,\eR)  \ar@{~>}[r] & B(M) \ar[d]^{p_B}\\& M
  }
\end{equation}
a structure of principal bundle\footnote{Much more details and proofs are given in \cite{Naber}.}. If $(\mU_{\alpha},\varphi_{\alpha})$ is a local coordinate chart on $M$, we define
		\begin{equation}
		\begin{aligned}
			\tilde{\varphi} \colon p_B^{-1}(\mU_{\alpha}) &\to \varphi_{\alpha}(\mU_{\alpha})\times GL(n,\eR)\
			b&\mapsto (\varphi_{\alpha}(x),A(b))
		\end{aligned}
	\end{equation}	
where $A(b)\in GL(n,\eR)$ is defined by the condition ${\bf b}_j=\bghd{A}{j}{i}\partial_i|_x$. The matrix $A(b)$ is the one which transforms the canonical basis (in the trivialization $\varphi_{\alpha}$) into $b\in B(M)_x$. That's for the principal bundle structure. 

The manifold structure of $B(M)$ is given by $\dpt{\Phi_{\alpha}}{p_B^{-1}(\mU_{\alpha})}{\mU_{\alpha}\times GL(\eR)}$,
\begin{equation}
\begin{split}
  \Phi(b)&=(\varphi_{\alpha}^{-1}\times \id|_{GL(n,\eR)})\circ\tilde{\varphi}(b)\\
         &=(x,A(b))\\
         &=(p_B(b),A(b)).
\end{split}
\end{equation}
It fulfils $A(b\cdot g)=A(b)\cdot g$. A section $\dpt{s}{\mU_{\alpha}}{B(M)}$ is sometimes called a \defe{moving frame}{moving frame}\index{frame!moving} over $\mU_{\alpha}$.

Frame bundle over $\eR^2$ is given as example in page \pageref{Pg_exempleRdeux}
% 
% \subsection{Space-time}\label{subsec:space_time}
% %----------------------
% 
% 
% We say that the basis $\{\gb_0,\ldots,\gb_3\}$ of $T_xM$ is oriented, time oriented and orthonormal if $g(\gb_i,\gb_j)=\eta_{ij}$ (pointwise) and if $\gb_0$ is time-like and future directed. All this is devoted to build the frame bundle
% 
% 
% \[
% \xymatrix{
%     \Lpf  \ar@{~>}[r] & L(M) \ar[d]^{p_L} \\
%     &M.
%   }
% \]
% 
% From discussion in \autoref{subsec:sym_nature}, we know that, in physics, the relevant group is \emph{not} the Lorentz group $\Lpf$, but $\SLdc$ . So we want to build a $\SLdc$-principal bundle which ``fit''{} as deeply as possible the bundle $L(M)$. It is done with a \defe{spin structure}{spin!structure} which is a principal bundle
% 
% \[
% \xymatrix{
%     \SLdc  \ar@{~>}[r] & S(M) \ar[d]^{p_S} \\
%     &M.
%   }
% \]
% with a map $\dpt{\lambda}{S(M)}{L(M)}$ such that $p_L\circ\lambda=p_S$ and $\lambda(\xi\cdot g)=\lambda(\xi)\cdot\mSpin(g)$. See \autoref{sec:spin_str} and \autoref{sec:incl_Lorentz}.

\subsection{Sections of principal bundle}
%----------------------------------------

A \defe{section}{section!of principal bundle} of a $G$-principal bundle is a smooth map $\dpt{s}{M}{P}$ such that $s(x)\in\pi^{-1}(x)$ for any $x\in M$. A trivialization $\phi_{\alpha}^P$ $P$ on $\mU_{\alpha}$ defines a section of $P$ over $\mU_{\alpha}$ by
\[
   \sigma_{\alpha}(x)=(\phi_{\alpha}^P)^{-1}(x,e)
\]
where $e$ is the neutral of the group. In the inverse sense, we have the following:

\begin{proposition}
If $\dpt{\sigma_{\alpha}}{\mU_{\alpha}}{P}$ is local section of $P$ over $\mU_{\alpha}\subset M$, then the definition $\phi_{\alpha}^P(\xi)=(x,a)$ if $\xi=\sigma_{\alpha}(x)\cdot a$ is a local trivialization.
\end{proposition}

\begin{proof}
The function $\phi_{\alpha}^P$ is well defined because $\xi\in\pi^{-1}(\mU_{\alpha})$ implies the existence of a $x\in\mU_{\alpha}$ such that $\xi\in\pi^{-1}(x)=\{\xi\cdot g\}\simeq G$. For this $x$, there exists a $g\in G$ such that $\xi=\sigma_{\alpha}(x)\cdot g$. 

Now we prove that the couple $(x,a)$ is unique in the sense that $s_{\alpha}(x)\cdot a=\sigma_{\alpha}(y)\cdot b$ implies $(x,a)=(y,b)$. The left hand side belongs to $\pi^{-1}(x)$ while the right one belongs to $\pi^{-1}(y)$. Then $x=y$. The condition $\pi^{-1}(x)\simeq G$ imposes the unicity of the $g$ making $\xi=\eta\cdot g$ for each couple, $\xi,\eta\in\pi^{-1}(x)$.
\end{proof}

If $\sigma$ and $\sigma'$ are two sections of the same principal bundle $P$, then there exists a differentiable map $\dpt{f}{M}{G}$ such that $\sigma'(x)=\sigma(x)\cdot f(x)$. So all the sections can be deduced from only one and multiplication by such a $f$.

\begin{theorem}
If $\dpt{\pi}{P(G,M)}{M}$ is a principal bundle, then the four following propositions are equivalent:
\begin{enumerate}
\item\label{enuymai} $P$ is trivial,
\item\label{enuymaii} $P$ has a global section,
\item\label{enuymaiii} there exists a differentiable map $\dpt{\gamma}{P}{G}$ such that $\gamma(\xi\cdot g)=g^{-1}\gamma(\xi)$ for all $\xi\in P$ and $g\in G$,
\item\label{enuymaiv} there exists a differentiable map $\dpt{\rho}{P}{G}$ such that $\rho(\xi\cdot g)=\rho(\xi)g$.
\end{enumerate}
\label{ThoYPrincBTrivSect}
\end{theorem}

\begin{proof}
\subdem{\ref{enuymai}$\Rightarrow$ \ref{enuymaii}} 
The diagram \eqref{diag:princ_triv} commutes and 
		\begin{equation}
		\begin{aligned}
			\tau \colon M &\to G\times M\
			x&\mapsto (e,x)
		\end{aligned}
	\end{equation}	
is a local section of $G\times M$. From it we build the following global section of $P$:
		\begin{equation}
		\begin{aligned}
			\sigma \colon M &\to P\
			x&\mapsto h(e,x).
		\end{aligned}
	\end{equation}	
 This is injective because $\pi\circ h=\pr_2$ and differentiable because this is a composition of $\xdp{x}{(e,x)}$ and $\xdp{(g,x)}{h(g,x)}$.
\subdem{\ref{enuymaii}$\Rightarrow$ \ref{enuymai}}
The principal bundle $P$ admits a global section $\dpt{\sigma}{M}{P}$. From it, we can build the differentiable map
		\begin{equation}
		\begin{aligned}
			h \colon G\times M &\to P\
			(g,x)&\mapsto \sigma(x)\cdot g
		\end{aligned}
	\end{equation}	
which satisfies $h(gh,x)=h(g,x)\cdot h$ and $\pi\circ(g,x)=x$. First we show that $h$ is a fiber homomorphism and an isomorphism between $P$ and $G\times M$ so that $P$ is trivial. For this remark that
\[
  g(gh,x)=g(g,x)\cdot h=\sigma(x)\cdot gh,
\]
hence equation \eqref{eq:def_princ_homo} reduces to $h( (g,x)\cdot h )=h(g,x)\cdot h_G(h)$ which is true with $h_G=\id$. Moreover $\dpt{h}{G\times M}{P}$ is bijective because $\sigma(\pi(\xi))$ belongs to the fiber of $\xi\in P$, therefore there is one and only one $\gamma(\xi)=u(\xi,\sigma(\pi(\xi)))$ such that $\xi\cdot\gamma(\xi)=(\sigma\circ\pi)\xi$. The inverse map is 
		\begin{equation}
		\begin{aligned}
			\theta \colon P &\to G\times M\
			\xi&\mapsto (\gamma(\xi),\pi(\xi))
		\end{aligned}
	\end{equation}	
which is differentiable because $\gamma$ and $\pi$ are. So far we see that $h$ and $h^{-1}$ are differentiable. Then $h$ is an isomorphism between $P$ and $G\times M$.

\subdem{\ref{enuymaii}$\Rightarrow$ \ref{enuymaiii}}
Let $\sigma$ be the global section and define
		\begin{equation}
		\begin{aligned}
			\gamma \colon P &\to G\
			\xi&\mapsto u(\xi,(\sigma\circ\pi)\xi)
		\end{aligned}
	\end{equation}	
where $\dpt{u}{R}{G}$ is the map defined by the condition $\xi\cdot(\xi,\eta)=\eta$. The map $\gamma$ is differentiable and we have to prove that $\gamma(\xi\cdot g)=g^{-1}\gamma(\xi)$. Since $\xi\cdot \gamma(\xi)=\sigma\circ\pi(\xi)$,
\[
  \gamma(\xi\cdot g)=u(\xi\cdot g,(\sigma\circ\pi)(\xi\cdot g))=u(\xi\cdot g,(\sigma\circ\pi)(\xi)).
\]
But $(\xi\cdot g)(g^{-1}\gamma(\xi))=\xi\cdot\gamma(\xi)=x$. So $\gamma(\xi\cdot g)=u(\xi\cdot g,x)$. Thus $\gamma(\xi\cdot g)=g^{-1}\gamma(\xi)$.

\subdem{\ref{enuymaiii}$\Rightarrow$ \ref{enuymaii}}
The given map $\gamma$ fulfils $\xi\cdot g\gamma(\xi\cdot g)=\xi\cdot(\xi)$, so 
		\begin{equation}
		\begin{aligned}
			\varphi \colon P &\to P\
			\xi&\mapsto \xi\cdot(\xi)
		\end{aligned}
	\end{equation}	
is just function of the class of $\xi$, thus we have a section $\dpt{\sigma'}{P/G}{P}$, but we know that $P/G$ and $M$ are isomorphic.

\subdem{\ref{enuymaiii}$\Rightarrow$ \ref{enuymaiv}}
Let us define $\dpt{\rho}{P}{G}$ by $\rho=J\circ\gamma$ with $J(g)=g^{-1}$, thus $\rho(\xi)=\gamma(\xi)^{-1}$ and
\[
  \rho(\xi\cdot g)=\gamma(\xi\cdot g)^{-1}=(g^{-1}\gamma(\xi))^{-1}=\gamma(\xi)^{-1} g=\rho(\xi)g.
\]
\subdem{\ref{enuymaiv}$\Rightarrow$ \ref{enuymaiii}} The proof is just the same with $\rho=J\circ\rho$.
\end{proof}

\begin{definition}  
A section $\psi\in\Gamma(P,TP)$ is \defe{$G$-equivariant}{equivariant!vector field on principal bundle} when
\[ 
  d\tau_{g}\psi(\xi)=\psi(\xi\cdot g).
\]
\label{DefEqVectPrinc}
\end{definition}
Be careful: this \emph{does not} define equivariant sections of the principal bundle.
\subsection{Equivalence of principal bundle}
%-------------------------------------------

Two principal bundles $\dpt{\pi}{P}{M}$ and $\dpt{\pi'}{P'}{M}$ are \defe{equivalent}{equivalence!of principal bundle} if there exists a diffeomorphism $\dpt{\varphi}{P}{P'}$ such that 

\begin{itemize}
\item  $\pi'\circ\varphi=\pi$
\item $\varphi(\xi\cdot g)=\varphi(\xi)\cdot g$.
\end{itemize}

If $\{\mU_{\alpha}\}_{\alpha\in I}$ is an open covering of $M$ on which we have trivializations $\phi_{\alpha}$ of $P$ and $\psi_{\alpha}$ of $P'$, the diffeomorphism $\varphi$ induces some functions $\dpt{\lambda}{\mU_{\alpha}}{G}$ by setting
\[
   (\phi_{\alpha}\circ\varphi^{-1}\circ\psi_{\alpha}^{-1})(x,a)=(x,\lambda_{\alpha}(x)a).
\]
This definition works because from the definitions of principal bundle and equivalence, one sees that $(\phi_{\alpha}\circ\varphi^{-1}\circ\psi_{\alpha}^{-1})(x,\cdot)=(x,\cdot)$.

\subsubsection{Transition functions}
%///////////////////////////////////

We have some transition functions for $P$ and $P'$ given by equations
\[
 \begin{aligned}
   (\phi_{\alpha}\circ\phi_{\beta}^{-1})(x,g)&=(x,g\bab(x)g)\\
   (\psi_{\alpha}\circ\psi_{\beta}^{-1})(x,g)&=(x,g'\bab(x)g).
 \end{aligned}
\]
Now, we want to know what is $g'\bab$ in function of $g\bab$. First remark that $(\psi_{\alpha}\circ\varphi\circ\phi_{\alpha}^{-1})(x,a)=(x,\lambda_{\alpha}(x)^{-1})a$, and next, compute
\begin{equation}
\begin{split}
  (x,g\bab(x)a)a&=(\psi_{\alpha}\circ\varphi\circ\phi_{\beta}^{-1}\circ\phi_{\beta}\circ\varphi^{-1}\circ\psi_{\beta}^{-1})(x,a)\\
                &=(\psi_{\alpha}\circ\varphi\circ\phi_{\beta}^{-1})(x,\lambda_{\beta}(x)a)\\
		&=(\psi_{\alpha}\circ\varphi\circ\phi_{\alpha}^{-1}\circ\phi_{\alpha}\circ\phi_{\beta}^{-1})(x,\lambda_{\beta}(x)a)\\
		&=(x,\lambda_{\alpha}(x)^{-1} g\bab(x)\lambda_{\beta}(x)a).
\end{split}
\end{equation}
Then 
\begin{equation}
    g\bab(x)=\lambda_{\alpha}(x)^{-1} g\bab(x)\lambda_{\beta}.
\end{equation}
One can show that if two principal bundle have transition functions whose fulfill this condition, they are equivalent. A $G$-principal bundle is \defe{trivial}{trivial!principal bundle} if it is equivalent to the one given by $\dpt{\pi_1}{M\times G}{M}$.

\subsection{Reduction of the structural group}
%---------------------------------------------

We say that a principal bundle $P(G,M)$ is \defe{reducible}{reducible!principal bundle} when there exists a principal bundle $P'(H,M)$ such that 

\begin{itemize}
\item $H$ is a subgroup of $G$,
\item there exists an homeomorphism $\dpt{h}{P'}{P}$ such that $\dpt{h_G}{H}{G}$ is an injective homomorphism.
\end{itemize}

In this case we say that $G$ is reducible to $H$ and that $P'$ is a reduced principal bundle.

\begin{theorem}
If $P$ is a principal bundle over $M$, the structural group $G$ is reducible to the Lie subgroup $H$ if and only if there exists an open covering $\{ \mU_i \}_{i\in I}$ of $M$ and transition functions $\varphi_{ij}$ taking their values in $H$.
\end{theorem}
\begin{proof}
No proof.
\end{proof}
The following comes from \cite{Dieu4}. Let us consider the principal bundle
\begin{equation}
\xymatrix{%
   G \ar@{~>}[r]		&	P\ar[d]^{\pi_P}\\
   				&	M
 }
\end{equation}
and $H$, a closed subgroup of $G$. We denote by $j\colon H\to G$ the inclusion map. The principal bundle
\begin{equation}
\xymatrix{%
   H \ar@{~>}[r]		&	Q\ar[d]^{\pi_Q}\\
   				&	M
 }
\end{equation}
is a \defe{reduction}{reduction of a principal bundle} of $P$ to the group $H$ if there exists a map $u\colon Q\to P$ such that $\pi_P\circ u=\pi_Q$ and $u(\xi\cdot h)=u(\xi)\cdot j(h)$. In this case, $u$ is an embedding\quext{plongement} of $Q$ in $P$ and the image is a closed submanifold of $P$.   

Let $M$ be a $n$-dimensional manifold and $B(M)$ be its frame bundle. This is a $\GL(n,\eR)$-principal bundle. If $G$ is a closed subgroup\footnote{Typically $\SO(p,q)$ or $\SO_0(p,q)$.} of $\GL(n,\eR)$, a \defe{$G$-structure}{$G$-structure} is a reduction of $B(M)$ to $G$.

%---------------------------------------------------------------------------------------------------------------------------
					\subsection{Density}
%---------------------------------------------------------------------------------------------------------------------------

A \defe{density}{density} on a $d$-dimensional manifold $M$ is a section of the principal bundle whose fiber $P_x$ over $x\in M$ is the space of homogeneous non vanishing maps
\begin{equation}
\rho\colon \Wedge^dT_xM\to \eR^*_+
\end{equation}
such that $\rho(\lambda v)=| \lambda |\rho(v)$ for every $\lambda\in\eR$ and $v\in\Wedge^d T_xM$.


\section{Associated bundle}  \index{bundle!associated}
%++++++++++++++++++++++++++

Let $\dpt{\pi}{P}{M}$ be a $G$-principal bundle and $\dpt{\rho}{G}{GL(V)}$, a representation of $G$ on a vector space $V$ (on $\eK=\eR$ or $\eC$) of dimension $r$.

The associated bundle $\dptvb{E=P\times_{\rho} V}{p}{M}$ is defined as following. On $P\times V$, we consider the equivalence relation 
\[
   (\xi,v)\sim(\xi\cdot g,\rho(g^{-1})v)
\]
for $g\in G$, $\xi\in P$ and $v\in V$. Then we define

\begin{itemize}
\item $E=P\times_{\rho} V:=P\times V/\sim$,
\item $p[(\xi,v)]=\pi(\xi)$
\end{itemize}
where $[(\xi,v)]$ is the class of $(\xi,v)$ in $P\times V$.

If $\phi^P_{\alpha}(\xi)=(\pi(\xi),a(\xi))$ is a  trivialization of $P$ on $\mU_{\alpha}$, then
\begin{equation}\label{eq:triv_P_E}
\phi^E[(\xi,v)]=(\pi(\xi),\rho(a)v) 
\end{equation}
is a trivialization of $E$.

In order to see that it is a good definition, let us consider $(\eta,w)\sim(\xi,v)$. It immediately gives the existence of a $g\in G$ such that $\eta=\xi\cdot g$ and $w=\rho(g^{-1})v$. Then $\phi^E[ (\xi\cdot g,\rho(g^{-1})v) ]=( \pi(\xi\cdot g),\rho(b)\rho(g^{-1})v )$.  From the definition of $\phi^E$, the vector $b$ is given by $\phi^P(\xi\cdot g)=( \pi(\xi\cdot g),b )$, and the definition of a principal bundle gives $b=\phi_{\pi(\xi)}(\xi\cdot g)=\phi_{\pi(\xi)}(\xi)\cdot g=ag$. The fact that $\rho$ is a homomorphism makes $\rho(ag)\rho(g^{-1})=\rho(a)v$ and $\phi^E$ is well defined.

Let $G$ be a Lie group, $\rho$ a representation of $G$ on $V$ and $M$, a manifold. We consider $\dptvb{P=M\times G}{\pr_1}{M}$, the trivial $G$-principal bundle on $M$. Then $\dptvb{E=P\times_{\rho}V}{p}{M}$ is \defe{trivial}{trivial!principal bundle}, i.e. we can build a $\dpt{\varphi}{P\times_{\rho} V}{M\times V}$ such that $\pr_1\circ\varphi=p$. It is rather easy: we define
\[
  \varphi\big[ \big((x,g),v\big) \big]=(x,\rho(g)v).
\]
It is easy to see that $(\pr_1\circ\varphi)[ (x,g),v ]=x$ and $p[(x,g),v]=\pr_1(x,g)=x$.

\subsection{Transition functions}
%--------------------------------

\begin{proposition}
Let $(\mU_{\alpha},\phi_{\alpha}^P)$ be a trivialization of $\dptvb{P}{\pi}{M}$ whose transition functions are $\dpt{g\bab}{\mU_{\alpha}\cap\mU_{\beta}}{G}$. Then $(\mU_{\alpha},\phi^E_{\alpha})$ given by \eqref{eq:triv_P_E} is a local trivialization of $\dptvb{E}{p}{M}$ whose transition functions $\dpt{g\bab^E}{\mU_{\alpha}\cap\mU_{\beta}}{GL(\dim V,\eK)}$ are given by
\[
   g\bab^E(x)=\rho(g\bab^P(x)).
\]

\end{proposition}

\begin{proof}
If we write $a:=\phi^E_{_{\beta} x}(\pi^{-1}(x))$, we have $\phi_{\beta}^P(\pi^{-1}(x))=(x,a)$ and $\phi_{\alpha}^E\circ(\phi_{\beta}^E)^{-1}(x,v)=\phi^E_{\alpha}[ (\pi^{-1}(x),\rho(a)^{-1} v) ]$. So,
\begin{equation}
\begin{split}
\phi_{\alpha}^E[ (\pi^{-1}(x),\rho(a)^{-1} v) ]&=\Big(x, 
                            \rho\big(\phi_{\alpha x}(\pi^{-1}(x))\big)\rho\big(\phi_{\beta x}(\pi^{-1}(x))\big)^{-1} v   \Big)\\
			&=\Big(   x,\rho\big(\phi_{\alpha x}( \pi^{-1}(x) )\phi_{\beta x}( \pi^{-1}(x) )        \big)    \Big).
\end{split}
\end{equation}
Then 
\begin{equation}
  g\bab^E=\rho\Big(  \phi_{\alpha x}( \pi^{-1}(x) )\phi_{\beta x}\big( \pi^{-1}(x) \big)  \Big)\\
         =\rho(g^P\bab(x)).
\end{equation}

\end{proof}

\subsection{Sections on associated bundle}  \label{sec_fnequiv}
%-----------------------------------------

\subsubsection{Equivariant functions}

We consider a bundle $\dptvb{E=P\times_{\rho} V}{p}{M}$ associated with the principal bundle $\dptvb{P}{\pi}{M}$ and a section $\dpt{\psi}{M}{E}$. 
\[
 \xymatrix{ P \ar[rd]_{\displaystyle\pi^P} &&
 E=P\times_{\rho} V \ar[ld]^{\displaystyle\pi^E} \\ & M }
\]
A \defe{section}{section!of an associated bundle} of $E$ is a map $\dpt{\psi}{M}{E}$ such that $\pi^E\circ\psi=id_M$. We define the function $\dpt{\hpsi}{P}{V}$ by
\begin{equation}\label{eq:equiv_psi}
   \psi(\pi(\xi))=[\xi,\hpsi(\xi)].
\end{equation}
Let us see the condition under which this equation well defines $\hpsi$. First, remark that a $\psi$ defined by this equation is a section because $p[\xi,v]=\pi(\xi)$, so that $(p\circ\psi)(\pi(\xi))=\pi(\xi)$. Now, consider a $\eta$ such that $\pi(\eta)=\pi(\xi)$. Then there exists a $g\in G$ for which $\eta\cdot g=\xi$. For any $g$ and for this one in particular,
\[
  \psi(\pi(\eta))=[\eta,\hpsi(\eta)]=[\eta\cdot g,\rho(g^{-1})\hpsi(\eta)].
\]
Then equation \eqref{eq:equiv_psi} defines $\hpsi$ from $\psi$ if and only if 
\begin{equation}\label{eq:equiv_psi_b}
  \hpsi(\xi\cdot g)=\rho(g^{-1})\hpsi(\xi).
\end{equation}
This condition is called the \defe{equivariance}{equivariant} of $\hpsi$. Reciprocally, any equivariant function $\hpsi$ defines a section of $E=P\times_{\rho} V$.

If $\eta=\xi\cdot g=\chi\cdot k$, one define a sum
\begin{equation}\label{eq:def:som_E}
  [\xi,v]+[\chi,w]=[\eta,\rho(g)v+\rho(k)w].
\end{equation}
If $\dpt{\psi,\eta}{M}{E}$ are two sections defined by the equivariant functions $\dpt{\hpsi,\hat\eta}{P}{V}$, then the section $\psi+\eta$ is defined by the equivariant function $\hat\psi+\hat \eta$.

\subsubsection{For the endomorphism of sections of \texorpdfstring{$E$}{E}}\label{equivendo}
%////////////////////////////////////////////////////

Let us now make a step backward, and take $A$ in $\End{\Gamma(E)}$. We will now see that $A$ defines (and is defined by) an equivariant function $\dpt{\hat A}{P}{\End{V}}$. Let $\dpt{\psi}{M}{E}$ be in $\Gamma(E)$. If $\psi(x)=[\xi,v]$, we define the new section $A\psi$ by
\[
         (A\psi)(x)=[\xi,\hat A(\xi)v]=[\xi,\hat A(\xi)\hat\psi(\xi)].
\]
In order for $A\psi$ to be well defined, the function $\hat A$ must satisfy
\begin{equation}
     \hat A(\xi\cdot g)=\rho(g^{-1})\hat A(\xi)\rho(g)                 \label{equivA}
\end{equation}
for all $g$ in $G$.

\subsubsection{Local expressions}
%////////////////////////////////

We consider a local trivialization $\dpt{\phi^P_{\alpha}}{\pi^{-1}(\mU_{\alpha})}{\mU_{\alpha}\times G}$ of $P$ on $\mU_{\alpha}$ and the corresponding section $\dpt{\sigma_{\alpha}}{\mU_{\alpha}}{P}$ given by 
\[
\sigma_{\alpha}(x)=(\phi^P_{\alpha})^{-1}(x,e).
\]
We saw at page \pageref{eq:triv_P_E} that a trivialization of $P$ gives a trivialization of the associated bundle $E=P\times_{\rho} V$; the definition is
\begin{equation}
  \phi^E_{\alpha}[(\xi,v)]=( \pi(\xi),\rho(a)v )
\end{equation}
if $\phi_{\alpha}^P(\xi)=(\pi(\xi),a)$. With $\xi=\sigma_{\alpha}(x)$, we find
\begin{equation}
   \phi^E_{\alpha}[(\sigma_{\alpha}(x),v)]=(  \pi(\sigma_{\alpha}(x)),\rho(a)v  )
                                 =(x,v).
\end{equation}

The section $\psi$ can also be seen with respect to the ``reference''{} sections $\sigma_{\alpha}$ by means of the definition
\begin{equation}\label{eq:def:psisa}
  \psi(x)=[\sigma_{\alpha}(x),\psisa(x)]
\end{equation}
for a function $\dpt{\psisa}{M}{V}$.

\begin{lemma}
Let $\dpt{\psi}{M}{E}$ be a section and $\dpt{\hpsi}{P}{V}$, the corresponding equivariant function. Then
\[
   \psisa(x)=\hpsi(\sigma_{\alpha}(x)).
\]
\end{lemma}

\begin{proof}
By definition, $\psi(x)=\psi(\pi(\xi))=[\xi,\hpsi(\xi)]$.  Thus if we consider in particular $\xi=\sigma_{\alpha}(x)$,
\begin{equation}
  \phi^E_{\alpha}(\psi(x))=\phi^E_{\alpha}[\xi,\hpsi(\xi)]
                        =\phi^E_{\alpha}[s_{\alpha}(x),\hpsi(\sigma_{\alpha}(x))]
                        =(x,\hpsi(\sigma_{\alpha}(x))).
\end{equation}

\end{proof}

Let us anticipate. A \defe{spinor}{spinor} is a section of an associated bundle $E=P\times_{\rho} V$ where $P$ is a Lorentz-principal bundle, $V=\eC^2$ and $\rho$ is the spinor representation of Lorentz on $\eC^2$. So a spinor $\dpt{\psi}{M}{E}$ is \emph{locally} described by a function $\dpt{\psi\bsa}{M}{\eC^2}$. The latter is the one that we are used to handle in physics. In this picture, the transformation law of $\psi$ under a Lorentz transformation comes naturally.

Let $\{e_i\}$ be a basis of V; we consider some ``reference'' sections $\gamai$ of the associated bundle $E=P\times_{\rho} V$ defined by
\begin{equation}\label{eq:def:gamai}
\gamai(x)=[\phi_{\alpha}^{-1}(x,e),e_i].
\end{equation}
A general section $\dpt{\psi}{M}{E}$ is defined by an equivariant function $\dpt{\hpsi}{P}{V}$ which can be written as $\hpsi(\xi)=a^i(\xi)e_i$. If $\eta=\phi_{\alpha}^{-1}(x,e)$ and $\xi=\eta\cdot g(\xi)$,
\begin{equation}
  \psi(x)=[\xi,a^ie_i]
         =a^i[\eta,\rho(g)e_i]
	 =a^i(\xi)\bghd{\rho(g(\xi))}{i}{j}[\eta,e_j]
	 =c^j(\xi)\gamaj(x).
\end{equation}
Since the left hand side of this equation just depends on $x$, the functions $c^j$ must actually not depend on the choice of $\xi\in\pi^{-1}(x)$. So we have $\dpt{c^j}{M}{\eR}$. Indeed, if we choose $\chi\in\pi^{-1}(x)$, 
\[
  \psi(x)=c^j(\xi)\gamaj(x)
         \stackrel{!}{=}[\xi,a^{i}(\chi)e_i]
	 =\ldots
	 =c^j(\chi)\gamaj(x),
\]
so that $c^j(\xi)=c^j(\chi)$. So any section $\dpt{\psi}{M}{E}$ can be decomposed (over the open set $\mU_{\alpha}$) as 
\begin{equation}
  \psi(x)=s_{\alpha}^i(x)\gamai(x).
\end{equation}


\subsection{Associated and vector bundle}
%----------------------------------------

\subsubsection{General construction}
%///////////////////////////////////

We are going to see that a vector bundle is an associated bundle. For this, we consider a vector bundle $\dpt{p}{F}{M}$ with a fiber $F_x=V$ of dimension $m$. Let $G=GL(V)$, $P$ be the trivial principal bundle $P=M\times G$ and $\rho$ be the definition representation of $G$ on $V$. We set $E=P\times_{\rho} V$. Our aim is to put a vector bundle structure on $E$ which is equivalent to the one of $F$. The bijection $\dpt{b}{F}{E}$ will clearly be 
\begin{equation}
   b(\phi^{-1}(x,v))=[(x,e),v].
\end{equation}
We define the projection $\dpt{q}{E}{M}$ by
\[
   q[(x,g),w]=x
\]
and we have to show that  $q^{-1}(x)=\{  [(x,g),w]\tq g\in G \textrm{ and } w\in V  \}$ is a vector space isomorphic to $V$. The following definitions define a vector space structure:
\begin{itemize}
\item multiplication by a scalar: $\lambda[(x,g),v]=[(x,g),\lambda v]$,
\item addition: $[(x,g),v]+[(x,h),w]=[(x,e),\rho(g)v+\rho(h)w]$.
\end{itemize}
As local trivialization map, we consider 
\begin{equation}
\begin{aligned}
 \chi\colon q^{-1}(\mU)&\to \mU\times V \\ 
    [(x,g),v]  &\mapsto(x,\rho(g)v).
\end{aligned}
\end{equation}
With this structure, the bijection $b$ is an equivalence because $b|_{F_x}$ is a vector space isomorphism and $q\circ b=p$.

\subsection{Equivariant functions for a vector field}	\label{equivvec}
%----------------------------------------------------

In order to define in the same way an equivariant function for a vector field $X\in\cvec(M)$, we need to see $TM$ as an associated bundle.

\begin{proposition}
If $M$ is a $n$ dimensional manifold, we have the following isomorphism:
\[
     \SO(M)\times_{\rhoM}\eR^m\simeq TM
\]
where $\dpt{\rhoM}{\SO(m)\times\eR^m}{\eR^m}$ is defined by $\rhoM(A)v=Av$.
\end{proposition}
\begin{proof}
Recall that an element $b\in \SO(M)_x$ is a map $\dpt{b}{\eR^m}{T_xM}$. The isomorphism is no difficult. It is $\dpt{\psi}{\SO(M)\times_{\rhoM}\eR^m}{TM}$ defined by
\[\psi[b,v]=b(v).\]
It prove no difficult to see that $\psi$ is well defined, injective and surjective.
\end{proof}

Now, let us consider $X\in\cvec(M)$. We can see it as an element of $\Gamma(\SO(M)\times_{\rhoM}\eR^m)$, and define an equivariant function $\dpt{\hX}{\SO(M)}{\eR^m}$.

Let us make it more explicit. A vector field $Y\in\cvec(M)$ is, for each $x$ in $M$, the data of a tangent vector $Y_x\in T_xM$. Hence the formula $b(v)=Y_x$ defines an element $[b,v]$ in $\SO(M)\times_{\rhoM}\eR^m$, and $Y$ defines a section $\tilde{Y}(x)=[b(x),v(x)]$ of $\SO(M)\times_{\rhoM}\eR^m$. The associated equivariant function is given by $\hY(b)=v$ if $b(v)=Y_x$. In other words, the equivariant function $\dpt{\hY}{\SO(M)}{\eR^m}$ associated with the vector field $Y\in\cvec(M)$ is given by
\begin{equation}\label{r1404e1}
  \hY(b)=b^{-1}(Y_x),
\end{equation}
 where $x=\pi(b)$.


\subsection{Gauge transformations}
%---------------------------------

A \defe{gauge transformation}{gauge!transformation!of principal bundle} of the $G$-principal bundle $\dpt{\pi}{P}{M}$ is a diffeomorphism $\dpt{\varphi}{P}{P}$ such that

\begin{itemize}
\item $\pi\circ\varphi=\pi$,
\item  $\varphi(\xi\cdot g)=\varphi(\xi)\cdot g$.
\end{itemize}

When we consider some local sections on $\dpt{\sigma_{\alpha}}{\mU_{\alpha}}{P}$, we can describe a gauge transformation with a function $\dpt{\tilde{\varphi}_{\alpha}}{M}{G}$ by requiring
\[
   \varphi(\sigma_{\alpha}(x))=\sigma_{\alpha}(x)\cdot\tilde{\varphi}_{\alpha}(x).
\]
This formula defines $\varphi$ from $\tilde{\varphi}$ as well as $\tilde{\varphi}$ from $\varphi$.

The group of gauge transformations\index{gauge!transformation!of section of associated bundle} has a natural action on the space of sections given by
\begin{subequations}
   \begin{align}
   (\varphi\cdot\psi)(x)&=[\varphi(\xi),v].
\intertext{if $\psi(x)=[\xi,v]=[\xi,\hat{\psi}(\xi)]$. This law can also be seen on the equivariant function $\hpsi$ which defines $\psi$. The rule is}
   \widehat{\varphi\cdot\psi}(\xi)&=\hpsi(\varphi^{-1}(\xi)).
   \end{align}
\end{subequations}
Indeed, in the same way as before we find $(\varphi\cdot\psi)(x)=[\xi,\widehat{\varphi\cdot\psi}(x)]\stackrel{!}{=}[\varphi(\xi),v]=[\varphi(\xi),\hpsi(\xi)]$. Taking $\xi\to\varphi^{-1}(\xi)$ as representative, $(\varphi\cdot\psi)(x)=[\xi,\hpsi\circ\varphi^{-1}(\xi)]$.


\section{Adjoint bundle}
%+++++++++++++++++++++++

Let $\pi\colon P\to M$ be a $G$-principal bundle. The \defe{adjoint bundle}{adjoint!bundle} is the associated bundle $\Ad(P)=P\times_{\Ad}\mG$. An element of that bundle is an equivalent class given by\nomenclature[D]{$\Ad(P)$}{Adjoint bundle of the principal bundle $P$}
\[ 
  [\xi,X]=[\xi\cdot g,\Ad(g^{-1})X]
\]
for every $g\in G$. Here $\xi\in P$ and $X\in\mG$.

\section{Connection on vector bundle: local description}\label{sec:conn_vect}
%+++++++++++++++++++++++++++++++++++++++++++++++++++++++

A \defe{connection}{connection!on vector bundle} on the vector bundle $p\colon E\to M$ is a bilinear map
		\begin{equation}
		\begin{aligned}
			\nabla \colon \cvec(M)\times\Gamma(E) &\to \Gamma(E),\
			(X,s)&\mapsto \nabla_Xs
		\end{aligned}
	\end{equation}	
 such that 
 \begin{itemize}
 \item $\nabla_{fX}s=f\nabla_Xs$,
 \item $\nabla_X(fs)=(X\cdot f)s+f\nabla_Xs$
 \end{itemize}
for all $X\in\cvec(M)$, $f\in\Cinf(M)$ and $s\in\Gamma(E)$. The operation $\nabla$ is often called a \defe{covariant derivative}{covariant!derivative!on vector bundle}.

An easy example is given on the trivial bundle $E=\pr_1\colon M\times\eC\to M$. For this bundle, $\Gamma(E)=\Cinf(M,\eC)$ and the common derivation is a covariant derivation: $\nabla_Xs=(ds)X$.

\begin{proposition}
The value of $(\nabla_Xs)(x)$ depends only on $X_x$ and $s$ on a neighbourhood of $x\in M$.
\end{proposition}

\begin{proof}
Let $X$, $Y\in\cvec(M)$ such that $Y_z=f(z)X_z$  with $f(x)=1$ and $f(z)\neq 1$ everywhere else. Then
\[
  (\nabla_Ys)(x)-(\nabla_Xs)(x)=(f(x)-1)(\nabla_Xs)(x)=0.
\]
Since it is true for any function, the linearity makes that it cannot depend on $X_z$ with $z\neq x$. If we consider now two sections $s$ and $s'$ which are equals on a neighbourhood of $x$, we can write $s'=fs$ for a certain function $f$ which is $1$ on the neighbourhood. Then
\[
  (\nabla_Xs')(x)-(\nabla_Xs)(x)=(f(x)-1)(\nabla_Xs)(x)+(Xf)s(x)
\]
which zero because on a neighbourhood of $x$, $f$ is the constant $1$.
\end{proof}

This proposition shows that it makes sense to consider only local descriptions of connections.  Let $\{e_1,\ldots,e_r\}$ be a basis of $V$ and consider the local sections $\dpt{\ovS_{\alpha i}}{\mU_{\alpha}}{E}$,
\[
  \ovS_{\alpha i}(x)=\phi_{\alpha}^{-1}(x,e_i).
\]
A local section $\dpt{s_{\alpha}}{\mU_{\alpha}}{V}$ can be decomposed as $s_{\alpha}(x)=s_{\alpha}^i(x)e_i$ with respect to this basis (up to an isomorphism between the different $V$ at each point). Then on $\mU_{\alpha}$,
\begin{equation}
  s_{\alpha}^i\ovS_{\alpha i}(x)=s_{\alpha}^i(x)\phi_{\alpha}^{-1}(x,e_i)
                              =\phi_{\alpha}^{-1}(x,s_{\alpha}^ie_i)
			      =\phi_{\alpha}^{-1}(x,s_{\alpha}(x))
			      =s(x).
\end{equation}
The first equality is the definition of the product $\eR\times F\to F$.

So any $s\in\Gamma(E)$ can be (locally !) written under the form\footnote{be careful on the fact that the ``coefficient'' $s_{\alpha}^i$ depends on $x$ : the right way to express this equation is $s(x)=s^i_{\alpha}(x)\ovS_{\alpha i}(x)$.} $s=s_{\alpha}^i\ovS_{\alpha i}$; in particular $\nabla_X(\ovS_{\alpha i})$ can. We define the coefficients $\theta$ by\nomenclature[D]{$(\theta_{\alpha})_i^j$}{Matrix associated with a connection}
\begin{equation}
 \nabla_X(\ovS_{\alpha i})=(\theta_{\alpha})^j_i(X)\ovS_{\alpha j}.
\end{equation}
where, for each $i$ and $j$, $(\theta_{\alpha})^j_i$ is a $1$-form on $\mU_{\alpha}$. We can consider $\theta_{\alpha}$ as a matrix-valued $1$-form on $\mU_{\alpha}$.

\begin{proposition}	\label{PropFormnabXthe}
The formula
\begin{equation}\label{eq:nab_theta}
   (\nabla_Xs)_{\alpha}=Xs_{\alpha}+\theta_{\alpha}(X)s_{\alpha}
\end{equation}
gives a local description of the connection. 
\label{prop:namba_theta_u}
\end{proposition}

\begin{proof}
For any $s\in\Gamma(E)$, we have
\[ 
\begin{split}
\nabla_Xs&=\nabla_X\big(  \sum_j s^j_{\alpha}\ovS_{\alpha j}  \big)\\
         &=\sum_j\Big(  (Xs_{\alpha}^j)\ovS_{\alpha j} + s^j_{\alpha}\nabla_X\ovS_{\alpha j}     \Big)\\
	 &=\sum_i\Big[   (Xs_{\alpha}^i)+s_{\alpha}^j(\theta_{\alpha})_j^i(X)  \Big]\ovS_{\alpha i}.
\end{split}
\]
\end{proof}

\subsection{Connection and transition functions}
%//////////////////////////////////////////////////

A connection determines some local matrix-valued $1$-forms $\theta_{\alpha}$ on the trivialization $\mU_{\alpha}$. Two natural questions raise. The first is the converse: does a matrix-valued $1$-form defines a connection~? The second is to know  what is $\theta_{\alpha}$ in function of $\theta_{\beta}$ on $\mU_{\alpha}\cap\mU_{\beta}$ ? The answer to the latter is  given by the following proposition :

\begin{proposition}
The $1$-form $\theta_{\alpha}$ relative to the trivialization $(\mU_{\alpha},\phi_{\alpha})$ is related to the $1$-form $\theta_{\beta}$ relative to the trivialization $(\mU_{\beta},\phi_{\beta})$ by
\begin{equation}\label{eq:theta_g}
  \theta_{\beta}=g\bab^{-1} dg\bab+g\bab^{-1}\theta_{\alpha} g\bab.
\end{equation}
\end{proposition}

This equation looks like something you know ? If you think to equation \eqref{trans_A} or \eqref{tr_de_A} or any physical equation of gauge transformation for the bosons, then you are almost right.
 
\begin{proof}
We can use equation \eqref{eq:tr_sec} pointwise on $(\nabla_X s)_{\alpha}$ :
\begin{equation}
\begin{split}
(\nabla_X s)_{\alpha}&=g\bab(\nabla_Xs)_{\beta}\\
                  &=g\bab\big(   Xs_{\beta}+\theta_{\beta}(X)s_{\beta}   \big) \\
                  &=g\bab\big(   X(g_{\alpha\beta} s_{\alpha})+\theta_{\beta}(X)g_{\alpha\beta} s_{\alpha}   \big).
\end{split}
\end{equation} 
We have to compare it with equation \eqref{eq:nab_theta}. Note that $g\bab$ and $\theta_{\alpha}(X)$ are matrices, then one cannot do
\[
   g\bab\theta_{\beta}(X)g_{\alpha\beta} =g\bab g_{\alpha\beta}\theta_{\beta}(X)=\theta_{\beta}(X)
\]
by using $g\bab g_{\alpha\beta}=\mtu$.  Taking carefully subscripts into account, one sees that the correct form is $(g\bab)^i_j\theta_{\beta}(X)^j_k(g_{\alpha\beta})^k_l$. Applying Leibnitz formula ($X(fg)=f(Xg)+(Xf)g$), and making the simplification $g\bab g_{\alpha\beta}=\mtu$ in the first term, we find
\[
  \theta_{\alpha}(X)s_{\alpha}=g\bab(Xg_{\alpha\beta})s_{\alpha}+g_{\alpha\beta}^{-1}\theta_{\beta}(X)g_{\alpha\beta} s_{\alpha}.
\]
The claim follows from the fact that $Xg_{\alpha\beta}=dg_{\alpha\beta}(X)$.
\end{proof}

Notice that formula \eqref{eq:theta_g} shows in particular that $\theta_{\alpha}$ takes its values in the Lie algebra $\gl(V)$, see for example subsection \ref{SubSecgmudg}.

The inverse is given in the
\begin{proposition}	\label{Propformconnve}
If we choose a family of $\gl(V)$-valued $1$-forms $\theta_{\alpha}$ on $\mU_{\alpha}$ satisfying \eqref{eq:theta_g},then the formula
\[
  (\nabla_Xs)_{\alpha}=Xs_{\alpha}+\theta_{\alpha}(X)s_{\alpha}
\]
defines a connection on $E$.\label{prop:thet_conn_F}
\end{proposition}

\begin{proof}
Note that $\theta$ is $\Cinf(M)$-linear, thus
\begin{equation}
  (\nabla_{fX}s)_{\alpha}=(fX)s_{\alpha}+\theta_{\alpha}(fX)s_{\alpha}
                        =f[ Xs_{\alpha}+\theta_{\alpha}(X)s_{\alpha} ]
			=f(\nabla_Xs)_{\alpha}.
\end{equation}
In expressions such that $\theta_{\alpha}(X)(fs_{\alpha})$, the product is a matrix times vector product between $\theta_{\alpha}(X)$ and $s_{\alpha}$; the position of the $f$ is not important. So we can check the second condition :
\begin{equation}
\begin{split}
(\nabla_X(fs))_{\alpha}&=X(fs_{\alpha})+\theta_{\alpha}(X)(fs_{\alpha}) \\
                     &=X(f)s_{\alpha}+f(Xs_{\alpha})+f\theta_{\alpha}(X)s_{\alpha}\\
		     &=df(X)s_{\alpha}+f(\nabla_Xs)_{\alpha}.
\end{split}
\end{equation}
This concludes the proof.
\end{proof}


\subsection{Torsion and curvature}
%----------------------------------

The map $\dpt{T^{\nabla}}{\cvec(X)\times\cvec(X)}{\cvec(X)}$ defined by
\begin{equation}
     T^{\nabla}(X,Y)=\nabla_XY-\nabla_YX-[X,Y]\label{deftorsion}
\end{equation}
is the \defe{torsion}{torsion!of a connection} of the connection $\nabla$. When $T^{\nabla}(X,Y)=0$ for every $X$ and $Y$ in $\cvec(X)$, we say that $\nabla$ is a \defe{torsion free}{torsion!free, connection} connection. Let $X$, $Y$ be in $\cvec(M)$, and consider the map $\dpt{R(X,Y)}{\Gamma(E)}{\Gamma(E)}$ defined by
		\begin{equation}
		\begin{aligned}
			R(X,Y) \colon \Gamma(E) &\to \Gamma(E)\
			s&\mapsto \nabla_X\nabla_Ys-\nabla_Y\nabla_Xs-\nabla_{[X,Y]}s.
		\end{aligned}
	\end{equation}	

For each $x\in M$, $R$ can be seen as a bilinear map $\dpt{R}{T_xM\times T_xM}{\End(E_x)}$. It is called the \defe{curvature}{curvature} of the connection $\nabla$. For every $f\in C^{\infty}(M)$, it satisfies
\[
 R(fX,Y)s=fR(X,Y)s=R(X,Y)fs.
\]
    

In a trivialization $(\mU_{\alpha},\phi_{\alpha})$, we have $(\nabla_Xs)_{\alpha}=Xs_{\alpha}+\theta_{\alpha}(X)s_{\alpha}$. In the expression of $(R(X,Y)s)_{\alpha}$, the terms coming from the $Xs_{\alpha}$ part of covariant derivative make
\[
  XYs_{\alpha}-YXs_{\alpha}-[X,Y]s_{\alpha}=0.
\]
The other terms are no more than matricial product, hence the formula
\begin{equation}
  (R(X,Y)s)_{\alpha}=\Omega_{\alpha}(X,Y)s_{\alpha}
\end{equation}
 defines a $2$-form $\Omega_{\alpha}$ which takes values in $GL(r,\eK)$. We can find an expression for $\Omega$ in terms of $\theta$ :
\[
  \Omega_{\alpha}(X,Y)=X\theta_{\alpha}(Y)-Y\theta_{\alpha}(X)-\theta_{\alpha}([X,Y])+\theta_{\alpha}(X)\theta_{\alpha}(Y)-\theta_{\alpha}(Y)\theta_{\alpha}(X);
\]
it is written as
\begin{equation}\label{eq:Omega_ttheta}
\Omega_{\alpha}=d\theta_{\alpha}+\theta_{\alpha}\wedge\theta_{\alpha}=d\theta_{\alpha}+\frac{1}{2}[\theta_{\alpha},\theta_{\alpha}]
\end{equation}
which is a notational shortcut for
\begin{equation}		\label{EaCurvdVVsq}
  \Omega_{\alpha}(X,Y)=d\theta_{\alpha}(X,Y)+[\theta_{\alpha}(X),\theta_{\alpha}(Y)].
\end{equation}
These equations are called \defe{structure equations}{structure!equations}. Pointwise, the second term is a matrix commutator; be careful on the fact that, when we will speak about principal bundle, the forms $\theta$'s will take their values in a Lie algebra. On $\mU_{\alpha}\cap\mU_{\beta}$, we have
\[
  \Omega_{\beta}(X,Y)=g\bab^{-1}\Omega_{\alpha}(X,Y)g\bab.
\]
The curvature and the connection fulfill the \defe{Bianchi identities}{Bianchi identities} :

\begin{lemma}
  \[
     d\Omega_{\alpha}+[\theta_{\alpha}\wedge\Omega_{\alpha}]=0.
  \]
\end{lemma}

\begin{proof}
For each matricial entry, $\theta_{\alpha}$ is a $1$-form on $\mU_{\alpha}$, then $\theta_{\alpha}(X)$ is a function which to $x\in M$ assign $\theta_{\alpha}(x)(X_x)\in\eR$. So we can apply $d$ and Leibnitz on the product $\theta_{\alpha}(X)\theta_{\alpha}(Y)$.
\[
 d\big(  \theta_{\alpha}(X)\theta_{\alpha}(Y)  \big)=\theta_{\alpha}(X)d\theta_{\alpha}(Y)+d\theta_{\alpha}(X)\theta_{\alpha}(Y).
\]
Differentiating equation \eqref{eq:Omega_ttheta}, $d\Omega_{\alpha}=d\theta_{\alpha}\wedge\theta_{\alpha}-\theta_{\alpha}\wedge d\theta_{\alpha}$. 
\end{proof}


%////////////////////////////////////////////////////////////////////////////////////////////////////////////////////////////
\subsection{Divergence, gradient and Laplacian}

We define the \defe{gradient}{gradient} of a function $f\in C^{\infty}(M)$, denoted by $\nabla f$\nomenclature[F]{$\nabla f$}{Gradient of the function $f$} as the vector field such that
\begin{equation}
	g(\nabla f,X)=X(f).
\end{equation}
The \defe{divergence}{divergence} of a vector field $X\in \Gamma(TM)$, is the function $\nabla\cdot X\in C^{\infty}(M)$\nomenclature[F]{$\nabla\cdot X$}{divergence of the vector field $X$} defined by
\begin{equation}
  (\nabla\cdot X)(x)=\tr\big( v\mapsto\nabla_vX \big)
\end{equation}
where the trace is the one of $v\mapsto\nabla_vX$ seen as an operator on $T_xM$. The \defe{Laplacian}{laplace operator} of the function $f$ is the function $\Delta f$\nomenclature[F]{$\Delta f$}{Laplace operator} given by
\begin{equation}
\Delta f=\nabla\cdot(\nabla f).
\end{equation}

\section{Connexion on vector bundle: algebraic view}
%+++++++++++++++++++++++++++++++++++++++++++++++++++

A \defe{connection}{connection!on vector bundle} on the vector bundle $\pi\colon E\to M$ is a linear map
\[ 
  \nabla\colon \Gamma^{\infty}(E)\to \Gamma^{\infty}(E)\otimes\Omega^1(M)
\]
which satisfies the Leibnitz rule
\begin{equation}
\nabla(\sigma f)=(\nabla\sigma)f+\sigma\otimes df
\end{equation}
for any section $\sigma\colon M\to E$ and function $f\colon M\to \eC$. If $\{ \sigma_i \}$ is a local basis of $E$, one can write $\sigma=\sigma_if^i$ and one defines the \defe{Christoffel symbols}{Christoffel symbol} $\Gamma_{i\mu}^{j}$ in this basis by
\begin{equation}
\nabla \sigma=\nabla (\sigma_if^i)
		=(\nabla \sigma_i)f^i+\sigma_i\otimes f(f^i)
		=f^i\Gamma_{i\mu}^{j}\sigma_j\otimes dx^{\mu}+\sigma_i\otimes d(f^i).
\end{equation}
The notations $d\sigma=\sigma_i\otimes d(f^i)$ and $\Gamma\sigma=f^i\Gamma_{i\mu}^{j}\sigma_j\otimes dx^{\mu}$ lead us to the compact usual form
\[ 
  \nabla\sigma=(d+\Gamma)\sigma.
\]

When $E=TM$ over a (pseudo)Riemannian manifold $M$, we know the Levi-Civita\index{connection!Levi-Civita} connection which is compatible with the metric:
\begin{equation}\label{eq_230605r1}
  g(\nabla X,Y)+g(X,\nabla Y)=d\big( g(X,Y) \big).
\end{equation}
One can see $g$ as acting on $\big( \cvec(M)\otimes\Omega^1(M) \big)\times\cvec(M)$ with
\[ 
 g\big( r^i_{\nu}\partial_i\otimes dx^{\nu},t^j\partial_j \big):=r^i_{\nu}j^jg(\partial_i,\partial_j)dx^{\nu}, 
\]
which at each point is a form. From condition \eqref{eq_230605r1}, we see $\nabla$ as a Levi-Civita connection on the bundle $E=T^*M$ which values in 
\[ 
 \Gamma^{\infty}(T^*M)\otimes\Omega^1(M)\simeq\Omega^1(M)\otimes\Omega^1(M).
\]
This is defined as follows. A $1$-form $\omega$ can always be written under the for $\omega=X^{\flat}:=g(X,.)$ for a certain $X\in\cvec(M)$. Then  \eqref{eq_230605r1} gives
\[ 
  (\nabla X)^{\flat}Y+\omega(\nabla Y)=d(\omega Y),
\]
and we put $\nabla\omega=(\nabla X)^{\flat}$, i.e
\begin{equation}
  (\nabla\omega)Y=d(\omega Y)-\omega(\nabla Y)
\end{equation}
for all $Y\in\cvec(M)$. When $\omega=dx^i$ and $Y=\partial_j$, we find
\begin{equation}
(\nabla dx^i)\partial_j=d(dx^i\partial_j)-dx^i(\nabla\partial_j)
		=d(\delta^i_j)-\Gamma_{jk}^{l}\partial_l\otimes dx^k
		=-\Gamma_{jk}^{l}\delta_l^i\otimes dx^k
		=-\Gamma_{jk}^{i}\,dx^k.
\end{equation}
So we get the local formula
\begin{equation}
\nabla dx^i=-\Gamma_{jk}^{i}\,dx^j\otimes dx^k.
\end{equation}
If the form writes locally $\omega=dx^if_i$,
\begin{equation}
  \nabla\omega=\nabla(dx^i)f_i+dx^i\otimes df_i
		=-f_i\Gamma_{jk}^{i}\,dx^j\otimes dx^k+d\omega
		=(d-\tilde\Gamma)\omega
\end{equation}
where we taken the notations $d\omega=dx^i\otimes df_i$ and $\tilde\Gamma\omega=f_i\Gamma_{jk}^{i}dx^j\otimes dx^k$.

%---------------------------------------------------------------------------------------------------------------------------
\subsection{Exterior derivative}

If $E$ is a $m$-dimensional vector bundle over $M$ and $s\colon M\to E$ is a section, we say that a \defe{exterior derivative}{exterior!derivative} is a map $D\colon \Gamma(E)\to \Gamma(E\otimes \Omega^1M)$ such that for every $f\in C^{\infty}(M)$ we have
\[ 
  D(fs)=s\otimes df+f(Ds).
\]
An exterior derivative can be extended to $D\colon \Gamma(E\otimes\Omega^kM)\to \Gamma(E\otimes\Omega^{k+1}M)$ imposing the condition
\begin{equation}		\label{EqExtExtDerrk}
D(\omega\wedge\alpha)=(D\omega)\wedge\alpha+(-1)^k\omega\wedge d\alpha
\end{equation}
for every $\omega\in\Gamma(E\otimes\Omega^kM)$ and $\alpha\in\Gamma(E\otimes\Omega^lM)$ . The result is an element of $\Gamma(E\otimes\Omega^{k+l+1}M)$.

Coordinatewise expressions are obtained when one choose a specific section $(e_i)$ of the frame bundle of $E$. In that case for each $i$, the derivative $e_i$ is an element of $\Gamma(E\otimes\Omega^1M)$ and we define $\omega_i^j\in\Omega^1(M)$\nomenclature{$\omega_i^j$}{Connection form} by
\begin{equation}
De_i=\sum_{j=1}^ke_j\otimes \omega_i^j.
\end{equation}
For each $i$ and $j$, we have an element $\omega_i^j\in\Omega^1(M)$, so that we say that $\omega\in\Omega^1(M,\gl(m))$. Now a section can be expressed as $s=s^ie_i$ where $s^i$ are functions, so we have
\begin{align}
  D(s)=D(s^ie_i)	&=e_i\otimes ds^i+s^iD(e_i)=e_i\otimes ds^i+s^ie_j\otimes \omega_i^j=e_i\otimes ds^i+e_i\otimes s^j\omega_j^i.
\end{align}
Expressed in component, we find $D(s)^i=ds^i+s^j\omega_j^i$, so that we often write
\begin{equation}
D=d+\omega.
\end{equation}
When a section $e$ is given, we write $s=s^i(e)e_i$, indicating the dependence of the functions $s^i$ in the choice of the frame $e$ :
\[ 
  D(s)=e_i\otimes ds^i(e)+e_i\otimes s^j(e)\omega(e)_j^i.
\]
When we apply both sides to a vector $X\in\Gamma(TM)$, we find
\begin{equation}
D_X(s)=e_i\otimes\Big( X(s^i)+s^j\omega^i_j(X) \Big).
\end{equation}

By convention we say that, when $f\in C^{\infty}(M)$, is a function, $D_X$ reduces to the action of the vector field $X$:
\begin{equation}
  D_X(f)=X(f).
\end{equation}


%----------------------------------------------------------------------------------------------------------------------------
\subsubsection{Covariant exterior derivative}

An important exterior derivative is the covariant exterior derivative. If the vector bundle $E$ is endowed by a covariant derivative $\nabla$, we define the corresponding \defe{covariant exterior derivative}{covariant!derivative!exterior } by the following :
{
\renewcommand{\theenumi}{\arabic{enumi}.}
\begin{enumerate}
\item for a section $s\colon M\to E$ (i.e. a $0$-form) we define
\begin{equation}
   (d_{\nabla}s)(X)=\nabla_Xs,
\end{equation}
\item and on the $1$-form $\sum_i(s_i\otimes\omega_i \big)\in\Gamma(E\otimes T^*M)$,
\begin{equation}
d_{\nabla}\big( \sum_is_i\otimes\omega_i \big)=\sum_i(d_{\nabla}s_i)\wedge\omega_i+\sum_is_i\otimes d\omega_i.
\end{equation}
\end{enumerate}
}		% Fin de la mise en nombre arabes pour la liste énumérée
The latter relation is the condition \eqref{EqExtExtDerrk} with $k=0$.

%\\\\\\\\\\\\\\\\\\\\\\\\\\\\\\\\\\\\\\\\\\\\\\\\\\\\\\\\\\\\\\\\\\\\\\\\\\\\\\\\\\\\\\\\\\\\\\\\\\\\\\\\\\\\\\\\\\\\\\\\\\\\
\subsubsection{Soldering form and torsion}


Let us particularize to the case where $E$ has the same dimension as the manifold. In that case, we can introduce a \defe{soldering form}{soldering form}, that is an element $\theta\in \Omega^1(M,E)$ such that for every $x\in M$ the map $\theta_x\colon T_xM\to E_x$ is a vector space isomorphism.
When a soldering form $\theta$ is given, the \defe{torsion}{torsion!of exterior derivative} is the exterior derivative $D$ is
\begin{equation}
	T=D\theta.
\end{equation}
Using a local frame $e$, we have forms $\theta^i(e)\in\Omega^1(M)$ such that
\[ 
  \theta(X)=\theta^i(X)e_i.
\]
We see $\theta$ as an element of $\Gamma(E\otimes \Omega^1(M))$ by identifying $\theta=e_i\otimes\theta^i$. Thus we have
\[ 
D\theta=D(e_i\otimes\theta^i)	=De_i\wedge\theta^i(e)+e_i\wedge d\theta^i(e)
				=(e_j\otimes^j_i)\wedge\theta^i(e)+e_i\wedge d\theta^i(e),
\]
or in coordinates :
\begin{equation}
  (D\theta)^i=\omega_j^i\wedge \theta^j(e)+d\theta^i(e).
\end{equation}
Notice that it provides the formula
\begin{equation}
T=d_{\omega}\theta
\end{equation}
for the torsion as exterior covariant derivative of the connection form.

%\\\\\\\\\\\\\\\\\\\\\\\\\\\\\\\\\\\\\\\\\\\\\\\\\\\\\\\\\\\\\\\\\\\\\\\\\\\\\\\\\\\\\\\\\\\\\\\\\\\\\\\\\\\\\\\\\\\\\\\\\\\\
\subsubsection{Example : Levi-Civita}

We consider the vector bundle $E=TM$ and the local basis $e_i=\partial_i$. An exterior derivative in this case is a map $D\colon \Gamma(TM)\to \Gamma\Big( TM\otimes\Omega^1M \Big)$. In that particular case, we denote by $\nabla_XY$ the vector field $D(Y)X$, and it is computed by first writing $D(X)_x=\sum_iZ_x^i\otimes\omega_x^i$ with $Z^i\in\Gamma(TM)$ and $\omega^i\in\Omega^1(M)$. The we have
\begin{equation}
D(X)_xY_x=\omega+x^i(Y_x)Z_x^i.
\end{equation}
A good choice of soldering form is $\theta_x=\id$ for every $x\in M$, or $\theta(X)=X$. In coordinates, that soldering form is given by $\theta^i(\partial_j)=\delta^i_j$. The \defe{Christoffel symbols}{Christoffel symbol} are defined by
\begin{equation}
\nabla_{\partial_i}\partial_j=\Gamma_{ij}^k\partial_k,
\end{equation}
and the covariant derivative reads
\begin{equation}		\label{EqCovDerGamChr}
\nabla_XY	= \nabla_{X^i\partial_i}(Y^j\partial_j)
		= X^i\Big( (\partial_iY^j)\partial_j+Y^j\nabla_{\partial_i}\partial_j \Big)
		= \Big( X(Y^k)+X^iY^j\Gamma_{ij}^k \Big) \partial_k.
\end{equation}

We can determine the Christoffel symbols in function of the connection form using the fact that on the one hand, $\nabla_{\partial_i}\partial_j=\Gamma_{ij}^k\partial_k$, and on the other hand,
\[ 
  \nabla_{\partial_i}\partial_j=D(\partial_j)(\partial_i)=\partial_k\otimes\omega_j^k(\omega_i), 
\]
so that
\begin{equation}
	\Gamma_{ij}^k=\omega_j^k(\partial_i)
\end{equation}
Now we can get the same result as equation \eqref{EqCovDerGamChr} using the exterior derivative formalism. First we have $DY=\partial_i\otimes dY^i+\partial_i\otimes X^j\omega_j^i$, so that
\[ 
  (DY)X=\partial_i\otimes dY^i(X)=\partial_i\otimes X^j\omega_j^i(X^k\partial_k),
\]
in which we use the relation $\omega_j^i(X^k\partial_k)=X^k\omega_j^i(\partial_k)=X^k\Gamma_{jk}^i$ to get
\[ 
  (DY)X=\big( X(Y^i)+X^jX^k\Gamma^i_{jk} \big)\partial_i.
\]
Notice that the anti-symmetric part of $\Gamma$ with respect to its two lower indices does not influence the covariant derivative. Let us compute the torsion in terms of $\Gamma$. For that remark that $d\theta^i=0$ because 
\[ 
  (d\theta^i)(X,Y)=X\theta^i(Y)-Y\theta^i(X)-\theta^i\big( [X,Y] \big)=X(Y^i)-Y(X^i)-[X,Y]^i=0.
\]
Thus we have
\begin{align*}
(D\theta)(\partial_k\otimes\partial_l)	&=\big( (D\partial_i)\partial_k \big)\theta^i(\partial_l)-\big( (D\partial_i)\partial_l \big)\theta^i(\partial_k)\\
					&=\delta_l^i\Gamma_{ik}^j\partial_j-\delta_k^i\Gamma_{il}^j\partial_j\\
					&=(\Gamma_{lk}^j-\Gamma^j_{kl})\partial_j.
\end{align*}
The connection $\nabla$ is moreover compatible with the metric because
\[ 
  \nabla_Z\big( g(X,Y) \big)=Z\big( \eta(eX,eY) \big)=\eta\big( \underbrace{D_Z(eX)}_{=e(\nabla_ZX)},eY \big)+\eta\big( eX,D_Z(eY) \big)=g(\nabla_ZX,Y)+g(X,\nabla_ZY).
\]




\section{Connection on principal bundle}  %\label{subsec_defconnprinc}
%------------------------------------------

\subsection{First definition:  \texorpdfstring{$1$}{1}-form}
%------------------------------

\label{pg_connpriic}
We consider a $G$-principal bundle 
\[
\xymatrix{%
   G \ar@{~>}[r]		&	P\ar[d]^{\pi}\\
   				&	  M
}
\]
and $\yG$, the Lie algebra of $G$. 

\begin{definition}
A \defe{connection}{connection!on principal bundle} on $P$  is a $1$-form $\omega\in\Omega(P,\yG)$ which fulfills

\begin{itemize}\label{pg:def:conne}
\item $\omega_{\xi}(A^*_{\xi})=A$,
\item $(R_g^*\omega)_{\xi}(\Sigma)=\Ad(g^{-1})(\omega_{\xi}(\Sigma))$,
\end{itemize}
for all $A\in\yG$, $g\in G$, $\xi\in P$ and $\Sigma\in T_{\xi} P$
\label{defconnform}
\end{definition}
Here, $R_g$ is the right action: $R_g\xi=\xi\cdot g$ and $A^*$ stands for the \defe{fundamental field}{fundamental!vector field} associated with $A$ for the action of $G$ on $P$:
\begin{equation} \label{defastar}
   A^*_{\xi}=\Dsdd{ \xi\cdot e^{-tA} }{t}{0},
\end{equation}
For each $\xi\in P$, we have $\dpt{\omega_{\xi}}{T_{\xi} P}{\yG}$. See section \ref{sec:fond_vec}.

If $\alpha$ is a connection $1$-form on $P$, we say that $\Sigma$ is an \defe{horizontal}{horizontal} vector field if $\alpha_{\xi}(\Sigma)=0$ for all $\xi\in P$. If $X_x\in T_xM$ and $\xi\in\pi^{-1}(x)$, there exists an unique\footnote{See \cite{kobayashi}, chapter II, proposition 1.2.} $\Sigma$ in $T_{\xi} P$ which is horizontal and such that $\pi_*(\Sigma)=X_x$. This $\Sigma$ is called the \defe{horizontal lift}{horizontal!lift} of $X_x$. We can also pointwise construct the horizontal lift of a vector field. The one of $X$ is often denoted by $\overline{X}$; it is an element of $\cvec(P)$.

\subsection[Horizontal space]{Second definition: horizontal space}
%-------------------------------------------------------------------

For each $\xi\in P$, we define the \defe{vertical space}{vertical space} $V_{\xi} P$ as the subspace of $T_{\xi} P$ whose vectors are tangent to the fibers: each $v\in V_{\xi} P$ fulfills $d\pi v=0$. Any such vector is given by a path contained in the fiber of $\xi$. So, $v\in V_{\xi} P$ if and only if there exists a path $g(t)\in G$ such that $v=\Dsdd{\xi\cdot g(t)}{t}{0}$.

A \defe{connection}{connection!on principal bundle} $\Gamma$ is a choice, for each $\xi\in P$, of an \defe{horizontal space}{horizontal!space} $H_{\xi} P$ such that

\begin{itemize}
\item $T_{\xi} P=V_{\xi} P\oplus H_{\xi} P$,
\item $H_{\xi\cdot g}=(dR_g)_{\xi} H_{\xi}$,
\item $H_{\xi} P$ depends on $\xi$ under a differentiable way.
\end{itemize}
The second condition means that the distribution $\xi\to H_{\xi}$ is invariant under $G$. Thanks to the first one, for each $X\in T_{\xi} P$, there exists only one choice of $Y\in H_{\xi} P$ and $Z\in V_{\xi} P$ such that $X=Y+Z$. These are denoted by $vX$ and $hX$ and are naturally named \emph{horizontal} and \emph{vertical components} of $X$. The third condition means that if $X$ is a differentiable vector field on $P$, then $vX$ and $hX$ are also differentiable vector fields. We will often write $V_{\xi}$ and $H_{\xi}$ instead of $V_{\xi} P$ and $V_{\xi} P$.

The word \emph{connection} probably comes from the fact that the horizontal space gives a way to jump from a fiber to the next one. 
When we consider a connection $\Gamma$, we can define a $\yG$-valued connection $1$-form by 
\[
   \omega(X)^*_{\xi}=vX_{\xi}.
\]
The existence is explained in section \ref{sec:fond_vec}. It is clear that $\omega(X)=0$ if and only if $X$ is horizontal. The theorem which connects the two definitions is the following.

\begin{theorem}
If $\Gamma$ is a connection on a $G$-principal bundle, and $\omega$ is its $1$-form, then

\begin{enumerate}
\item\label{enuyai} for any $A\in\yG$, we have $\omega(A^*)=A$,
\item\label{enuyaii} $(R_g)^*\omega=\Ad(g^{-1})\omega$, i.e. for any $X\in T_{\xi} P$, $g\in G$ and $\xi\in M$,
\[
    \omega((dR_g)_{\xi} X)=\Ad(g^{-1})\omega_{\xi}(X)
\]
\end{enumerate}
Conversely, if one has a $\yG$-valued $1$-form on $P$ which fulfills these two requirement, then one has one and only one connection on $P$ whose associated $1$-form is $\omega$.

\end{theorem}

\begin{proof}
\ref{enuyai} The definition of $\omega$ is $\omega(X)^*_{\xi}=vX$. Then $\omega(A^*)^*_{\xi}=vA^*_{\xi}=A^*_{\xi}$ because $A^*$ is vertical. From lemma \ref{lem:As_Bs_A_B}, $\omega(A^*)=A$.

\ref{enuyaii} Let $X\in\cvec(P)$. If $X$ is horizontal, the definition of a connection makes $dR_d X$ also horizontal, then the claim becomes $0=0$ which is true. If $X$ is vertical, there exists a $A\in\yG$ such that $X=A^*$ and a lemma shows that $dR_gX$ is then the fundamental field of $\Ad(g^{-1})A$. Using the properties of a connection,
\begin{equation}
  (R^*_g\omega)_{\xi}(X)=\omega_{\xi\cdot g}(dR_g X)=\Ad(g^{-1})A=\Ad(g^{-1})\omega_{\xi}(X).
\end{equation}

Now we turn our attention to the inverse sense: we consider a $1$-form which fulfills the two conditions and we define
\begin{equation}
   H_{\xi}=\{X\in T_{\xi} P\tq \omega(X)=0\}.
\end{equation}
We are going to show that this prescription is a connection. First consider a $X\in V_{\xi}$, then $X=A^*$ and $\omega(X)=A$. So $H_{\xi}\cap V_{\xi}=0$. Now we consider $X\in T_{\xi} P$ and we decompose it as
\[
   X=A^*+(X-A^*)
\]
where $A^*$ is the vertical component of $X$. If $\omega(dR_g X)=0$ for all $g\in G$, then $\omega(X)=0$, then a vector $X\in H_{\xi}$ fulfills at most $\dim G$ independent constraints $\omega(dR_g X)=0$ and $\dim H_{\xi}$ is at least $\dim P-\dim G$. On the other hand, $\dim V_{\xi}=\dim G$; then
\[
  \dim V_{\xi}+\dim H_{\xi}\geq\dim G+\dim P-\dim G.
\]
Then the equality must holds and $V_{\xi}\oplus H_{\xi}=T_{\xi} P$.

We have now to prove that $\omega$ is the connection form of $H_{\xi}$, i.e. that $\omega(X)$ is the unique $A\in\yG$ such that $A^*_{\xi}$ is the vertical component of $X$. Indeed if $X\in T_{\xi} P$, it can be decomposed as into $A^*\in V_{\xi}$ and $Y\in H_{\xi}$ and
\[
   \omega(X)=\omega(A^*+Y)=\omega(A^*)=A.   
\]

It remains to be proved that the horizontal space $H_{\xi}$ of any connection $\Gamma$ is related to the corresponding $1$-form $\omega$ by $H_{\xi}=\{X\in T_{\xi} P\tq\omega_{\xi}(X)=0\}$. From the connection $\Gamma$, the $1$-form is defined by the requirement that $\omega(X)^*_{\xi}=vX_{\xi}$. For $X\in H_{\xi}$, it is clear that $vX=0$, so that $\omega(X)^*=0$. This implies $\omega(X)=0$ because we suppose that the action of $G$ is effective.

\end{proof}

The projection $\dpt{\pi}{P}{M}$ induces a linear map $\dpt{d\pi}{T_{\xi} P}{T_xM}$. We will see that, when a connection is given, it is an isomorphism between $H_{\xi}$ and $T_xM$ (if $x=\pi(\xi)$). The \defe{horizontal lift}{horizontal!lift}\index{lift!horizontal} of $X\in\cvec(M)$ is the unique horizontal vector field (i.e. it is pointwise horizontal) such that $d\pi(\ovX_{\xi})=X_{\pi(\xi)}$. The proposition which allows this definition is the following.

\begin{proposition}
For a given connection on the $G$-principal bundle $P$ and a vector field $X$ on $M$, there exists an unique horizontal lift of $X$. Moreover, for any $g\in G$, the horizontal lift is invariant under $dR_g$.

The inverse implication is also true: any horizontal field on $P$ which is invariant under $dR_g$ for all $g$ is the horizontal lift of a vector field on $M$. 
\end{proposition}

This proposition comes from \cite{kobayashi}, chapter II, proposition 1.2.

\begin{proof}
We consider the restriction $\dpt{d\pi}{H_{\xi}}{T_{\pi(\xi)}M}$. It is injective because $d\pi(X-Y)$ vanishes only when $X-Y$ is vertical or zero. Then it is zero. It is cleat that $\dpt{d\pi}{T_{\xi} P}{T_{\pi(\xi)}M}$ is surjective. But $d\pi X=0$ if $X$ is vertical, then $d\pi$ is surjective from only $H_{\xi}$.

So we have existence and unicity of an horizontal lift. Now we turn our attention to the invariance. The vector $dR_g\ovX_{\xi}$ is a vector at $\xi\cdot g$. From the definition of a connection, $dR_g H_{\xi}=H_{\xi\cdot g}$, then $dR_g\ovX_{\xi}$ is the unique horizontal vector at $\xi\cdot g$ which is sent to $X_x$ by $d\pi$. Thus it is $\ovX_{\xi\cdot g}$.

For the inverse sense, we consider $\ovX$, an horizontal invariant vector field on $P$. If $x\in M$, we choose $\xi\in\pi^{-1}(x)$ and we define $X_x=d\pi(\ovX_{\xi})$. This construction is independent of the choice of $\xi$ because for $\xi'=\xi\cdot g$, we have
\[
   d\pi(\ovX_{\xi'})=\pi(dR_g\ovX_{\xi})=\pi(\ovX_{\xi}).
\]
\end{proof}

An other way to see the invariance is the following formula:
\[
   \ovX_{\xi\cdot g}=(dR_g)_{\xi} \ovX_{\xi}.
\]
By definition, $\ovX_{\xi\cdot g}$ is the unique vector of $T_{\xi\cdot g}P$ which fulfils $d\pi\ovX_{\xi\cdot g}=X_x$ if $\xi\pi^{-1}(x)$, so the following computation proves the formula:
\begin{equation}
  (d\pi)_{\xi\cdot g}((dR_g)_{\xi}\ovX_{\xi})=d(\pi\circ R_g)_{\xi}\ovX_{\xi}
                                         =d\pi_{\xi}\ovX_{\xi}
					 =X_x.
\end{equation}

\subsection{Curvature}
%////////////////////////

The curvature of a vector or associated bundle satisfies $\Omega_{\alpha}=d\theta_{\alpha}+\theta_{\alpha}\wedge\theta_{\alpha}$. So we naturally define the \defe{curvature}{curvature!on principal bundle} of the connection $\omega$ on a principal bundle as the $\yG$-valued $2$-form
\begin{equation} 
  \Omega=d\omega+\omega\wedge\omega.
\end{equation}
When we consider a local section\label{PgLocSecCurv} $\dpt{\sigma_{\alpha}}{\mU_{\alpha}}{P}$ on $\mU_{\alpha}\subset M$, we can express the curvature with a $2$-form on $M$ instead of $P$ by the formula \label{pg:curv_princ}
\[ 
 F\bsa=\sigma_{\alpha}^*\Omega,
\]
or, more explicitly, by $F\bsa_x(X,Y)=\Omega_{\sigma_{\alpha}(x)}(d\sigma_{\alpha} X,d\sigma_{\alpha} Y))$. Note that if $\yG$ is abelian, $\Omega=d\omega$ and $d\Omega=0$.

\section{Exterior covariant derivative and Bianchi identity} 
%--------------------------------------------------------------

Let $\omega\in\Omega^1(P,\mG)$ be a connection $1$-form on the $G$-principal bundle $P$. Using the operation $[.\wedge .]$ defined in section \ref{SecLiaAlgformval}, we define the \defe{exterior covariant derivative}{exterior!covariant derivative} by\nomenclature{$d_{\omega}$}{Exterior covariant derivative associated with the connection form $\omega$}  %\index{exterior covariant derivative}
\begin{align}
d_{\omega}\alpha&=d\alpha+\frac{ 1 }{2}[\omega\wedge\alpha]&\text{when $\alpha\in\Omega^1(P,\mG)$},\\
d_{\omega}\beta&=d\beta+[\omega\wedge\beta]&\text{when $\beta\in\Omega^2(P,\mG)$},
\end{align}

The \defe{curvature}{curvature!form} is the $2$-form defined by
\begin{equation}
\Omega=d_{\omega}\omega=d\omega+\omega\wedge\omega
\end{equation}
where $d_{\omega}$ is the exterior covariant derivative associated with the connection form $\omega$, and the wedge has to be understood as in equation \eqref{EqAbuswesgeomom}.

\begin{proposition}
The curvature form satisfies the identity
\begin{equation}
d_{\omega}\Omega=0
\end{equation}
which is the Bianchi identity\index{Bianchi identities}
\end{proposition}

\begin{proof}
taking the differential of $\Omega=d\omega+\omega\wedge\omega$, we find
\[ 
  d\Omega=d^2\omega+d\omega\wedge\omega-\omega\wedge d\omega
\]
in which $d^2\omega=0$ and we replace $d\omega$ by $\Omega-\omega\wedge\omega$, so that
\[ 
  d\Omega=\Omega\wedge\omega-\omega\wedge\Omega,
\]
which becomes the Bianchi identity using the definition of $d_{\omega}$ and the notation \eqref{EqDefCrochwedgedeux}.
\end{proof}
Remark that the Bianchi identity reads $d_{\omega}^2\omega=0$, but that in general $d_{\omega}$ does not square to zero. 

\section{Covariant derivative on associated bundle}
%-----------------------------------------------------

Now we consider a general $G$-principal bundle $\dpt{\pi}{P}{M}$ and an associated bundle $E=P\times_{\rho} V$. We define a product $\eR\times E\to E$ by
\begin{equation}\label{eq:def:REE}
  \lambda[\xi,v]=[\xi,\lambda v].
\end{equation}
It is clear that the equivariant function $\widehat{\lambda \psi}$ defines the section $\lambda\psi$. A \defe{covariant derivative}{covariant!derivative!on associated bundle} is a map 
		\begin{equation}
		\begin{aligned}
			\nabla \colon \cvec(M)\times \Gamma(M,E) &\to \Gamma(M,E)\
			(X,\psi)&\mapsto \nabla_X\psi
		\end{aligned}
	\end{equation}	
such that 
\begin{subequations}
\begin{align}
\nabla_{fX}\psi&=f\nabla_X\psi,   \\
\nabla_X(f\psi)&=(X\cdot f)\psi+f\nabla_X\psi                                      \label{eq:def:der_covii}
\end{align}
\end{subequations}
where products have to be understood by formula \eqref{eq:def:REE}.

\begin{theorem}  
A connection on a principal bundle gives rise to a covariant derivative on any associated bundle by the formula
\begin{equation}
  \widehat{\nabla^E_X\psi}(\xi)=\ovX_{\xi}(\hpsi)
\end{equation}
where $\dpt{\hpsi}{P}{V}$ is the function associated with the section $\dpt{\psi}{M}{E}$. 
\label{tho_dercovassoequiv}
\end{theorem}
We have to prove that it is a good definition: the function $\widehat{\nabla^E_X\psi}$ must define a section $\dpt{\nabla_X^R\psi}{M}{E}$ and the association $\psi\to\nabla^E_X\psi$ must be a covariant derivative.

With the discussion of page \pageref{pg:vecto_vecto} about the application of a tangent vector on a map between manifolds, we have $(d\varphi X)f=X(f\circ\varphi)$. By using this equality in the case of $\ovX$ with $\hpsi$ and $R_g$, we find $(dR_g\ovX)(\hpsi)=\ovX(\hpsi\circ R_g)$ and thus
\[
   \ovX_{\xi\cdot g}(\hpsi)=\ovX_{\xi}(dR_g\hpsi).
\]
We prove the theorem step by step.

\begin{proposition}
The function $\widehat{\nabla_X^E\psi}$ defines a section of $P$.
\end{proposition}

\begin{proof}
We have to see that $\widehat{\nabla_X^E\psi}$ is an equivariant function. The equivariance of $\hpsi$ gives $\hpsi\circ R_g=\rho(g^{-1})\hpsi$, thus
\begin{equation}
\widehat{ \nabla_X^E\psi }(\xi\cdot g)=\ovX_{\xi\cdot g}(\hpsi)\\
                                      =\big( (dR_g)_{\xi}\ovX_{\xi}\big)(\hpsi)\\
				      =\ovX_{\xi}(\hpsi\circ R_g)\\
				      =\ovX_{\xi}( \rho(g^{-1})\hpsi )\\
				      =\rho(g^{-1})\ovX_{\xi}(\hpsi).
\end{equation}
The last equality comes from the fact that the product $\rho(g^{-1})\hpsi$ is a linear product ``matrix times vector''{} and that $\ovX_{\xi}$ is linear.
\end{proof}

\begin{theorem}  
The definition
\[
   \widehat{\nabla^E_X\psi}(\xi)=\ovX_{\xi}(\hpsi)
\]
defines a covariant derivative.
\label{tho_nablaE}
\end{theorem}

\begin{proof}
We have to check the two conditions given on page \pageref{sec:conn_vect}.

\subdem{First condition}
By definition, $\widehat{\nabla_{fX}^E\psi}(\xi)=\overline{fX_{\xi}}(\hpsi)$. Now we prove that 
\begin{equation}\label{eq:fXhpsi}
  \overline{fX_{\xi}}(\hpsi)=(f\circ\pi)(\xi)\ovX_{\xi}(\hpsi).
\end{equation}
 This formula is coherent because $\ovX_{\xi}(\hpsi)\in V$ and $(f\circ\pi)(\xi)\in\eR$. By definition of the horizontal lift, $\overline{fX}_{\xi}$ is the unique vector such that

 \begin{itemize}
 \item $d\pi_{\xi}(\overline{fX}_{\xi})=(fX)_x=f(x)d\pi\ovX_{\xi}=(f\circ\pi)(\xi)d\pi\ovX_{\xi}$,
 \item $\omega_{\xi}(\overline{fX}_{\xi})=0$.
 \end{itemize}
We check that $(f\circ\pi)(\xi)\ovX_{\xi}$ also fulfills these two conditions because $d\pi$ and $\omega$ are $\Cinf(P)$-linear. Equation \eqref{eq:fXhpsi} immediately gives 
\begin{equation}
\widehat{\nabla_{fX}^E\psi}(\xi)=(f\circ\pi)(\xi)\widehat{\nabla_X^E\psi}(\xi).
\end{equation}
Now we show that $\widehat{ f\nabla_X^E\psi }$ is the same. The section $\dpt{f\nabla_X^E\psi}{M}{E}$ is given by  $(f\nabla_X^E\psi)(x)=f(x)(\nabla_X^E\psi)(x)$, and by definition of the associated equivariant function,
\[
  f(x)(\nabla_X^E\psi)(x)=[ \xi,f(x)\widehat{\nabla_X^E\psi}(\xi) ].
\]
Then
\begin{equation}
  \widehat{f\nabla_X^E\psi}(\xi)=f(x)\widehat{\nabla_X^E\psi}(\xi)=(f\circ\pi)(\xi)\widehat{\nabla_X^E\psi}(\xi).
\end{equation}
All this shows that
$  \nabla_{fX}^E\psi=f\nabla_X^E\psi$.
\subdem{Second condition}
This is a computation using the Leibnitz rule:
\begin{equation}
\begin{split}
  \widehat{\nabla_X^E(f\psi)}(\xi)&=\ovX_{\xi}( \widehat{f\psi} )
                                  	\stackrel{(a)}{=}\ovX_{\xi}((\pi\circ f)\hpsi)\\
				  &\stackrel{(b)}{=}\ovX_{\xi}(\pi^*f)\hpsi(\xi)+(\pi^*f)(\xi)\ovX_{\xi}\hpsi
				  =d(f\circ\pi)_{\xi}\ovX_{\xi}\hpsi(\xi)+f\widehat{\nabla_X^E\psi}(x)\\
				  &=df_{\pi(\xi)}d\pi_{\xi}\ovX_{\xi}\hpsi(\xi)+f\widehat{\nabla_X^E\psi}(x)
				  =X_x(f)\hpsi(\xi)+f\widehat{\nabla_X^E\psi}(x)\\
				  &=\widehat{(Xf)\psi}(\xi)+\widehat{f\nabla_X^E\psi}(\xi)
\end{split}
\end{equation}
where (a) is because $\widehat{f\psi}=\pi^*f\hat{\psi}$, and (b) is an application of the Leibnitz rule.
\end{proof}




\begin{theorem}
Using the local coordinates related to the sections $\dpt{\sigma_{\alpha}}{\mU_{\alpha}}{P}$, the covariant derivatives reads:
\begin{equation}\label{eq:nabla_coord}
(\nabla_X\psi)\bsa(x)=X_x\psi\bsa-\rho_*(\sigma_{\alpha}^*\omega_x(X))\psi\bsa(x)
\end{equation}
where $\dpt{\rho_*}{\yG}{\End(V)}$ is defined by
\begin{equation}  \label{eq:def_rho_s}
  \rho_*(A)=\Dsdd{\rho(e^{tA})}{t}{0}
\end{equation}

\end{theorem}

\begin{proof}
The problem reduces to the search of $\ovX$ because
\[
   (\nabla_X\psi)\bsa(x)=\widehat{\nabla_X\psi}(\sigma_{\alpha}(x))=\ovX_{\sigma_{\alpha}(x)}(\hpsi).
\]
We claim that $\ovX_{\sigma_{\alpha}(x)}=d\sigma_{\alpha} X_x-\omega(d\sigma_{\alpha} X_x)^*$. We have to check that $d\pi\ovX=X$ and $\omega(\ovX)=0$. The latter comes easily from the fact that $\omega(A^*)=A$. For the first one, remark that $s_{\alpha}$ is a section, then $d(\pi\circ s_{\alpha})=\id$, and $d\pi(ds_{\alpha} X_x)=X_x$, while
\begin{equation}
  d\pi(A^*\bxi)=d\pi\Dsdd{\xi\cdot e^{-tA}}{t}{0}
               =\Dsdd{\pi(\xi)}{t}{0}
               =0.
\end{equation}
Since the horizontal lift is unique, we deduce
\begin{equation}
  (\nabla_X\psi)\bsa(x)=\big(  d\sigma_{\alpha} X_x-\omega(d\sigma_{\alpha} X_x)^*  \big)\hpsi.
\end{equation}
From the definition of a fundamental vector field,
\begin{equation}
\begin{aligned}
    \omega(d\sigma_{\alpha} X_x)^*_{\sigma_{\alpha}(x)}\hpsi
           &=\Dsdd{\hpsi\big(\sigma_{\alpha}(x)\cdot e^{-t\omega(d\sigma_{\alpha} X_x)}  \big) }{t}{0}\\
           &=\Dsdd{\rho(e^{t\omega(d\sigma_{\alpha} X_x)})\hpsi(\sigma_{\alpha}(x))}{t}{0}&&\text{from \eqref{eq:equiv_psi_b}}\\
           &=(d\rho)_e(\omega\circ d\sigma_{\alpha})X_x(\hpsi\circ\sigma_{\alpha})(x)\\
           &=\rho_*\big( (\sigma_{\alpha}^*\omega)(X_x) \big)\psi\bsa(x)    &&\text{by \eqref{eq:def_rho_s}}
\end{aligned}
\end{equation}

\end{proof}
 
We can express the covariant derivative by means of some maps $\dpt{\theta_{\alpha}}{\cvec(M)\times M}{\End(V)}$ given by
 \begin{equation}
\nabla_X\gamai=\bghd{\theta_{\alpha}(X)}{i}{j}\gamaj.
 \end{equation}
where the $\gamai$'s were given in equation \eqref{eq:def:gamai}. By the definition \eqref{eq:def:der_covii},
\[
\begin{split}
  (\nabla_X\psi)(x)&=(X\cdot s^i_{\alpha})_x\gamai(x)+s^i_{\alpha}(x)(\nabla_X\gamai)(x)\\
                   &=(X\cdot s^i_{\alpha})_x\gamai(x)+s^i_{\alpha}(x)\bghd{\theta_{\alpha}(X)}{i}{j}\gamaj(x).
\end{split}
\]
On the othre hand with the notations of equation \eqref{eq:def:psisa}, $\gamsai=e_i$ and $X_x\gamsai=0$. Then equation \eqref{eq:nabla_coord} gives $\theta_{\alpha}(X)=\rho_*(\sigma_{\alpha}^*\omega_x(X))$, or
\begin{equation}
\theta_{\alpha}=\rho_*(\sigma^*_{\alpha}\omega_x).
\end{equation}

\subsection{Curvature on associated bundle}

From the definition \eqref{eq:def:som_E}, it makes sense to define the curvature $2$-form by
\[
  R(X,Y)\psi=\nabla_X\nabla_Y\psi-\nabla_Y\nabla_X\psi-\nabla_{[X,Y]}\psi.
\]
It is also clear that $\psisa$ defines a section of the trivial vector bundle $F=M\times V$ by $x\to (x,\psisa(x))$, so one can define  $\dpt{\Omega_{\alpha}(X,Y)}{\Gamma(M,E)}{\Gamma(M,E)}$ by
\[
  \big(  R(X,Y)\psi  \big)\bsa=\Omega_{\alpha}(X,Y)\psisa
\]
and take back all the work around Bianchi because of the relation \eqref{eq:nabla_coord} which can be written as $(\nabla_X\psi)\bsa(x)=X_x\psisa+\theta_{\alpha}(X)\psisa(x)$ and which is the same as in proposition \ref{prop:thet_conn_F}.

\subsection{Connection on frame bundle}\index{frame!bundle}
%//////////////////////////////////////

\subsubsection{General framework}
%////////////////////////////

The frame bundle was defined at page \pageref{pg:frame_bundle}. Let $\dptvb{F}{p}{M}$ be a $\eK$-vector bundle with some local trivialization $(\mU_{\alpha},\phi^E_{\alpha})$ and the corresponding transition functions $\dpt{g\bab}{\mU_{\alpha}\cap\mU_{\beta}}{GL(r,\eK)}$. We consider $\dpt{\pi}{P}{M}$, the frame bundle of $F$; it is a $GL(r,\eK)$-principal bundle. Let $\nabla$ be a covariant derivative on $F$ and $\theta_{\alpha}$, the associated matrices $1$-form. The frame bundle is
\[
  P=\bigcup_{x\in M}(\text{frame of $F_x$}).
\]
A connection is a $\yG$-valued $1$-form; in our case it is a map
\[
  \dpt{\omega^{\alpha}\bxi}{T\bxi\big(\pi^{-1}(\mU_{\alpha})\big)}{\gl(r,\eK)}.
\]
We define our connection by, for $g\in GL(r,\eK)$, $x\in \mU_{\alpha}$, $X_x\in T_xM$ and $A\in \gl(r,\eK)$,
\begin{equation}\label{eq:def_omega_frame}
  \omega_{S_{\alpha}(x)\cdot g}^{\alpha}
         \big(   {R_g}_*s_{\alpha}(x)_* X_x + A^*_{S_{\alpha}(x)\cdot g}  \big):=A+\Ad(g^{-1})\theta_{\alpha}(X_x).
\end{equation}
where $\dpt{S_{\alpha}}{\mU_{\alpha}}{P}$ is the section defined by the trivialization $\phi^P_{\alpha}$:
\[
   S_{\alpha}(x)=\{ \ovv_{\alpha}={\phi^E_{\alpha}}^{-1}(x,e_i) \}_{i=1,\ldots,r}.
\]
Since $\theta_{\alpha}(X_x)\in\End(\eK^r)\subset\gl(r,\eK)$, the second term of \eqref{eq:def_omega_frame} makes sense. This formula is a good definition of $\omega$ because of the following lemma:

\begin{lemma}
If $\xi=S_{\alpha}(x)\cdot g$ and $\Sigma\in T\bxi P$, there exists a choice of $A\in\yG$, and $X_x\in T_xM$ such that
\begin{equation}\label{eq:geneSigma}
  \Sigma={R_g}_*{s_{\alpha}(X)}_* X_x+A^*_{S_{\alpha}(x)\cdot g}.
\end{equation}
\end{lemma}

\begin{proof}
If $\xi\in P$ is a basis of $E$ at $y$, there exists only one choice of $x\in M$ and $g\in G$ such that $\xi=S_{\alpha}(x)\cdot g$.

Let us consider a general path $\dpt{c}{\eR}{P}$ under the form $c(t)=s_{\alpha}(x(t))\cdot g(t)$ where $x$ and $g$ are path in $M$ and $GL(r,\eK)$. The frame $c(t)$ is the one of $F_{x(t)}$ obtained by the transformation $g(t)$ from $s_{\alpha}(x(t))$. It is a set of $r$ vectors, and each of them can be written as a combination of the vectors of $s_{\alpha}(x(t))$, so we write
\begin{equation}
  c^i(t)=s_{\alpha}^j(x(t))g_j^i(t)
\end{equation}
where $s_{\alpha}^j(x(t))\in F_{x(t)}$ and $g_j^i(t)\in\eK$. We compute $\Sigma=c'(0)$ by using the Leibnitz rule and we denote $x'(0)=X_x$, $x(0)=x$ and $g^i_j(0)=g^i_j$ (the matrix of $g$):
\begin{equation}
\begin{split}
  \Sigma^i&=\Dsdd{  s^j_{\alpha}(x(t))  }{t}{0}g^i_j+s^j_{\alpha}(x)\Dsdd{g^i_j(t)}{t}{0}\\
          &=(ds_{\alpha}^j)_xX_xg^i_j+{g^i_j}'(0)s^j_{\alpha}(x).
\end{split}
\end{equation}
Going to more compact matrix form, it gives
\[
  \Sigma=(ds_{\alpha})_xX_x\cdot g+s_{\alpha}(x)g'(0).
\]
The second term, $s_{\alpha}^j(x)g'^i_j(0)$, is a general vector tangent to a fiber. So it can be written as a fundamental field $A^*\bxi$.

\end{proof}

\begin{lemma}
On $\mU_{\alpha}\cap\mU_{\beta}$, the form fulfills $\omega^{\alpha}=\omega\hbeta$.
\end{lemma}

\begin{proof}
Let $\dpt{\gamma}{\eR}{M}$ be a path whose derivative is $X_x$. Then
\begin{equation}
\begin{split}
   (R_g)_*s_{\alpha}(x)_*X_x&=\Dsdd{s_{\alpha}(\gamma_t)\cdot g}{t}{0}
                          =\Dsdd{  s_{\beta}(\gamma_t)g_{\alpha\beta}(\gamma_t)\cdot g  }{t}{0}\\
                          &=\Dsdd{ s_{\alpha}(\gamma_t)g_{\alpha\beta}(x)\cdot g }{t}{0}
                          +\Dsdd{ s_{\beta}(x)\cdot g_{\alpha\beta}(\gamma_t)\cdot g }{t}{0}.
\end{split}
\end{equation}
What is in the derivative of the first term is $R_{g_{\alpha\beta}(x)g}(s_{\beta}(\gamma_t))$. Taking the derivative, we find the expected ${R_{g_{\alpha\beta}(x)g}}_*{s_{\beta}}_*X_x$.

For the second term, we note $r:=s_{\beta}(x)\cdot g_{\alpha\beta}(g)g$, and we have to compute the following, using equation \eqref{eq:rdotht},
\begin{equation}
\begin{aligned}
 \Dsdd{  r\cdot\Ad_{g^{-1}}( g_{\alpha\beta}^{-1}(x)&g_{\alpha\beta}(\gamma_t) ) }{t}{0}\\
                  &=\Dsdd{ r\cdot\exp t\big(    (d\AD_{g^{-1}})_e( g_{\alpha\beta}^{-1}(x)(dg_{\alpha\beta})_xX_x )   \big)}{t}{0}\\
                  &=\Dsdd{r\cdot\exp t\big( \Ad_{g^{-1}}g_{\alpha\beta}^{-1}(x)dg_{\alpha\beta}}{t}{0}\\
                  &=\left(  \Ad_{g^{-1}}g^{-1}_{\alpha\beta}(x)dg_{\alpha\beta} X_x    \right)^*_r.
\end{aligned}
\end{equation}
Using this, we can perform the computation:
\begin{equation}
\begin{aligned}
\omega\hbeta_{S_{\alpha}(x)\cdot g}\big(  {R_g}_*{s_{\alpha}(x)}_*X_x+A^*_{S_{\alpha}(x)\cdot g}  \big)
                          &=\omega\hbeta_{S_{\beta}(x)\cdot g_{\alpha\beta}(x)g}\Big(  {R_{g_{\alpha\beta}}(x)g}_*{s_{\beta}(x)}_*X_x \\
                              &\qquad + (  \Ad_{g^{-1}}g^{-1}_{\alpha\beta}(x)dg_{\alpha\beta} X_x  )^*_r +A^* \Big)\\
                          &=\Ad_{(g_{\alpha\beta}(x)g)^{-1}}\theta_{\beta}(X_x)\\
			&\quad+\Ad_{g^{-1}}g_{\alpha\beta}^{-1}(x)dg_{\alpha\beta}(X_x)+A\\
                          &=\Ad_{g^{-1}}\big(  (g_{\alpha\beta}^{-1}\theta_{\beta} g_{\alpha\beta}+g_{\alpha\beta}^{-1} dg_{\alpha\beta})(X_x)  \big)+A\\
                          &=\omega^{\alpha}_{S_{\alpha}(x)g}\big( {R_g}_*{s_{\alpha}(x)}_*X_x+A^A_{S_{\alpha}(x)\cdot g} \big).
\end{aligned}
\end{equation}

\end{proof}

\begin{proposition}
The $\omega$ defined by formula \eqref{eq:def_omega_frame} is a connection $1$-form.
\label{prop_omconfrom}
\end{proposition}

\begin{proof}
The first condition, $\omega(A^*\bxi)=A$, is immediate from the definition. The lemma \ref{lem:dRgAstar} gives the second condition in the case $\Sigma=A^*\bxi$. It remains to be checked that $\omega(dR_g\Sigma)=\Ad(g^{-1})\omega(\Sigma)$ in the case $\Sigma=dR_hds_{\alpha} X_x$. This is obtained using the fact that $\Ad$ is a homomorphism.
\end{proof}

\subsubsection{Levi-Civita connection}\label{subsubsec_levi}
%---------------------------------------------

Let $(M,g)$ be a Riemannian manifold. We look at a connection $1$-form $\alpha\in\Omega^1(\SO(M),so(\eR^m))$ on $\SO(M)$, and we define a covariant derivative $\dpt{\nabla^{\alpha}}{\cvec(M)\times T(M)}{T(M)}$, where $T(M)$ is the tensor bundle on $M$ by (cf. theorem \eqref{tho_nablaE})
\begin{eqnarray}\label{r2804e1}
 \widehat{\nabla^{\alpha}_X s}=\overline{X}\hat{s},
\end{eqnarray}
for any $s\in T(M)$.  Our purpose now is to prove that an automatic property of this connection is $\nabla^{\alpha} g=0$. The unique such connection which is torsion-free is the \defe{Levi-Civita}{Levi-Civita connection} one.

The metric $g$ is a section of the tensor bundle $T^*M\otimes T^*M$. So we have, in order to find $\hg$ and to use equation \eqref{r2804e1}, to see $T^*M\otimes T^*M$ as an associated bundle. As done in \ref{equivvec}, we see that
\[
 T^*M\otimes T^*M\simeq \SO(M)\times_{\rho}(V^*\otimes V^*),
\]
with the following definitions:
\begin{itemize}
\item The isomorphism is given by $\psi[b,\alpha\otimes\beta](X\otimes Y)=\alpha(b^{-1} X)\beta(b^{-1} Y)$,
\item $\rho(A)\alpha=\alpha\circ A$,
\item $b\cdot A=b\circ A$.
\end{itemize}
Here, $V=\eR^m$; $\dpt{b}{V}{T_xM}$; $\alpha,\beta\in V^*$; $X$, $Y\in T_xM$ and $A\in \SO(m)$ is seen as $\dpt{A}{V}{V}$.

The following shows that $\psi$ is well defined:
\begin{equation}
\begin{split}
 \psi[b\cdot A,\rho(A^{-1})\alpha\otimes\beta](X\otimes Y)&=(\alpha\circ A)
                            (A^{-1}\circ b^{-1} X)(\beta\circ A)(A^{-1}\circ b^{-1} Y)\\
                                           &=\psi[b,\alpha\otimes\beta](X\otimes Y)
\end{split}
\end{equation}

\begin{proposition}
The function $\hg$ is given by
\[
 \hg(b)(v\otimes w)=g_x(b(v)\otimes b(w))=v\cdot w.
\]
\end{proposition}

\begin{proof}
The second equality is just the fact that $\dpt{b}{(\eR^m,\cdot)}{(T_xM,g_x)}$ is isometric. On the other hand, if $\hg(b)=\alpha\otimes\beta$, we have:
\begin{equation}
\begin{split}
 g_x(X\otimes Y)&=\psi[b,\alpha\otimes\beta](X\otimes Y)
                =\alpha(b^{-1} X)\beta(b^{-1} Y)\\
                &=\alpha\otimes\beta(b^{-1} X\otimes b^{-1} Y)
                =\hg(b)(b^{-1} X\otimes b^{-1} Y).
\end{split}
\end{equation}

Since $b$ is bijective, $X$ and $Y$ can be written as $bv$ and $bw$ respectively for some $v$, $w\in V$, so that
\[g_x(bv\otimes bw)=\hg(b)v\otimes w.\]
\end{proof}

It is now easy to see that $\oX\hg=0$. As $\hg$ takes its values in $V^*\otimes V^*$, $\oX\hg$ belongs to this space and can be applied on $v\otimes w\in V\otimes V$. Let $\oX(t)$ be a path in $\SO(M)$ which defines $\oX$; if $\oX\in T_b\SO(M)$, $\oX(0)=b$. We have
\begin{equation}
 \oX\hg(v\otimes w)	=\dsdd{\hg(\oX(t))v\otimes w}{t}{0}
			=\dsdd{v\cdot w}{t}{0},
\end{equation}
which is obviously zero.

\subsection{Holonomy}
%--------------------

Let the principal bundle 
\begin{equation}
\xymatrix{%
   G \ar@{~>}[r]		&	P\ar[d]^{\pi}\\
   				&	M
 }
\end{equation}
  and $\omega$ a connection on $G$. Let $\gamma\colon [0,1]\to M$, a closed curve piecewise smooth; $\gamma(0)=\gamma(1)=x$. For each $p\in\pi^{-1}(x)$, there exists one and only one horizontal lift $\tilde\gamma\colon [0,1]\to P$ such that $\tilde\gamma(0)=p$. There exists of course an element $g\in G$ such that $\tilde\gamma(1)=p\cdot g$.

We define the following equivariance relation on $P$: we say that $p\sim q$ if and only if $p$ and $q$ can be joined by a piecewise smooth path. The \defe{holonomy group}{holonomy group} at the point $p$ is 
\[ 
  \Hol_p(\omega)=\{ g\in G\tq p\sim p\cdot g \}.
\]
\subsection{Connection and gauge transformation}
%-----------------------------------------------


\begin{proposition}
If $\omega$ is a connection on a $G$-principal bundle and $\varphi$, a gauge transformation, the form $\beta=\varphi^*\omega$ is a connection $1$-form too.
\label{prop:vp_conn}
\end{proposition}

\begin{proof}
It is rather easy to see that $\varphi_*A^*\bxi=A^*_{\varphi(x)}$:
\[
  \varphi_*A^*\bxi=\Dsdd{ \varphi(\xi e^{-tA})  }{t}{0}=\Dsdd{ \varphi(\xi)e^{-tA}  }{t}{0}=A^*_{\varphi(x)}.
\]
The same kind of reasoning leads to $\varphi_*{R_g}_*={R_g}_*\varphi_*$. From here, it is easy to see that
\[
  (\varphi^*\omega)\bxi(A^*\bxi)=\omega_{\varphi(\xi)}(\varphi_*A^*\bxi)=A,
\]
and
\[
\big( R^*_g(\varphi^*\omega)\bxi \big)(\Sigma)
                =(R_g^*\omega)_{\varphi(\xi)}(\varphi_*\Omega)=\Ad(g^{-1})\big( (\varphi^*\omega)\bxi(\Sigma) \big).
\]
\end{proof}
So, the ``gauge transformed'' of a connection is still a connection. It is hopeful in order to define gauge invariants objects (Lagrangian) from connections (electromagnetic fields).

\subsubsection{Local description}

Let $\dpt{\pi}{P}{M}$ be a $G$-principal bundle given with some trivializations $\dpt{\phi^P_{\alpha}}{\pi^{-1}(\mU_{\alpha})}{\mU_{\alpha}\times G}$ over $\mU_{\alpha}\subset M$ and $\dpt{s_{\alpha}}{\mU_{\alpha}}{\pi^{-1}(\mU_{\alpha})}$, a section. In front of that, we consider an associated bundle $\dpt{p}{E=P\times_{\rho} V}{M}$ with a trivialization $\dpt{\phi^E_{\alpha}}{E}{\mU_{\alpha}\times V}$. One can choose a section $s_{\alpha}$ compatible with the trivialization in the sense that $\phi^P_{\alpha}(s_{\alpha}(x)\cdot g)=(x,g)$; the same can be done with $E$ by choosing $\phi^E_{\alpha}([s_{\alpha}(x),v])=(x,v)$.

A section\index{section!local description} $\dpt{\psi}{M}{E}$ is described by a function $\dpt{\psi_{\alpha}}{\mU_{\alpha}}{V}$ defined by $\phi^E_{\alpha}(\psi(x))=(x,\psi_{\alpha}(x))$.  In the inverse sense, $\psi$ is defined (on $\mU_{\alpha}$) from $\psi_{\alpha}$ by
$\psi(x)=[s_{\alpha}(x),\psi_{\alpha}(x)]$.
In the same way, a gauge transformation\index{gauge!transformation!local description} $\dpt{\varphi}{P}{P}$ is described by functions $\dpt{\tilde{\varphi}_{\alpha}}{\mU_{\alpha}}{G}$,
\begin{equation}
  \varphi(s_{\alpha}(x))=s_{\alpha}(x)\cdot\tilde{\varphi}_{\alpha}(x).
\end{equation}
The function $\tilde{\varphi}_{\alpha}$ also fulfil
\begin{equation}
  (\phi_{\alpha}^P\circ\varphi\circ{\phi_{\alpha}^P}^{-1})(x,g)=(x,\tilde{\varphi}(x)\cdot g) 
\end{equation}
because
\begin{equation}
\begin{split}
  (\phi_{\alpha}^P\circ\varphi\circ{\phi_{\alpha}^P}^{-1})(x,g)&=(\phi_{\alpha}^P\circ\varphi)(s_{\alpha}(x)\cdot g)\\
                                                      &=\phi_{\alpha}^P( \varphi(s_{\alpha}(x))\cdot g )\\
                                                      &=\phi_{\alpha}^P( s_{\alpha}(c)\cdot\tilde{\varphi}_{\alpha}(x)g)\\
                                                      &=(x,\tilde{\varphi}_{\alpha}(x)g).
\end{split}
\end{equation}

We know that a connection on $P$ is given by its $1$-form $\omega$. Moreover we have the following:
\begin{proposition}
A connection on $P$ is completely determined on $\pi^{-1}(\mU_{\alpha})$ from the data of the $\yG$-valued $1$-form $\sigma_{\alpha}^*\omega$ on $\mU_{\alpha}$.
\end{proposition}

\begin{proof}
We consider a $1$-form $\omega$ which fulfils the two conditions of page \pageref{pg:def:conne}. Our purpose is to find back $\omega\bxi(\Sigma)$, $\forall\xi\in P,\Sigma\in T\bxi P$ from the data of $\sigma_{\alpha}^*\omega$ alone. For any $\xi$, there exists a $g$ such that $\xi=\sigma_{\alpha}(x)\cdot g$. We have
\begin{equation}
  \Ad_{g^{-1}}(\omega_{\sigma_{\alpha}(x)\Sigma})=(R^*_g\omega)_{\sigma_{\alpha}(x)}(\Sigma)
             =\omega_{\sigma_{\alpha}(x)\cdot g}\big( (dR_g)_{\sigma_{\alpha}(x)}\Sigma \big).
\end{equation}
If we know $s_{\alpha}^*\omega$, then we know $\omega\big(  (ds_{\alpha})_xv  \big)$ for any $v\in T_xM$. So 
\[
   \omega_{\sigma_{\alpha}(x)\cdot g}\big( (dR_g)_{\sigma_{\alpha}(x)}\Sigma\big)
\]
is given from $\sigma_{\alpha}^*\omega$ for every $\Sigma$ of the form $\Sigma=(d\sigma_{\alpha})_xv$. From the form \eqref{eq:geneSigma} of a vector in $T\bxi P$, it just remains to express $\omega_{\sigma_{\alpha}(x)\cdot g}(A^*_{\sigma_{\alpha}(x)\cdot g})$ in terms of $s_{\alpha}^*$. The definition of a connection makes that it is simply $A$.

\end{proof}

\subsubsection{Covariant derivative}

If we have a connection on $P$, we can define a covariant derivative on the associated bundle $E$ by 
\[
  (\nabla_X\psi)\bsa(x)=X_x(\psi_{\alpha})+\rho_*( s_{\alpha}^*\omega_x(X) )\psi\bsa(x),
\]
the matricial $1$-form being given by $\theta_{\alpha}=\rho_*\sigma^*_{\alpha}\omega$. The gauge transformation $\varphi$ acts on the connection $\omega$ by defining $\omega^{\varphi}:=\varphi^*\omega$.

\begin{proposition}
If $\beta=\varphi^*\omega$, then
\[
   s^*_{\alpha}(\beta)=\Ad_{\tilde{\varphi}_{\alpha}(x)^{-1}}s^*_{\alpha}(\omega)+\tilde{\varphi}_{\alpha}(x)^{-1} d\tilde{\varphi}_{\alpha}.
\]
\end{proposition}

\begin{proof}
Let $\dpt{\gamma}{\eR}{M}$ be a path such that $\gamma(0)=x$ and $\gamma'(0)=X_x$. We have to compute the following:
\begin{equation}\label{eq:ppu}
  (s_{\alpha}^*\beta)(X_x)=(s_{\alpha}^*\varphi^*\omega)(X_x)=\omega_{(\varphi\circ s_{\alpha})(x)}\big(  (\varphi\circ s_{\alpha})_*X_x  \big).
\end{equation}
What lies in the derivative is:
\begin{equation}
\begin{split}
  (\varphi\circ s_{\alpha})_*(X_x)&=\Dsdd{ (\varphi\circ s_{\alpha}\circ\gamma)(t) }{t}{0}\\
                            &=\Dsdd{  s_{\alpha}(\gamma(t))\cdot\tilde{\varphi}_{\alpha}(\gamma(t))  }{t}{0}\\
                            &=\Dsdd{ s_{\alpha}(\gamma(t))\cdot\tilde{\varphi}_{\alpha}(\gamma(0)) }{t}{0}
                             +\Dsdd{ s_{\alpha}(\gamma(0))\cdot\tilde{\varphi}_{\alpha}(\gamma(t)) }{t}{0}\\
                            &={R_{\tilde{\varphi}_{\alpha}(x)}}_*{s_{\alpha}}_*X_x
                             +\Dsdd{  s_{\alpha}(x)\cdot\tilde{\varphi}_{\alpha}(x)e^{ t\tilde{\varphi}_{\alpha}(x)^{-1}(d\tilde{\varphi}_{\alpha})_x\gamma'(0)}}{t}{0}.
\end{split}
\end{equation}
A justification of the remplacement $\tilde{\varphi}_{\alpha}(\gamma(t))\to \tilde{\varphi}_{\alpha}(x)e^{t\tilde{\varphi}_{\alpha}(x)^{-1}(d\tilde{\varphi}_{\alpha})_x\gamma'(0)}$ is given in the corresponding proof at page \pageref{pg:justif_s}.
If we put this expression into equation \eqref{eq:ppu}, the first term becomes
\[
\begin{split}
   \omega_{(\varphi\circ s_{\alpha})(x)}\big(  {{R_{\tilde{\varphi}_{\alpha}(x)}}_*{s_{\alpha}}_*X_x}   \big)
           &=(R^*_{\tilde{\varphi}_{\alpha}(x)}\omega)_{s_{\alpha}(x)}({s_{\alpha}}_*X_x)\\
           &=\Ad_{\tilde{\varphi}_{\alpha}(x)^{-1}} \big(\omega_{s_{\alpha}(x)} ( {s_{\alpha}}_*X_x ) \big)\\
           &=\Ad_{\tilde{\varphi}_{\alpha}(x)^{-1}}  (s_{\alpha}^*\omega)(X_x).
\end{split}
\]
The second term is the case of a connection applied to a fundamental vector field.

\end{proof}

\section{Product of principal bundle}\label{sec:produit_bundle}
%++++++++++++++++++++++++++++++++++++

In this section, we build a $G_1\times G_2$-principal bundle from the data of a $G_1$ and a $G_2$-principal bundle. The physical motivation is clear: as far as electromagnetism is concerned, particles are sections of $U(1)$-principal bundle while the relativistic invariance must be expressed by means of a $\SLdc$-associated bundle. So the physical fields must be sections of something as the product of the two bundles. See subsection \ref{subsec:incl_Lorentz}.

\subsection{Putting together principal bundle}
%------------------------------------

Let us consider two principal bundle over the same base space
\[
\xymatrix{
    G_1  \ar@{~>}[r] & P_1 \ar[r]^{p_1}& M,}
\]
and
\[
\xymatrix{
    G_2  \ar@{~>}[r] & P_2 \ar[r]^{p_2}& M.
  }
\]
First we define the set
\begin{equation}
  P_1\circ P_2=\{   (\xi_1,\xi_2)\in P_1\times P_2\tq p_1(\xi_1)=p_2(\xi_2)    \}
\end{equation}
which will be the total space of our new bundle. The projection $\dpt{p}{P_1\circ P_2}{M}$ is naturally defined by
\[
  p(\xi_1,\xi_2)=p_1(\xi_1)=p_2(\xi_2),
\]
while the right action of $G_1\times G_2$ on $P_1\circ P_2$ is given by
\[
  (\xi_1,\xi_2)\cdot(g_1,g_2)=(\xi_1\cdot g_1,\xi_2\cdot g_2)
\]
With all these definitions,
\[
\xymatrix{
    G_1\times G_2  \ar@{~>}[r] & P_1\circ P_2 \ar[d]^{p}\\& M\\
  }
\]
is a $G_1\times G_2$-principal bundle over $M$. We define the natural projections 
		\begin{equation}
		\begin{aligned}
			\pi_i \colon P_1\times P_2 &\to P_i\
			(\xi_1, \xi_2)&\mapsto \xi_i,
		\end{aligned}
	\end{equation}	
%
and if $e_i$ denotes the identity element of $G_i$, we can identify $G_1$ to $G_1\times \{e_2\}$ and $G_2$ to $G_2\times \{e_1\}$; in the same way, $\yG_1=\yG_1\times\{0\}\subset\yG_1\times\yG_2$. So we get the following principal bundles:
\[
\xymatrix{
    G_2  \ar@{~>}[r] & P_1\circ P_2 \ar[r]^{\pi_1}& P_1\\
    G_1  \ar@{~>}[r] & P_1\circ P_2 \ar[r]^{\pi_2}& P_2.
  }
\]
It is clear that the following diagram commutes:
\[
\xymatrix{
    P_1  \ar[rd]_{P_1} & P_1\circ P_2\ar[r]^{\pi_2} \ar[l]_{\pi_1}\ar[d]_p& P_2 \ar[ld]^{p_2}\\
    &M
  }
\]

\subsection{Connections}
%-----------------------

Let $\omega_i$ be a connection on the bundle $\dpt{p_i}{P_i}{M}$. Using the identifications, $\pi_1^*\omega_1$ is a connection on $\dpt{\pi_2}{P_1\circ P_2}{P_2}$ (the same is true for $1\leftrightarrow 2$), and $\pi_1^*\omega_1\oplus\pi_2^*\omega_2$ is a connection on $\dpt{p}{P_1\circ P_2}{M}$. Let us prove the first claim.

Let $A\in\yG_1$. We first have to prove that $\pi_1^*\omega_1(A^*)=A$. For this, remark that $A=(A,0)\in\yG_1\oplus\yG_2$ and
\begin{equation}
   A^*\bxi=\Dsdd{\xi\cdot e^{-t(A,0)}}{t}{0}
          =\Dsdd{(\xi_1,\xi_2)\cdot(e^{-tA},e_2)}{t}{0}
          =\Dsdd{(\xi_1\cdot e^{-tA},\xi_2)}{t}{0},
\end{equation}
so $d\pi_1A^*=\Dsdd{\pi_1(\ldots)}{t}{0}=\omega_1(A^*)=A$. Let now $\Sigma\in T_{(\xi_1,\xi_2)}(P_1\circ P_2)$ be given by the path $(\xi_1(t),\xi_2(t))$. In this case we have
\begin{equation}
\begin{split}
   \big(  R^*_{(g,e_2)\pi_1^*\omega_1}  \big)_{(\xi_1,\xi_2)}\Sigma&=
      (\pi_1^*\omega_1)(dR_{(g,e_2)}\Sigma)\\
     &=\omega_1( \Dsdd{\xi_1(t)\cdot g}{t}{0}  )\\
     &=\omega_1( dR_g\Dsdd{\xi_1(t)}{t}{0} )\\
     &=\Ad(g^{-1})\pi_1^*\omega_1( \Dsdd{( \xi_1(t),\xi_2(t) )}{t}{0} )\\
     &=\Ad(g^{-1})\pi_1^*\omega_1\Sigma.
\end{split}
\end{equation}

\subsection{Representations}
%----------------------------

Let $V$ be a vector space and $\dpt{\rho_i}{G_i}{GL(V)}$ be some representations such that 
\begin{equation}\label{eq:cond_reprez}
   [\rho_1(g_1),\rho_2(g_2)]=0
\end{equation}
for all $g_1\in G_1$ and $g_2\in G_2$ (in the sense of commutators of matrices). In this case, one can define the representation $\dpt{\rho_1\times\rho_2}{G_1\times G_2}{GL(V)}$ by
\begin{equation}
   (\rho_1\times\rho_2)(g_1,g_2)=\rho_1(g_1)\circ\rho_2(g_2)=\rho_2(g_2)\circ\rho_1(g_1).
\end{equation}
The relation \eqref{eq:cond_reprez} is needed in order for $\rho_1\times\rho_2$ to be a representation, as one can check by writing down explicitly the requirement
\[
  (\rho_1\times \rho_2)\big(  (g_1,g_2)(g'_1,g'_2)  \big)=(\rho_1\times \rho_2)(g_1g'_1,g_2g'_2) 
\]

%+++++++++++++++++++++++++++++++++++++++++++++++++++++++++++++++++++++++++++++++++++++++++++++++++++++++++++++++++++++++++++
\section{Hodge decomposition theorem and harmonic forms}
%+++++++++++++++++++++++++++++++++++++++++++++++++++++++++++++++++++++++++++++++++++++++++++++++++++++++++++++++++++++++++++

Among other sources for Hodge decomposition and harmonic forms, we have \cite{JohnsonHodge,CohoHarBound,UndergradDeRham}. Some parts of the wikipedia article \wikipedia{en}{Hodge_dual}{Hodge\_dual} are interesting.

Let $E$ be an oriented Euclidian space of dimension $m=2n$. We define the operation $*$ by
\begin{equation}		\label{EqGradWedge}
	\begin{aligned}
		*\colon \Wedge E&\to \Wedge E \\
		e_{i_1}\wedge e_{i_2}\wedge\cdots\wedge e_{i_k}&\mapsto e_{i_{k+1}}\wedge\cdots\wedge e_{i_m} 
	\end{aligned}
\end{equation}
when $\{ e_i \}$ is an oriented basis of $E$ and $\{ i_k \}$ is an even permutation of $\{ 1,2,\cdots,m \}$. If it is impossible to build an even permutation, then we add a minus sign. We have $**\omega=(-1)^{p(m-p)}\omega$ belongs to $\omega\in\Wedge^pE$.
\begin{example}
	If we consider the space $\eR^4$ with the coordinates $(x,y,z,t)$,
	\begin{equation}
		*(dx\wedge dz\wedge dt)=dy
	\end{equation}
	because $(x,z,t,y)$ is an even permutation of $(x,y,z,t)$. Now, $*dy=dz\wedge dx\wedge dt=-(dx\wedge dz\wedge dt)$ because $(y,z,x,t)$ is an even permutation of $(x,y,z,t)$.

	More generally, if we have a differential $p$-form $\omega$ on a $m$ dimensional space, we have
	\begin{equation}
		*(e_{\sigma(1)}\wedge e_{\sigma(2)}\wedge\ldots\wedge e_{\sigma(p)})=e_{\sigma(p+1)}\wedge \ldots\wedge e_{\sigma(m)}
	\end{equation}
	In order to compute $*(e_{\sigma(p+1)}\wedge \ldots\wedge e_{\sigma(m)})$, we need a permutation of $\big( \sigma(1),\ldots, \sigma(m)\big)$ which \emph{begins} by $\sigma(p+1)\ldots\sigma(m)$. This reduces to permute the $m-p$ elements $\sigma(p+1),\ldots,\sigma(m)$ with the $p$ first elements. Thus we have
	\begin{equation}
		**\omega=(-1)^{p(m-p)}\omega.
	\end{equation}
\end{example}

Let $V$ be a compact, oriented manifold. Each of the spaces of sections $ C^{\infty}\big( V, \Wedge^k_{\eC}(T^*V)\big)$ is endowed with a $2$-form
\begin{equation}		\label{EqProdWedgeHOfge}
	\langle \omega_1, \omega_2\rangle =\int_V\omega_1\wedge *\omega_2.
\end{equation}

\begin{lemma}
	The \defe{codifferential}{codifferential} $\delta$ defined by
	\begin{equation}
		\begin{aligned}[]
			\delta	\colon \Wedge^k_{\eC}(T^*V)&\to \Wedge_{\eC}^{k-1}(T^*V)\\
			\beta				&\mapsto (-1)^{mk+m+1}*d*\beta
		\end{aligned}
	\end{equation}
	is a formal adjoint of $d$ for the product \eqref{EqProdWedgeHOfge}.
\end{lemma}

\begin{proof}
	If $\beta\in\Wedge_{\eC}^k(T^*V)$ and $\alpha\in\Wedge_{\eC}^{k-1}(T^*V)$, we have
	\begin{equation}
		\begin{aligned}[]
			*\delta\beta&=(-1)^{mk+m+1}* *\big( d*\beta \big)\\
			&=(-1)^{mk+m+1}(-1)^{(m-k+1)(m-m+k-1)}d*\beta\\
			&=(-1)^kd*\beta.
		\end{aligned}
	\end{equation}
	Using that formula we find
	\begin{equation}
		\begin{aligned}[]
			\langle d\alpha, \beta\rangle -\langle \alpha, \delta\beta\rangle &=\int_V d\alpha\wedge *\beta-\alpha\wedge *\delta\beta\\
			&=\int_Vd\alpha\wedge *\beta-(-1)^k\alpha\wedge d*\beta\\
			&=\int_Vd\alpha\wedge *\beta+(-1)^{k+1}\alpha\wedge d*\beta\\
			&=\int_Vd(\alpha\wedge *\beta)\\
			&=\int_{\partial V}\alpha\wedge *\beta\\
			&=0.
		\end{aligned}
	\end{equation}
	This proves that $\langle d\alpha, \beta\rangle =\langle \alpha, \delta\beta\rangle$.
	\begin{probleme}
		I do not understand why the integral in the boundary is zero. 
	\end{probleme}	
\end{proof}

Now we define the \defe{Laplace-Beltrami operator}{Laplace-Beltrami operator}\index{operator!Laplace-Beltrami} by\nomenclature[D]{$\Delta$}{Laplace-Beltrami operator}
\begin{equation}
	\Delta=\delta d+d\delta
\end{equation}
and the space of \defe{harmonic forms}{harmonic form}
\begin{equation}
	H^k=\{ \omega\in\Omega^k\tq\Delta\omega=0 \}.
\end{equation}

\begin{lemma}
	If $M$ is a closed manifold, a $k$-form is harmonic if and only if $d\omega=\delta\omega=0$.
\end{lemma}

\begin{proof}
	No proof.
\end{proof}

\begin{theorem}[Hodge decomposition theorem]\index{theorem!Hodge decomposition}
	For every integer $0\leq k\leq m$, the space $H^p$ is finite dimensional and $\Omega^k(M)$ has the orthogonal decomposition
	\begin{equation}
		\Omega^k(M)=H^k\oplus\Delta\big( \Omega^k(M) \big),
	\end{equation}
	i.e. the space splits into the kernel of $\Delta$ and its image.
\end{theorem}

\begin{theorem}
	Let $M$ be a compact orientable manifold of dimension $m$. Any exterior differential $k$-form can be written as a unique sum of an exact form, a coexact form and an harmonic form :
	\begin{equation}
		\omega=d\alpha+\delta\beta+\gamma.
	\end{equation}
	with $\omega\in\Omega^k(M)$, $\alpha\in\Omega^{k-1}(M)$, $\beta\in\Omega^{k+1}(M)$ and $\gamma\in\Omega^k(M)$ harmonic.
\end{theorem}

The operator $\Delta$ commutes with the differential $d$ and we have $d\Delta=\Delta d= d\delta d$ since
\begin{equation}
	d\Delta \omega=dd\delta\omega+d\delta d\omega=d\delta d\omega,
\end{equation}
because $d^2=0$, while
\begin{equation}
	\Delta d\omega=d\delta d\omega+\Delta d d\omega=d\delta d\omega.
\end{equation}

\begin{lemma}
	On a close manifold, $\Delta\omega=0$ if and only if $\delta\omega=d\omega=0$.
\end{lemma}

In the case of a closed manifold, a form is harmonic if and only if is belongs to the kernel of $d+\delta$. Moreover, a form in the image of $d+\delta$ is orthogonal to the harmonic forms :
\begin{equation}
	\langle d\alpha^{k-1}+\delta\beta^{k+1}, \gamma\rangle =0
\end{equation}
whenever $\gamma$ is harmonic on a closed manifold.


%\chapter{Examples of groups and representations}        \label{ChapThoComsGroupes}
%\input{iwasawa}

\chapter{From Clifford algebras to Dirac operator}
Bibliography for Clifford algebras, spin group and related topics are \cite{memP,Michelson,Witkowski,mellor,ResEtaDiracType}. More algebraic point of view  can be found in \cite{Fult,Chevalley}. More details about ``square rooting'' second order differential operators are in \cite{Bronn}. For physical concerns, the reader should refer to \cite{Weinberg,Peskin,schwabl}. 

\section{Invitation : Clifford algebra in quantum field theory}\index{quantum!field theory}
%++++++++++++++++++++++++++++++++++++++++++++++++

\label{Secqft}
\subsection{Schrödinger, Klein-Gordon and Dirac}
%----------------------------------

The origin of the Klein-Gordon equation is almost the same as the one of the Schrödinger: one replace physical functions by operators. For a free particle, the correspondence are
\begin{align*}
 \textrm{energy}&& E&\rightarrow i\hbar\dsd{}{t},\\
 \textrm{momentum}&& \overline{p}&\rightarrow -i\hbar\overline{ \nabla }.
\end{align*}
The Schrödinger equation\index{equation!Schrödinger} (which is the non relativistic quantum wave equation) comes from replacement in the non non relativistic expression of the Hamiltonian
\[
E=\frac{\overline{p}^2}{2m}\longrightarrow\left(\partial_t-\frac{i\hbar}{2m}\nabla^2\right)\psi=0,
\]
while the Klein-Gordon\index{equation!Klein-Gordon} one (which is the relativistic quantum wave equation) comes from the relativistic corresponding equation:
\[
E^2=\overline{p}^2c^2+m^2c^4\longrightarrow\left(\partial\hmu\partial_{\mu}+(\frac{mc}{\hbar})^2\right)\psi=0.
\]

This is a second order differential equation; there are however no ``law of nature''{} which forbid a first order equation. We try
\[
 i\hbar\dsd{\psi}{t}
 =\left(\frac{\hbar c}{i}\alpha^k\partial_k+\beta mc^2\right)\psi\equiv \hat{H}\psi.
\]

There are some physical constraints on the coefficients $\alpha^k$ and $\beta$. We will study one of them: we want the components of $\psi$ to satisfy the Klein-Gordon equation, so that the plane waves fulfill the fundamental relation $E^2=p^2c^2+m^2c^4$.

In order to see the implications of this constraint on the coefficients, we apply two times the operator $\hat{H}$, and we compare the result with the Klein-Gordon equation. We find:
\begin{subequations}
\begin{align}
 \alpha^i\alpha^j+\alpha^j\alpha^i&=2\delta^{ij}\mtu,\\
 \alpha^i\beta+\beta\alpha^i&=0,\\
 (\alpha^i)^2=\beta^2&=\mtu.
\end{align}
\end{subequations}
%
If we define $\gamma^0=\beta$ and $\gamma^i=\beta\alpha^i$, we find that the matrices $\gamma^{\mu}$ have to give a representation of the Clifford algebra\footnote{Don't be afraid with the extra minus sign: the quantum field theory is most written with the metric $(+,-,-,-)$ instead of $(-,+,+,+)$.}:
\begin{equation}\label{cliffphys}
	\gamma\hmu\gamma\hnu+\gamma\hnu\gamma\hmu=2\eta^{\mu\nu}\mtu.
\end{equation}
The Dirac equation reads
\[
\left(-i\gamma\hmu\partial_{\mu}+\frac{mc}{\hbar}\right)\psi=0.
\]
If we want to perform some computation with the quantum field theory, we need an explicit form for the $\gamma$'s; that's the reason why we study representations of the Clifford algebra. The \defe{Dirac operator}{dirac!operator} $\Dir$ is the operator which lies in the Dirac equation:
\begin{equation}
 \label{dirflat}\Dir=\sum_{\mu=0}^3\gamma\hmu\dsd{}{x\hmu}.
\end{equation}

\subsection{Lorentz algebra}
%---------------------------

There is an other physical reason (which is in fact the same, but differently presented) justifying the study of the Clifford algebra. The quantum field theory need representation of the Lorentz algebra\footnote{When one think to real infinitesimal rotation matrices, the presence of $i$ seems not natural, but one redefines $J\to iJ$ for formalism reasons.}\index{lorentz!algebra}
\[
 [J^{\mu\nu},J^{\rho\sigma}]=i(\eta^{\nu\rho}J^{\mu\sigma}-\eta^{\mu\rho}J^{\nu\sigma}
 -\eta^{\nu\sigma}J^{\mu\rho}+\eta^{\mu\sigma}J^{\nu\rho}).
\]
A proof of these relations is given in lemma \ref{LemCommsopqAlg}. Dirac had a trick to find such $J$ matrices from a representation of the Clifford algebra. If we have $n\times n$ matrices $\gamma_{\mu}$ such that
\[
	\gamma^{\mu}\gamma^{\nu}+\gamma^{\nu}\gamma^{\mu}=2\eta^{\mu\nu}\mtu_{n\times n},
\]
a $n$-dimensional representation of the Lorenz algebra is obtained by
\[
	S^{\mu\nu}=\frac{i}{4}\left[\gamma^{\mu},\gamma^{\nu}\right].
\]

\begin{lemma}		\label{LemCommsopqAlg}
The matrices of $\so(p,q)$ satisfy the definition relation
\begin{equation}
	M^t\eta+\eta M=0,
\end{equation}
and if $M^{ab}$ is the ``rotation'' in the place of directions $a$ and $b$ (i.e. a trigonometric or an hyperbolic rotation following that $a$ and $b$ are of the same type or not), then the action on $\eR^{(p,q)}$ is given by $(x')^{\mu}=(M^{ab})^{\mu}_{\nu}x^{\nu}$ with
\begin{equation}
	(M^{ab})^{\mu}_{\nu}=\eta^{a\mu}\delta^b_{\nu}-\eta^{b\mu}\delta^a_{\nu}.
\end{equation}
The commutation relations are given by
\begin{equation}
	[M^{ab},M^{cd}]=-\eta^{ac}M^{bd}+\eta^{ad}M^{bc}+\eta^{bc}M^{ad}-\eta^{bd}M^{ac}.
\end{equation}
Notice that $M^{ab}=-M^{ba}$. 
\end{lemma}
See section 12.5 of \cite{Schomblond_em}. By a simple redefinition $J=iM$, one obtains 
\begin{equation}			\label{EqJJietaJcomm}
	[J,J]=i\eta J
\end{equation}
instead of $[M,M]=\eta M$, and the matrices $J$ are Hermitian. Here $\eta$ is the matrix $\eta=diag(\underbrace{+,\ldots,+}_{\text{$p$ times}},\underbrace{-,\ldots,-}_{\text{$q$ times}})$. As convention, we say that a direction corresponding to a \emph{positive} entry in the metric is a \emph{time} direction, while the spatial directions are negative. That corresponds to the convention of page \pageref{PgDefsGenre} to say that a \emph{time-like} vector has positive norm.

\section{Clifford algebra}
%++++++++++++++++++++++++++

\subsection{Definition and universal problem}
%------------------------------------------------------


\begin{definition}
Let $V$ be a (finite dimensional) vector space and $q$, a bilinear quadratic form over $V$. The \defe{Clifford algebra}{Clifford!algebra}\index{algebra!Clifford} $\Cliff(V,q)$ is the unital associative algebra generated by $V$ subject to the relation
\begin{equation}\label{501r1}
       v\cdot v=q(v)
\end{equation}
for all $v$ in $\Cliff(V,q)$. Here the dot denotes the algebra product and $q(v)$ means $q(v,v)$.
\end{definition}
Theorem \ref{tho_Cliffunif} proves unicity of such an algebra, so that it makes sense.

\begin{remark}
The relation \eqref{501r1} is no more a restriction for the elements in $\Cliff(V,q)$ than a restriction on the choice of the algebra product.
\end{remark}


\begin{theorem}\index{Clifford!algebra!Universal property of}
Let $E$ be an unital associative algebra and $\dpt{j}{V}{E}$ a linear map such that
\begin{equation}
    j(v)\cdot j(v)=q(v)1.        \label{102r1}
\end{equation}
Then we have an unique extension of $j$ to a homomorphism $\dpt{\tilde\jmath}{\Cliff(V,q)}{E}$. Moreover, $\Cliff(V,q)$ is the unique associative algebra which have this property for all such~$E$.
\[
\xymatrix{
    \Cliff(V,q) \ar@{^{(}->}[d]_{\displaystyle i} \ar[rd]^{\displaystyle\tilde{j}} &  \\
    V \ar[r]_{\displaystyle j} & D
  }
\]
\label{tho_Cliffunif}
\end{theorem}
This theorem can be seen as a definition of $\Cliff(V,q)$.

\begin{proof}
The proof shall belongs two parts: the first one will show how to extend $j$ and why it is unique, and the second one will prove the unicity of $\Cliff(V,q)$.

We begin by define the extension of $j$. First note that any linear map $\dpt{f}{V}{E}$ can be extended to an algebra homomorphism $\dpt{\overline{f}}{T(V)}{E}$ in only one way. Indeed, the homomorphism condition require that $\overline{f}(v\otimes w)=f(v)\cdot f(w)$.  The whole map $\overline{f}$ is then well defined by the data of $f$ alone.

As far as the map $j$ is concerned, we have the relation \eqref{102r1} which says that $\overline{j}(\mI)=0$. Indeed,
\begin{equation}
 \ovj(v\otimes v-q(v)\cdot(1))=\ovj(v)\cdot\ovj(v)-q(v)\ovj(1)
                              =j(v)\cdot j(v)-q(v)1
                              =0.
\end{equation}
Thus $\dpt{\ovj}{T(V)}{E}$ is a class map for $\mI$, and we can descent $\ovj$ from $T(V)$ to $\Cliff(V,q)$, We define
$\dpt{\tilde\jmath}{\Cliff(V,q)}{E}$ by
\begin{equation}
         \tilde\jmath[x]=\ovj(x)
\end{equation}
where $[x]$ is the class of $x$. That's for the existence part.

The unicity is clear: $f_1=f_2$ on $V$ implies that $\overline{f_1}=\overline{f_2}$ on $T(V)$. Thus $\tilde{f_1}=\tilde{f_2}$ on $\Cliff(V,q)$.

We turn now our attention to the unicity of $\Cliff(C,q)$. Let $D$ be an unital associative algebra such that
\begin{enumerate}
\item $V\subset D$,
\item For any unital associative algebra $E$ and for any $\dpt{f}{D}{E}$ such that $f(v)\cdot f(v)=-q(v)1$, there exists only one homomorphic map $\dpt{\tilde{f}}{D}{E}$ which extend $f$.
\end{enumerate}
We should find a homomorphic map $\dpt{\tilde{k}}{D}{\Cliff(V,q)}$. Let $i$ be the canonical injection $\dpt{i}{V}{D}$. Clearly, we have a homomorphism $V\rightarrow i(V)$. Now, as a space $E$, we can take $\Cliff(V,q)$; $i$ can be seen as a linear map $\dpt{i}{V}{\Cliff(V,q)}$ such that $i(v)\cdot i(v)=q(v)1$. The assumptions say that $i$ can be extended (in only one way) to a homomorphic map $\dpt{\tilde{i}}{D}{\Cliff(V,q)}$.

The Clifford algebra is thus unique up to a homomorphism.

\end{proof}

What we proved is the following: if for any $E$ and for any $\dpt{j}{V}{E}$ such that $j(v)\cdot j(v)=q(v)1$, there exist an unique $\dpt{\tilde{j}}{D}{E}$ which extend $j$, then $D=\Cliff(V,q)$ up to a homomorphism. One ays that $\Cliff(V,q)$ solve an \defe{universal problem}{universal!problem}.

An explicit construction of $\Cliff(V,q)$ can be achieved in the following way. We consider the tensor algebra $T(V)=\bigoplus_{n\geq 0}\left(\otimes^nV\right)=\eC\oplus V\oplus(V\otimes V)\oplus\ldots$ over $V$ the two-sided ideal $\mI$ generated by elements of the form $v\otimes v-q(v)1$. The  \defe{Clifford algebra}{Clifford!algebra}\index{algebra!Clifford} for $(V,q)$ is given by\nomenclature[G]{$\Cliff(p,q)$}{Clifford algebra of $\eR^{1,3}$}
\begin{equation}	\label{defI}
	\Cliff(p,q):=T(V)/\mI
\end{equation}
in which product of $\Cliff(V,q)$ is naturally defined by $[a]\otimes[b]=[a\otimes b]$ if $[a]$ is the class of $a\in T(V)$.

Let us now fix some notations more adapted to what we want to do. Let $V=\eR^{p,q}$ the vector space $\eR^{p+q}$ endowed with a diagonal metric which contains $p$ plus sign and $q$ minus signs. For $v$, $w\in V$, the inner product with respect to the metric $\eta$ of $v$ by $w$ will be denoted by $\eta(v,w)$.  The norm on $V$ will be defined by $\|v\|^2=-\eta(v,v)$. It is neither positive defined, nor negative defined. The explanation of the minus sign will come soon. The Clifford algebra is the quotient $\Cliff(p,q):=T(V)/\mI$ of the tensor algebra by the two-sided ideal $\mI$ generated by elements of the form
\[
	(v\otimes w)\oplus (w\otimes v)\oplus 2\eta(v,w)1
\]
 for $v,w$ in $V$. Depending on the context, we will often use the notations $\Cliff(\eta)$ or $\Cliff(V)$ or $\Cliff(p,q)$. The algebra product is $[x]\cdot[y]=[x\otimes y]$, $x$, $y\in T(V)$.  As long as $z\in V\subset\Cliff(p,q)$, the expression $\eta(z,z)$ is meaningful. The definition of $\Cliff$ is such that $z\cdot z=-\eta(z,z)$. This leads to the somewhat surprising formula  $z^2=\|z\|^2=-\eta(z,z)$.

\subsection{First representation}
%----------------------------------

Let $(V,g)$ be a metric vector space and $\Cliff(V,g)$ its Clifford algebra. For each $v\in V$, we define the two following elements of $\End_{\eR}(\Wedge V)$:
\begin{subequations}
\begin{align}
\epsilon(v)\big( u_1\wedge\cdots\wedge u_k \big)&=v\wedge u_1\wedge\cdots\wedge u_k\\
\iota(v)\big( u_1\wedge\cdots\wedge u_k \big)&=\sum_{j=1}^k(-1)^{j-1}g(u,u_j)u_1\wedge\cdots\wedge\hat u_j\wedge\cdots\wedge u_j.
\end{align}
\end{subequations}
One has $\epsilon(v)^2=0$ and $\iota(v)^2=0$ because $v\wedge v=0$. In order to understand the latter, we wonder what are the terms with $g(v,u_i)g(v,u_j)$ are in
\[ 
  \iota(v)^2\big( u_1\wedge\cdots\wedge u_k \big)=\sum_{l=1}^k(-1)^{j-1}g(v,u_j)\sum_{l=1}^{k-1}(-1)^{l-1}g(v,u_l)u_1\wedge\hat u_l\wedge\hat u_j\wedge\cdots\wedge u_k.
\]
Let's suppose $i<j$. The first term comes when the first $\iota(v)$ acts on $u_j$, its sign is given by $(-1)^{j-1}(-1)^{i-1}$. The second term has the same $(-1)^{i-1}$, but in this term, $u_j$ is on the position $j-1$ because $u_i$ has disappeared.

Now we use $c(v)=\epsilon(v)+\iota(v)$ which fulfils for all $u$, $v\in V$:
\[ 
\begin{split}
c(v)^2&=g(v,v,)1\\
c(u)v(v)+c(v)c(u)&=2g(u,v)1.
\end{split}  
\]
Therefore $c$ can be extended to a representation $c\colon \Cliff(V,g)\to \End(\Wedge V)$. If $\{ e_0,\cdots e_n \}$ is an orthonormal basis of $V$ (i.e. $g(e_i,e_j)=\eta_{ij}$); in this case the $c(e_j)$ are anticommuting and a basis of $\Cliff(V,g)$ is given by 
\begin{equation}
 \{ c(e_{k_1})\cdots c(e_{k_r})\tq 1\leq k_1<\cdots<k_r\leq n \}.
\end{equation}

\subsection{Some consequences of the universal property}
%---------------------------------------------------

The map $-\id|_V$ extends to $\alpha\in\Aut\big( \Cliff(V) \big)$,
\[ 
  \alpha(v_1\cdots v_r)=(-1)^rv_1\cdots v_r
\]
($v_i\in V$) and provides a graduation
\[ 
  \Cliff(V)=\Cliff^0(V)\oplus\Cliff^1(V).
\]
The map $\tau\colon \Cliff(V)\to \Cliff(V)$ extends $\id|_V$ to an anti-homomorphism:
\begin{equation}
\tau(v_1\cdots v_r)=v_r\cdots v_1.
\end{equation}

The \defe{complexification}{complexification!of Clifford algebra} of $\Cliff(V,g)$ is 
\[ 
  \CCliff(V,g):=\Cliff(V,g)\otimes_{\eR}\eC\simeq\Cliff(V^{\eC},g^{\eC}),
\]
the isomorphism being a $\eC$-algebra isomorphism. The $\eR$-linear operator $v\mapsto \overline{ v }$ in $V^{\eC}$ of complex conjugation extends to a $\eR$-linear automorphism $a\mapsto \overline{ a }$. We define the \defe{adjoint}{adjoint!in Clifford algebra} by
\begin{equation}
  a^*=\tau(\overline{ a })
\end{equation}

\subsection{Trace}
%---------------------

\begin{theorem}
There exists one an only one trace $\tr\colon \CCliff(V)\to \eC$ such that
\begin{enumerate}
\item $\tr(1)=1$,
\item $\tr(a)=0$ when $a$ is odd.
\end{enumerate}

\end{theorem}

\begin{proof}

Let $\{ e_1,\cdots,e_n \}$ be an orthonormal basis of $(V,g)$ and $a\in\CCliff(V)$. When one decomposes $a$ into the basis of $e_i$, one finds a lot of terms of each order. Since $\tr$ is a trace, when the $k_i$ are all different,
\[ 
\tr(e_{k_1}\cdots e_{k_{2r}})	=\tr(-e_{k_2}\cdots e_{k_{2r}} e_{k_1}
				=\tr(-e_{k_1}\cdots e_{k_{2r}})\\
\]
So the trace of any even element is zero. We decompose $a$ into
\[ 
  a=\sum_K a_K\prod_{i\in K}e_i
\]
where the sum is taken on the subsets of $\{ 1,\ldots, n \}$. A trace which fulfils the conditions must vanishes on even (but non zero) elements as well as on odd elements, so the only possible form is
\[ 
  \tr a=a_{\emptyset}.
\]
Notice that in order to get this precise form, we used $\tr(1)=1$ and linearity. This proves unicity and existence. Now we have to prove that this is a good definition in the sense that an other choice of basis gives the same result. So we take a new orthonormal basis
\[ 
  e'_j=\sum_{k=1}^nH_{jk}e_k
\]
with $H^tH=\mtu_{n\times n}$. Now we have
\[ 
  a=\sum_{K}^{}a_K\prod_{i\in K}e_i=\sum_{K}^{}a'_K\prod_{i\in K}e'_i,
\]
and we will prove that $a_{\emptyset}=a'_{\emptyset}$. Let's compute a lot:
\[ 
\begin{split}
   e_i'e_j'&=\sum_{k}^{}\sum_{l}^{}H_{ik}H_{jl}e_ke_l\\
		&=\sum_{k=l}^{}H_{ik}H_{jl}e_{k}e_{l}+\sum_{k\neq l}^{}H_{ik}H_{jl}e_{k}e_{l}\\
		&=\sum_{k}^{}H_{ik}H_{jk}1+\sum_{k\neq l}^{}H_{ik}H_{jl}e_{k}e_{l}\\
		&=(HH^t)_{ij}1+\sum_{k\neq l}^{}H_{ik}H_{jl}e_{k}e_{l}.
\end{split}  
\]
The sense of this formula is that when $i\neq j$, the product $e'_{i}e'_{j}$ has no term of order zero. In other terms, as long as we only have terms of order zero, one and two, a change $e\to e'$ does not change the term of order zero. We are now going to an induction proof: we want to prove that $e'_{j_{1}}\ldots e'_{j_{2r}}e'_{l}e'_{k}$ has no scalar term assuming that no even combination has scalar terms up to $2(r-1)$. It reads
\[ 
  \sum_{\text{$K$ even}}^{}a_{K}\prod_{i\in K} e_{i}e'_{l}e'_{k},
\]
therefore we just have to look at terms of the form
\[ 
  e_{j_{1}}\ldots e_{j_{2r}}\Big( (HH)^t_{kl}1-\sum_{i\neq j}^{}C_{kl}^{ij}e_{i}e_{j} \Big)
\]
where the $e_{j_{l}}$ are all different. The first term cannot produce a scalar term. In order to find a scalar term in $e'_{j_{1}}\ldots e'_{j_{2r}}e_{k}e_{l}$, we begin to look at terms whose decomposition of $e'_{j_{1}}\ldots e'_{j_{2r}}$ ends by $e_{l}e_{k}$, i.e.
\[ 
  H_{j_{2r-2}l}H_{j_{2r-1}k}e'_{j_{1}}\ldots e'_{2r-3}e_{l}e_{k}e_{k}e_{l}.
\]
The induction assumption says that there are no scalar term in $e'_{2r-3}e_{l}e_{k}e_{k}e_{l}$.

\end{proof}
One can prove that $\CCliff(C)$ is a Hilbert space with the scalar product 
\begin{equation}
\langle a |b\rangle=\tr(a^*b).
\end{equation}


Let $v\in V$ with $g(v,v)=1$ (thus in $\Cliff(V)$, we have $v^2=1$); since $v=\overline{ v }$, we have
\[ 
  a^*v=vv^*=v^2=1.
\]

\begin{lemma}
The maps $a\mapsto ua$ and $a\mapsto au$ are unitary if and only if $uu^*=u^*u=1$.
\end{lemma}

\begin{proof}
We pick $\lambda\in U(1)$ and $w=\lambda v\in V^{\eC}$ which fulfils $w^*w=1$. This is the most general element such that $ww^*=w^*w=1$. Now for an arbitrary $a$, $b\in\CCliff(V)$, we compute the two followings:
\[ 
  \langle wa|wb\rangle=\tr\big( (wa)^*wb \big)
		=\tr\big( a^*w^*wb \big)
		=\tr(a^*b)
		=\langle a|b\rangle,
\]
and
\[ 
  \langle aw|bw\rangle=\tr\big( w^*a^*bw \big)
		=\tr(ww^*a^*b)
		=\tr(a^*b)
		=\langle a|b\rangle.
\]
This proves that $a\mapsto wa$ and $a\mapsto aw$ are two unitary operators on the Hilbert space $\CCliff(V)$.

For the converse, we impose for all $a$, $b\in\CCliff(V)$:
\[ 
  \langle ua|ub\rangle=\tr(ba^*u^*u)\stackrel{!}{=}\tr(ba^*).
\]
In particular with $a^*b=1$, $\tr(u^*u)=\tr(1)=1$, thus the scalar part of $u^*u$ is $1$. So we write $u^*u=1+f$ where $f$ is non scalar, and for any $x\in\CCliff(V)$ , we have
\[
   \tr(x)=\tr(xu^*u)=\tr(x)+\tr(xf).
\]
We conclude that $\tr(xf)=0$, and therefore that $f=0$.
\end{proof}


\section{Spinor representation}
%+++++++++++++++++++++++++++++

For the spinor representation, we restrict ourself to the even case $p+q=2n$.

The aim of this subsection is to find some faithful representations of the complex Clifford algebra $\Cliff^{\eC}(p,q)$. In order to achieve this, we first consider $V^{\eC}$, the complex vector space of $V$ with an orthonormal basis $\{ e_1,\cdots,e_{p-1},e_p,\cdots,e_q  \}$. The metric is $\eta(e_k,e_k)=1$ and $\eta(e_{p+k},e_{p+k})=-1$ for $k=0,\cdots,p-1$. We use the following basis:
\begin{align}
f_k&=\frac{1}{2}(e_k+e_{p+k}),& g_k&=\frac{1}{2}(e_k-e_{p+k}),\\
 f_{p+s}&=\frac{1}{2}(e_{2p+2s}+ie_{2p+2s+1}),& g_{p+s}&=\frac{1}{2}(e_{2p+2s}-ie_{2p+2s})
\end{align}
for $k=0,\cdots,p-1$.
We note that $\{f_0,g_0\}$ spans a $\eC^2$-space which is $\eta$-orthogonal to the one which is spanned by $\{f_1,g_1\}$. The following two  spaces will prove to be useful:
\begin{subequations}
\begin{align}
  W           &=\Span_{\eC}\{f_0,f_1\}\simeq\eC^2,\\
 \underline{W}&=\Span_{\eC}\{g_0,g_1\}\simeq\eC^2.
\end{align}
\end{subequations}
\nomenclature{$W,\underline{W}$}{Totally isotropic subspace}
It is easy to compute the various products; among others we find
\begin{equation}
 \eta(f_0,f_0)=0,\quad
 \eta(f_1,f_0)=0,\quad
  \eta(f_1,f_1)=0;
\end{equation}
so that for any $w\in W$, we have $\scal{w}{w}=0$; for this reason, we say that $W$ is a \defe{completely isotropic}{isotropic!subspace!completely} subspace of $(V^{\eC},\eta^{\eC})$. The space $\underline{W}$ has the same property.

\begin{proposition}
We have
\begin{equation}
  \underline{W}\simeq W^*,
\end{equation}
where $W^*$ is the dual space of $W$. By $\simeq$ we mean that there exists a linear bijective map $\dpt{\psi}{\underline{W}}{W^*}$.
\end{proposition}
\begin{proof}
For each $\uw\in\uW$, we define $\dpt{\psi(\uw)}{W}{\eC}$ by
\[
   \psi(\uw)(w)=\eta(w,\uw).
\]
We first show that the map $\psi$ is injective. Let $\uw\in\uW$ be so that $\psi(\uw)=0$. Thus for all $v\in W$, we have
\begin{eqnarray}
   \label{3101r1}\psi(\uw)v=\eta(\uw,v)=0.
\end{eqnarray}
By decomposing $\uw=ag_0+bg_1$ and taking successively $v=f_0$ and $v=f_1$, we see that $a=b=0$.

The next step is to see that the map $\psi$ is surjective. We know that $dim_{\eC}\uW=dim_{\eC}W^*=2$ and that $\psi(g_0)\neq 0$. Let's prove that $\{\psi(g_0),\psi(g_1)\}$ is a basis of $W^*$. It is clear by linearity that $\{\psi(ag_0):a\in\eC\}=\Span\{\psi(g_0)\}$. The fact that $\psi$ is injective  imposes that $\psi(g_1)$ doesn't belong to $\Span\{\psi(g_0)\}$. So $\{\psi(g_0),\psi(g_1)\}$ is a two-dimensional free subset of $W^*$, and therefore is a basis of $W^*$.
\end{proof}

We turn our attention to the exterior algebra $\Lambda W=\eC\oplus W\oplus(W\wedge W)\oplus\cdots\oplus\wedge^{p+q}W$\nomenclature{$\Lambda W$}{Space of spinor representation} of $W$. 
\nomenclature{$\dpt{\tilde\rho}{(\eR^{1+3})^{\eC}}{\End(\Lambda W )}$}{Spinor representation}

\begin{definition}
\index{endomorphism!of $\Lambda W$}
We define the homomorphism $\dpt{\tilde\rho}{V^{\eC}}{\End(\Lambda W)}$ by 
\begin{equation}
\begin{split}
 \tilde\rho(f_i)\alpha&=f_i\wedge\alpha,\\
 \tilde\rho(g_i)\alpha&=-\iota(g_i)\alpha
\end{split}
\end{equation}
($v\in V^{\eC}$, $\alpha\in\Lambda W$) where $\iota$ denotes the interior product defined in page \pageref{pg_DefProdExt}.\
\label{defrt}
\end{definition}
  More explicitly, for all $z\in\eC$ and for all $w,w'\in W$, we have
\begin{subequations}
\begin{align}
 \tilde\rho(f_i)z&=zf_i,&\tilde\rho(g_i)z&=0,\\
 \tilde\rho(f_i)w&=f_i\wedge w,&\tilde\rho(g_i)w&=-\eta(g_i,w)1,\\
 \tilde\rho(f_i)(w\wedge w')&=0,&\tilde\rho(g_i)(w\wedge w')&=-\eta(g_i,w)w'+\eta(g_i,w')w.
\end{align}
\end{subequations}
We will see that, \emph{via} some manipulations, $\tilde\rho$ provides a faithful representation of the Clifford algebra, the \defe{spinor representation}{spinor!representation}.
\index{representation!of Clifford algebra}

\begin{remark}
By ``endomorphism of $\Lambda W$'', we mean an endomorphism for the \emph{linear} structure of $\Lambda W$. We obviously not have $\tilde\rho(x)(\alpha\wedge\beta)=\tilde\rho(x)\alpha\wedge\tilde\rho(x)\beta$. 
\end{remark}

\begin{proposition}
The map $\tilde\rho$ is injective.
\end{proposition}

\begin{proof}
We have to show that $\tilde\rho(v)=0$ ($v$ in $V^{\eC}$) implies $v=0$. Any $v\in V^{\eC}$ can be written as
$v=a^if_i+b^ig_i$ with a sum over $i$. We first have that
\[
 \tilde\rho(a^if_i+b^ig_i)z=za^if_i=0,
\]
 but the $f_i$ are independents and then $a^i=0$. We can also write
\[
 \tilde\rho(b^0g_0+b^1g_1)f_1=-b^0\eta(g_0,f_1)-b^1\eta(g_1,f_1)=-\frac{b^1}{2}=0,
\]
 then $b^1=0$. The same with $f_0$ proves that $b^0=0$.
 \end{proof}

The homomorphism $\tilde\rho$ extends to the whole the tensor algebra of $V^{\eC}$ by the following definitions:
\begin{subequations}
\begin{align}
 \tilde\rho(1)              &=\id_{\Lambda W},\\
 \tilde\rho(e_k)            &=\tilde\rho(e_k),\\
 \tilde\rho(e_{k_1}\otimes\ldots\otimes e_{k_r})&=
                      \tilde\rho(e_{k_1})\circ\ldots\circ\tilde\rho(e_{k_r}).\label{eq:3101r2}
\end{align}
\end{subequations}
So we get $\dpt{\tilde\rho}{T(V^{\eC})}{\End(\Lambda W)}$.  The following proposition will allow us to descent $\tilde\rho$ to a representation of the Clifford algebra.

\begin{proposition}
The homomorphism $\tilde\rho$ maps $\mI$ to $0$: $\tilde\rho(\mI)=0$.
\end{proposition}

\begin{probleme}
This proposition is wrong: there is a double covering.

Moreover, there is a sign problem in the proof: the sign in the first lines is not the one used in the definition of the Clifford algebra.
\end{probleme}


\begin{proof}
We have to check the following:
\[\tilde\rho(v\otimes w\oplus w\otimes v-2\eta(v,w)1)=0\]
for any choice of
 $v,w$ in $\{e_0,e_1,e_2,e_3\}$.
  Here we will just check it explicitly for $v=e_0$ and $w=e_1$. The computation uses the definition \eqref{eq:3101r2}:
\begin{equation}
\begin{split}
\tilde\rho(e_0\otimes e_1\oplus e_1\otimes
             e_0-2\eta(e_0,e_1)&=\tilde\rho(e_0)\circ\tilde\rho(e_1)+\tilde\rho(e_1)\circ\tilde\rho(e_0)\\
                               &=2\left[\tilde\rho(f_0)^2-\tilde\rho(g_0)^2\right].
\end{split}
\end{equation}
It is easy to see that $\tilde\rho(f_0)^2=0$:
\begin{equation}
 \tilde\rho(f_0)^2\left[z\oplus w\oplus w_1\wedge w_2 \right]=\tilde\rho(f_0)[zf_0\oplus f_0\wedge w]
                                                   =zf_0\wedge f_0,
							=0.
\end{equation}
 The proof that $\tilde\rho(g_0)^2=0$ is almost the same:
\[ 
 \tilde\rho(g_0)^2\left[z\oplus w\oplus w_1\wedge w_2 \right]
 =\tilde\rho(g_0)[-\eta(g_0,w)1\oplus-\eta(g_0,w_1)w_2\oplus\eta(g_0,w_2)w_1].
\]

\end{proof}

We can now see $\tilde\rho$ as a map $\dpt{\tilde\rho}{\Cliff^{\eC}(p,q)}{\End(\Lambda W )}$. By construction, it is a homomorphism and, thus, is a representation of $\Cliff^{\eC}(p,q)$ on $\Lambda W$. For compactness, we use the notation \index{dirac!matrices}\nomenclature{$\gamma_i$}{Abstract definition of Dirac matrices}
\begin{equation}
 \label{defgamma}\gamma_a:=\sqrt{2}\tilde\rho(e_a). 
\end{equation}

\begin{lemma}
The $\gamma$'s operators satisfy the following relation:
\begin{equation}\label{3101r3}
  \gamma_a\gamma_b+\gamma_b\gamma_a=-2\eta_{ab}\mtu.
\end{equation}
\label{3101l1}
\end{lemma}

\begin{proof}
We have to check this equality on any element of $\Lambda W$. If we choose
$w_1=af_0+bf_1$ and $w_2=a'f_0+b'f_1$,
 we find $w_1\wedge w_2=(ab'-ba')f_0\wedge f_1$.

For example, we will explicitly check \eqref{3101r3} with $a=b=0$, i.e. $\tilde\rho(e_0)\circ\tilde\rho(e_0)=\frac{1}{2}\id$, which proves that $\gamma_0\circ\gamma_0=\id$.
 \begin{equation}
\begin{split}
   \tilde\rho(e_0)^2[z\oplus w\oplus(ab'-ba') f_0\wedge f_1]&=\tilde\rho(f_0+g_0)^2[z\oplus w\oplus(ab'-ba') f_0\wedge f_1]\\ 
                                              &=\tilde\rho(f_0+g_0)\Big[zf_0\oplus f_0\wedge w\oplus-\eta(g_0,w)1\\
                                                            &\qquad-(ab'-ba')\eta(g_0,f_0)f_1\\
                                                             &\qquad+(ab'-ba')\eta(g_0,f_1)f_0\Big]\\
                                              &=\frac{1}{2}(z\oplus w\oplus(ab'-ba') f_0\wedge f_1).
\end{split}				      
\end{equation}
\end{proof}

\begin{lemma}
For any sequence $i_0,\ldots i_3$ of $0$ and $1$ (with at least one of them equals to $1$), we have
\begin{equation}
 \tr(\gamma_0^{i_0}\cdots \gamma^{i_{2n-1}}_{2n-1})=0.
\end{equation}
We take the convention that $\gamma_a^0=\mtu$.
\label{3101l2}
\end{lemma}

\begin{proof}
 If the number of nonzero $i_k$ is even (say $2m$), we have:
\[
	\tr(\gamma_{a_1}\ldots\gamma_{a_{2m}})=\tr(\gamma_{a_{2n}}\gamma_{a_1}\ldots\gamma_{a_{2m-1}})
\] 
because the trace is invariant under cyclic permutations. But we can also permute $\gamma_{a_{2m}}$ with the $2m-1$ other $\gamma$'s.  $\tr(\gamma_{a_1}\ldots\gamma_{a_{2m}})=(-1)^{2n-1}\tr(\gamma_{a_{2m}}\gamma_{a_1}\ldots\gamma_{a_{2m-1}})$ because each permutation gives an extra minus sign (\hbox{lemma \ref{3101l1}}). Then the trace is zero.

If the number of nonzero $i_k$ is odd (say $2m-1$). Let $i_a=0$ (we restrict ourself to the even dimensional case). We have $\tr(A)=-\eta_{aa}\tr(A\gamma_a\gamma_a)$. Using once again the cyclic invariance of the trace, $\tr(\gamma_{a_1}\ldots\gamma_{a_{2m-1}}\gamma_a\gamma_a)=\tr(\gamma_a\gamma_{a_1}\ldots\gamma_{a_{2m-1}}\gamma_a)$. But, if we permute the \emph{first} $\gamma_a$ with the $2m-1$ first $\gamma$'s, we find \hbox{$\tr(\gamma_{a_1}\ldots\gamma_{a_{2m-1}}\gamma_a\gamma_a)=-\tr(\gamma_a\gamma_{a_1}\ldots\gamma_{a_{2m-1}}\gamma_a)$ }, and the trace is zero again.
\end{proof}

\begin{proposition}
The subset
\[
	\left\{\mtu,\dga{a}{b}\,(a<b),\tga{a}{b}{c}\,(a<b<c), \cdots, \gamma_{0}\cdots\gamma_{2n} \right\}
\]
 is free in $\End(\Lambda W)$.
\end{proposition}

\begin{proof}
We consider a general linear combination of these operators:
\[
 E=\lambda\mtu+\sum_a\lambda_a\gamma_a+\sum_{a<b}\lambda_{ab}\dga{a}{b}+\ldots+
 \sum_{a<b<c<d}\lambda_{abcd}\qga{a}{b}{c}{d}.
\]
The claim is that if $E=0$, then all the coefficients $\lambda_{(\ldots)}$ must be zero. First note that $Tr(E)=0=\lambda$ by lemma \ref{3101l2}. It is also clear that $Tr(\gamma_iE)=0=\lambda_i$. In order to see that $\lambda_{ij}=0$, we compute $Tr(\gamma_j\gamma_iE)=0=\lambda_{ij}$. And so on.
\end{proof}

How many operators does we have in this free system ? Any operators in this system can be written as $\gamma_{0}^{i_{0}},\cdots\gamma^{i_{2n-1}}_{2n-1}$ with $i_k$ equal to zero or one. Thus we have $2^{2n}$ operators. On the other hand, we know that $dim_{\eC}\Lambda W=2p+2$, and then that $dim_{\eC}\End(\Lambda W)=4^2=16$. The result is that $\{ \gamma_{0}^{i_{0}},\cdots\gamma^{i_{2n-1}}_{2n-1}  \tq i_k=0\,or\,1\}$ is a basis of $\End(\Lambda W)$. In other words (if we suppose a suitable ordering), the image by $\tilde\rho$ of 
\[ 
B=\{1,e_a,e_a\otimes e_b,e_a\otimes e_b\otimes e_c,e_a\otimes e_b\otimes e_c\otimes e_d\}
\]
 is a basis of $\End(\Lambda W)$.

If $B$ is a basis of $C^{\eC}_{(p,q)}$, then $\tilde\rho$ is bijective and thus isomorphic.  Therefore, we expect $\dpt{\tilde\rho}{C^{\eC}_{(p,q)}}{\End(\Lambda W)}$ to be a faithful representation\index{representation!of Clifford algebra}. It is not difficult to see that $B$ is indeed a basis thanks to the equivalence relation.

\subsection{Explicit representation}
%------------------------------

First, we choose a basis for $\Lambda W$:
\begin{eqnarray}\label{102r2} 1=\left(\begin{matrix}
1 \\
0 \\
0 \\
0
\end{matrix}\right),\quad
f_0=\left(\begin{matrix}
0 \\
1 \\
0 \\
0
\end{matrix}\right),\quad
f_1=\left(\begin{matrix}
0 \\
0 \\
1 \\
0
\end{matrix}\right),\quad
f_0\wedge f_1=\left(\begin{matrix}
0 \\
0 \\
0 \\
1
\end{matrix}\right).
\end{eqnarray}
 Here is the explicit computation for the matrix $\gamma_0$ in this basis. First remark that $\tilde\rho(e_0)1=f_0$, $\tilde\rho(e_0)f_0=\frac{1}{2}$, $\tilde\rho(e_0)f_1=f_0\wedge f_1$, $\tilde\rho(e_0)(f_0\wedge f_1)=\frac{1}{2} f_1$. Then
\begin{equation}
\begin{split}
 \gamma_0\left(\begin{matrix}1 \\0 \\0 \\0\end{matrix}\right)=\sqrt{2}\left(\begin{matrix}0 \\1 \\0  \\0\end{matrix}\right),\quad
 \gamma_0\left(\begin{matrix}0 \\1 \\0 \\0\end{matrix}\right)=\sqrt{2}\left(\begin{matrix}\frac{1}{2} \\0 \\0 \\0\end{matrix}\right),\\
 \gamma_0\left(\begin{matrix}0 \\0 \\1 \\0\end{matrix}\right)=\sqrt{2}\left(\begin{matrix}0 \\0 \\0 \\1\end{matrix}\right),\quad
 \gamma_0\left(\begin{matrix}1 \\0 \\0 \\0\end{matrix}\right)=\sqrt{2}\left(\begin{matrix}0 \\0 \\\frac{1}{2} \\0\end{matrix}\right).
\end{split}
\end{equation}
This allows us to write down $\gamma_0$; the same computation gives the other matrices.\index{dirac!matrices}\nomenclature{$\gamma_i$}{Explicit form of gamma matrices}
\begin{equation}
\begin{split}
\gamma_0=\sqrt{2}\begin{pmatrix}
0 & \frac{1}{2} & 0 & 0 \\
1 & 0 & 0 & 0 \\
0 & 0 & 0 & \frac{1}{2} \\
0 & 0 & 1 & 0
\end{pmatrix}, \qquad
\gamma_1=\sqrt{2}\left(\begin{matrix}
0 & -\frac{1}{2} & 0 & 0 \\
1 & 0 & 0 & 0 \\
0 & 0 & 0 & -\frac{1}{2} \\
0 & 0 & 1 & 0
\end{matrix}\right),\\
\gamma_2=\sqrt{2}\begin{pmatrix}
0 & 0 & -\frac{1}{2} & 0 \\
0 & 0 & 0 & \frac{1}{2} \\
1 & 0 & 0 & 0 \\
0 & -1 & 0 & 0
\end{pmatrix},\qquad
\gamma_3=\sqrt{2}\begin{pmatrix}
0 & 0 & -\frac{i}{2} & 0 \\
0 & 0 & 0 & \frac{i}{2} \\
-i & 0 & 0 & 0 \\
0 & i & 0 & 0
\end{pmatrix}.
\end{split}
\end{equation}
It is easy to check that these matrices satisfies \eqref{3101r3}. 

Notice that, up to a suitable change of basis in $\Lambda W $, these are the usual Dirac matrices\index{dirac!matrices}. Indeed we actually solved the physical problem to find a representation of the algebra \eqref{cliffphys}.  We understand by the way why do physicists work with $4$-components spinors: the $\gamma$'s are operators on the four-dimensional space $\Lambda W$; hence the Dirac operator will naturally acts on four-components objects.

The main result of this section is an explicit faithful representation of $\CCliff(p,q)$. This allows us to write a \defe{Dirac operator}{dirac!operator!on $\protect\eR^{1,3}$} which solve (see the invitation \ref{Secqft} and \cite{Bronn}) the problem  to find a ``square root'' of the d'Alembert operator: the differential operator $\Dir=\gamma\hmu\partial_{\mu}$ satisfies $\Dir^2=\Box$.

\subsection{A remark}  % C'est ce qui arrive quand je ne sais pas quel titre donner.
%---------------------

 Let us compare the two faithful representations
\[ 
\begin{split}
  c\colon \Cliff(V)&\to \End_{\eR}(\wedge V)\\
\tilde\rho\colon \CCliff&\to  \End_{\eR}(\wedge W).
\end{split}  
\]
They obviously comes from the same ideas. One common point is that 
\[ 
  c(e_1)(e_1\wedge e_2)=2\tilde\rho (e_1)(e_1\wedge e_2)=e_2,
\]
but they are different:
\[ 
\begin{split}
  \tilde\rho(e_3)(e_0\wedge e_2)&=0\\
c(e_3)(e_0\wedge e_2)&=e_3\wedge e_0\wedge e_1.
\end{split}  
\]

\subsection{General two dimensional Clifford algebra}
%---------------------------------------------------

The Clifford algebra for the metric
\[ 
  g=\begin{pmatrix}
\alpha&\delta\\\delta&\beta
\end{pmatrix}
\]
is realised by matrices
\[ 
  \gamma_1=
\epsilon\begin{pmatrix}
\sqrt{\alpha}\\ & -\sqrt{\alpha}
\end{pmatrix},\quad
\gamma_2=
\epsilon\begin{pmatrix}
\delta/\sqrt{\alpha}& \beta-\delta^2/| \alpha |\\
1		& -\delta/\sqrt{\alpha}	
\end{pmatrix}
\]
where $\epsilon=\pm 1$ is chosen in such a way that $\epsilon| \alpha |=\alpha$. 
 
\section{Spin group}
%+++++++++++++++++++

We will not immediately go on with Dirac operators on Riemannian manifolds because we still have to build some theory about the Clifford algebra itself. In particular, we have to define the spin group which will play a central role in the definition of the Dirac operator. Almost all --and (too ?) much more-- the concepts we will introduce in this section can be found in \cite{Chevalley}; a more physical oriented but useful approach can be found in \cite{Preparation}.

Let define the map $\dpt{\chi}{\Gamma(p,q)}{GL(\eR^{1,3})}$ by 
\begin{equation}
                \chi(x)y=\alpha(x)\cdot y\cdot x^{-1}.
\end{equation}
\nomenclature{$\chi$}{A representation of $\Gamma(p,q)$}\nomenclature[G]{$\Gamma(p,q)$}{Clifford group}
Let
\[
 \Gamma(p,q)=\{x\in\Cliff(p,q)\tq\textrm{$x$ is invertible and }  \chi(x)y  \in V\textrm{ for all $y\in V$}\}.
\]
It should be remarked that this definition comes back to the real Clifford algebra. The Clifford algebra product gives this subset a group structure which is called the \defe{Clifford group}{Clifford!group}. Any $x\in V$ is invertible since $x\cdot x=-\eta(x,x)1$, the inverse of $x$ is given by $x^{-1}=x/\|x\|^2$.

\index{Clifford!algebra!grading of}
The subset $\Cliff(p,q)^+$ (resp. $\Cliff(p,q)^-$) of $\Cliff(p,q)$ is the image of even (resp. odd) tensors of $T(V)$ by the canonical projection $T(V)\to\Cliff(p,q)$. With these definitions, we have a natural grading of $\Cliff$:
\begin{equation}
 \label{directC}\Cliff(p,q)=\Cliff(p,q)^+\oplus\Cliff(p,q)^-,
 \end{equation}
and the subgroups
\begin{align}
\label{defgplus}
\Gamma(p,q)^+&=\Gamma(p,q)\cap\Cliff(p,q)^+,&\Gamma(p,q)^-&=\Gamma(p,q)\cap\Cliff(p,q)^-.
\end{align}\nomenclature[G]{$\Cliff(p,q)^{\pm}$}{Grading of Clifford algebra}

For $x_1,\ldots,x_n\in V$, we have $\tau(x_1\cdots x_n)=x_n\cdots x_1$. \nomenclature[G]{$\Spin(p,q)$}{Spin group of $\eR^{1,3}$} The \defe{spin group}{spin!group!on $\protect\eR^{1,3}$} is
\begin{equation}   \label{defSpinun}
 \Spin(p,q)=\{x\in\Gamma(p,q)^+\vert\tau(x)=x^{-1}\}
\end{equation}
while the \defe{spin norm}{spin!norm}\nomenclature{$\dpt{N}{\Gamma(p,q)}{\Gamma(p,q)}$}{Spin norm} is the map $\dpt{N}{\Gamma(p,q)}{\Gamma(p,q)}$ defined by
\[
 N(x)=x\tau(\alpha(x)).
\]
We will see in proposition \ref{proppourN} that $N$  actually takes its values in $\eR$ and is therefore a homomorphism $N\colon \Gamma(p,q)\to \eR$

\begin{remark}
The elements of $\Spin(p,q)$ are spin-normed at $1$. Indeed, take a $s$ in $\Spin(p,q)$. We have $N(s)=s\cdot \tau(s)=1$ because $\alpha(s)=s$ and $\tau(s)=s^{-1}$. In particular $\Spin(p,q)\cap\eR=\eZ_{2}$.
\label{rem:spin_norm_u}
\end{remark}

\subsection{Studying the group structure}
%--------------------------------

\begin{proposition}
The set $\Gamma(p,q)$ admits a Lie group structure.
\end{proposition}
\begin{proof}

During this proof, $\mu$ denotes the Clifford multiplication: $\mu(x,y)=x\cdot y$. We know that $\Cliff^{\eC}(p,q)$ is isomorphic to $\End(\Lambda W)$ in which the multiplication is a continuous map. Thus $\mu$ is continuous on $\Cliff^{\eC}(p,q)$. But $\Cliff(p,q)$ is a closed subset of $\Cliff^{\eC}(p,q)$, so $\mu$ is a continuous map in $\Cliff(p,q)$. This proves that  $\chi$ seen as a map from $\Gamma(p,q)\times V$ to $V$ is a continuous map.

The space $V$ is closed in $\Cliff(p,q)$, thus $\sigma^{-1}(V)$ is also closed. But $\sigma^{-1}(V)=\Gamma(p,q)\times\Cliff(p,q)$. So $\Gamma(p,q)$ is closed in $\Cliff(p,q)$.

Now the result is just a consequence of theorems \ref{Helgason2.3} and \ref{Helgason4.2}. Indeed, let us study the subset $\mI$ which appears in the definitions of the Clifford algebra. It makes no difficult to convince ourself that it is a closed subgroup of $T(V)$. The theorem \ref{Helgason4.2} thus makes $\Cliff(p,q)=T(V)/\mI$ a Lie group. But we just say that $\Gamma(p,q)$ is closed in $\Cliff(p,q)$, and the fact that $\Gamma(p,q)$ is a subgroup of $\Cliff(p,q)$ is clear. By theorem \ref{Helgason2.3} we conclude that there exists a Lie group structure on $\Gamma(p,q)$.
\end{proof}

\begin{lemma}
The map $\chi$ is a homomorphism, in other words $\chi$ is a representation of $\Gamma(p,q)$.
\index{representation!of $\Gamma(p,q)$}
\end{lemma}

\begin{proof}
The following computation uses the fact that $\alpha$ is a homomorphism:
\[
\begin{split}
\chi(a\cdot b)y&=\alpha(a\cdot b)\cdot y\cdot (a\cdot b)^{-1}
               =\alpha(a)\cdot\alpha(b)y\cdot b^{-1}\cdot a^{-1}\\
               &=\alpha(a)\cdot\chi(b)y\cdot a^{-1}
               =\chi(a)\chi(b)y.
\end{split}
\]
\end{proof}
Let $y\in\Gamma(p,q)^-$ and $v\in V$. Where is $y\cdot v$ ? First note that $(y\cdot v)^{-1}=v^{-1}\cdot y^{-1}$, so that
\begin{equation}
\begin{split}
  \alpha(y\cdot v)\cdot w\cdot(y\cdot v)^{-1}&=-\alpha(y)\cdot v\cdot w\cdot v^{-1}\cdot y^{-1}\\
                                            &=-\alpha(y)\big( 2\eta(v,w)-w\cdot v \big)\cdot v^{-1}\cdot y^{-1}\\
					    &=-2\eta(v,w)\alpha(y)\cdot v^{-1}\cdot y+\alpha(y)\cdot w\cdot y^{-1}
\end{split}
\end{equation}
which belongs to $V$ because $y\in\Gamma(p,q)$. This reasoning shows that (apart for $0$), $y\cdot v\in\Gamma(p,q)^+$ if and only if $y\in\Gamma(p,q)^-$.

\begin{lemma}
If $x\in V$ is non-isotropic (i.e. $\eta(x,x)\neq 0$), the automorphism $\chi(x)$  is the orthogonal symmetry with respect to $x^{\perp}$.
\end{lemma}

We recall that\nomenclature{$x^{\perp}$}{Space orthogonal to $x$}
\[ 
  x^{\perp}=\{ y\in V\tq\eta(x,y)=0  \}.
\]
We will denote by $\sigma^x$ the orthogonal symmetry with respect to $x^{\perp}$.

\begin{proof}
When the operator $\sigma^x$ acts on $y$, it just change the sign of the ``$x$-part''\ of $y$. So we can write $\sigma^x y=y-2\eta(x,y) 1_x$, where $1_x:=x/\|x\|$. It should be checked if
$\chi(x)y=\alpha(x)\cdot y\cdot x^{-1}$ is equal to $y-2\eta(x,y) 1_x$ or not. We know that $x\cdot x=\eta(x,x)1=-\|x\|$. It follows that
\[
  x\cdot y+y\cdot x=2\eta(x,y)\frac{x\cdot x}{\|x\|}.
  \]
If we multiply this at right by $x^{-1}$, using the fact that $\alpha(x)=-x$, we find
\[
-\alpha(x)\cdot y\cdot x^{-1}=-y+2\eta(x,y) 1_x,
\]
which is precisely the identity we wanted to check.
\end{proof}

The following result will help us to identify subgroups of Clifford group as isometry groups.
\begin{theorem}[Cartan-Dieudonné theorem]
\index{Cartan-Dieudonné theorem}\index{theorem!Cartan-Dieudonné}
Each $\sigma$ in $O(1,3)$ can be written as
\hbox{$\sigma=\tau_1\circ\ldots\circ\tau_m$}, where the $\tau$'s are orthogonal symmetries with respect to hyperplanes which are orthogonal to non-isotropic vectors.
\label{CartanDieu}
\end{theorem}

\begin{proposition}
\[
              \chi(\Gamma(p,q))=O(p,q).
\]
\label{prop1001t1}
\end{proposition}

\begin{proof}
In order to show that $\chi(\Gamma(p,q))\subset O(p,q)$ take $z\in V$ and $x\in\Gamma(p,q)$. Since $\alpha(x)\cdot z\cdot x^{-1}$ lies in $V$, we can write:
\[
\alpha(x)\cdot z\cdot x^{-1}=-\alpha\left(\alpha(x)\cdot z\cdot x^{-1}\right)
=-x\cdot\alpha(z)\cdot\alpha(x^{-1})=x\cdot z\cdot\alpha(x^{-1}).
\]
In order to see that $\chi(x)\in O(p,q)$, we have to prove that $\left\|\chi(x)y\right\|_{(p,q)}^2=\|y\|_{(p,q)}^2$. This is achieved by the following computation:
\begin{equation}
\begin{split}
 \left\|\chi(x)y\right\|_{(p,q)}^2&=-\left(\alpha(x)\cdot y\cdot x^{-1}\right)^2
                                  =\left(\alpha(x)\cdot y\cdot x^{-1}\right)\left(x\cdot y\cdot\alpha(x^{-1})\right)\\
                                  &=-\alpha(x)\cdot y^2\cdot\alpha(x^{-1})
                                  =\|y\|^2_{(p,q)}.
\end{split}
\end{equation}
The last step is simply the fact that $y^2\in\eR$ and therefore commutes with anything. We now know that $\chi(x)\in O(p,q)$ for all $x\in\Gamma(p,q)$. Thus $\chi(\Gamma(p,q))\subset O(p,q)$.

For the second part, let $\sigma$ be in $O(p,q)$. The Cartan-Dieudonné theorem\index{Cartan-Dieudonné theorem}\index{theorem!Cartan-Dieudonné}(theorem \ref{CartanDieu}) says that $\sigma=\sigma^{x_1}\circ\ldots\circ\sigma^{x_r}$ for some $x_1,\ldots, x_r$ in $V$. Thus $\sigma=\chi(x_1\cdots x_r)$, \hbox{and $O(p,q)\subset\chi(\Gamma(p,q))$}.
\end{proof}

\begin{proposition}
\begin{equation}
      \ker\chi=\eR\invtible
\end{equation}
where the right hand side is the set of invertible elements of $\eR$.
\label{prop1001p1}
\end{proposition}
\nomenclature{$\cA\invtible$}{The set of invertible elements of the algebra $\cA$; for example $\eR\invtible=\eR\setminus\{ 0 \}$}

\begin{proof}
Before beginning the proof, we want to insist on the fact that $x\in \ker\chi$ does not mean that $\chi(x)y=0$ for all $y$ in $V$. The ``zero''\ of an algebra is the element $e$ which satisfies $e\cdot y=y\cdot e=y$ for all $y$ in the algebra. In other words, $x$ is in the kernel of $\chi$ if and only if $\chi(x)=\id$.

First we show that $\eR_0\subset \ker\chi$. If $x\in \eR$, then $\alpha(x)=x$. Therefore, when $x\neq 0$,
\[
\chi(x)y=\alpha(x)\cdot y\cdot x^{-1}=y,
\]
because the algebra product $\cdot$ between an element of $\Cliff(p,q)$ and a real is commutative. Note that this does not work with $x=0$.

We are now going to show that $\ker\chi\subset\eR$. Let $z\in\ker\chi$. We decompose (definitions \eqref{defgplus}) it into his odd and even part: $z=z^++z^-$, with $z^{\pm}\in\Gamma(p,q)^{\pm}$. These two can be written as $z^+=e_{j_1}\cdots e_{j_{2r}}$ and $z^-=e_{i_1}\cdots e_{i_{2r-1}}$ with no two $i_k$ or $j_k$ equals. This is almost the general form of elements in even and odd part of $\Gamma(p,q)$: the only other possibility is $z$ in $\eR$. Obviously $\alpha(z^{\pm})=\pm z^{\pm}$. Multiplying the condition $\chi(z)y=y$ at right by $(z^++z^-)$, we find \[(z^+-z^-)y=y(z^++z^-).\] Thanks to equation \eqref{directC}, we can split this condition into even and odd parts:
\begin{align}
 z^+y&=yz^+,
 &z^-y&=-yz^-.
\end{align}
The first equation with $y=e_{j_1}$ gives $e_{j_1}\cdots e_{j_{2r}}\cdot e_{j_1}=e_{j_1}e_{j_1}\cdots e_{j_{2r}}$. In the left hand side, permute the last $e_{j_1}$ from last to second position. So we find the right hand side, with an extra minus sign. This means that $z^+=0$. In the same way, the second equation gives $z^-=0$. We are left with the last possibility: $z\in\eR$.
\end{proof}

\begin{corollary}
For any $s\in\Gamma(p,q)$, there exists some non-isotropic vectors $x_1,\ldots,x_r$, and $c\in\eR$ such that $s=cx_1\cdots x_r$.
\label{602c1}
\end{corollary}

\begin{proof}
Let us take a $s\in\Gamma(p,q)$; we just saw (theorem \ref{prop1001t1}) that $\chi(s)$ is an element of $O(p,q)$. It can be written $\chi(s)=\sigma_1\circ\ldots\circ\sigma_m$. But we had shown that $\sigma_i=\chi(x_i)$ for any $x_i$ normal to the hyperplane defining $\sigma_i$. We thus have
\[
     \chi(s)=\chi(x_1\cdots x_m),
\] 
where $s$ belongs to $\Gamma(p,q)$ and is therefore invertible. This leads us to write $\id=\chi(s^{-1}\cdot x_1\cdots x_m)$. But the kernel of $\chi$ is $\eR$ (proposition \ref{prop1001p1}); so one can find a $r\in\eR$ such that $s^{-1}\cdot x_1\cdots x_m=r$. The claim follows.
\end{proof}

\begin{lemma}
If $v\in V$,
\begin{equation}
                 \det\chi(v)=-1.
\end{equation}
\end{lemma}

\begin{proof}
We already know that $det\chi(v)=\pm 1$. To check that the right sign is plus, take the following basis of $V$: $\{v,v_i^{\perp}\}$ where $\{v_i^{\perp}\}$ is a basis of $v^{\perp}$. Calculating the action of $\chi(v)$ on this basis, we find:
\begin{equation}
\begin{split}
 \chi(v)v&=-v\cdot v\cdot v^{-1}=-v,\\
 \chi(v)v_i^{\perp}&=-v\cdot v_i^{\perp}\cdot v^{-1}
                   =v_i^{\perp}\cdot v\cdot v^{-1}
                   =v_i^{\perp}.
\end{split}
\end{equation}
In this computation\nomenclature{$\sQ$}{A subgroup of $\sG$}, we used the relation $v\cdot w=-w\cdot v-2\brak{v}{w}$ which is true for all $v$, $w$ in $V$. The action of $\chi(v)$ on this basis is thus to let unchanged three vectors and to change the sign of the fourth. This proves the claim.
\end{proof}

\begin{theorem}
\begin{equation}
                   \chi(\Gamma(p,q)^+)=\SO(p,q).
\end{equation}
\label{2102p1}
\end{theorem}

\begin{proof}
From corollary \ref{602c1}, and definition \ref{defgplus}, an element $s\in\Gamma(p,q)^+$ reads $s=cv_1\cdots v_{2r}$. Thus
\begin{equation}
 \det\chi(s)=\det\chi(v_1\cdots v_{2r})
            =\det\left[\chi(v_1)\ldots\chi(v_{2r})\right].
\end{equation}
 But we know that, for all $v_i$ in $V$, $det\chi(v_i)=-1$. So $\det\chi(s)=1$ and $\chi(\Gamma(p,q)^+)\subseteq \SO(p,q)$. As set, 
\[
  \Gamma(p,q)=\Gamma(p,q)^+\cup\Gamma(p,q)^-,
\]
but the lemma shows that $\det\chi(\Gamma(p,q)^-)=-1$ so, from theorem \ref{prop1001t1}, $\chi(\Gamma(p,q)^+)$ must be the whole $\SO(p,q)$.
\end{proof}

\begin{proposition}
The map $N$ takes values in $\eR$ and the formula
\begin{equation}
             N(x\cdot y)=N(x)N(y),
\end{equation}
holds for all $x$, $y\in\Gamma(p,q)$.
\label{proppourN}
\end{proposition}

\begin{proof}
We write as usual $x\in\Gamma(p,q)$ as $x=cv_1\cdots v_r$. So,
\begin{equation}
 N(x)=cv_1\cdots v_r\tau(\alpha(cv_1\cdots v_r))
     =\me{r}c^2v_1\cdots v_r\cdot v_r\cdots v_1.
\end{equation}
The first equality comes from the fact that $\alpha(cv_1\cdots v_r)=\me{r}cv_1\cdots v_r$. Now $N(x)\in\eR$ because $v_i\cdot v_i=-\brak{v_i}{v_i}\in\eR$ for all $i$. Hence the following hold:
\begin{equation}
\begin{split}
 N(x\cdot y)&=v\cdot y\cdot\tau(\alpha(v\cdot y))\\
            &=v\cdot y\cdot\tau(\alpha(y))\cdot\tau(\alpha(v))\\
            &=v\cdot N(y)\tau(\alpha(v))\\
            &=N(y)N(x).
\end{split}
\end{equation}
This is the claim.
\end{proof}


\begin{theorem}
We have the following isomorphism of groups
\[ 
  \Spin(p,q)=\SO_0(p,q).
\]
provided by the map $\chi$.
\end{theorem}

\begin{probleme}
	This result is wrong because of a double covering issue. The real proposition is the next one. I should try to merge the proofs.
\end{probleme}

\begin{proof}
Let $\{ e_1,\cdots,e_p,f_1,\cdots,f_p \}$ be a basis of $\eR^{p+q}$ where the $e_i$'s are time-like and the $f_j$'s are space-like.
Following the discussion at page \pageref{PgDisGeoConnSO}, we have
\[ 
  \SO(p,q)=\SO_0(p,q)\cup\xi \SO_0(p,q)
\]
where $\xi$ is defined as follows: $\xi e_1=-e_1$, $\xi f_1=-f_1$ and $\xi e_k=e_k$, $\xi f_k=f_k$ for $k\neq 1$. This element can be implemented as $\xi=\chi(g)$ for $g=e_1f_1$. It is easy to see that $g^{-1}=-f_1e_1$ and that $\tau(g)=f_1e_1$, so that $g\notin\Spin(p,q)$.

Is it possible to find another $h\in\Gamma(p,q)$ such that $\chi(h)=\xi$ ? If $\chi(a)=\chi(b)$ for $a$, $b\in\Gamma(p,q)$, then $a=rb$ for a certain $r\in\eR$. So we find that $h=g^{-1}/r$ is the general form of an element in $\Gamma(p,q)$ such that $\chi(h)=\xi$. This is an element of $\Spin(p,q)$ if and only if $\tau(h)=h^{-1}$, or $-e_1f_1/r=re_1f_1$ which has no solutions. We conclude that no element of $\Spin(p,q)$ is send on $\xi$ by $\chi$. So
\[ 
  \chi\big( \Spin(p,q) \big)\subset SO_0(p,q).
\]

\begin{probleme}
	I still have to prove the surjectivity of $\chi$ from $\Spin(p,q)$ to $\SO(p,q)$.
\end{probleme}

\end{proof}
\begin{theorem}
\begin{equation}	\label{EqchiSpinSO}	
             \chi(\Spin(p,q))=\SO_0(p,q)
\end{equation} 
where the index $0$ means the identity component.
\end{theorem}

\begin{proof}
Proposition \ref{prop1001p1}, theorem \ref{2102p1} and remark \ref{rem:spin_norm_u} show that the map $\dpt{\chi}{\Spin(p,q)}{\SO(p,q)}$ is a homomorphism with $\mathbb{Z}_2$ as kernel. We begin to prove that $\dpt{\chi}{\Spin(p,q)}{\SO_0(p,q)}$ is surjective. On the one hand, elements of $\Spin(p,q)$ satisfy one more condition than the ones of $\Gamma(p,q)^+$. Thus the algebra $\Spin(p,q)$ has codimension one in $\Gamma(p,q)^+$.

On the other hand, we know that $\SO(p,q)=\SO_0(p,q)\cup h\SO_0(p,q)$ where $h$ is the matrix such that $he_i=-e_i$ for $i=0,\ldots,3$. Since $\Spin(p,q)$ has codimension one in $\Gamma(p,q)^+$, there is at most one more generator in $\chi(\Gamma(p,q)^+)$ than in $\chi(\Spin(p,q))$ (because $\chi$ is a homomorphism). In order to prove this theorem, we just need to show that elements of $\chi(\Gamma(p,q)^+)$ which do not belong to $\chi(\Spin(p,q))$ is $h$.

Is is no difficult to see that $\chi(e_0\cdot e_1\cdot e_2\cdot e_3)e_i=-e_i$ for $i=0\ldots 3$: just write
$\chi(e_0\cdot e_1\cdot e_2\cdot e_3)e_i=e_0\cdot e_1\cdot e_2\cdot e_3\cdot e_i\cdot e_3^{-1}\cdot e_2^{-1}\cdot e_1^{-1}\cdot e_0^{-1}$ and use the commutation relations. An easy computation gives
$N(e_0\cdot e_1\cdot e_2\cdot e_3)=-1$; then this is not in $\Spin(p,q)$ by remark \ref{rem:spin_norm_u}.
\end{proof}

We write it by the exact sequence
\begin{equation}
 \xymatrix{
    \eZ_2  \ar@{^{(}->}[r] & \Sppq \ar[r]^{\chi} & \SO_0(p,q)
  } 
\end{equation}
we say that the group $\Spin(p,q)$ is a \defe{double covering}{double covering!of $\SO_{0}(p,q)$} of $\SO_0(p,q)$.

\begin{lemma}
If $\dpt{\pi}{\tX}{X}$ is a covering which satisfies

\begin{enumerate}
\item $X$ is path connected,
\item $\forall x\in X$, $\tX_x:=\pi^{-1}(x)$ is path connected in $\tX$ \emph{i.e.} for all $a$, $b\in \tX$, there exist a path in $\tX$ which joins $a$ and $b$,
\end{enumerate}
then $\tX$ is path connected.
\label{lem_cov_path_con}
\end{lemma}
\begin{proof}
If $\tx$ and $\ty$ are in $\tX$, we can suppose that $\pi(\tx)\neq\pi(\ty)$ (because if $\pi(\tx)=\pi(\ty)$, the second assumption gives the thesis). We define $x$ and $y$ as their projections: $x=\pi(\tx)$ and $y=\pi(\ty)$. Let $\gamma$ be a path such that $\gamma(0)=x$ and $\gamma(1)=y$, and $\tgamma$ be the lift of $\gamma$ in $\tX$ which contains $\tx$: $\tgamma(0)=\tx$ and $\pi(\tgamma(1))=\gamma(1)=y$. Then $\tgamma(1)$ lies in $\tX_y$. Therefore, we can consider $\gamma'$ which joins $\tgamma(1)$ and $\ty$.

So, $\gamma'\circ\tgamma$ is a path which contains $\tx$ and $\ty$.
\end{proof}


\begin{proposition}
 The group $\Spin(p,q)$ is connected.
\end{proposition}

\begin{proof}
We will prove that the covering $\dpt{\chi}{\Spin(p,q)}{\SO_0(p,q)}$ fulfils lemma \ref{lem_cov_path_con}. We just have to show that $\Spin(p,q)$ fulfills the second assumption of the lemma. First note that $\chi(\tx)=\chi(\ty)$ implies $\chi(\tx\ty^{-1})=e$, and then $\tx=\pm\ty$ because of proposition \ref{prop1001p1}. Since the other case is trivial, we can suppose $\tx=-\ty$.

It remains to prove that for every $g\in\Spin(p,q)$, there is a path in $\Spin(p,q)$ which joins $g$ and $-g$. The answer is given by the path $t\mapsto \gamma(t)g$ where
\[
\gamma(t)=\exp(te_1\cdot e_2)=\cos(t)(-1)+\sin(t)e_1\cdot e_2
\]
which satisfies $\gamma(0)=1$ and $\gamma(\pi)=-1$. 
\end{proof}

\begin{proposition}

The homomorphism $\tilde\rho$ restricts to a homomorphism $\tilde\rho\colon \Spin(p,q)\to \GL(\Lambda^+W)$.
\end{proposition}

\begin{proof}
An element in $\Spin(p,q)$ reads $s=cv_1\cdots v_{2r}$ and its image by $\tilde\rho$ is
\[ 
  \tilde\rho(s)=c\tilde\rho(v_1)\circ \cdots \circ\tilde\rho(v_{2r}).
\]
When one applies $\tilde\rho(v_1)$ to an element $\alpha\in\Lambda^kW$, one obtains a linear combination of an element of $\Lambda^{k-1}W$ and one of $\Lambda^{k+1}W$. The element $\tilde\rho(s)$ being an even composition of such maps, its transforms an element of $\Lambda^+W$ into an element of $\Lambda^+W$. 
\end{proof}

Notice that an element of $V$ ---no $V^{\eC}$--- is represented on $\Lambda^+W$ by complex matrices. This is not a problem. In the case of $\eR^{1,3}$, we have $\dim\Lambda^+W=2$ and thus 
\[ 
  \tilde\rho\big( \Spin(1,3) \big)\subset \GL(2,\eC).
\]
The following is the lemma 8.5 (page 57) of \cite{Michelson}.

\begin{lemma}
Let $\rho\colon \Cl(p,q)\to \Hom_{\eC}(E,E)$ be a representation of the Clifford algebra on a vector space $E$. If $p+q\geq 2$, then for all $s\in\Spin(p-1,q)\subset \Cl(p,q) $,
\[ 
  \det{}_{\eC}\big( \rho(s) \big)=\pm 1.
\]

\end{lemma}
\begin{proof}
No proof.
\end{proof}

\begin{theorem}
The representation $\tilde\rho$ provides a group isomorphism 
\[ 
  \Spin(1,3)\simeq \SL(2,\eC)
\]

\end{theorem}

\begin{proof}
In the case $p=2$, $q=3$, the lemma assures us that for each $s$ in the spin group, $\det\tilde\rho(s)=1$. Since $\Spin(1,3)$ is connected and the determinant function is continuous, we deduce that $\det\tilde\rho(s)\equiv 1$. This proves that $\tilde\rho\big( \Spin(1,3) \big)\subset \SL(2,\eC)$. The proposition \ref{PropUssGpGenere} thus implies that
\[ 
  \tilde\rho\big( \Spin(1,3) \big)=\SL(2,\eC),
\]
 but from $\Cl(1,3)$, the representation $\tilde\rho$ is yet injective. \emph{A forciori}, the representation $\tilde\rho$ is injective from $\Spin(1,3)$. This finishes the proof.
\end{proof}

\subsection{Redefinition of \texorpdfstring{$\Spin(V)$}{Spin(V)}}
%----------------------------------------------------------------

As it, this new definition only holds when $g$ is positive defined.

\begin{probleme}
	When we work with a signature $(p,q)$, maybe we only get the connected part. To be checked.
\end{probleme}

Let us take $v$, $x\in V$ with $g(v,v)=1$. We have 
\[ 
  -vxv^{-1}=-vxv=-2g(x,v)v+xv^2
		=x-2g(x,v)v\in V.
\]
The effect was to reverse the $v$ component of $x$; the map $x\mapsto -vxv^{-1}$ is $\sigma^v$. Now, when $\lambda\in U(1)$ and $w=\lambda v$, we also have that $x\mapsto -wxw^{-1}$ is $\sigma^v$. Now we look at $\chi(a)\colon x\mapsto \alpha(a)xa^{-1}$ with $a=w_{1}\ldots w_{r}$, a product of unitary vectors in $V^{\eC}$. Explicitly,
\[ 
  \chi(a)x=(-1)^{r} w_{1}\ldots w_{r}xw_{r}^{-1}\ldots w_{1}^{-1},
\]
a composition of reflexions in $V$. When $r$ is even, it is a rotation. We conclude that when $a$ is an even product of unitary vectors in $V^{\eC}$, then $\chi(a)\in \SO(V)$. Theorem \ref{CartanDieu} states that any rotation of $V$ is a composition of reflexions. So we define\nomenclature[G]{$\Spin^{c}(V)$}{A group related to $\Spin$}
\begin{equation}
\Spin^{c}(V)=\{ w_{1}\ldots w_{2k}\tq w_{j}\in V^{\eC},\,w_{j}^*w_{j}=1 \}\subset \CCliff^{0}(V),
\end{equation}
and $\chi\colon \Spin^{c}(V)\to \SO(V)$ is a surjective group homomorphism. The inverse in $\Spin^{c}(V)$ is given by
\[ 
  (w_{1}\ldots w_{2k})^{-1}=w_{2k}^*\ldots w_{1}^*=\overline{ w_{2k} }\ldots\overline{ w_{1} }.
\]
In the real case, proposition \ref{prop1001p1} says that $\ker\chi=\eR\invtible$. In the complex case we get  $\ker\chi=\eC\invtible$ and, when we look at $\ker\chi|_{\Spin^{c}(V)}$, we find
\begin{equation}
\ker\chi=U(1).
\end{equation}
Then we find the short exact sequence 
\begin{equation}
\xymatrix{%
   1 \ar[r]^-{\id}&U(1) \ar[r]^-{\id}&\Spin^{c}(V) \ar[r]^-{\chi}&\SO(V)\ar[r]^-{\id}&1.
}
\end{equation}
Let $u=w_{1}\ldots w_{2k}\in\Spin^{c}(V)$ with $w_{j}=\lambda_{j}v_{j}$ and $\lambda_{j}\in V$, so $\tau(u)=w_{2k}\ldots w_{1}$ and
\[ 
  \tau(u)u=w_{2k}\ldots w_{1}w_{1}\ldots w_{2k}
		=\lambda_{1}^{2}\ldots \lambda_{2k}^{2}\in U(1).
\]
This proves that $\tau(u)u$ is central in $\Spin^{c}(V)$. We define the homomorphism
\begin{equation}
\begin{aligned}
\nu \colon \Spin^{c}(V)&\to U(1) \\ 
u&\mapsto \tau(u)u. 
\end{aligned}
\end{equation}
This is a homomorphism because
\[ 
\begin{split}
  \nu(u_{1}u_{2})&=\tau(u_{1}u_{2})u_{1}u_{2}
		=\tau(u_{2})\underbrace{\tau(u_{1})u_{1}}_{\text{central}}u_{2}
		=\tau(u_{2})u_{2}\tau(u_{1})u_{1}\\
		&=\nu(u_{2})\nu(u_{1})
		=\nu(u_{1})\nu(u_{2}).
\end{split}  
\]
The map $\nu$ naturally restricts to $U(1)$ as
\[ 
  \nu(\lambda)=\lambda^{2}.
\]
The combined map $(\chi,\nu)\colon \Spin^{c}(V)\to \SO(V)\times U(1)$ has kernel $\{ \pm 1 \}$. We define\nomenclature[G]{$\Spin(V)$}{The spin group}
\begin{equation}  \label{eq_defSpindeux}
\Spin(V)=\ker\nu|_{\Spin^{c}(V)}.
\end{equation}

\begin{lemma}
This group is the same as the one defined in equation \eqref{defSpinun}. 
\end{lemma}

\begin{proof}
Let $u\in\Spin(V)$ (in the sense of equation \eqref{eq_defSpindeux}). The fact for $u$ to belongs to $\Spin(V)$ implies the two following:
\begin{enumerate}
\item $u\in\Spin^{c}(V)\Rightarrow u^*u=1$,
\item $u\in\ker\nu\Rightarrow \tau(u)u=1$.
\end{enumerate}
The second point says that $u^{-1}=\tau(u)$, which is a first good point to fit the first definition of $\Spin(V)$. Now we have to prove that $u\in\Gamma^{+}(V)$: $u$ must be invertible and $\chi(u)x$ must belongs to $V$ for all $x\in V$. These two points are contained in the definition of $\Spin^{c}(V)$.
\end{proof}
Let us see in the new definition how is $\chi\colon \Spin(V)\to \SO(V)$. On $\Spin^{c}(V)$, we have $\ker\chi=U(1)$, but on $\Spin(V)$ we require moreover $\tau(u)u=1$, thus an element of $\ker\chi$ in $\Spin(V)$ fulfils $\tau(\lambda)\lambda=1$, so that $\lambda=\{ \pm1 \}$. We conclude that $\ker\chi|_{\Spin(V)}=\{ \pm 1 \}$, and then that $\Spin(V)$ is a double covering of $\SO(V)$.\index{double covering!of $\SO(V)$}


\subsection{A few about Lie algebra}
%----------------------------------

\nomenclature[G]{$\spin(p,q)$}{Lie algebra of the group $\Spin(p,q)$}
\begin{proposition}
We have an isomorphism
\[ 
                    \spin(p,q)\simeq\so(p,q)
\]
between the Lie algebras of $\Spin(p,q)$ and $\SO(p,q)$.
\label{prop:spin_so}   
\end{proposition}

\begin{proof}
Using the second part of lemma \ref{Helgason5.1}, with the map $\dpt{\chi}{\Spin(p,q)}{\SO(p,q)}$, we find that $d\chi_e(\spin(p,q))=\so(p,q)$. Then we know (lemma \ref{1203r1}) that 
\[
	\so(p,q)=\spin(p,q)/\ker\,d\chi_e.
\]
On the other hand, the first part of the same lemma gives us that $\chi^{-1}(e)$ is a Lie subgroup of $\Spin(p,q)$ whose Lie algebra is $\ker\,d\chi_e$. But $\chi^{-1}(e)=\eZ_2$, so $\ker\,d\chi_e=\{0\}$.
\end{proof}

Let us now shortly speak about the Lie algebra of $\Gamma(p,q)^+$. A basis of $\Cliff(p,q)^+$ is \[\{1,\gamma_0\cdot\gamma_1,\gamma_0\cdot\gamma_1 ,\gamma_0\cdot\gamma_3
,\gamma_0\cdot\gamma_1\cdot\gamma_2\cdot\gamma_3  \}.\] Thanks to the anticommutation relations, we don't need $\gamma_1\cdot\gamma_2$ in the basis.

Remember that $\Gamma(p,q)^+$ is the set of the $x\in\Cliff^+(p,q)$ such that $x\cdot v\cdot\alpha(x^{-1})$ lies in $V$ for all $v\in V$. Let $x(t)$ be a path in $\Gamma(p,q)^+$ such that $x(0)=e$ and $\dot{x}(0)=X$. Differentiating the definition relation, we find
 \[
 \dot{x}\cdot v\cdot\alpha(x^{-1})|_0+x\cdot v\cdot(-)\alpha(\dot{x})|_0=X\cdot v-v\cdot X,
 \]
 therefore\nomenclature[G]{$\Lie{\Gamma(p,q)^+}$}{Algèbre de $\Gamma(p,q)^+$}
\[
  \Lie{\Gamma(p,q)^+}=\left\{X\in\Cliff^+(p,q)\textrm{ such that } X\cdot v-v\cdot X\in V,\,\forall v\in V\right\}.
\]

It is clear that $\eC$ is a subset of $\Lie{\Gamma(p,q)^+}$, and that $V$ is not. The following computation shows that $V\cdot V$ is a subset $\Lie{\Gamma(p,q)^+}$:
\[
         a\cdot b\cdot v-v\cdot a\cdot b=2\eta(v,a)b-2\eta(v,b)a.
\]
 We can also check that $V\cdot V\cdot V\cdot V\cap\Lie{\Gamma(p,q)^+}=\emptyset$. A basis of $\Lie{\Gamma(p,q)^+}$ is
\[
	\{ 1,e_{\alpha}\cdot e_{\beta}\tq \alpha<\beta \}
\]

 We know that $\ker[\dpt{\chi}{\Gamma(p,q)^+}{\SO(p,q)}]=\eR_0$. So the kernel of the restriction of $d\chi_e$ to $\Lie{\Gamma(p,q)^+}$ is the Lie algebra of $\eR_0$ (see lemma \ref{Helgason5.1}), which is $\eR$. Therefore, a basis of $\spin(p,q)$ is 
\[
	\{e_{\alpha}\cdot e_{\beta}\tq \alpha<\beta\}.
\]

\subsection{Grading \texorpdfstring{$\Lambda W$}{LW}}
%-------------------------------

We already know that $\Lambda W=\eC\oplus W\oplus\Lambda^2W$. This space can be written as \[\Lambda W =\Lambda W^+\oplus\Lambda W^-,\] with $\Lambda W^+=W$ and $\Lambda W^-=\eC\oplus\Lambda^2W$. The interest of such a decomposition lies in the definition of an action of $\Cliff^+(p,q)$ on $\Lambda W $. This action will be defined by $\dpt{\bullet}{\Cliff^+(p,q)\times\Lambda W }{\Lambda W }$, 
  \[
 x\bullet\alpha=\tilde\rho(x)\alpha
 \]
for any $x$ in $\Cliff^+(p,q)$ and any $\alpha$ in $\Lambda W $ (see definition \ref{defrt}).

\begin{proposition}
This action preserves the grading of $\Lambda W $:
\begin{equation}
\begin{split}
 \Cliff^+(p,q)\bullet\Lambda W^+&=\Lambda W^+\\
 \Cliff^+(p,q)\bullet\Lambda W^-&=\Lambda W^-.
\end{split}
\end{equation}

\end{proposition}
\begin{proof}
For $x\in\eC$, theses equalities are obvious. We have to check it for $x=e_i\cdot e_j$. Here, we will just check that $(e_1\cdot e_0)\bullet(v\wedge w)\in\Lambda W^+$. This follows from a simple computation:
\begin{equation}
\begin{split}
\tilde\rho(e_1)\tilde\rho(f_0+g_0)(v\wedge w)&=
                         \tilde\rho(f_1+g_1)\left[-\eta(g_0,v)w+\eta(g_0,w)v\right]\\
                    &=-\eta(g_0,v)f_1\wedge w+\eta(g_0,w)f_1\wedge v\\
                    &\quad+\eta(g_0,v)\eta(g_1,w)-\eta(g_0,w)\eta(g_1,v).
\end{split}
\end{equation}
\end{proof}

Since $\Spin(p,q)$ is a subgroup of $\Cliff^+(p,q)$, we can construct two new representation of $\Spin(p,q)$. These are $\dpt{\rho^{\pm}}{\Spin(p,q)\times\Lambda W ^{\pm}}{\Lambda W ^{\pm}}$,
\begin{equation}
\begin{split}
 \rho^-(s)w^-&=\tilde\rho(s)w^-,\\
 \rho^+(s)w^+&=\tilde\rho(s)w^+,
\end{split}
\end{equation}
for $w^{\pm}$ in $\Lambda W ^{\pm}$. This is no more than the fact that $\tilde\rho$ is reducible and that two invariant subspaces are $\Lambda W^+$ and $\Lambda W^-$.
\subsection{Clifford algebra for \texorpdfstring{$V=\eR^2$}{V=R2}}\label{cliffR2}
%----------------------------------------------

\subsubsection{General definitions}
%/////////////////////////////////

The whole construction can also be applied to $V=\eR^2$ with the Euclidean metric. This is our business now. We take the complex vector space $V^{\eC}$ and an orthonormal basis $\{e_1,e_2\}$. As before, we define
 \[
f_1=\frac{1}{2}(e_1+ie_2),\qquad g_1=\frac{1}{2}(e_1-ie_2).
\]
There are no difficulties to see that $Span(f_1)$ is a completely isotropic subspace\index{isotropic!subspace!in $\eR^{2}$} of $V^{\eC}$. Thus we define $W=\eC f_1$, $\Lambda W =\eC\oplus W$, $\Lambda W^+=\eC$, and $\Lambda W^-=W$\nomenclature{$\Lambda W^{\pm}$}{Spinor space}. The homomorphism $\dpt{\tilde\rho}{V^{\eC}}{\End(\Lambda W )}$\nomenclature{$\dpt{\tilde\rho}{(\eR^2)^{\eC}}{\End(\Lambda W )}$}{Spinor representation} in $\Lambda W $ is defined by
\begin{equation}
\begin{split}
 \tilde\rho(f_1)\alpha&=f_1\wedge\alpha,\\
 \tilde\rho(g_1)\alpha&=-i(g_1)\alpha,
\end{split}
\end{equation}
where $\alpha$ is any element of $\Lambda W $. In the basis $1=\begin{pmatrix}
1 \\
0
\end{pmatrix} $ and $f_1=\begin{pmatrix}
0 \\
1
\end{pmatrix} $, we easily find that
\[
 \tilde\rho(e_1)=\begin{pmatrix}
 0 & -\frac{1}{2} \\
 1 & 0
 \end{pmatrix},\quad\tilde\rho(e_2)=\begin{pmatrix}
 0 & -\frac{i}{2} \\
 -i & 0
 \end{pmatrix}.\]
For $c\in\eR$ we	 also have $\tilde\rho(c)f_1=cf_1$ and $\tilde\rho(c)1=c$, thus we assign the matrix $\begin{pmatrix}
c & 0 \\
0 & c
\end{pmatrix}$ to the number $c$.

As before, we define $\gamma_i=\sqrt{2}\tilde\rho(e_i)$. We immediately have $\gamma_1\gamma_2+\gamma_2\gamma_1=0$ and $\gamma_i\gamma_i=-2\mtu$, so that the $\gamma$'s satisfy the Clifford algebra for the euclidian metric.

For notational conveniences, it proves useful to make a change of basis so that we get
\begin{equation}\label{gammaR2}
\gamma_1=\begin{pmatrix}
0 & -1 \\
1 & 0
\end{pmatrix},\quad\gamma_2=-\begin{pmatrix}
0 & i \\
i & 0
\end{pmatrix}.
\end{equation}

The algebra $\Cliff(2)$\nomenclature[G]{$\Cliff(2)$}{Clifford algebra of $\eR^2$} is isomorphic to the algebra which is generated by direct sum $\Cliff(2)\simeq\eR\oplus\gamma_1\oplus\gamma_2\oplus\eR\gamma_1\gamma_2$. A general element of $\Cliff(2)$ can be written as $x\gamma_1+y\gamma_2+x'\eR+y'\gamma_1\gamma_2$. In the representation of $\tilde\rho$, a general element of $\Cliff(2)$ is therefore
\[\begin{pmatrix}
x'+iy' & x+iy \\
-x+iy & x'-iy'
\end{pmatrix},\] so that we can write the Clifford algebra of $\eR^2$ as\index{algebra!Clifford}
\[
\Cliff(2)=\left\{\begin{pmatrix}
 \alpha & \beta \\
 -\obeta & \oalpha
 \end{pmatrix}\,:\,\alpha,\beta\in\eC\right\}.
\]
The following four matrices provide a basis:
\begin{align}\label{pauli} 
1&=\begin{pmatrix}
1 & 0 \\
0 & 1
\end{pmatrix}, &i&=\begin{pmatrix}
-i & 0 \\
0 & i
\end{pmatrix},&j&=\begin{pmatrix}
0 & i \\
i & 0
\end{pmatrix},&k&=\begin{pmatrix}
0 & 1 \\
-1 & 0
\end{pmatrix}.
\end{align}
We can check that these matrices satisfies the quaternionic algebra\index{quaternion!algebra}\index{algebra!quaternion} :
\begin{equation}
\begin{split}
i^2&=j^2=k^2=-1\\
ij &=-ji=k,\\
jk &=-kj=i,\\
ki &=-ik=j.
\end{split}
\end{equation}
The algebra $\Cliff(2)=\eH$\nomenclature{$\eH$}{quaternionic algebra} is represented by $\tilde\rho$ on $\eC^2$ by the \defe{Pauli matrices}{pauli matrices} $1,i,j,k$ which are given by \eqref{pauli}.

\subsubsection{The maps \texorpdfstring{$\alpha$}{a} and \texorpdfstring{$\tau$}{t}}
%///////////////////////////////////////////////

What are the matrices which represent $V$ ? These are $\tilde\rho(e_1)$ and $\tilde\rho(e_2)$. Thus we can write $V=\Span_{\eR}\{\gamma_1,\gamma_2\}=\Span_{\eR}\{j,k\}$, or
\[
 V=\left\{\begin{pmatrix}
 0 & \xi \\
 -\oxi & 0
 \end{pmatrix}\,:\,\xi\in\eC\right\}.
\]

As before, $\alpha$ is the unique homomorphic extension to $\Cliff(2)$ of $-\id$ on $V$. From the definitions, we get $\alpha(j)=-j$, $\alpha(k)=-k$.
The extension present no difficult. For example: $\alpha(i)=\alpha(jk)=\alpha(j)\alpha(k)=jk=i$, but $\alpha(jk)=\alpha(i)$; then $\alpha(i)=i$. The same gives $\alpha(1)=1$.

The case of $\tau$ is treated in similar way. We find: $\tau(j)=j$, $\tau(k)=k$, $\tau(i)=-i$, $\tau(1)=1$.

Now, we can find the group $\gud$. The condition for $x\in\Cliff(2)$ to be in $\gud$ is $\alpha(x)yx^{-1}$ to belongs to $V$ for all $y\in V$. We put
\[ x=\begin{pmatrix}
\alpha & \beta \\
-\obeta & \oalpha
\end{pmatrix},\qquad\alpha(x)=\begin{pmatrix}
\alpha & -\beta \\
\obeta & \oalpha
\end{pmatrix}.\]
A typical $y$ in $V$ is
\[
 y=\begin{pmatrix}
 0 & \eta \\
 -\oeta & 0
 \end{pmatrix}.
\]
A few computation gives:
\[
 \alpha(x)yx^{-1}=\us{|\alpha|^2+|\beta|^2}\begin{pmatrix}
 \alpha\eta\obeta+\beta\oeta\oalpha & \alpha\alpha\eta-\beta\beta\oeta \\
 \obeta\obeta\eta-\oalpha\oalpha\eta & \eta\alpha\obeta+\oalpha\oeta\beta
 \end{pmatrix}.
\]
If we impose it to be of the form $\begin{pmatrix}
0 & \xi \\
-\oxi & 0
\end{pmatrix} $ for all $\eta\in\eC$, we get, for all $\eta\in\eC$, 
 $\Reel(\oalpha\beta\oeta)=0$, which implies $\oalpha\beta=0$. So we conclude:
\[
 \gud=\left\{\begin{pmatrix}
 \alpha & 0 \\
 0 & \oalpha
 \end{pmatrix}, \begin{pmatrix}
 0 & \beta \\
 -\obeta & 0
 \end{pmatrix}\,:\,\alpha,\beta\in\eC\textrm{ not both equals zero}\right\}.
\]
Be careful on a point: $\gud$ is the \emph{multiplicative} group generated by these two matrices, not the additive one.

\subsubsection{The spin group}
%////////////////////////////

It present no difficult to find that
\begin{equation}
 \gud^+=\left\{\begin{pmatrix}
 \alpha & 0 \\
 0 & \oalpha
 \end{pmatrix}\,:\,\alpha\neq 0\right\}.
\end{equation}
The \defe{spin group}{spin!group!on $\protect\eR^2$} is made of elements of $\gud^+$ which satisfy $\tau(x)=x^{-1}$. We know that
$\tau\begin{pmatrix}
\alpha & 0 \\
0 & \oalpha
\end{pmatrix} =\begin{pmatrix}
\oalpha & 0 \\
0 & \alpha
\end{pmatrix}$ and that $\begin{pmatrix}
\alpha & 0 \\
0 & \oalpha
\end{pmatrix}^{-1} =\us{\displaystyle\alpha\oalpha}\begin{pmatrix}
\oalpha & 0 \\
0 & \alpha
\end{pmatrix}$. Thus the condition \hbox{$\tau(x)=x^{-1}$} becomes $|\alpha|^2=1$. The first conclusion is that
\begin{equation}
                    \Spin(2)=U(1).
\end{equation}
A typical $s$ in $\Spin(2)$ is
\[s=e^{i\theta}=\begin{pmatrix}
e^{i\theta} & 0 \\
0 & e^{-i\theta}
\end{pmatrix}.\]

The next point is to see the action of $\Spin(2)$ on $V$.\index{action!of $\Spin(2)$ on $\eR^2$} The action of $s\in\Spin(2)$ on a vector $v\in V$ is still defined by $s\bullet v=\chi(s)v=\alpha(s)\cdot v\cdot s^{-1}$. More explicitly:
\begin{equation}
 \chi(s)v=\begin{pmatrix}
 e^{i\theta} & 0 \\
 0 & e^{-i\theta}
 \end{pmatrix} \begin{pmatrix}
 0 & z \\
 -\overline{z} & 0
 \end{pmatrix} \begin{pmatrix}
 e^{-i\theta} & 0 \\
 0 & e^{i\theta}
 \end{pmatrix}=\begin{pmatrix}
 0 & e^{2i\theta}z  \\
 -e^{-2i\theta}\overline{z} & 0
 \end{pmatrix},
\end{equation}
where the  matrix $\begin{pmatrix}
0 & z \\
\overline{z} & 0
\end{pmatrix} $ denotes the representation of the vector $v$ of $V$. This equality can be written $e^{i\theta}\cdot v=e^{2i\theta}v$. If we note $v=v_1+iv_2=\begin{pmatrix}
v_1 \\
v_2
\end{pmatrix} $, we get
\[ e^{2i\theta}\bullet v=\begin{pmatrix}
\cos 2\theta & -\sin 2\theta \\
\sin 2\theta & \cos 2\theta
\end{pmatrix}\begin{pmatrix}
v_1 \\
v_2
\end{pmatrix}. \]
Therefore, we can write
\[\chi(e^{i\theta})=\begin{pmatrix}
\cos 2\theta & -\sin 2\theta \\
\sin 2\theta & \cos 2\theta
\end{pmatrix}.\]

So $\chi$ projects $U(1)$ into $\SO(2)$ with a kernel $\eZ_2$, for this reason, we say that $U(1)$ is a \defe{double covering}{double covering!of $\SO(2)$} of $\SO(2)$. We note it
\begin{equation}
            \eZ_2\rightarrow U(1)\stackrel{\chi}{\rightarrow}\SO(2).
\end{equation}

\section{Clifford modules}  \label{susec_Cliffmodule}\index{Clifford!module}
%---------------------------

References: \cite{ResEtaDiracType,mellor}.

Let $M$ be a manifold. We denote by $\CCliff(M)$ the bundle whose fibre at $x\in M$ is the complex Clifford algebra of the metric $g_x$ : $\CCliff(M)_x=\CCliff(g_x)$. We define the important map
\begin{equation}
\begin{aligned}
 \gamma\colon \Gamma(M,\CCliff(M))&\to \oB(\hH) \\ 
\gamma(dx^{\mu})&\mapsto \gamma^{\mu}(x)  
\end{aligned}
\end{equation}
which can be extended to the whole Clifford algebra.

Let $V$ be a vector space endowed with a bilinear symmetric form. We consider $\Cliff(V)$, the corresponding Clifford algebra. A \defe{Clifford module}{Clifford!module} is a real vector space $E$ with a $\eZ_2$-graduation and a morphism 
\[ 
  \rho_E\colon \Cliff(V)\to \End(E)
\]
of $\eZ_2$-graded vector spaces. It is defined by a linear map $\rho_E\colon V\to \End(V)$ such that
\begin{equation}
\rho_E(v)\rho_E(w)+\rho_E(w)\rho_E(v)=B(v,w)\id
\end{equation}
for every $v$, $w\in E$. The element $\rho_E(x)v$ will often be denoted by $x\cdot v$ and the operation $\rho_E$ is the \defe{Clifford multiplication}{Clifford!multiplication}. The \defe{dual module}{dual module} $E^*$ is defined by $\rho_{E^*}(x)=\rho_E(x^t)^*$, i.e.
\begin{equation}
\langle \rho_{E^*}(x)\psi,v \rangle =(-1)^{| \psi | |x |}\langle \psi, \rho_E\big( \tau(x) \big)v\rangle 
\end{equation}
for every $\psi\in E^*$ and $v\in E$. Here

Let $\cA$ be a $\eZ_2$-graded subalgebra of $\Cliff(V)$ and $E_1$, a $\cA$-module. Then the space\nomenclature{$\Ind_{\cA}^{\Cliff(V)}(E_1)$}{Induced Clifford module}
\[ 
  E=\Ind_{\cA}^{\Cliff(V)}(E_1)=\Cliff(V)\otimes_{\cA}E_1
\]
has a structure of Clifford module, the \defe{induced module}{induced!Clifford module}. The tensor product $\otimes_{\cA}$ is the usual one modulo the subspace spanned by elements of the form 
\[ 
  x\otimes a\cdot y-xa\otimes y
\]
for every $x$, $a\in\Cliff(V)$ and $y\in E_1$. In a similar way, if $E$ is a complex vector space we have a notion of $\CCliff(V)$-module. 

Let $x\in\Cliff(V)$ be such that $x^2=1$. In that case the Clifford multiplication $\rho_E(x)$ decomposes $E$ in eigenspaces
\[ 
  E^{\pm}=\frac{ 1 }{2}\big( 1\pm\rho_E(x) \big)E.
\]

If $V$ is a $n$-dimensional vector space with an oriented orthonormal basis $\{ e_1,\ldots, e_n \}$, the algebra $\Cliff(V)$ has a \defe{volume element}{volume!element} $\omega=e_1e_2\ldots e_n$ which does not depend on the choice of the basis. The volume element squares to
\begin{equation}
\omega^2=(-1)^{n(n+1)/2}.
\end{equation}
In the complex case, we consider the complex vector space $V^{\eC}$ and the complex Clifford algebra $\CCliff(V)=\Cliff(V)\otimes_{\eR}\eC$, and the volume element is defined as
\begin{equation}
\omega_{\eC}=i^{[ (n+1)/2 ]}\omega.
\end{equation}
where $[x]$ is denotes the integer part of $x$. Performing a separate computation for $n$ even or odd, it is easy to see that in both case,
\begin{equation}
\omega_{\eC}^2=1.
\end{equation}
So in the complex case we always have an element in $\Cliff(V)$ which squares to $1$, and a $\CCliff(V)$-module $W$ always accepts a decomposition as $W^{\pm}=\frac{ 1 }{2}(1+\omega_{\eC})W$.

One says that a representation\index{representation!of Clifford algebra} $\rho$ of $\Cliff(V)$ on $W$ is \defe{reducible}{reducible!representation of Clifford} if there exists a splitting $W=W_1\oplus W_2$ such that $\rho(\Cliff(V))W_i\subset W_i$. If the representation is not reducible, it is said to be irreducible. Two representations $\rho_j\colon \Cliff(V)\to \End(W_j)$ are \defe{equivalent}{equivalence!representation of Clifford} if there exists a linear isomorphism $F\colon W_1\to W_2$ such that $F\circ\rho_1(x)\circ F^{-1}=\rho_2(x)$ for every $x\in\Cliff(V)$.

\begin{proposition}
The real Clifford algebra has
\[ 
 \begin{cases}
2&\text{if $n+1=0\mod 4$}\\
1&\text{otherwise}
\end{cases} 
\]
inequivalent irreducible representations. The complex Clifford algebra $\CCliff(V)$ has
\[ 
 \begin{cases}
2&\text{if $n$ is odd}\\
1&\text{if $n$ is even}
\end{cases} 
\]
inequivalent irreducible representations. 
\end{proposition}
\begin{proof}
No proof.
\end{proof}


If $M$ is a manifold, we denote by $\Cliff(M)=\Cliff(TM)$ the bundle whose fiber at $x$ is the Clifford algebras of $T_xM$. We consider an orthonormal basis $\{ e_i \}$ and if $\Sigma$ is a multi-index $\{ 1\leq\sigma_1,\ldots,\leq\sigma_t\leq m \}$, we pose $e_{\Sigma}=e_{\sigma_1}\ldots e_{\sigma_t}\in\Cliff(M)$. By convention, $e_{\emptyset}=1$. Since the elements $e_i$ are ordered, they provide an orientation:
\begin{equation}
d\vol=e_1\wedge\ldots\wedge e_m\in\Wedge^m(M).
\end{equation}
Since the map $e_{\sigma_1}\wedge\ldots\wedge e_{\sigma_t}\mapsto e_{\sigma_1\ldots e_{\sigma_t}}$ is an isomorphism between $\Cliff(M)$ and $\Wedge(M)$, we say that $d\vol\in\Cliff(M)$. Now we define 
\[ 
  \kappa=i^{-[(m+1)/2]}d\vol,
\]
which is nothing else that the volume form normalised in such a way that $\kappa^2=1$. If $m$ is even, it anti-commutes with $TM$, and if $m$ is odd, it commutes with $TM$.

Let $V$ be a $m$-dimensional real vector space, and $\CCliff(V)$, the corresponding complex Clifford algebra.
\begin{lemma}
Every $\CCliff(V)$-module accepts an unique decomposition as sum of irreducible representations as follows
\begin{enumerate}
\item if $m=2n$, there exists one and only one irreducible $\CCliff(V)$-module $\Delta$ and $\dim(\Delta)=2n$,
\item if $m=2n+1$, we have two inequivalent irreducible modules $\Delta_{\pm}$ with $\gamma(\kappa)=\pm 1$ on $\Delta_{\pm}$ and $\dim(\Delta_{\pm})=2^n$.
\end{enumerate}
\end{lemma}
\begin{proof}
No proof.
\end{proof}

Let $V$ be a vector bundle over $M$. A structure of $\Cliff(M)$-module over $V$ is a morphism of unital algebra $\gamma\colon \Cliff(M)\to \End(V)$. When one has a basis $\{ e_i \}$ of $V$, we pose $\gamma_i=\gamma(e_i)$. The following lemma is the lemma 1.2 of \cite{ResEtaDiracType}. 
\begin{lemma}			\label{LemGammaBaseConstant}
Let $V$ be a $\Cliff(V)$-module and $\{ e_i \}$, an orthonormal basis for $TM$ on a contractible open set $V$. Then there exists a local frame for $V$ such that the matrices $\gamma(e_i)$ are constant.
\end{lemma}
We also define $\gamma^i=\gamma(dx^i)=g^{ij}\gamma_j$. One easily proves that
\begin{equation}
\gamma^i\gamma^j+\gamma^j\gamma^i=-2g^{ij}
\end{equation}
where $(g^{ij})$ is the inverse matrix of $(g_{ij})$. If the endomorphisms $\gamma_i$ are constant in the basis $\{ e_i \}$, then the endomorphisms $\gamma^i$ are constant in the basis $\{ f_i=g_{ki}e_k \}$.



\section{Spin structure}	\label{sec:spin_str}
%++++++++++++++++++++++++

We consider a (pseudo-)Riemannian manifold $(M,g)$ with metric signature $(p,q)$, and $\SO(M)$, its frame bundle; it admits a $\SO(p,q)$-principal fibre bundle structure which is well defined by the metric $g$ (see \ref{subsubsecframebundle}).

\begin{definition}
We say that $(M,g)$ is a  \defe{spin manifold}{spin!manifold} if there exists a $\Sppq$-principal bundle $P$ over $M$ and a principal bundle homomorphism $\dpt{\varphi}{P}{\SO(M)}$ which induced covering for the structure groups is $\chi$, i.e.
$\varphi(\xi\cdot s)=\varphi(\xi)\cdot\chi(s)$. A choice of $P$ and $\varphi$ is a \defe{spin structure}{spin!structure} on $M$.
\label{defvarspin}
\end{definition}
\[
\xymatrix{ \Sppq \ar@{~>}[r]	& P \ar[rr]^-{\displaystyle\varphi} \ar[rd]_{\displaystyle\pi} && \SO(M) \ar[ld]^{\displaystyle p}&\SO(p,q) \ar@{~>}[l]  \\& & M }
\]
The wavy arrows mean ``structural group of''.

\begin{remark}
When we will use the concept of spin structure in the physical oriented chapters, we will naturally use $\SLdc$ as group instead of $\Sppq$. The isomorphism $\SLdc\simeq\Sput$ is proved in \cite{Michelson}. A physical motivation of such a structure is given at page \pageref{pg_spinenphyz}.
\end{remark}

\subsection{Example: spin structure on the sphere \texorpdfstring{$S^2$}{S2}}
%----------------------------------------------------------------

It is no difficult to see that $\SO(S^2)\simeq \SO(3)$. Indeed, each element of $\SO(S^2)$ is described by three orthonormal vectors: one which point to an element $x$ of $S^2$ and two which gives a basis of $T_xS^2$. The action $\SO(3)\times S^2\to S^2$ is transitive, and the stabilizer of any element is $\SO(2)$.

We define $\dpt{\alpha}{\SO(3)/\SO(2)}{S^2}$ by $\alpha(g\SO(2))=g$. Proposition \ref{propHelgason4.3} shows that $\alpha$ is a diffeomorphism. Then
\[
                S^2=\frac{\SO(3)}{\SO(2)}.
\]

On the other hand,  we know that
\begin{eqnarray}\label{explsu2} T_eSU(2)=su(2)=\left\{\begin{pmatrix}
ix & \xi \\
-\oxi & -ix
\end{pmatrix}\,:\,\xi\in\eC,x\in\eR\right\}.
\end{eqnarray}
 It is a classical result that $\mathfrak{su}(2)\simeq\eR^3$ not only as set but also as metric space with the identification
\[
\langle X,Y\rangle=-\frac{1}{2}\tr(XY),
\]
 for all $X$, $Y\in su(2)$. As we are in matrix groups, we know (see \cite{Lie} to get more details) that $Ad_xY=xYx^{-1}$. In our case, this gives the formula
\[
                \langle Ad(g)X,Ad(g)Y\rangle=\langle X,Y\rangle.
\]
We can now state the result for $S^2$.

 \begin{proposition}
The manifold $S^2$ with the usual metric induced from $\eR^3$ admits the following spin structure:
\begin{eqnarray}\label{spins2}
\xymatrix{ \Spin(2)\ar@{~>}[r] &SU(2) \ar[rr]^{\displaystyle \varphi=Ad} \ar[rd]_{\displaystyle U(1)}^{\displaystyle\pi} && \SO(3) \ar[ld]^{\displaystyle \SO(2)}_{\displaystyle p} \\& & S^2 },
\end{eqnarray} where the arrow
$\xymatrix{X \ar[r]^{f}_G & Y }$ means that $G$ is the kernel of the map $\dpt{f}{X}{Y}$.
\end{proposition}
\begin{proof}
 First, let us precise the concept of frame bundle for $S^2$, and how it is well described by $\SO(3)$. Let $\{e_1,e_2,e_3\}$ be the canonical basis of $\eR^3$. To $A\in \SO(3)$, we make correspond the basis $\{Ae_2,Ae_3\}$ at the point $Ae_1$ of $S^2$. The projection $\dpt{p}{\SO(3)}{S^2}$ is then defined by $p(A)=Ae_1$. It is clear that we will  define the map $\dpt{\pi}{SU(2)}{S^2}$ in the same way: $\pi(U)=p(Ad(U))$.

For the rest of the demonstration, we will use the ``$su(2)$ description''\ of $\eR^3$ given by \eqref{explsu2} with $\xi=y+iz$.

Now, let us show that $\dpt{\pi}{SU(2)}{S^2}$ is a $\Spin(2)$-principal bundle. Since we had already shown that $\Spin(2)\simeq U(1)$, we define the right action of $\Spin(2)$ on $SU(2)$ by right multiplication: $U\cdot s=Us$ with $s=\begin{pmatrix}
e^{i\theta} & 0 \\
0 & e^{-i\theta}
\end{pmatrix}$. It is clear that $\pi(Us)=\pi(U)$:
\begin{equation}
 Ad(Us)e_1=(Us)\begin{pmatrix}
 1 \\
 0 \\
 0
 \end{pmatrix}s^{-1} U^{-1}=Us
 \begin{pmatrix}
 i&0\\
 0&-i
 \end{pmatrix}s^{-1} U^{-1},
\end{equation}
because $\begin{pmatrix}
i&0\\
0&-i
\end{pmatrix}$ is the vector $e_1$ in the ``$su(2)$ description''\ of $\eR^3$.

In order for $\dpt{\pi}{SU(2)}{S^2}$ to be a $\Spin(2)$-principal bundle, we still need to show that for all $x\in S^2$,
\[
   \pi^{-1}(x)=\left\{\xi\cdot g\tq g\in\Spin(2)\,\forall\xi\in\pi^{-1}(x)\right\}.
\]
Take $A$, $B\in\pi^{-1}(x)$, i.e. $Ae_1=Be_1=x$. We need to find a $s\in\Spin(2)$ such that
\begin{eqnarray}
 \label{1603r3} A=B\cdot s.
\end{eqnarray}
The matrices $A$ and $B$ are such that
\begin{eqnarray}\label{1603r1}
 B^{-1} A\begin{pmatrix}
 i&0\\
 0&-i
         \end{pmatrix}A^{-1} B=\begin{pmatrix}
 i&0\\
 0&-i
         \end{pmatrix}.
\end{eqnarray}
This implies that $B^{-1} A\in\Spin(2)$. As $Ad$ is surjective from $SU(2)$ into $\SO(3)$, a general $C$ in $\SO(3)$ which acts on $e_1$ can be written $Ue_1U^{-1}$ for $U\in SU(2)$ such that $Ad(U)=C$. The condition \eqref{1603r1} becomes
\[
\begin{pmatrix}
\alpha&\beta\\
-\obeta&\oalpha
\end{pmatrix}
\begin{pmatrix}
i&0\\
0&-i
\end{pmatrix}
\begin{pmatrix}
\oalpha&-\beta\\
\obeta&\alpha
\end{pmatrix}=
\begin{pmatrix}
i&0\\
0&-i
\end{pmatrix},
\]
which implies $\alpha=e^{i\theta}$, $\beta=0$. Then $B^{-1} A$ belongs to $\Spin(2)$, and $s=B^{-1} A$ fulfills the condition \eqref{1603r3}.

What about the induced covering for the structural groups ? The structural group of $\dpt{\pi}{SU(2)}{S^2}$ is $\Spin(2)$, while the one of $\dpt{p}{\SO(3)}{S^2}$ is $\SO(2)$. Indeed, for each $x\in S^2$, $\SO(2)$ acts on $T_xS^2$, leaving $x$ unchanged. We have the following associations:
\[
         U\in SU(2)\stackrel{\varphi}{\longrightarrow}A\in \SO(3),
\]
the matrix $A$ being defined by $A\cdot X=UXU^{-1}$. For $s\in\Spin(2)$ we of course also have
\[
         Us\in SU(2)\stackrel{\varphi}{\longrightarrow}As\in \SO(3),
\]
with $As\cdot X=UsXs^{-1} U^{-1}$. As we act by $\Spin(2)$ on $SU(2)$, in the fibres of $\SO(3)$, the action of $\Spin(2)$ is --via $\varphi$-- the composition with $X\to sXs^{-1}$. But this is exactly $\chi(s)X$ because $\alpha(s)=s$, since $s\in\Spin(2)$.
\end{proof}

\subsection{Spinor bundle}
%--------------------------

Let us take once again the spin structure on the (pseudo-)Riemannian manifold $(M,g)$:
\[
  \xymatrix{ \Sppq \ar@{~>}[r]& P \ar[rr]^-{\displaystyle\varphi}
   \ar[rd]_{\displaystyle\pi} && \SO(M) \ar[ld]^{\displaystyle p}&\SO(p,q) \ar@{~>}[l]
   \\& &   M }
\]
with $\varphi(\xi\cdot g)=\varphi(\xi)\cdot\chi(g)$.

Let us define $S=\Lambda W $, and $\mS=P\times_{\rho}S$. Take $\dpt{\rho}{\Sppq\times\mS}{\mS}$, $\rho(g,s)=\tilde\rho(g)s$, where $\tilde\rho$ is the spinor representation of $\Sppq$ on $S$. We also have
$\dpt{\chi}{\Sppq}{\SO_0(p,q)}$, $\chi(g)v=\alpha(g)\cdot v\cdot g^{-1}$, with $\alpha(g)=g$ for $g\in\Sppq$.

The \defe{spinor bundle}{spinor!bundle} is the associated bundle
\begin{equation}
                   \mS=P\times_{\rho}S\to M
\end{equation}
A \defe{spinor field}{spinor!field} is an element of $\Gamma(\mS)$, the space of section of the spinor bundle.

On $\SO(M)$, we look at a connection $1$-form $\alpha\in\Omega^1(\SO(M),so(\eR^m))$,
and, if $T(M)$ is the tensor bundle over $M$, we define a covariant derivative $\dpt{\nabla^{\alpha}}{\cvec(M)\times T(M)}{T(M)}$ by
 \[
             \widehat{\nabla^{\alpha}_X s}=\overline{X}\hat{s},
\]
 for any $s\in T(M)$. See theorem \ref{tho_nablaE}, and the fact that $T(M)$ can be see as an associated bundle; it is explicitly done for $\cvec(M)$ at page \pageref{equivvec}.

As seen in point \ref{subsubsec_levi}, an automatic property of this connection is $\nabla^{\alpha} g=0$ if $g$ is the metric of $M$. The \defe{Levi-Civita connection}{connection!Levi-Civita}\index{connection!Levi-Civita} is the unique\footnote{We will not prove unicity.} such connection which is torsion-free: $T^{\nabla^{\alpha}}=0$.


\begin{proposition}
The $1$-form $\talpha=\varphi^*\alpha\in\Omega^1(P,so(\eR^{m}))$ defines a connection on $P$. See definition \ref{defconnform} and theorem \ref{tho_nablaE}.
\end{proposition}

\begin{proof}
Let us denote by $R_g$ the right action of $g\in\Sppq$ on $P$ (\emph{id est} $R_g\xi=\xi\cdot g$), and by $R_u^{\SO(M)}$ the right action of $u\in\Sopq$ on $\SO(M)$.
We  have to check the usual two conditions of a connection.

\subdem{First condition}
The first one is:
\[
            (R_g^*\talpha)_{\xi}(\Sigma)=Ad(g^{-1})(\talpha_{\xi}(\Sigma)),
\]
for all $\xi\in P$, and $\Sigma\in T\bxi P$. In order to check this, we first remark that $\varphi\circ R_g=R_{\chi(g)}^{\SO(M)}\circ\varphi$. Indeed, for all $\xi\in P$, definition \ref{defvarspin} gives us $\varphi(R_g\xi)=\varphi(\xi\cdot g)=\varphi(\xi)\cdot\chi(g)$.  With this, we can make the following computation:
\begin{equation}\label{1603r4}
\begin{aligned}
R_g^*\talpha&=R_g^*\varphi^*\alpha=(\varphi\circ   R_g)^*\alpha	=(R_{\chi(g)}^{\SO(M)}\circ\varphi)^*\alpha\\
            &=\varphi^*R_{\chi(g)}^{\SO(M)*}\alpha=\varphi^*(Ad(\chi(g)^{-1})\circ\alpha).
\end{aligned}
\end{equation}
The last equality comes from the fact that $\alpha$ is a connection $1$-form. As we are in matrix groups, we have $Ad(g)x=gxg^{-1}$, so
\begin{equation}
   [Ad(\chi(g))x]v=[\chi(g) x \chi(g)^{-1}]v
                  =\chi(g)[xg^{-1} vg]
                  =gxg^{-1}.
\end{equation}
In the first line, the product is the usual matrix product which can be seen as operator composition.

But $(Ad(g)x)v=gxg^{-1} v$. Then $Ad(g)=Ad(\chi(g))$, if we identify $\sppq\simeq\sopq$ by proposition \ref{prop:spin_so}. Moreover, the action of $Ad$ is linear, so it commutes with $\varphi^*$. With these remarks, we can continue the computation \eqref{1603r4}:
\begin{equation}
 \varphi^*(Ad(\chi(g)^{-1})\circ\alpha)=\varphi^*(Ad(g^{-1})\circ\alpha)
                                  =Ad(g^{-1})\circ\varphi^*\alpha
                                  =Ad(g^{-1})\circ\talpha.
\end{equation}
This proves the first condition.

\subdem{Second condition}
The second one is $\talpha(A^*\bxi)=-A$ with the definition \eqref{defastar}. This is also a computation. First remark
\[
 \talpha\bxi(A\bxi^*)=(\varphi^*\alpha)\bxi(A^*\bxi)=\alpha_{\varphi(\xi)}(\varphi_{*\xi}A^*\bxi).
\]
 We compute $\varphi_{*\xi}A^*$ with lemma \ref{lemsur5d}:
\begin{equation}
\begin{split}
 \varphi_{*\xi}A^*&=\dsdd{\varphi(\xi\cdot\exp -tA)}{t}{0} =\dsdd{(R_{\chi(\exp -tA)}^{\SO(M)}\circ\varphi)(\xi)}{t}{0}\\
              &=\dsdd{\varphi(\xi)\cdot\chi(\exp -tA)}{t}{0}=\dsdd{\varphi(\xi)\cdot\exp(-td\chi_eA)}{t}{0}=(d\chi_eA)^*_{\varphi(\xi)}.
\end{split}
\end{equation}
But $d\chi_e=\id_{so(p,q)}$, thus $\varphi_{*\xi}A^*=A^*_{\varphi(\xi)}$. The whole makes that:
\[
\talpha\bxi(A^*\bxi)=\alpha_{\varphi(\xi)}(\varphi_{*\xi}A^*\bxi)=\alpha_{\varphi(\xi)}(A^*_{\varphi(\xi)})=-A.
\] 
This completes the proof.
\end{proof}

\begin{definition}
This connection $1$-form on $P$ is called the \defe{spinor connection}{spinor!connection}. It gives us a covariant derivative on any associated bundle and in particular on the spinor bundle, $\dpt{\tnab}{\cvec(M)\times\Gamma(\mS)}{\Gamma(\mS)}$.
 \label{spinconn}
\end{definition}
\nomenclature[D]{$\dpt{\tnab}{\cvec(M)\times\Gamma(\mS)}{\Gamma(\mS)}$}{Covariant derivative for the spinor connection}

\begin{proposition}
If $X$, $Y\in\cvec(M)$ are such that $X_x=Y_x$, then for all $s\in\Gamma(\mS)$,
\[
              (\tnab_Xs)(x)=(\tnab_Ys)(x).
\]
 \label{2303p1}
\end{proposition}
\begin{proof}
We just have to show that for all vector field $Z$ such that $Z_x=0$, $(\tnab_Zs)(x)=0$. Such a $Z$ can be written as $Z=fZ'$ for a function $f$ on $M$ which satisfies $f(x)=0$. We have:
\[
\tnab_Zs=\tnab_{fZ'}s=f\tnab_{Z'}s,
\]
which is obviously zero at $x$.
\end{proof}

Let $x\in M$ and $\{e_{\alpha\,x}\}$ be an orthonormal basis of $T_xM$. We can extend it to $\{e_{\alpha}\}$, a local basis field around $x$ such that $e_{\alpha}$ is a section of the frame bundle (in other words, we ask the extension to be smooth). The claim of proposition \ref{2303p1} is that $\tnab_{e_{\alpha}}(x)$ is an element of $\mS_x$ which doesn't depend on the extension.

\section[Dirac operator]{Dirac operator\protect\quad{\Huge\Smiley}}		\label{applgamma}
%---------------------------------------------------------------------

\subsection{Preliminary definition}
%----------------------------------

Let $M$ be a $m$-dimensional (pseudo)Riemannian manifold with its spin structure 
\[
\xymatrix{ \Sppq \ar@{~>}[r]& P \ar[rr]^{\displaystyle\varphi} \ar[rd]_{\displaystyle\pi} && \SO(M) \ar[ld]^{\displaystyle p}&\SO(p,q) \ar@{~>}[l]  \\& & M }
\]
where $\varphi$  satisfies $\varphi(\xi\cdot g)=\varphi(\xi)\cdot\chi(g)$.

Recall that for any vector space, one can see $\End{V}=V^*\otimes V$ with the definition $(v^*\otimes v)w=(v^*w)v$. This allows us to define an action of $\Sppq$ on $\End{S}$ by defining an action of $\Sppq$ on $S$ and $S^*$ separately. We know the action 
\begin{equation}
\begin{aligned}
 \Spin(p,q)\times S&\to S \\ 
(g,v)&\mapsto \tilde\rho(g)v,
\end{aligned}
\end{equation}
and as action on $S^*$, we take the dual one
\begin{equation}
\begin{aligned}
 \Spin(p,q)\times S^*&\to S^* \\ 
 g\cdot\alpha&= \alpha\circ\tilde\rho(g^{-1}) 
\end{aligned}
\end{equation}
for all $g\in\Spin(p,q)$ and $\alpha\in S^*$.

Now we can make the following computation with $g\in\Sppq$, $\alpha\in S^*$ and $v\in S$, using the fact that $\tilde\rho$ is linear:
\begin{equation}
\begin{split}
[g\cdot(\alpha\otimes v)]w&=[(\alpha\circ\tilde\rho(g^{-1}))w]\tilde\rho(g)v\\
                          &=\tilde\rho\left([(\alpha\circ\tilde\rho(g^{-1}))w]g\right)v\\
                          &=\big[\tilde\rho(g)\circ(\alpha\otimes v)\circ\tilde\rho(g^{-1})\big]w.
\end{split}
\end{equation}
Then we write the action of $\Sppq$ on $\End{S}$\index{action!of $\Sppq$ on $\End{S}$} by ($A\in\End S$)
\begin{equation}
     g\cdot A=\tilde\rho(g)\circ A\circ\tilde\rho(g^{-1}).                          \label{actspin}
\end{equation}
Notice that this definition is the one required in condition \eqref{equivA}.

The tangent bundle $T_xM$ is given with a metric $g_x$. As usual, we build $S_x=\Lambda W _x$, a completely isotropic subspace of $T_xM$ with respect to the metric $g_x$, and a representation
\[ 
\tilde \rho_x\colon T_xM\to \End(\Lambda W_x)
\]
The first step in the definition of $\gamma(X)$ is to build $\dpt{\ha_X}{P}{\End(\Lambda W )}$ setting\footnote{See subsection \ref{equivvec} for the definition of $\hX$.} $\ha_X(p)=\tilde\rho(\hX_{\varphi(p)})$.

\begin{lemma}
The function $\hat a$ is equivariant, i.e. it satisfies
\begin{equation}
     \ha_X(p\cdot g)=g^{-1}\cdot\ha_X(p)                             \label{equivaX}
\end{equation}
for all $g\in\Sppq$.
\end{lemma}

\begin{proof}
It is no more than a simple computation using the equivariance of $\hX$. Indeed:
\begin{equation}
\begin{split}
 \ha_X(p\cdot g)&=\tilde\rho(\hX_{\varphi(p\cdot g)})=\tilde\rho(\hX_{\varphi(p)\chi(g)})=\tilde\rho(\chi(g^{-1})\cdot\hX_{\varphi(p)})\\
		&=\tilde\rho(g^{-1}\cdot\hX_{\varphi(p)}\cdot g)=\tilde\rho(g^{-1})\circ\tilde\rho(\hX_{\varphi(p)})\circ\tilde\rho(g)\\
                &=g^{-1}\cdot\ha_X(p).
\end{split}
\end{equation}
In the fourth line, the dots mean the Clifford product, and the last equality comes from the definition of the action \eqref{actspin} of $\Sppq$ on $\End{S}$.
\end{proof}

From the discussion of section \ref{sec_fnequiv}, the function $\dpt{\ha_X}{P}{\End{S}}$ defines a section $\dpt{a_X}{M}{\End{\mS}}$. We define $\dpt{\gamma}{\cvec(M)}{\End{\Gamma(\mS)}}$ by
\nomenclature[D]{$\dpt{\gamma}{\cvec(M)}{\End{\Gamma(\mS)}}$}{A key ingredient for Dirac operator}
\begin{equation}		\label{EqDefgammax}
                       \gamma(X)=a_X.
\end{equation}
We immediately have
\[
                     \widehat{\gamma(X)}(p)=\tilde\rho(\hX_{\varphi(p)})
\]
for any $p\in P$. If we define
\begin{equation}\label{3103r1}
  \widehat{\gamma\cdot a_X}(p)=\widehat{\gamma(X)}(p),
\end{equation}
the map $\gamma$ can be seen as an action on the section of $\mS$. Indeed, $\widehat{\gamma\cdot s}_X$ is an equivariant function:
\begin{equation}
\begin{split}
 \hat{\gamma}(p\cdot g)(\ha_X(p\cdot g))
                  &=\rho(g)^{-1}\hat{\gamma}(p)\rho(g)\rho(g^{-1})\ha_X(p)\\
                  &=\rho(g)^{-1}\hat{\gamma}(p)\ha_X(p)\\
                  &=\rho(g^{-1})\widehat{\gamma\cdot a_X}(p),
\end{split}
\end{equation}
 so that
\[
    \widehat{\gamma\cdot a_X}(p)=\rho(g^{-1})\widehat{\gamma\cdot a_X}(p).
\]

The map $\dpt{\widehat{\gamma\cdot a_X}}{P}{\End{\Lambda W }}$ defined by \eqref{3103r1} is equivariant, and thus defines a section
$\gamma\cdot a_X\in\Gamma(\mS)$, as seen in the section \ref{sec_fnequiv}.

\subsection{Definition of Dirac}
%-------------------------------

If we consider a basis $\{e_{\alpha}\}$ of $TM$, \emph{i.e.} $m$ sections $\dpt{e_{\alpha}}{M}{TM}$ such that for all $x$ in $M$, the set $\{e_{\alpha x}\}$ is a basis of $T_xM$, we note $\gamma^{\alpha}:=\gamma(e_{\alpha})\in\End(\mS)$.

\begin{remark}  
This is not always globally possible. The example of the sphere is given in subsection \ref{subsec_DimofModule}. 
\label{rem_secnoglobal}
\end{remark}

For any $s\in\Gamma(\mS)$, we consider the local\footnote{Extensions of $e_{\alpha}$ do not always globally exist, see remark \ref{rem_secnoglobal}.} section $\psi$ of $\mS$ given by
\[
    \psi(x)=
   \sum_{\alpha\beta}g_x(e_{\alpha},e_{\beta})\gamma_x\hbeta(\tnab_{e_{\alpha}}s)(x).
\]

For each $x\in M$, take a $A_x$ in\footnote{By $\SO(g_x)$, we mean the set of all the matrix $A$ such that $A^tg_xA=g$; $A_x$ is an isometry of $(T_xM,g_x)$. In other words, we consider $A$ as a section of what we could call the ``isometry bundle''.} $\SO(g_x)$, and consider the new basis $e'_{\alpha}=A_{\alpha}^{\phantom{\alpha}\beta}e_{\beta}$. As $A$ is an isometry, $g_x(e'_{\alpha},e'_{\beta})=g_x(e_{\alpha},e_{\beta})$; and since $\tilde\rho$ is linear, $\gamma_x'^{\alpha}=\tilde\rho_x(e'_{\alpha x})=A_{\alpha}^{\phantom{\alpha}\beta}\tilde\rho(e_{\beta x})=A_{\alpha}^{\phantom{\alpha}\beta}\gamma_x\hbeta$. In the new basis, the section reads:
\begin{equation}
\begin{split}
   \psi(x)&=\sum_{\alpha\beta\eta\sigma}g_x(e_{\alpha},e_{\beta})
                A_{\beta}^{\phantom{\beta}\sigma}\gamma_x^{\sigma}
                (\tnab_{A_{\alpha}^{\phantom{\alpha}\eta}e_{\eta}}s)(x)\\
          &=\sum_{\alpha\beta\eta\sigma}(A^t)\heta_{\phantom{\eta}\alpha}
                  g_{\alpha\beta}(x)A_{\beta}^{\phantom{\beta}\sigma}
                  \gamma_x^{\sigma}(\tnab_{e_{\eta}}s)(x)\\
          &=\sum_{\eta\sigma}g_x(e_{\eta},e_{\sigma})\gamma_x^{\sigma}(\tnab_{e_{\eta}}s)(x),
\end{split}
\end{equation}
where we used the fact that $A^tgA=g$ and that all the $A_{\alpha}^{\phantom{\alpha}\beta}$ are $C^{\infty}$ functions on $M$, so that
$\tnab_{A_{\alpha}^{\phantom{\alpha}\beta}X}=A_{\alpha}^{\phantom{\alpha}\beta}\tnab_X$.  This shows that $\psi(x)$ doesn't depend on the choice of the basis, so it defines a section from the data of $s$ alone.


The \defe{Dirac operator}{dirac!operator!on $(M,g)$, a spin manifold}\nomenclature[D]{$\Dir$}{Dirac operator} $\Dir\colon \Gamma(\mS)\to \Gamma(\mS)$ acting on a spinor field is defined by 
\begin{equation}\label{dirac}
(\Dir s)(x)=g_x(e_{\alpha},e_{\beta})\gamma_x\hbeta(\tnab_{e_{\alpha}}s)(x).
\end{equation}

\begin{proposition}
If the field of basis $e_{\alpha}\in\cvec(M)$ is everywhere an orthonormal basis, the Dirac operator reads
\begin{equation}
(\Dir s)(x)=g_{\alpha\beta}\gamma^{\alpha}(\tilde\nabla_{e_{\beta}}s)(x)
\end{equation}
where $\gamma^{\alpha}$ is a constant numeric matrix acting on $\Lambda W$.
\end{proposition}

\begin{proof}
The building of the Dirac operator begins by considering the vector space $T_xM$ endowed with the metric $g_x$; then the spinor representation $\tilde\rho_x\colon T_xM\to \End(\Lambda W_x)$ where $\Lambda W_x$ is build from isotropic vectors of $T_xM$ is defined. If the vector fields $e_{\alpha}\in\cvec(M)$ are everywhere orthonormal for the metric $g$, then we have the matricial equality
\begin{equation}
	\tilde\rho_x\big( (e_{\alpha})_x \big)_{ij}=\tilde\rho(v_{\alpha})_{ij}
\end{equation}
where the left hand side describe the matrix component of a linear operator acting on $\Lambda W_x$ while in the right hand side we have the matrix component of a linear operator acting on $\Lambda W$ and $v_{\alpha}$ is a basis on $\eR^n$ with respect to which the metric is the same as the metric $g_x$ in the basis $(e_{\alpha})_x$. Let $\hat{\psi}\colon P\to \Lambda W$ be an equivariant function; from definition \eqref{EqDefgammax} of $\gamma$ we have
\[ 
  \big( \gamma(e_{\alpha}\hat{\psi}) \big)(\xi)=(a_{\alpha}\hat{\psi})(\xi)
\]
where $a_{\alpha}(\xi)=\tilde\rho\Big( \tilde{e}_{\alpha}\big( \phi(\xi) \big) \Big)$. In this expression, $\tilde{e}_{\alpha}$ is the equivariant function associated with the vector field $e_{\alpha}\in\cvec(M)$. It is defined in subsection \ref{equivvec} as
\begin{equation}
\begin{aligned}
 \tilde{e}_{\alpha}\colon \SO(M)&\to \eR^m \\ 
b&\mapsto b^{-1}\big( (e_{\alpha})_{\pi(b)} \big). 
\end{aligned}
\end{equation}
So we have $\hat a_{\alpha}\colon P\to \End(\Lambda W)$ defined by
\[ 
  \hat a_{\alpha}(\xi)=\tilde\rho\big( \varphi(\xi)^{-1}e_{\alpha}(x) \big)
\]
with $x=\pi(\xi)$. Now if $\xi$ is any element of $\pi^{-1}(x)$, we have
\begin{align*}
\big( \gamma(e_{\alpha})\psi \big)(x)&=(a_{\alpha}\psi)(X)=\big[ \xi,\hat a_{\alpha}(\xi)\hat\psi(\xi) \big]
		=\big[ \xi,\tilde\rho\big( \varphi(\xi)^{-1}e_{\alpha}(x) \big)\hat{\psi}(\xi) \big].
\end{align*}
There exists a $g\in\Spin(p,q)$ such that $\varphi(\xi\cdot g)=\mtu$; taking this element and using equivariance of the latter expression,
\begin{align}
  \big( \gamma(e_{\alpha})\psi \big)(x)=\big[ \xi\cdot g,\tilde\rho\big( e_{\alpha}(x) \big)\hat{\psi}(\xi\cdot g) \big]
		=\big[ \xi\cdot g,\gamma^{\alpha}\hat{\psi}(\xi) \big]
		=[\xi,\gamma^{\alpha}\hat{\psi}(\xi)].
\end{align}
What we proved is that $\big( \gamma e_{\alpha}\psi \big)(x)=\gamma^{\alpha}\psi(x)$ is the sense that
\begin{equation}
	\widehat{\gamma(e_{\alpha})\psi}=\gamma^{\alpha}\hat{\psi}.
\end{equation}
Hence the Dirac operator reads
\[ 
  (\Dir s)(x)=g_{\alpha\beta}\gamma^{\alpha}\big( \tilde\nabla_{e_{\beta}}s \big)(x)
\]
in the sense that
\begin{equation}
\widehat{\Dir s}=g_{\alpha\beta}\gamma^{\alpha}\widehat{  \tilde\nabla_{e_{\beta}}s }.
\end{equation}

\end{proof}


An often more convenient way to write the Dirac operator is to consider an orthonormal basis (so that the metric $g$ and the matrices $\gamma$ are constant) and to consider the equivariant functions:
\[ 
  \widehat{\Dir\psi}=g_{\alpha\beta}\gamma^{\alpha}\widehat{\nabla_{e_{\alpha}}\psi}.
\]
This formulation is typically used when one search for Dirac operator on Lie groups. In this case, we choose left invariant vector fields generated by an orthonormal basis of the Lie algebra. The resulting field of basis is everywhere Killing-orthonormal.

Acting on a function $\dpt{f}{M}{\eR}$, it is defined by $\dpt{\Dir}{C^{\infty}(M)}{C^{\infty}(M)}$\index{dirac!operator!on functions},
\begin{equation} \label{eq_defDirac_f}
(\Dir f)(x)=g_x(e_{\alpha},e_{\beta})\gamma\hbeta_x(e_{\alpha x}\cdot f).
\end{equation}
With these definitions, one has
\[(\Dir(fs))(x)=(f\Dir s)(x)+(\Dir f)(x).\] Indeed,
\begin{equation}
\begin{split}
   (\Dir(fs))(x)&=g_{\alpha\beta}\gamma_x\hbeta(\tnab_{e_{\alpha}}fs)(s)\\
                &=g_{\alpha\beta}\Big((e_{\alpha}\cdot f)s(x)+f(x)(\tnab_{e_{\alpha}}s)(x)\Big)\\
                &=f(x)(\Dir s)(x)+g_{\alpha\beta}\gamma_x\hbeta(e_{\alpha x}\cdot f)\\
                &=(f\Dir s)(x)+(\Dir f)(x).
\end{split}
\end{equation}
With that definition, the Dirac operator becomes a derivation of the spinor bundle.


\section[Dirac operator on  \texorpdfstring{$\eR^2$}{R2}]{Example: Dirac operator on \texorpdfstring{$\eR^2$}{R2} with the euclidian metric}\label{Pg_exempleRdeux}\index{dirac!operator!on $\eR^2$}
%---------------------------------------------------

%\subsubsection{Example: tangent bundle}
%--------------------------

Since the frame bundle $B(M)$ is a principal bundle (see subsection \ref{subsec_frbundle}), one can consider some associated bundles on it. We are now going to see that the one given by the definition representation $\dpt{\rho}{GL(n,\eR)}{GL(n,\eR)}$ on $\eR^n$ is the tangent bundle. So we study $B(M)\times_{\rho} \eR^n$. By choosing a basis on each point of $M$, we identify each $T_xM$ to $\eR^n$. An element of $B(M)\times \eR^n$ is a pair $(b,v)$ with $b=(\overline{b}_1,\ldots,\overline{b}_n)$ and $v=(v^1,\ldots,v^n)$. We can identify $v$ to the element of $T_xM$ given by $v=v^i\overline{b}_i$.

In order to build the associated bundle, we make the identifications
\[
  (b,v)\cdot g\sim(b\cdot g,g^{-1} v).
\]
Here, by $gv$ we mean the vector whose components are given by $(gv)^i=v^j\bghd{g}{j}{i}$. The tangent vector given by $(b\cdot g,g^{-1} v)$ is $(g^{-1} v)^i(b\cdot g)_i=v^j\bghd{(g^{-1})}{j}{i}\bghd{g}{i}{k}\overline{b}_k=v^k\overline{b}_k$ So the identification map $\dpt{\psi}{B(M)\times_{\rho}\eR^n}{TM}$ given by
\[
  \psi([b,v])=v^i\overline{b}_i
\]
is well defined.

\index{spin!structure!on $\eR^2$}
The following step is to consider the following spin structure:
 \[\xymatrix{
    \Spin(2)  \ar@{~>}[r]&  \eR^2\times \SO(2) \ar[r]^-{\displaystyle\varphi} & \SO(\eR^2)  & \SO(2) \ar@{~>}[l].
  }\]

We have to define the two actions and $\varphi$. One of the main result of section~\ref{cliffR2} is that $\dpt{\chi}{\Spin(2)=U(1)}{\SO(2)}$ is surjective. So, we can define the action of $\Spin(2)$ on $P$ by
\[(x,b)\cdot s=(x,\chi(s)^{-1} b).\]

On the other hand, an element $A$ in $\SO(\eR^2)$ can be written as $A=\baz{a}{x}$ where $e_i$ is the canonical basis of $T_x\eR^2$, and $a$ is a matrix of $\SO(2)$. See subsection \ref{subsec_frbundle}. For $g\in \SO(2)$, we define
\begin{eqnarray}
 \label{r1504d2}A\cdot g=\{g^{-1} ae_i\}_x.
\end{eqnarray}
and  $\dpt{\varphi}{\eR^2\times \SO(2)}{\SO(\eR^2)}$ by
\[
\varphi(x,b)=\{be_i\}_x.
\]
The following shows that these definitions give a spin structure:
\begin{equation}
   \varphi((x,b)\cdot s)=\varphi(x,\chi(s)^{-1} b)
                    =\{\chi(s)^{-1} be_i\}_x
                    =\{be_i\}_x\cdot\chi(s)
                    =\varphi(x,b)\cdot\chi(s).
\end{equation}


\subsection{Connection on \texorpdfstring{$\SO(\eR^2)$}{SO(R2)}}\index{connection!on $\SO(\eR^2)$}
%///////////////////////////////////////

We are searching for a torsion-free connection on the simplest metric space: the euclidian $\eR^2$. Thus we will try the simplest choice of horizontal space: we want an horizontal vector to be tangent to a curve of the form $X(t)=\baz{b}{x(t)}$. For this reason, we want to define the connection $1$-form by $\omega(X)=b'(0)$. For technical reasons which will soon be apparent, we will not exactly proceed in this manner. For $X(t)=\baz{b}{x(t)}$, we define
\begin{equation}
                       \omega(X)=-(b(t)b(0)^{-1})'(0).
\end{equation}
We of course have $\omega(X)=0$ if and only if $b'(0)=0$: this choice of $\omega$ follows our first idea. In order for $\omega$ to be a connection form, we have to verify the two conditions of definition \ref{defconnform}.

\begin{proposition}
The $1$-form defined by
\[
              \omega(X)=-(b(t)b(0)^{-1})'(0)
\]
for $X=\displaystyle\dsdd{\baz{b(t)}{x(t)}}{t}{0}$ is a connection $1$-form.
\end{proposition}

\begin{proof}
Let $A\in \SO(2)$. If $u=\baz{b}{x}$, equation \eqref{r1504d2} gives:
\[
   A^*_u=\dsdd{\baz{e^{-tA}b}{x}}{t}{0},
\]
 so that $\omega(A^*_u)=-(e^{-tA}bb^{-1})'(0)=A$. This checks the first condition. For the second, one remarks that the path in $\SO(\eR^2)$ which defines the vector $R_{g*}X$ is $(R_{g*}X)(t)=\baz{g^{-1} b(t)}{x}$. It follows that
\begin{equation}
\begin{split}
\omega(R_{g*}X)&=-(g^{-1} b(t)b(0)^{-1} g)'(0)\\
               &=-\left(\AD_{g^{-1}}(b(t)b(0)^{-1})\right)'(0)\\
               &=-Ad_{g^{-1}}(b(t)b(0)^{-1})'(0)\\
               &=Ad_{g^{-1}}\omega(X).
\end{split}
\end{equation}

\end{proof}

\begin{proposition}
The covariant derivative induced on $M$ by this connection is
\begin{equation}\label{derrcovexplicite}
                 \nabla_XY=X(Y).
\end{equation}
\end{proposition}

\begin{proof}
In this demonstration, we will use the equivariant functions defined in \ref{equivvec}. In order to compute $(\nabla_XY)_x$, we have to use the definition of theorem \ref{tho_nablaE}. We first have to compute the horizontal lift of $X$. It is no difficult to see that $\oX_{\baz{b}{x}}$ is given by the path
\[\oX(t)=\baz{b}{X(t)}\]
if the vector field $X$ is given by the path $X(t)$ in $M$. Indeed, it is trivial that $\omega(\oX)=0$, and
\[d\pi_*\oX=\dsdd{\pi\baz{b}{X(t)}}{t}{0}=\dsdd{X(t)}{t}{0}=X.\]

Now, we compute $(\oX\hs)(b)$ for $b=\{Se_i\}_x$. We begin using the basic definitions and notations:
\[
(\oX\hs)(b)=\oX_b\hs=\dsdd{\hs(\oX_b(t))}{t}{0}=\dsdd{\hs(\baz{S}{X(t)})}{t}{0}.
\]
We can rewrite it with $\hY$ instead of $\hs$. By construction (see \eqref{r1404e1}), if $b=\baz{S}{x}$, $\hY(b)=S^{-1}(Y_x)$. Thus
\[
(\oX\hY)(b)=\dsdd{S^{-1}(Y_{X(t)})}{t}{0},
\]
where, if $\{\oui\}$ is a basis of $\eR^m$, then $S$ is 
\begin{equation}
\begin{aligned}
 S\colon\eR^m&\to T_{X(t)}M \\ 
 v^i\oui &\mapsto S^i_jv^j(\partial_j)_{X(t)} 
\end{aligned}
\end{equation}
So if we write $Y_x=Y^i(x)\partial_i$, we have
\[
S^{-1}(Y_{X(t)})=(S^{-1})^i_jY^j(X(t))\oui
\]
and
\[
\dsdd{S^{-1}(Y_{X(t)})}{t}{0}=(S^{-1})^i_j\dsdd{Y^j(X(t))}{t}{0}\oui=(S^{-1})^i_jX(Y^j)\oui.
\]
Since $b$ is an isomorphism, we can apply $b$ on both side of $\hX(b)=b^{-1}(X_x)$, and take $\nabla_XY$ instead of $X$:
\begin{equation}
 (\nabla_XY)(x)=b\big((S^{-1})^i_jX(Y^j)\oui\big)
               =S^k_i(S^{-1})^i_jX(Y^j)(\partial_k)_x
               =X(Y^j)(\partial_j)_x
               =X(Y)_x.
\end{equation}
\end{proof}

From this and definition \ref{deftorsion}, we immediately conclude that our connection is torsion-free. In a certain manner, one can say that our covariant derivative is the usual one.

\subsection{Construction of \texorpdfstring{$\gamma$}{g}}
%//////////////////////////////////////

Now, we construct the map $\gamma$ of subsection \ref{applgamma}. The first step is to define $\dpt{\ha_X}{P}{\End{(\Lambda W )}}$ by
\[
\ha_X(p)=\tilde\rho(\hX_{\varphi(p)}).
\]
Here, $\Lambda W $ is the completely isotropic subspace of $(\eR^2)^{\eC}$ with euclidian metric; thus we can use the result of section \ref{cliffR2}. In particular, we know the representation $\tilde\rho$.

To see it more explicitly, we need the expression of $\hX$. It is given in subsection \ref{equivvec}: if $b$ is the basis $\baz{b}{x}$, $\hY(b)=b^{-1}(Y_x)$. As $\varphi(b,x)=\baz{b}{x}$, we have
\[
\ha_X(b,x)=\tilde\rho(b^{-1}(X_x)).
\]

The subsection \ref{equivendo} explains how to explicitly get $\gamma(X)$ with the definition $\gamma(X)=a_X$. If $\psi$ is a section of $\mS$ and $\psi(x)=[\xi,v]$, the general definition gives us $(a_X\psi)(x)=[\xi,\ha_X(\xi)v]$ and in our particular case, if $\xi=(b,x)$, we get:
\begin{eqnarray}
 \label{gammaX}(\gamma(X)\psi)(x)=[\xi,\tilde\rho(b^{-1}(X_x))v].
\end{eqnarray}

\subsection{Covariant derivative on \texorpdfstring{$\protect\Gamma(\mS)$}{S}}
%///////////////////////////////////////////

Remember the spin structure of $\SO(\eR^2)$: $\varphi(x,S)=\{Se_i\}_x$. We now construct the connection on $P=\eR^2\times \SO(2)$. It is defined by the $1$-form $\tomega=\varphi^*\omega$. If $v$ is a vector of $P$, it is described by a path $v(t)=(x(t),b(t))$, then the path of $d\varphi(v)$ is $\{b(t)e_i\}_{x(t)}$ and $\tomega(v)=\omega(d\varphi(v))=-(b(t)b(0)^{-1})'(0)$.

The next step defining the Dirac operator is to find out an explicit form for the map $\dpt{\tnab}{\cvec(M)\times\Gamma(\mS)}{\Gamma(\mS)}$. A section $s\in\Gamma(\mS)$ is a map $\dpt{s}{M}{\mS=(\eR^2\times \SO(2))\times_{\rho}\Lambda W }$; it is defined by an equivariant function $\dpt{\hs}{P}{\Lambda W }$. In order to find the value of $(\tnab_Xs)(x)$ for $X\in\cvec(M)$, we use the definition
\[
 \widehat{\tnab_Xs}(\xi)=\oX\bxi(\hs)
\]
where $\oX$ is the horizontal lift in the sense of $\tomega$. For the same reason as in the proof of proposition \ref{derrcovexplicite}, $\oX_{(b,x)}$ is given by the path $\oX(t)=(b,X(t))$ where $X(t)$ is the path which defines $X$. So we have
\[
 \widehat{\tnab_Xs}(\xi)=\oX_{(b,x)}(\hs)=\dsdd{\hs(b,X(t))}{t}{0}.
\]
Remark that $\Lambda W $ is a vector space; so for every $\alpha\in\Lambda W $, the identification $T_{\alpha}\Lambda W =\Lambda W $ is correct.

Our first form of $\tnab$ is
\[
(\tnab_Xs)(x)=\Big[\xi,\dsdd{\hs(b,X(t))}{t}{0}\Big],
\]
but we can modify this in order to get simpler expressions. Remark that we have an equivalence class, so that we can always choose the element of the class such that $\xi=(\mtu,x)$. We define $\dpt{\os}{\eR^2}{\Lambda W }$, $\os(v)=\hs(\mtu,v)$. Our second and final form for $\tnab$ is:
\begin{subequations}
 \begin{align}
 (\tnab_Xs)(x)&=\Big[(\mtu,x),\dsdd{\os(X(t))}{t}{0}\Big]\\\label{nabs}
              &=[(\mtu,x),X(\os)],
\end{align}
\end{subequations}
where $X(\os)$ is well defined because $\os$ is a map from $\eR^2$ into a vector space (namely: $\Lambda W $).

\subsection{Dirac operator on the euclidian \texorpdfstring{$\eR^2$}{R2}}
%///////////////////////////////////////////////////
\index{dirac!operator!on $\eR^2$}

We continue to write explicitly the definition \eqref{dirac}. Putting together \eqref{gammaX} and \eqref{nabs}, one finds
\begin{equation}
 \gamma^{\alpha}_x(\tnab_{e_{\beta}}s)(x)	=\gamma(e_{\alpha x})[\xi,e_{\beta}(\os)]
                                     		=[\xi,\tilde\rho(b^{-1}(e_{\alpha x}))e_{\beta}(\os)].
\end{equation}
Here, $e_{\beta}=\partial_{\beta}$ and $b=\mtu$, then
\[
 \gamma^{\alpha}_x(\tnab_{e_{\beta}}s)(x)=[(\mtu,x),\tilde\rho(e_{\alpha})\partial_{\beta}\os].
\]
Now, the Dirac operator reads
\[
 (\Dir s)(x)=[(\mtu,x),\gamma^{\alpha}\partial_{\alpha}\os].
\]

We can obtain a more compact expression by defining ``$Ys$''\ and ``$As$'' when $s\in\Gamma(\mS)$, $Y\in\cvec(\eR^2)$ and $A\in\End{\Lambda W }$. The definitions are
\begin{align*}
(Ys)(x)&=[(\mtu,x),(Y\os)(x)],\\
(As)(x)&=[(\mtu,x),A\os(x)].
\end{align*}
With these conventions, one writes:
\[
(\Dir s)(x)=\gamma^{\alpha}(\partial_{\alpha} s)(x).
\]
This justifies the expression \eqref{dirflat}: $\Dir=\gamma^{\alpha}\partial_{\alpha}$ on flat spaces. With a good choice of basis of $\Lambda W $, the matrices $\gamma^{\alpha}$ are given by \eqref{gammaR2}, and
\[
\gamma^{\alpha}\partial_{\alpha}=
\begin{pmatrix}
0 & -1 \\
1 & 0
\end{pmatrix}\partial_x-
\begin{pmatrix}
0 & i \\
i & 0
\end{pmatrix}\partial_y.
\] 
If we identify $\eR^2$ with $\eC$ we have the following definitions:
\[
\partial_z=\frac{1}{2}(\partial_x-i\partial_y),\qquad\partial_{\overline{z}}=\frac{1}{2}(\partial_x+i\partial_y),\]
so that
\[\Dir=\begin{pmatrix}
0 & -\partial_{\overline{z}} \\
\partial_z & 0
\end{pmatrix}.
\]
\section{Clifford algebras and Morita equivalence}
%++++++++++++++++++++++++++++++++++++++++++++++++

Let $\cA$ be an algebra. An algebra $\cB$ is said to be \defe{Morita equivalent}{Morita equivalence}\label{PgMoritaEq} to $\cA$ if $\cB=\End_{\cA}(\modE)$ for some finite projective module $\modE$ over $\cA$. The algebra $\cA$ is Morita equivalent to itself taking the trivial module $\modE=\cA$.

We consider a manifold $M$ of dimension $n=2m$.

\begin{probleme}
	The two following statements are imprecise.
\end{probleme}

\begin{proposition}
A module which implement a Morita equivalence between two $C^*$-algebras is finite projective.
\end{proposition}

\begin{theorem}[Serre-Swan]
If one of the two Morita equivalent is the continuous function space over a manifold $\cA=C(M)$, then the module which gives the Morita equivalence is the section of continuous sections of a vector bundle over $M$, $\modE=\Gamma(E)$.
\end{theorem}
Furthermore, if $\cA=C(M)$ and $\cB=\Gamma(\Cl(M))$, we have $\End E\simeq \Cl(M)$ as isomorphism of vector bundle. Since $\Cl M$ is of rank $2^n$, $\End E$ has same rank and $E_x$ has dimension $\sqrt{2^n}=2^m$. So it is possible to choose the Clifford action in such a way that $\Gamma(E)$ is an irreducible Clifford module.

We often look at an anti-linear map $J\colon \Gamma(E)\to \Gamma(E)$ such that for all $\psi\in\Gamma(E)$
\begin{enumerate}
\item $J(\psi f)=(J\psi)\overline{ f }$ for all $f\in C(M)$,
\item $J(a\psi)=\epsilon(a)a J\psi$ for all $a\in\Gamma^{\infty}(\Cl M)$.
\end{enumerate}
How to define $a\psi$ ? We consider $\cA=C(M)$, $\cB=\Gamma(\Cl M)$ and we define $\Gamma(E)$ is such a way that it implements a Morita equivalence between $\cA$ and $\cB$; hence $\Gamma(E)$ is a $C(M)$-module. From dimensional considerations, we can define on $\Gamma(E)$ a Clifford module structure, i.e. a $C(M)$-linear
\begin{equation}
  c\colon \Gamma(\Cl M)\to \End(\Gamma E),
\end{equation}
hence $a\psi$ makes sense for any $a\in\Gamma^{\infty}(\Cl M)$ and $\psi\in\Gamma(E)$ with definition
\begin{equation}
 (a\psi)(x)=\big( c(a)\psi \big)(x)
		=c(a(x))\psi(x)
\end{equation}

\begin{theorem}
Let $(M,S,J)$ be a spin manifold of dimension $n$. There exists an unique connection 
\[ 
  \nabla^S\colon \Gamma^{\infty}(S)\to \Gamma^{\infty}(S)\otimes\Omega^1(S)
\]
such that
\begin{enumerate}
\item $\scalp{ \nabla^S\psi }{ \phi }+\scalp{ \psi }{ \nabla^S\phi }=d\scalp{ \psi }{ \phi }$,
\item $[\nabla^S,J]=0$,
\item $\nabla^S\big( c(a)\psi \big)=c\big( \nabla a \big)\psi+c(a)\nabla^S\psi$ for all $a\in\Cl(M)$ and $\psi\in\Gamma^{\infty}(S)$.
\end{enumerate}
In the latter, the action of $\Gamma^{\infty}(\Cl M)$ on $\Gamma^{\infty}(S)$ is induced from the action $c\colon \Cl(T^*_xM)\to \End S$. The $\nabla$ which acts on $a$ is the connection extended to $\Gamma^{\infty}(\Cl M)$ by virtue of Leibnitz rule $\nabla(uv)=\nabla(u)v+u\nabla(c)$.

\end{theorem}

\begin{proof}
No proof
\end{proof}


In this setting, we define 
\begin{equation}
\begin{aligned}
 \hat c\colon\Gamma^{\infty}(S)\otimes\Gamma^{\infty}(\Cl M)&\to \Gamma^{\infty}(S) \\ 
 \psi\otimes a&\mapsto c(a)\psi. 
\end{aligned}
\end{equation}
Then we define the \defe{Dirac operator}{dirac!operator} $\Dir\colon \Gamma^{\infty}(S)\to \Gamma^{\infty}(S)$,
\begin{equation}
  \Dir=-i(\hat c\circ\nabla^S).
\end{equation}

\subsection{Example: quantum field theory}
%-----------------------------------------

Let us show how does this operator gives back the usual Dirac operator of quantum field theory. Let $M$ be a manifold and with two local basis $\{ \partial_u \}$ and $\{ \partial_{\alpha} \}$ of $T_xM$. The first one is the ``natural'' basis: $g(\partial_u,\partial_v)=g_{uv}$ has no particular properties while the second one is orthonormal $g(\partial_{\alpha},\partial_{\beta})=\delta_{\alpha\beta}$. The first dual basis is defined by $dx^{\alpha}\partial_{\beta}=\delta^{\alpha}_{\beta}$. 

We write $\partial_{\alpha}=e_{\alpha}^u\partial_u$ and for the dual basis, $dx^{\alpha}=e^{\alpha}_u\,dx^u$. In order these definition to be coherent, we impose $dx^{\alpha}\partial_{\beta}=\delta^{\alpha}_{\beta}$ :
\begin{equation}
  dx^{\alpha}\partial_{\beta}=e_u^{\alpha}dx^u\big( e^v_{\beta}\partial_v \big)
		=e^{\alpha}_ue^v_{\beta}\delta^u_v
		=e^{\alpha}_ue^u_{\beta}.
\end{equation}
We conclude that the \defe{vielbein}{vielbein} $(e^{\alpha}_u)$ is the inverse of $(e^v_{\beta})$: $e^{\alpha}_ue^u_{\beta}=\delta^{\alpha}_{\beta}$. The vielbein are eventually complexes.



\subsection{An other definition of the Dirac operator}
%-----------------------------------------------------

Let us consider an orthonormal basis $\{ e_a \}$ of $M$, i.e. on each $x\in M$, 
\[ 
  g_x(e_a(x),e_b(x))=\eta_{ab}.
\]
This basis is related to a ``natural'' basis $\{ \partial_{\mu} \}$ by
\begin{equation} 
  e_a=e_a^{\mu}\partial_{\mu}
\end{equation}
where  $e_a^{\mu}$ is called \defe{vielbein}{vielbein} (here, they are more precisely $n$-beins). As far as metric is concerned we have
\begin{subequations}
\begin{align}
	g^{\mu\nu}&=e_a^{\mu}e_b^{\nu}\eta_{ab}\\
	\eta_{ab}&=e_a^{\mu}e_b^{\nu}g_{\mu\nu}.
\end{align}
\end{subequations}
If $\nabla$ is the covariant derivative associated with $g$, we define the coefficients $\omega_{\mu a}^b$ by
\begin{equation}
\nabla_{\mu}e_a=\omega_{\mu a}^be_b.
\end{equation}
On the other hand, $\nabla$ is related to the Christoffel symbols by
\begin{equation}
\nabla_{\mu}\partial_{\nu}=\Gamma_{\mu\nu}^{\sigma}\partial_{\sigma}.
\end{equation}
Let $\Cl(M)$ be the Clifford module whose fibre is the Clifford complex algebra $\Cl(T^*_xM)^{\eC}$. We consider $\Gamma(\Cl(M))$, the module of corresponding sections. It gives an algebra morphism
\begin{equation}
\begin{aligned}
 \gamma\colon\Gamma(\Cl(M))&\to \opB(\hH) \\ 
dx^{\mu}&\mapsto \gamma^{\mu}(x)=\gamma^ae_a^{\mu} 
\end{aligned}
\end{equation}
which can be extended to the whole Clifford algebra. One can choose matrices $\gamma^{\mu}(x)$ and $\gamma^a$ to be hermitian; they satisfy
\begin{subequations}
\begin{align}
\gamma^{\mu}(x)\gamma^{\nu}(x)+\gamma^{\nu}(x)\gamma^{\mu}(x)&=-2g(dx^{\mu},dx^{\nu})=-2g^{\mu\nu}\\
\gamma^a\gamma^b+\gamma^b\gamma^a&=-2\eta^{ab}.
\end{align}
\end{subequations}
All this allow us to lift the Levi-Civita connection from the tangent bundle to the spinor bundle by defining
\begin{equation}
\nabla_{\mu}^S=\partial_{\mu}+\omega^S_{\mu}=\partial_{\mu}+\frac{ 1 }{2}\omega_{\mu ab}\gamma^a\gamma^b.
\end{equation}
The \defe{Dirac operator}{operator!Dirac} is then given by
\[ 
  \Dir=\gamma\circ\nabla
\]
and can locally be written under the form
\begin{equation}  \label{eq_Dirac_deux}
\Dir=\gamma^{\mu}(x)(\partial_{\mu}+\omega_{\mu}^S)
	=\gamma^ae_a^{\mu}(\partial_{\mu}+\omega^S_{\mu}).
\end{equation}

\section{Dirac operator on \texorpdfstring{$AdS_2$}{AdS2}}
%-------------------------------------

Why to compute Dirac operator on anti de Sitter spaces ? Let $M=AdS_2$ and $R=AN$ acts on $M$. Let $\mO$ be an open orbit of $R\times M\to M$. In the specific case of $AdS_2$, we have $R=\mO=R\cdot\mfo$. In larger dimensions, there is a $\SO(1,n)$ which causes that the orbit is not exactly the acting group. It is
\[ 
  \mO=\frac{ R }{ R\cap \SO(1,n) }.
\]

\subsection{Clifford algebra and spin group}
%/////////////////////////////////////


As definition, we retain
\begin{equation}
\begin{split}
  AdS_2&\equiv t^2+u^2-x^2=1\\
	&=\frac{ \SO(2,1) }{ \SO(1,1) }.
\end{split}
\end{equation}
Let $V=\eR^{1,1}$ and $e_0$, $e_1$ an orthonormal basis. We pose 
\begin{align*}
  f_0=\frac{ 1 }{2}(e_0+e_1)\quad g_0=\frac{ 1 }{2}(e_0-e_1)
\end{align*}
and we define $\tilde\rho$ by
\begin{subequations}
\begin{align}
  \tilde\rho(f_0)\alpha&=f_0\wedge\alpha\\
	\tilde\rho(g_0)\alpha&=-i(g_0)\alpha
\end{align}
\end{subequations}
where $\alpha\in\Lambda W$, $W$ being the space spanned by $f_0$. More explicitly we have :
\begin{subequations}
\begin{align}
  \tilde\rho (f_0)1&=f_0&\tilde\rho (f_0)f_0&=0\\
\tilde\rho (g_0)1&=0&\tilde\rho (g_0)f_0&=-\eta(f_0,g_0).
\end{align}
\end{subequations}
As element of $\Lambda W$, $f_0$ stands for $\eta(f_0,.)$. If we choose the basis
\[ 
  1=
\begin{pmatrix}
1\\0
\end{pmatrix},
\quad
f_0=
\begin{pmatrix}
0\\1
\end{pmatrix},
\]
 the matrices of $\tilde\rho$ are given by
\[ 
  \tilde\rho(e_0)=
\begin{pmatrix}
0&-1/2\\
1&0
\end{pmatrix},
\quad
\tilde\rho(e_1)=
\begin{pmatrix}
0&1/2\\
1&0
\end{pmatrix}.
\]
Up to a change of basis,
\[ 
  \gamma_0=
\begin{pmatrix}
0&1\\1&0
\end{pmatrix},
\quad
\gamma_1=
\begin{pmatrix}
0&-1\\1&0
\end{pmatrix}
\quad
\gamma_0\gamma_1=
\begin{pmatrix}
1&0\\0&-1
\end{pmatrix},
\]
and a general element of $Cl_{(1,1)}$ reads
\[ 
  x\gamma_0+y\gamma_1+u\eR+v\gamma_0\gamma_1=
\begin{pmatrix}
u+v&x-y\\
x+y&u-v
\end{pmatrix}.
\]
With the change of basis $e_1\to ie_1$, we write it under a more simple form :
\begin{equation}
Cl_{(1,1)}\leadsto 
\begin{pmatrix}
\alpha&\beta\\
\bar\beta&\bar\alpha
\end{pmatrix}
\end{equation}
with $\alpha,\beta\in\eC$. In particular, an element of $V$, i.e. a combination of $\gamma_0$ and $\gamma_1$ is
\begin{equation}
V\leadsto
\begin{pmatrix}
0&\xi\\
\overline{\xi}&0
\end{pmatrix}.
\end{equation}
Let
\[ 
  1=
\begin{pmatrix}
1&0\\0&1
\end{pmatrix},
\quad
a=
\begin{pmatrix}
i&0\\0&-i
\end{pmatrix},
\quad b=
\begin{pmatrix}
0&1\\1&0
\end{pmatrix},
\quad
c=
\begin{pmatrix}
0&i\\-i&0
\end{pmatrix}
\]
Let us now determine $\alpha$, the extension of $-\id|_V$ into an automorphism and $\tau$, the extension of $\id|_V$ into an anti-automorphism. We have $\alpha(b)=-b$, $\alpha(c)=-c$, $\tau(b)=b$ and $\tau(c)=c$. We find the others by virtue of relations $bc=-a$ and $b^2=1$. Finally
\begin{equation}
\begin{aligned}
  \alpha(1)&=1&\alpha(a)&=a\\
\alpha(b)&=-b&\alpha(c)&=-c
\end{aligned}
\end{equation}
and
\begin{equation}
\begin{aligned}
  \tau(1)&=1&\tau(a)&=-a\\
\tau(b)&=b&\tau(c)&=c.
\end{aligned}
\end{equation}

The condition for $s\in Cl_{(1,1)}$ to belongs to $\Gamma_{(1,1)}$ is that $\alpha(s)v s^{-1}\in V$ for all $v\in V$. If we consider $s=
\begin{pmatrix}
\alpha&\beta\\\bar\beta&\bar\alpha
\end{pmatrix}$,
we have
\[ 
  \alpha(s)=
\begin{pmatrix}
\alpha&-\beta\\-\bar\beta&\bar\alpha
\end{pmatrix},
\text{ and }
s^{-1}=\frac{1}{ | \alpha |^2-| \beta |^2 }
\begin{pmatrix}
\bar\alpha&-\beta\\-\bar\beta&\alpha
\end{pmatrix}.
\]
If we impose $\alpha(s)v s^{-1}$ to be of the form $\begin{pmatrix}
0&\eta\\\bar\eta&0
\end{pmatrix}$ for all $v$ of the form $\begin{pmatrix}
0&\xi\\\bar\xi&0
\end{pmatrix}$, we find $\Reel(\bar\alpha\beta\bar\xi)=0$ and the $\bar\alpha\beta=0$. So generators of $\Gamma_{(1,1)}$ are
\begin{equation}
\Gamma_{(1,1)}\leadsto
\begin{pmatrix}
\alpha&0\\0&\bar\alpha
\end{pmatrix},\quad
\begin{pmatrix}
0&\beta\\\bar\beta&0
\end{pmatrix}.
\end{equation}
Elements of $Spin_{(1,1)}$ are elements of $\Gamma_{(1,1)}^+$ such that $\tau(s)=s^{-1}$. So
\begin{equation}
Spin_{(1,1)}\leadsto
\begin{pmatrix}
\alpha&0\\0&\bar\alpha
\end{pmatrix},\text{ such that }| \alpha |^2=1.
\end{equation}
We recognize $Spin_{(1,1)}=U(1)$.

\begin{probleme}
This is wrong: in fact $\Spin(1,1)\neq U(1)$.
\end{probleme}

\subsection{Relation between \texorpdfstring{$SU(1,1)$}{SU(1,1)} and \texorpdfstring{$\SO(2,1)$}{SO(2,1)}}
%//////////////////////////////////

A general matrix of $SU(1,1)$ is
\[ 
  U=\begin{pmatrix}
\alpha&\beta\\\bar\beta&\bar\alpha
\end{pmatrix},
\quad
U^{-1}=\begin{pmatrix}
\bar\alpha&-\beta\\-\bar\beta&\alpha
\end{pmatrix}
\]
with $| \alpha |^2 - | \beta |^2=1$. They are matrices which fulfil $\det U=1$ and $U^+g=gU^{-1}$.  If we denote by $V$ the space of matrices of the form $(r,z)=\begin{pmatrix}
r&\bar z\\z&r
\end{pmatrix}$ with $r\in\eR$ and $z\in\eC$, we have a bijection $\psi\colon \eR^{1,1}\to V$ given by
\[ 
  \begin{pmatrix}
u\\t\\x
\end{pmatrix}\mapsto
\begin{pmatrix}
x&t-iu\\ t+iu&x
\end{pmatrix}.  
\]
It becomes an isometry if we pose $\| (r,z) \|=z\bar z-r^2=-\det(r,z)$. The group $SU(1,1)$ has an isometric action on $V$ given by
\[ 
  Uv=UvU^{\dag}.
\]
We immediately remark that $Uv=(-U)v$. We define
\begin{equation}
\begin{aligned}
 T\colon SU(1,1)&\to \SO(2,1) \\ 
T(U)\begin{pmatrix}
u\\t\\x
\end{pmatrix}&= \psi^{-1}\big( U\psi\begin{pmatrix}
u\\t\\x
\end{pmatrix}U^{\dag} \big). 
\end{aligned}
\end{equation}

Now we want to know when $T(U)=T(\tilde U)$. Using the fact that $U^{-1}=gU^{\dag}g$ in the condition $UvU^{\dag}=\tilde Uv\tilde U^{\dag}$, we find
\[ 
  VvV^{\dag}=v
\]
with $V=\tilde U^{-1}U$. Then imposing
\[ 
  \begin{pmatrix}
r&\bar z\\z&r
\end{pmatrix}=
\begin{pmatrix}
\alpha&\beta\\\bar\beta&\bar\alpha
\end{pmatrix}
\begin{pmatrix}
r&\bar z\\z&r
\end{pmatrix}
\begin{pmatrix}
\bar \alpha&\beta\\\bar\beta&\alpha
\end{pmatrix},
\]
we find $T(U)=T(\tilde U)\Leftrightarrow \tilde U=\pm U$. We have
\[ 
  T\begin{pmatrix}
i\\&-i
\end{pmatrix}
=
\begin{pmatrix}
-1\\&-1\\&&1
\end{pmatrix}.
\]
The map $T\colon SU(1,1)\to \SO(2,1)$ is a double covering.

We are now going to explicitly compute the map $T$. First :
\[ 
 \begin{split}
\begin{pmatrix}
\alpha&\beta\\
 \bar\beta&\bar \alpha
\end{pmatrix}
\begin{pmatrix}
r&\bar z\\
z&r
\end{pmatrix}
&
\begin{pmatrix}
\bar\alpha&\beta\\
\bar\beta&\alpha
\end{pmatrix}=\\
&\begin{pmatrix}
\bar\alpha(\alpha r+\beta z)+\bar\beta(\alpha\bar z+\beta r)& \beta(\alpha r+\beta z)+\alpha(\alpha\bar z+\beta r)\\
\bar\alpha(\bar\beta+\bar\alpha z)+\bar\beta(\bar\beta\bar z+\bar\alpha r)&\beta(\bar\beta r+\bar\alpha z)+\alpha(\bar\beta\bar z+\bar\alpha r)
\end{pmatrix}.
\end{split} 
\]
When we pose $z=0$ and $r=1$, i.e., when we look at $\begin{pmatrix}
0\\0\\1
\end{pmatrix}$, we find
\[ 
  \begin{pmatrix}
\bar\alpha\alpha+\bar\beta\beta&2\beta\alpha\\
2\bar\alpha\bar\beta& \beta\bar\beta+\alpha\bar\alpha
\end{pmatrix}
\]
which corresponds to $x=\bar\alpha\alpha$, $t-iu=2\alpha\beta$ and $t+iu=2\bar\alpha\bar\beta$. We conclude that
\[ 
  T\begin{pmatrix}
\alpha&\beta\\\bar\beta&\bar\alpha
\end{pmatrix}
=
\begin{pmatrix}
.&.&i(\alpha\beta-\bar\alpha\bar\beta)\\
.&.&\alpha\beta+\bar\alpha\bar\beta\\
.&.&\alpha\bar\alpha+\beta\bar\beta
\end{pmatrix}.
\]
Similar computations lead to
\begin{equation}
T\begin{pmatrix}
\alpha&\beta\\
\bar\beta&\bar\alpha
\end{pmatrix}
=
\begin{pmatrix}
\frac{ \bar\alpha^2+\alpha^2-\beta^2-\bar\beta^2 }{2} &		\frac{ i }{2}(\alpha^2-\bar\alpha^2+\beta^2-\bar\beta^2) &	i(\alpha\beta-\bar\alpha\bar\beta)\\
\frac{ i }{2}(\beta^2-\bar\beta^2+\bar\alpha^2-\alpha^2) &	\frac{ 1 }{2}(\alpha^2+\bar\alpha^2+\beta^2+\bar\beta^2) &	\alpha\beta+\bar\alpha\bar\beta\\
i(\bar\alpha\beta-\bar\beta\alpha)			&	\bar\alpha\beta+\bar\beta\alpha 			&	\alpha\bar\alpha+\beta\bar\beta
\end{pmatrix}
\end{equation}

\subsection{Spin structure on \texorpdfstring{$AdS_2$}{AdS2}}
%/////////////////////////////////////////

We are going to build elements of the following spin structure:
\[
\xymatrix{ \Spin(1,1) \ar@{~>}[r]& SU{(1,1)} \ar[rr]^{\displaystyle\varphi} \ar[rd]_{\displaystyle\pi} && \SO(AdS_2) \ar[ld]^{\displaystyle p}&\SO(2,1) \ar@{~>}[l]  \\& & AdS_2 }
\]
First let $\{ e_t,e_u,e_x \}$ be a basis of $\eR^{1,1}$ with $e_t\in AdS_2$, $e_t\cdot e_t=e_u\cdot e_u=-e_x\cdot e_x=-1$. We suppose that $e_u$ and $e_x$ span tangent space at $e_t$. Let $T$ be a representation of $\SO(1,1)$ on $\eR^{2,1}$  which leaves $e_t$ unchanged: $T(A)e_t=e_t$ for all $A\in \SO(1,1)$. To each element $B\in \SO(AdS_2)$, one can associate an element of $B'\in \SO(AdS_2)$ such that $B$ has the form
\begin{equation} \label{eq_blaseAdS}
  B=\{ B'e_u,B'e_x \}_{B'e_t}.
\end{equation}
We define
\[ 
  p(B)=B'e_t.
\]
Now the action of $A\in \SO(1,1)$ on $B\in \SO(AdS_2)$ is defined, if $B$ has the form \eqref{eq_blaseAdS}, by
\begin{equation}
  B\cdot A=\{ T(A)B'e_u,T(A)B'e_x \}_{B'e_t}.
\end{equation}
The map $\varphi\colon SU(1,1)\to \SO(AdS_2)$ is given by
\begin{equation}
\big( \varphi(U) \big)'=(T\circ S)(U),
\end{equation}
and the projection $\pi\colon SU(1,1)\to AdS_2$, 
  $\pi=p\circ\varphi$.

The group $\Spin(1,1)$ must act on $SU(1,1)$; we define
\begin{equation}
U\cdot s=\varphi^{-1}\big( \varphi(U)\cdot \chi(s) \big).
\end{equation}
We have $\pi(U\cdot s)=\pi(U)$ because
\[ 
  \pi(U\cdot s)=p\big( \varphi(U)\cdot\chi(s) \big),
		=\big[ \varphi(U)\cdot \chi(s) \big]'e_t
		=\varphi(U)\circ T(\chi(s))e_t
		=\varphi(U)'e_t
		=\pi(U).
\]
We have used the fact that  $\chi(s)\in \SO(1,1)$ and that, therefore, $(T\circ\chi)(s)e_t=e_t$.


\subsection{Spinor bundle and connection}
%///////////////////////////////////////

We define $S=\Lambda W$ where $W$ is the (one dimensional) space spanned by $f_0$ and we define 
\begin{equation} \label{eq_mSSUrho}
  \mS= SU(1,1)\times_{\rho}S
\end{equation}
where $\rho\colon Spin_{(1,1)}\times\Lambda W\to \Lambda W$ is the representation of $Spin_{(1,1)}$ on $SU(1,1)$ given by
\begin{equation}
  \rho(s,\alpha)=\tilde\rho(s)\alpha.
\end{equation}
Recall that $\alpha$ is either a scalar either a multiple of $f_0$. The equivalence relation which arises in equation \eqref{eq_mSSUrho} is
\begin{equation}
  (U,\alpha)\sim(U\cdot s,\rho(s^{-1})\alpha).
\end{equation}
The projection is
\[ 
  \pi_{\mS}[(U,\alpha)]=\pi(U).
\]

For the connection on $\SO(AdS_2)$, we want that horizontal vector are tangent vectors to curves formed by parallel transport. In other word, a path 
\[
B(s)=\{ B'(s)e_u,B'(s)e_x \}_{B'(s)e_t}
\]
 has horizontal tangent vector if $B'(s)e_i$ ($i=u,x$) is a parallel transport of $B'(0)e_i$ along the curve $B'(s)e_e$ on $AdS_2$. Here, $B'(s)$ denotes the matrix of $\SO(2,1)$ associated with the basis $B(s)$ : the prime doesn't denotes a derivation. Let us define the $so(1,1)$ valued connection $1$-form which corresponds to this intuition. We consider $b_i(s)$ the parallel transported along the curve $B'(s)e_t$ of $B'(0)e_i$, and $A(s)$, the matrix of $\SO(1,1)$ such that $A(s)B'(s)e_i=b_i(s)$ ($i=u,x$). The definition is
\[ 
  \omega(\dot B)=\Dsdd{ A(s) }{s}{0}.
\]


\begin{proposition}
It is a connection $1$-form.
\end{proposition}

\begin{proof}
First we consider a fundamental vector field
\[ 
  X^*_B=\Dsdd{ B\cdot e^{-tX} }{t}{0}=\Dsdd{ \{ T( e^{-tX})B'e_u,T( e^{-tX})B'e_x \}_{B'e_t} }{t}{0}.
\]
The path in $AdS_2$ on which this path in $\SO(AdS_2)$ is build is constant: it is $B'e_t$. So the parallel transport is constant and the path $A(s)$ is given by
\[ 
  A(s)T( e^{-tX})B'e_u=B'e_u
\]
and $\omega(X^*_B)=X$.

It remains to be proved that for all $B\in \SO(AdS_2)$, $g\in \SO(1,1)$ and $X\in T_B\SO(AdS_2)$,
\begin{equation}
\omega\big( (dR_g)_BX \big)=\Ad(g^{-1})\omega_B(X).
\end{equation}
We give $X$ by the path 
\[ 
  X(s)=\{ B'(s)e_u,B'(s)e_x \}_{B'(s)e_t}.
\]
 The differential $dR_g$ gives rise to the new path
\[ 
  (dR_gX)(s)=\{ gB'(s)e_u,gB'(s)e_x \}_{B'(s)e_t}.
\]
Let $b_i(s)$ be the parallel transport of $B'(0)e_i$ ($i=u,x$) along the path $B'(s)e_t$. We have to compute $\omega_B(X)$ with $A(s)$ defined by $A(s)B'(s)e_i=b_i$. The parallel transport of $gB'(0)e_i$ is given by $gb_i$. Therefore $\omega(dR_gX)$ is given by the path $A^g(s)$ which satisfies $A^g(s)gB'(s)e_i=gA(s)B'(s)e_i$. So
\[ 
  A^g(s)=gA(s)g^{-1}
\]
and
\[ 
  \Dsdd{ A^g(s) }{s}{0}=\Ad(g)\omega(X).
\]

\end{proof}

\subsection{Clifford algebra \texorpdfstring{$(1,1)$}{(1,1)}}
%-------------------------------------

We consider the left invariant vector fields
\begin{subequations}
\begin{align}
  \tilde{e}_J(r_0)&=\Dsdd{ r_0 e^{-sJ} }{s}{0}=-r_0J\\
\tilde{e}_L(r_0)&=\Dsdd{ r_0 e^{-sL} }{s}{0}=-r_0L.
\end{align}
\end{subequations}
More precisely, we consider the vectors given by action of these matrices on the ``base point'' $\begin{pmatrix}
0\\1\\0
\end{pmatrix}$. Hence
\begin{align}
\tilde{e}_J(r_0)=-r_0\begin{pmatrix}
0\\0\\1
\end{pmatrix},\quad
\tilde{e}_L(r_0)=-r_0\begin{pmatrix}
1\\0\\0
\end{pmatrix}
\end{align}
and 
\[ 
  g=\begin{pmatrix}
1\\&-1
\end{pmatrix}.
\]
Remark that this metric is constant (it does not depend on $r_0$) because $r_0$ is an isometry. For this reason, we now turn our attention to Clifford algebra and spin group for $V=\eR^{1,1}$. Following matrices fulfill relation \eqref{3101r3}
\[ 
  \gamma_J=\begin{pmatrix}
0&1\\1&0
\end{pmatrix},\quad
\gamma_L=\begin{pmatrix}
0&-1\\1&0
\end{pmatrix}. 
\]
The complete Clifford algebra has the following matrices too :
\[ 
  1=\begin{pmatrix}
1&0\\0&1
\end{pmatrix},\quad
\gamma_{JL}=\begin{pmatrix}
-1\\&1
\end{pmatrix}.
\]
The Clifford algebra is nothing else than $GL(2,\eR)$, the set of all real $2\times 2$ matrices. From definitions, one can check that
\[ 
 \begin{aligned} 
\alpha(J)&=-J	&\tau(J)&=J\\
\alpha(L)&=-L	&\tau(L)&=L\\
\alpha(JL)&=JL	&\tau(JL)&=-JL\\
\alpha(1)&=1	&\tau(1)&=1
\end{aligned}
\]
Inverse and $\alpha$ of general element in $\Cl(1,1)$ are given by
\[ 
  \begin{pmatrix}
p&q\\r&s
\end{pmatrix}^{-1}=
\frac{1}{ ps-qr }\begin{pmatrix}
s&-q\\-r&p
\end{pmatrix},\quad
\alpha\begin{pmatrix}
p&q\\r&s
\end{pmatrix}=
\begin{pmatrix}
p&-q\\-r&s
\end{pmatrix}.
\]
A general element in $\eR^{1,1}$ is $\begin{pmatrix}
0&\alpha\\\beta&0
\end{pmatrix}$ with $\alpha,\beta\in\eR$, so the condition to belongs to $\Gamma(1,1)$ is that 
\[ 
  \frac{1}{ ps-qr }\begin{pmatrix}
p&-q\\-r&s
\end{pmatrix}
\begin{pmatrix}
0&\alpha\\\beta&0
\end{pmatrix}
\begin{pmatrix}
s&-q\\-r&p
\end{pmatrix}
\]
belongs to $\eR^{1,1}$ for all $\alpha$ and $\beta$. It requires, among others, that $qs\beta-rp\alpha=0$ for all $\alpha$ and $\beta$. Hence $qs=rp=0$, but the alternatives $p=r=0$ and $q=s=0$ are ruled out because we want the determinant $ps-qr$ to be non zero. Therefore, $\Gamma(1,1)$ is generated by
\[ 
  \Gamma(1,1)\leadsto
\begin{pmatrix}
p&0\\0&s
\end{pmatrix},\begin{pmatrix}
0&q\\r&0
\end{pmatrix}.
\]
The latter belongs to $\eR^{1,1}$, so
\[ 
  \Gamma^+(1,1)\leadsto\begin{pmatrix}
z+c&0\\0&z-c
\end{pmatrix}=z\mtu+c\gamma_{JL}.
\]
From 
\[ 
  \tau(z\mtu+c\gamma_{JL})=z\mtu-c\gamma_{JL},
\]
elements in $\Spin(1,1)$ are subject to the relation
\[ 
  \begin{pmatrix}
s&0\\0&p
\end{pmatrix}=
\frac{1}{ ps }\begin{pmatrix}
s&0\\0&p
\end{pmatrix}.
\]
As consequence, we find
\begin{equation}
\Spin(1,1)=\eR_0\leadsto
\begin{pmatrix}
1/p\\&p
\end{pmatrix}.
\end{equation}
 If we put (see decomposition \eqref{eq:expo_ANK})
\[ 
  A=\begin{pmatrix}
 e^{a}\\& e^{-a}
\end{pmatrix},
\]
we have $\Spin(1,1)=A\times\eZ_2$. Let us check that $\Spin(1,1)$ is a double covering of $\SO_0(1,1)$. We know that 
\[ 
  \SO(1,1)=\begin{pmatrix}
\cosh\xi&\sinh\xi\\\sinh\xi&\cosh\xi
\end{pmatrix}\times\eZ_2
\]
while $\SO_0(1,1)$ is 
\[ 
  \SO_0(1,1)=\begin{pmatrix}
\cosh\xi&\sinh\xi\\\sinh\xi&\cosh\xi
\end{pmatrix}=\eR.
\]
This structure of $\SO(1,1)$ comes from the fact (true for all $\SO(1,n)$) that $| \Lambda^0_0 |\geq1$ when $\Lambda$ is a Lorentz transformation. So $\mtu$ and $-\mtu$ cannot belong to the same connected component. Note that $\cosh\xi\geq 1$.
We see intuitively how to cover two times $\eR$ with $\eR_0$. Let us see how the map $\chi$ does that. From definition, $\chi(x)y=\alpha(x)yx^{-1}$, so it is easy to see that 
\[ 
  \chi(1)=\chi(-1)=\id|_{\eR^{1,1}}
\]

\subsection{Parallel transport}
%------------------------------

We have a connection on the frame bundle of $AdS_2$ and we wan to lift the vectors $\tilde{e}_J$ and $\tilde{e}_L$, i.e. we consider a point 
\[ 
  \xi_0=(r_0,v_1,v_2)\in \SO(AdS_2)
\]
where $v_1$ and $v_2$ form an orthonormal (in the sense of $g$) basis of $T_{r_0}AdS_2$. Then we have to find a path $s\to\xi(s)$ in $\SO(AdS_2)$ such that $\xi(0)=\xi_0$, $\omega(\xi'(0))=0$ and $dp\xi'(0)=\tilde{e}_a$. The latter condition allows us to compute $r(s)$ in the expression
\[ 
  \xi(s)=(r(s),v_1(s),v_2(s)),
\]
namely, $r(s)$ is the path of $\tilde{e}_a$. The condition to be horizontal imposes that vectors $v_i(s)$ are parallel transport of $v_i$ along $\tilde{e}_a$. So we have to compute the different $T_a(\tilde{e}_b)(s)$ which is the parallel transported of $\tilde{e}_b$ along the path of $\tilde{e}_a$ at a distance $s$; this is an element of $T_{\tilde{e}_a(s)}AdS_2$. It will be decomposed in the basis 
\[ 
 \begin{split}
  \tilde{e}_J\big( \tilde{e}_a(s) \big)&=-r_0 e^{-s a}J\\
\tilde{e}_L\big( \tilde{e}_a(s) \big)&=-r_0 e^{-sJ}L.
\end{split} 
\]
where we imply the action on the base point $\begin{pmatrix}
0\\1\\0
\end{pmatrix}$. For notational simplicity, from now we write $a(s)$ instead of $\tilde{e}_a(s)$. Various products are easy to compute; for example
\[ 
  \tilde{e}_J(J(s))\cdot\tilde{e}_L(J(s))
		=r_0 e^{-sJ}J\cdot r_0 e^{-sJ}L\\
		=J\cdot L\\
		=\begin{pmatrix}
0\\0\\1
\end{pmatrix}\cdot\begin{pmatrix}
1\\0\\0
\end{pmatrix}\\
		=0
\]
because $r_0 e^{-sJ}$ is an isometry.
In general :
\[ 
  \tilde{e}_a\big( c(s) \big)\cdot \tilde{e}_b\big( c(s) \big)=a\cdot b
\]
Now we pose in general
\[ 
  T_a(\tilde{e}_b)(s)=\alpha(s)\tilde{e}_b\big( a(s) \big)+\beta(s)\tilde{e}_L\big( a(s) \big),
\]
and we want to find the (real valued) functions $\alpha$ and $\beta$. Parallel transport fulfils two conditions: the norm and the angle with the path are constant. This leads us to two conditions :
\begin{subequations}
\begin{align}
  T_a\big( \tilde{e}_b(s) \big)\cdot T_a\big( \tilde{e}_b(s) \big)&=b\cdot b\\
 T_a\big( \tilde{e}_b(s) \big)\cdot \tilde{e}_a\big( a(s) \big)&=b\cdot a.
\end{align}
\end{subequations}
These equations extends to
\begin{subequations}
\begin{align}
  \beta(s)^2-\alpha(s)^2&=b\cdot b\\
\alpha(s)J\cdot a+\beta(s)L\cdot a&=b\cdot a.
\end{align}
\end{subequations}
There are four cases to be considered following that $a=J,L$ and $b=J,L$. The result is that
\begin{equation}
T_a\big(\tilde{e}_b \big)=\tilde{e}_b,
\end{equation}
in other terms, the vectors $\tilde{e}_J$ and $\tilde{e}_L$ are not only parallel vector fields, but each is parallel along the path of the other.

\subsection{Covariant derivative}
%--------------------------------

We will give the horizontal lift of $\tilde{e}_a$ at point
\[ 
  \xi(0)=\{ B_1^b\tilde{e}_b,B_2^c\tilde{e}_c \}_{r_0e_t}
\]
under the form of the path
\[ 
  \xi(s)=\{ B_1^b\tilde{e}_b\big( a(s) \big),B_2^c\tilde{e}_c\big( a(s) \big) \}_{\tilde{e}_a(s)}.
\]
We create a connection on the spinor bundle from the connexion via the formula
\[ 
  \widehat{\nabla_X^E\psi}(\xi)=\overline{ X }_{\xi}(\hat \psi).
\]
In our case, we take $\psi\colon M\to \mS$, or $\hat{\psi}\colon SU(1,1)\to \Lambda W$ such that
\[ 
  \hat{\psi}(U\cdot g)=\rho(g^{-1})\hat{\psi}(U).
\]
Since $\tilde\omega=\varphi^*\omega$, we have $\tilde\omega(X)=\omega(d\varphi X)$ and
\[ 
  \overline{ e }_a{}_{\xi_0}=\varphi^{-1}\big( \tilde{e}_a(s),\ldots \big).
\]
Therefore
\begin{equation}  \label{eq_whidpsinabla}
\widehat{\nabla_a\psi}(\xi_0)=\Dsdd{ (\hat{\psi}\circ\varphi^{-1})\{ B_i^c\tilde{e}_c\big( a(s) \big) \}_{\tilde{e}_a(s)} }{s}{0}
\end{equation}
where $\varphi$ is defined by
\[ 
  \varphi(U)=\{ U\tilde{e}_J,U\tilde{e}_L \}_{Ur_0e_t}.
\]
We have to find
\begin{equation}  \label{eq_varpBic}
  \varphi^{-1}\{ B_i^c\tilde{e}_c\big( a(s) \big) \}_{\tilde{e}_a}.
\end{equation}
Before to write down the inverse of $\varphi$, let us perform some computations.
\[ 
  J=\begin{pmatrix}
&0\\
0&0&1\\
&1
\end{pmatrix},\quad
L=\begin{pmatrix}
0&1&1\\
-1\\
1
\end{pmatrix},
\]
and as far as we only wants to compute derivatives, we can write the exponentials as
\begin{align} 
 e^{sJ}&=\mtu+sJ=\begin{pmatrix}
1\\
&1&s\\
&s&1
\end{pmatrix}\\
 e^{sL}&=\mtu+sL=\begin{pmatrix}
1&s&s\\
-s&1&0\\
s&0&1
\end{pmatrix}.
\end{align}
The path are given by
\begin{equation}
\tilde{e}_a(s)=r_0 e^{-sa}\begin{pmatrix}
0\\1\\0
\end{pmatrix},
\end{equation}
in particular
 \begin{align}
\tilde{e}_J(s)&=r_0\begin{pmatrix}
0\\1\\-s
\end{pmatrix},
&\tilde{e}_L(s)&=r_0\begin{pmatrix}
-s\\1\\0
\end{pmatrix}.
\end{align} 
For the various $\tilde{e}_b\big( a(s) \big)$, we have
\begin{equation}
\tilde{e}_b\big( a(s) \big)=\Dsdd{ a(s) e^{-tb} }{t}{0}\begin{pmatrix}
0\\1\\0
\end{pmatrix}\\
	=\Dsdd{ r_0 e^{-sa} e^{-tb} }{t}{0}\begin{pmatrix}
0\\1\\0
\end{pmatrix}\\
	=-r_0 e^{-sa}b\begin{pmatrix}
0\\1\\0
\end{pmatrix}.
\end{equation}
Results are
\begin{subequations}
\begin{align}
  \tilde{e}_J\big( J(s) \big)&=-r_0\begin{pmatrix}
0\\-s\\1
\end{pmatrix}
&
\tilde{e}_J\big( L(s) \big)&=-r_0\begin{pmatrix}
-s\\0\\1
\end{pmatrix}\\
\tilde{e}_L\big( J(s) \big)&=-r_0\begin{pmatrix}
1\\0\\0
\end{pmatrix}
&
\tilde{e}_L\big( L(s) \big)&=-r_0\begin{pmatrix}
1\\s\\-s
\end{pmatrix}.
\end{align}
\end{subequations}
We finally have to know that
\[ 
  B^c\tilde{e}_c\big( J(s) \big)=-r_0\begin{pmatrix}
B^L\\-sB^J\\B^J
\end{pmatrix},
\quad
B^c\tilde{e}_c\big( L(s) \big)=-r_0\begin{pmatrix}
-sB^J+B^L\\sB^L\\B^J-sB^L
\end{pmatrix}.
\]

Following equation \eqref{eq_varpBic}, in order to write down $\widehat{\nabla_a\psi}$, we have to find $U(s)\in SU(1,1)$ such that
\begin{enumerate}
\item $Ur_0e_t=\tilde{e}_a(s)$,
\item $U\tilde{e}_J=B^c_1\tilde{e}_c\big( a(s) \big)$,
\item $U\tilde{e}_L=B^c_2\tilde{e}_c\big( a(s) \big)$.
\end{enumerate}
If $\overline{ f }$ and $\overline{ g }$ are vectors, solutions in $U$ of equation $Ur_0\overline{ f }=r_0\overline{ g }$ are $U=\AD(r_0)B$ where $B$ fulfils $B\overline{ f }=\overline{ g }$. In the case of $a=J$, the three conditions successively give
\begin{subequations}
\begin{align}
U&=\AD(r_0)\begin{pmatrix}
.&0&.\\
.&1&.\\
.&-s&.
\end{pmatrix}\\
U&=\AD(r_0)\begin{pmatrix}
.&.&B_1^L\\
.&.&-sB_1^J\\
.&.&B_1^J
\end{pmatrix}\\
U&=\AD(r_0)\begin{pmatrix}
-B^L_2&.&.\\
sB_2^J&.&.\\
B_2^J&.&.
\end{pmatrix}.
\end{align}
\end{subequations}
Putting all together in equation \eqref{eq_whidpsinabla} we find
\begin{equation} \label{eq_nabJmoi}
\begin{split}
  \widehat{\nabla_J\psi}(\xi_0)&=\frac{ d }{ ds }\hat{\psi}\AD(r_0)
\begin{pmatrix}
-B_2^L	&	0	&	B_1^L\\
sB_2^J	&	1	&	-sB_1^J\\
B_2^J	&	-s	&	B_1^J
\end{pmatrix}\\
	&=d\hat{\psi}\AD(r_0)\begin{pmatrix}
0&0&0\\
B_2^J&0&-B_1^J\\
0&-1&0
\end{pmatrix}.
\end{split}
\end{equation}
The same with $L$ instead of $J$ leads to
\begin{equation} \label{eq_nabLmoi}
\widehat{\nabla_L \psi}(\xi_0)=d\hat{\psi}\AD(r_0)\begin{pmatrix}
-B_2^J&-1&B_1^J\\
B_2^L&0&B_1^L\\
-B_2^L&0&-B_1^L
\end{pmatrix}.
\end{equation}

However it should be shocking to get $3\times 3$ matrices in $SU(1,1)$ : we had abused between $\SO(2,1)$ and $SU(1,1)$.

\subsection{Another way to write a section (wrong way to do)}
%------------------------------------------------------------

The equivariant function $\hat{\psi}\colon SU(1,1)\to \Lambda W$ fulfills 
\[ 
  \hat{\psi}(U\cdot g)=\rho(g^{-1})\hat{\psi}(U)
\]
for all $g\in\Spin(1,1)$; in particular with $g=-\mtu$,
\begin{equation}
 \hat{\psi}(-U)=-\hat{\psi}(U).
\end{equation}
This gives the idea that it is not impossible to define $\hat{\psi}$ from its projection on $\SO(2,1)$ : we want to get $\tilde{\psi}\colon \SO(2,1)\to \Lambda W$ and define 
\[ 
  \hat{\psi}(U)=\tilde{\psi}\big( T(U) \big).
\]
More precisely, we parametrize $SU(1,1)$ by $\alpha$ and $\beta$ such that $| \alpha |^2-| \beta |^2=1$. Then we divide $SU(1,1)$ into two parts: $\alpha=x+iy$ is green when $x>0$ and when $x=0$, $y<0$; $\alpha$ is red when $x<0$ and when $x=0$, $y>0$. When $\alpha=0$, we classify following $\beta$ in the same way. The result is that $U$ is green if and only if $-U$ is red. For a map $\tilde{\psi}\colon \SO(2,1)\to \Lambda W$, we define
\begin{equation}
\hat{\psi}(U)=
\begin{cases}
\tilde{\psi}\big( T(U) \big)&\text{if $U$ is green}\\
-\tilde{\psi}\big(T(U)\big)&\text{if $U$ is red}
\end{cases}
\end{equation}
We define $T^{-1}\colon \SO(2,1)\to SU(1,1)$ as follows: $T^{-1}(A)$ is the green element of $SU(1,1)$ whose image by $T$ is $A$. In any cases we have
\[ 
\hat{\psi}\circ T^{-1}=\tilde{\psi}.  
\]
The meaning of equations \eqref{eq_nabJmoi} and \eqref{eq_nabLmoi} is that $\AD(r_0)$ is a matrix whose inverse image by $T$ should be given to $\hat{\psi}$; the difficulty is to know which of the two. When $U_0$ is green,
\[ 
 \begin{split}
\widehat{\nabla_a\psi}(U_0)&=\Dsdd{ (\hat{\psi}\circ T^{-1})  \AD(r_0)\Big( \cdots \Big)   }{s}{0} \\
		&=\Dsdd{ \tilde{\psi} \AD(r_0)\Big( \cdots \Big)}{s}{0},
\end{split} 
\]
while when $U_0$ is red,
\[ 
  \widehat{\nabla_a\psi}(U_0)=-\Dsdd{ \tilde{\psi}\AD(r_0)\Big( \cdots \Big) }{s}{0}.
\]
These two show that
\begin{equation}
\widetilde{\nabla_a\psi}\big( T(U_0) \big)=\Dsdd{   \tilde{\psi}\AD(r_0)\Big( \cdots \Big)    }{s}{0}
\end{equation}
All this is only proved in the interior of the green and red regions so that the path $U(s)$ keeps on only one region.

\subsection{Once again}  \label{pg_DiracADsdeux}
%-------------------------

We see $AdS_2$ as\footnote{Here, $G=SL(2,\eR)$} $\mO=\Ad(G)H$ and we consider a base point $o=\Ad(k_0)H$ with $G=SL(2,\eR)=ANK$. Let the principal bundle
\[ 
\xymatrix{%
   A \ar@{~>}[r]^{R}		&	G\ar[d]^{\pi}\\
    &	\mO
}
\]
with $A$ acting on $G$ by $(a,g)\mapsto ga$ and the projection
\begin{equation}
\pi(rk_0a)=\Ad(rk_0a)H.
\end{equation}
where $r\in R$ and $a\in A$. More precisely, the principal bundle we look at is
\begin{equation}
\xymatrix{%
   A \ar@{~>}[r]^{R}		&	\mU_G\ar[d]^{\pi}\\
   	&	\mU
}
\end{equation}
where $\mU_G=Rk_0A$ and $\mU=\pi(\mU_G)=\Ad(Rk_0A)H=\Ad(Rk_0)H=\Ad(R)o$. The $\mU_G$ is so defined in order to be the $\pi^{-1}$ of an orbit $\mU=\Ad(R)o$.

We have a manifold isomorphism $R\simeq\mU$ given by
\[ 
  \phi\colon r\to \Ad(r)o.
\]
How to see a left invariant vector field on $R$ \emph{via} this identification ? 
\[ 
  d\phi\tilde X_r=d\phi\Dsdd{ r e^{tX} }{t}{0}
		=\Dsdd{ \Ad(r)\Ad( e^{tX})o }{t}{0}.
\]
This leads us to consider the following field for $X\in\sR$. We define $\xi_X(rk_0a)\in T_{rk_0a}\mU_G$,
\begin{equation}
  \xi_X(rk_0a)=\Dsdd{ r e^{tX}k_0a }{t}{0}.
\end{equation}
Let's see the projection :
\[ 
\begin{split}
d\pi_{rk_0a}\xi_X(rk_0a)&=\Dsdd{ \pi(r e^{tX}k_0a) }{t}{0}\\
		&=\Dsdd{ \Ad(r e^{tX}k_0a)H }{t}{0}\\
		&=\Dsdd{ \Ad(r e^{tX})o }{t}{0}.
\end{split} 
\]
This gives us the idea to define $X^{\sharp}\in T_{\Ad(rk_0a)H}\mU=T_{\pi(rk_0a)}\mU$ by
\begin{equation}
X^{\sharp}_{rk_0a}=\Dsdd{ \Ad(re^{tX})o }{t}{0},
\end{equation}
which is a good definition because $\pi(rk_0a)=\pi(r'k_0a')$ only when $r=r'$. We put the following connection on $\mU_G$ :
\begin{equation}
\alpha_{rk_0a}(\Sigma)=\left[     \big( dL_{rk_0a}^{-1} \big)_{rk_0a}\Sigma    \right]_{\sA}.
\end{equation}
We hope $\xi_X$ to be the horizontal lift\footnote{We will see in proposition \ref{prop_horliftXdiz} that it is not the case, but for the moment, we hope it.} of $X^{\sharp}$; by construction $d\pi\xi_X=X^{\sharp}$. We have
\[ 
 \begin{split}
\alpha_{rk_0a}(\xi_X)&=\left[ dL_{(rk_0a)^{-1}}\xi_X(rk_0a) \right]_{\sA}\\
		&=\Dsdd{ a^{-1}k_0^{-1}r^{-1}r e^{tX}k_0a }{t}{0}^{\sA}\\
		&=\Dsdd{ a^{-1}\AD(k_0^{-1}) e^{tX}a}{t}{0}^{\sA}\\
		&=\left[ \Ad(a^{-1}k_0^{-1})X \right]_{\sA}.
\end{split} 
\]
One can, by brute force computation\footnote{Or by remarking that $\sA$ is abelian.}, show that the difference $\Ad(ak_0)X-\Ad(k_0)X$ is skew-diagonal when
\[ 
  X=\begin{pmatrix}
a'&n\\0&-a'
\end{pmatrix},
\quad k_0=\begin{pmatrix}
\cos a&\sin a\\\sin a&\cos a
\end{pmatrix},
\quad
a=\begin{pmatrix}
a&0\\0&1/p
\end{pmatrix}.
\]
So $\Ad(a)$ does not change the $\sA$-component of $\Ad(k_0)X$. We conclude that
\begin{equation}
  \alpha(\xi_X)=\left[ \Ad(k_0^{-1})X \right]_{\sA}.
\end{equation}
When $X\in\sR$, we consider $\tilde X_g=(dL_g)_eX$;
\begin{equation} 						 \label{eq_defXtilde}
  \tilde X_{rk_0a}=dL_{rk_0a}X,
\end{equation}
in particular, $\tilde X_r=\Dsdd{ r e^{tX} }{t}{0}$.
We denote by $\tau$ the action
\begin{equation}
\begin{aligned}
 \tau_g\colon\mO&\to \mO \\ 
\tau_g\Ad(r)H&= \Ad(gr)H 
\end{aligned}
\end{equation}
In particular
\[ 
 \begin{split}
d\pi_gdL_g Y&=\Dsdd{ \pi(g e^{tY}) }{t}{0}\\
		&=\Dsdd{ \ad(g e^{tY}H) }{t}{0}\\
		&=\Dsdd{ \tau_g\Ad( e^{tY})H }{t}{0}\\
		&=(d\tau_g)_Hd\pi_e Y,
\end{split} 
\]
thus
\begin{equation}
d\pi\circ dL = d\tau\circ d\pi.
\end{equation}
With definition \eqref{eq_defXtilde}, we have $\alpha(\tilde X)=X_{\sA}$ because
 \[ 
\alpha\big( dL_{rk_0a}X \big)=\left[ dL_{(rk_0a)^{-1}}dL_{rk_0a}X \right]_{\sA}
		=X_{\sA}.
\]
We are now able to find some horizontal lift. 
\begin{proposition}  \label{prop_horliftXdiz}
The horizontal lift of $X^{\sharp}$ is
\[
  \overline{ X^{\sharp} }=\xi_X-\widetilde{  [\Ad(k_0^{-1})X]_{\sA}  }.
\]

\end{proposition}

\begin{proof}
First, we have
\[ 
 \begin{split}
d\pi\overline{ X^{\sharp} }|_{rk_0a}&=d\pi\xi_X-d\pi(dL_{rk_0a})_e[\Ad(k_0^{-1})X]_{\sA}\\
				&=\Dsdd{ \Ad(r e^{tX}o) }{t}{0}-d\tau\,d\pi[\Ad(k_0^{-1})X]_{\sA}.
\end{split} 
\]
The first term is $X^{\sharp}$ while the second is zero because if $A\in\sA$, 
\[ 
 \begin{split}
d\pi A&=\Dsdd{ \pi( e^{tA}) }{t}{0}\\
		&=\Dsdd{ \Ad( e^{tA})H }{t}{0}\\
		&=0.
\end{split} 
\]
On the other hand,
\[ 
\alpha(\overline{ X^{\sharp} })=\alpha(\xi_X)-[\Ad(k_0)^{-1}X]_{\sA}=0.
\]
\end{proof}

Now we prove that the function $\overline{ X^{\sharp} }\cdot\hat{\psi}$ is equivariant, and therefore that the definition
\[ 
  \widehat{\nabla_{X^{\sharp}}\psi}=\overline{ X^{\sharp} }\cdot\hat{\psi}
\]
works. Using equivariance of $\hat{\psi}$, 
\begin{equation}
\begin{aligned}
  \overline{ X^{\sharp} }\cdot \hat{\psi}(ga_1)&=\Dsdd{ \hat{\psi}(\xi_X(t) }{t}{0}-\hat{\psi}\big( dL_{ga_1}[\Ad(k_0^{-1})X]_{\sA} \big)\\
		&=\Dsdd{ \hat{\psi}\big( r e^{tX}k_0aa_1 \big) }{t}{0}-\Dsdd{ \hat{\psi}\big( ga_1 e^{t[\Ad(k_0^{-1})X]_{\sA}} \big) }{t}{0}\\
		&=\Dsdd{ \rho(a_1)\hat{\psi}\big( r e^{tX}k_0a \big) }{t}{0}-\Dsdd{ \rho(a_1)\hat{\psi}\big( g e^{t[\Ad(k_0^{-1})X]_{\sA}} \big) }{t}{0},\\
		&=\rho(a_1)\hat{\psi}(\xi_X)-\Dsdd{ \rho(a_1)\hat{\psi}\big( \widetilde{[\Ad(k_0^{-1})X]_{\sA} }|_g \big) }{t}{0}\\
		&=\rho(a_1)\hat{\psi}(\xi_X)-\rho(a_1)\widetilde{ [\Ad(k_0^{-1})X]_{\sA}  }\hat{\psi}(g)\\
		&=\rho(a_1)(\overline{ X^{\sharp} }\cdot\hat{\psi})(g).
\end{aligned}
\end{equation}
for the third line, we used the fact that $\sA$ is abelian
\subsubsection{Clifford algebra for \texorpdfstring{$AdS_2$}{AdS2}}
%----------------------------------------------------------------

Our basis of $\sA\oplus\sN$ is
\[ 
  H=\begin{pmatrix}
1&0\\0&-1
\end{pmatrix}\quad\text{and}\quad
E=\begin{pmatrix}
0&1\\0&0
\end{pmatrix}
\]
and we choose
\[ 
  o=\Ad(k_0)H=\cos(2k_0)H+\sin(2k_0)(E+F).
\]
Since (at first order in $t$) $\Ad( e^{tH})o=\cos(2k_0)\mtu+\sin(2k_0)\big( E+2tE+F-2tF \big)$,
\[ 
  H^{\sharp}_{rk_0a}=2\sin(2k_0)\Ad(r)(E-F),
\]
and
\[ 
  E^{\sharp}_{rk_0a}=\Ad(r)\big( -2\cos(2k_0)E+\sin(2k_0)H \big).
\]
We have to compute the metric matrix for this basis; we know from equation \eqref{eq_KillAdinvariant}, the Killing form is $\Ad$-invariant and $(\mathfrak{sl}(2,\eR),B)\simeq(\eR^3,\eta_{21})$. So the $\Ad(r)$ disappears in the computation of $B(X^{\sharp},Y^{\sharp})$. We get
\[
\begin{split}
B(H^{\sharp},H^{\sharp})&=4\sin^2(2k_0)B(E-F,E-F)\\
		&=-32\sin^2(2k_0)\\
B(E^{\sharp},E^{\sharp})&=\sin^2(2k_0)B(H,H)\\
		&=8\sin^2(2k_0)\\
B(E^{\sharp},H^{\sharp})&=-4\sin(2k_0)\cos(2k_0)B(E,E-F)+2\sin^2(2k_0)B(H,E-F)\\
		&=16\sin^2(2k_0)\cos(2k_0).
\end{split}
\]
So the metric is in the basis $\{ H^{\sharp},E^{\sharp} \}$
\begin{equation}
g=
\begin{pmatrix}
-32\sin^2(2k_0) & 16\sin(2k_0)\cos(2k_0)\\
16\sin(2k_0)\cos(2k_0) & 8\sin^2(2k_0)
\end{pmatrix}.
\end{equation}
When we consider the orbit of $E+F$, we choose $o=E+F$, i.e. $\cos(2k_0)=0$, $\sin(2k_0)=1$ so that
\begin{equation}
H^{\sharp}_{rk_0a}=2\Ad(r)(E-F),\quad E^{\sharp}_{rk_0a}=\Ad(r)H,
\end{equation}
and
\[ 
  g=\begin{pmatrix}
-32&0\\0&8
\end{pmatrix};
\]
in the case of the orbit of $-(E+F)$, we get the same. The negative vector is $H^{\sharp}$ and the positive one is $E^{\sharp}$.

\subsubsection{Identification \texorpdfstring{$\sQ\leftrightarrow\Lambda W$}{QW}}
%////////////////////////////////////////////////////////////////////////////////

We want a linear bijection $\phi\colon \sQ\to \Lambda W$ such that
\[ 
  \rho(s)\phi(X)=\phi\big( \rho(s)X \big)
\]
where the left hand side action of $\Spin$ is the usual on $\Lambda W$ while the right hand side one remains to be defined. The implementation of this is easy: we can take any bijection between $\sQ$ and $\Lambda W$ and define
\begin{equation}
  \rho(s)X=\phi^{-1}\big( \rho(s)\phi(X) \big).
\end{equation}

Spinors on $AdS_2$ are given by equivariant functions $\hat{\psi}\colon \mU_G\to \Lambda W$ which are now replaced by $\tilde{\psi}\colon R\to \sQ\simeq\Lambda W$ by
\[ 
  \hat{\psi}(rk_0a)=\rho(a^{-1})\tilde{\psi}(r).
\]
So the set of sections of the spinor bundle over $\mU$ is
\[ 
  \Gamma_{\mU}\simeq  C^{\infty}(R,\Lambda W).
\]

\subsubsection{Covariant derivative}
%////////////////////////////////

The aim is now to compute 
\[ 
\begin{split}
   \widetilde{\nabla_{X^{\sharp}}\psi}(r)&=\widehat{\nabla_{X^{\sharp}}\psi}(rk_0)\\
		&=\overline{ X^{\sharp} }\cdot \hat{\psi}|_{rk_0}\\
		&=\big( \xi_X-\widetilde{ [\Ad(k_0^{-1})X]_{\sA}  } \big)\cdot\hat{\psi}|_{rk_0}\\
		&=\Dsdd{ \tilde{\psi}(r e^{tX}) }{t}{0}-\Dsdd{ \rho\big(  e^{t[\Ad(k_0^{-1})X]_{\sA}} \big) }{t}{0}\tilde{\psi}(r)\\
		&=\tilde X_r\tilde{\psi}(r)-d\rho_e\big( [\Ad(k_0^{-1})X]_{\sA} \big)\tilde{\psi}(r).
\end{split}  
\]
Our final formula for the covariant derivative is
\begin{equation}
  \widetilde{\nabla_{X^{\sharp}}\psi}(r)=\tilde X_r\tilde{\psi}-d\rho\big( [\Ad(k_0^{-1})X]_{\sA} \big)\tilde{\psi}.
\end{equation}
The Dirac operator will be a linear combination of vectors of the form
\[ 
  \tilde X+d\rho\big( [\Ad(k_0^{-1})X]_{\sA}  \big).  
\]
Notice that $\tilde X$ is left invariant and the second term is even independent of the point, so the whole is left invariant.

\section{Dirac operator on \texorpdfstring{$AdS_{3}$}{AdS3}}  \label{PgDiracAdSTrois}
%--------------------------------------------------------------

The definition is 
\[ 
  AdS_{3}=\frac{ \SO(2,2) }{ \SO(1,2) },
\]
and the group which acts is the $AN$ of $\SO(2,2)$. The Lie algebra is given by
\[ 
\begin{split}
  \sA&=\{ J_{1},J_{2} \}\\
  \sN&=\{ M,L \}
\end{split}  
\]
which has dimension $4$. So there is a stabiliser. One can prove that for the open orbit of $u=\begin{pmatrix}
0&1\\-1&0
\end{pmatrix}$, the stabiliser is $\{  e^{aJ_{2}} \}$, i.e.
\begin{equation}
  [ e^{aJ_{2}}u]=[u].
\end{equation}
For the spin group, we find 
\[ 
  \Spin(2,1)\simeq SL_2^*(\eR),
\]
the group of $2\times 2$ matrices with determinant equals to $\pm 1$ (cf \cite{Michelson}). Let us recall that the isomorphism $AdS_{3}\simeq SL(2,\eR)$ is given by
\[ 
  SL(2,\eR)=\begin{pmatrix}
t+x&y-u\\
y+u&t-x
\end{pmatrix}
\]
with $u^{2}+t^{2}-x^{2}-y^{2}=1$. For sake of simplicity, we denote $SL(2,\eR)$ by $G$. It is explained in \cite{Clement} that the map
\begin{equation}
\begin{aligned}
 \psi\colon (G\times G)\times AdS_{3}&\to AdS_3 \\ 
(g_{1},g_{2})x&= g_{1}xg_{2}^{-1} 
\end{aligned}
\end{equation}
provides a local isomorphism $G\times G\simeq O(2,2)$. Moreover we have locally :
\[ 
  \frac{ G\times G }{ \eZ_{2} }\simeq \SO(2,2).
\]
At each point $x\in AdS_3$, we have an isomorphism
\[ 
  \SO(2,2)_{x}\simeq \SO(2,1)
\]
where $\SO(2,2)_{x}$ is the stabiliser of $x$ in $\SO(2,2)$. So we define the isomorphism
\[ 
  \chi_{x}\colon \Spin(2,1)\to \SO(2,2)_{x}
\]
which is a double covering. If $d\psi\colon \mG\oplus\mG\to \mathfrak{so}(2,2)$ is the isomorphism of \cite{Clement}, we define $\psi\colon G\times G\to \SO(2,2)$ by
\[ 
  \psi( e^{X})= e^{d\psi X},
\]
which is a good definition because the exponential is surjective on $G\times G$. For each $x\in AdS_3$, we consider the isomorphism
\[ 
  \phi_{x}\colon \SO(2,1)\to \SO(2,2)_{x}
\]
such that $\phi_{x}\big( \SO(2,1) \big)=\SO(2,2)_{x}$.

We define $\chi(s)_i\colon \Spin(2,1)\to \SO(2,2)$ by
\[ 
  \chi(s)=\chi(s)_1v\chi(s)_2.
\]
The choice of $\chi(s)_i$ is not unique. So we define the action of $\Spin(2,1)$ on $G\times G$ by
\begin{equation}
(g,h)\cdot s=\big( \chi(s)_1g,\chi(s)_2^{-1}h \big).
\end{equation}
Therefore we have
\[ 
\begin{split}
  \psi\big( (g,h)\cdot s \big)x&=\chi(s)_1gxh^{-1}\chi(s)_2\\
		&=\chi(s)\big( gxh^{-1} \big)\\
		&=\chi(s)\big(\psi(g,h)x\big).
\end{split}  
\]
\subsection{Spin structure on \texorpdfstring{$AdS_3$}{AdS3}}
%+++++++++++++++++++++++++++++++++++

From previous considerations, the first choice should be
\[ 
  P=\frac{ AN }{ S }\times\Spin(2,1),
\]
but it is easy to remark that $\sR'=\{ J_{1},M,L \}$ is a Lie algebra. So we use the corresponding Lie group $R$ instead of the homogeneous space $AN/S$ (these two are isomorphic). Thus the choice is
\begin{equation}
P=R'\times\Spin(2,1),
\end{equation}
with the projection $\pi\colon P\to AdS_3$,
\[ 
  \pi\big( r',s \big)=\left[ r'\begin{pmatrix}
0&1\\-1&0
\end{pmatrix} \right]
\]
 We consider
\begin{equation}
\begin{aligned}
 \theta\colon R'&\to \mU=Ro \\ 
r'&\mapsto ro=\left[ r'\begin{pmatrix}
0&1\\-1&0
\end{pmatrix} \right].
\end{aligned}
\end{equation}
The projection $\pi\colon P\to \mU$ reads $\pi=\theta\circ\pr_{1}$,
\[ 
  \pi\big( r',s \big)=[ro].
\]
This definition works because for all $a$, there exists a $h\in H$ such that
\[ 
   e^{aJ_{2}}\begin{pmatrix}
0&1\\-1&0
\end{pmatrix}=
\begin{pmatrix}
0&1\\-1&0
\end{pmatrix}h,
\]
from construction of $S$. Then we look at the following :
\[ 
\xymatrix{%
   \Spin(2,1) \ar@{~>}[r]&R'\times\Spin(2,1)\ar[rr]^{\displaystyle\varphi}\ar[rd]_{\displaystyle\pi}&&\SO\big( \mU \big)\ar[ld]&\SO(2,1)\ar@{~>}[l]	  \\
  &&\mU
}
\]
The action of $\Spin(2,1)$ on $P$ is
\[ 
  \big( r',s' \big)\cdot s=\big( r',s's \big).
\]
In order to define $\varphi$, we consider the isomorphism
\[ 
  \phi_{x}\colon \SO(2,1)\to \SO(2,2)_{x}
\]
between $\SO(2,1)$ and the stabiliser of $x$ in $\SO(2,2)$. This extends to an automorphism
\[ 
  \phi_{x}\colon \SO(2,2)\to \SO(2,2),
\]

\begin{probleme}
	I'm not sure of that extension, but we do not use it here.
\end{probleme}


and we define the action of $\SO(2,1)$ on $\SO\big( AdS_3 \big)$ by
\begin{equation}
\{ b_{i} \}_{x}\cdot g=\{ \phi_{x}(g)b_{i} \}_{x}.
\end{equation}
Then we define
\begin{equation}
  \varphi\big( r',s \big)=\{ \phi_{\pi r'}\big( \chi(s) \big)b_{i} \}_{\pi r'}
\end{equation}
 if $\{ b_{i} \}$ is a reference basis at $\pi[r]$. So this construction implies the choice of a section of $\SO\big( AdS_3 \big)$. Now, using the fact that both $\phi_{x}$ and $\chi$ are morphisms, we find 
\begin{equation}
\begin{split}
\varphi\big( (r',s)\cdot s' \big)&=\left\{ \phi_{\pi r'}\big( \chi(ss') \big) \right\}_{\pi r'}\\
			&=\left\{ \phi_{\pi r'}\big( \chi(s) \big)b_{i} \right\}\cdot\chi(s)\\
			&=\varphi(r',s)\cdot\chi(s').
\end{split}
\end{equation}
This proves that the construction gives a spin structure.

\subsection{Connection on the spinor bundle}
%-----------------------------------------------

A left invariant vector on $\mU$ is of the form
\[ 
  \tilde X_{xo}=\Dsdd{ x e^{tX}o }{t}{0}
		=\Dsdd{ \pi\big( x e^{tX},s \big) }{t}{0}
\]
for any $s\in\Spin(2,1)$. On $AdS_3$ (in fact on $\mU$) we consider the left invariant vector field
\begin{equation}
X^{\sharp}_{[x]}=\Dsdd{ [x e^{tX}] }{t}{0}
\end{equation}
which leads us to consider the following field on $P$ :
\begin{equation}
\xi_{X}\big( r',s \big)=\Dsdd{ r' e^{tX},s }{t}{0}\in T_{( r',s )}P.
\end{equation}
This defines a field which projects to the left invariant field on $\mU$ :
\begin{equation}   \label{eq_xiXprojXsharp}
d\pi\xi_{X}(r',s)=X^{\sharp}_{r'}.
\end{equation}  


\begin{lemma}
On the general vector 
\begin{equation}   \label{eq_gebevectSig}
  \Sigma=\Dsdd{ r'(t),s(t) }{t}{0},
\end{equation}
the formula
\begin{equation}
\alpha_{(r',s_{0})}\Sigma=-\Dsdd{ s_{0}^{-1}s(t) }{t}{0}\in\spin(2,1)
\end{equation}
where $s_{0}=s(0)$ defines a connection form.

\end{lemma}

\begin{proof}
First let $A\in\spin(2,1)$ and
\[ 
  A^*_{\xi}=\Dsdd{ \xi\cdot e^{-tA} }{t}{0}.
\]
We have
\[ 
   \alpha\big( A^*_{(r'),s_{0}} \big)=\alpha \Dsdd{ (r',s_{0})\cdot e^{-tA}  }{t}{0}
		=\alpha\Dsdd{ (r',s_{0} e^{-tA}) }{t}{0}
		=-\Dsdd{ s_{0}^{-1}s_{0} e^{-tA}} {t}{0}
		=A.
\]
Now we take back the vector $\Sigma$ of equation \eqref{eq_gebevectSig}, an element $a\in\Spin(2,1)$ and we compute
\[ 
\begin{split}
  (dR_{a}\alpha)\Sigma&=\alpha\Dsdd{ \big( r'(t),s(t) \big)\cdot a }{t}{0}\\
		&=\alpha\Dsdd{ \big( r'(t),s(t)a \big) }{t}{0}\\
		&=-\Dsdd{ a^{-1}s(0)^{-1}s(t)a }{t}{0}\\
		&=-\Ad(a^{-1})\Dsdd{ s(0)^{-1}s(t) }{t}{0}\\
		&=\Ad(a^{-1})\alpha(\Sigma).
\end{split}  
\]

\end{proof}
Thus that is a connection. This is however not the spin connection. Let $\beta$ be the Levi-Civita connection on the frame bundle $\SO(AdS_3)$. If 
\[ 
  \Sigma=\Dsdd{ r'(t),s(t) }{t}{0},
\]
we have
\begin{equation} \label{eQbetadphiSigma}
  \beta d\phi\Sigma=\left. \phi_{r'}\big( \chi(s_{0}) \big)^{-1}\Dsdd{ \phi_{r'(t)}\big( \chi(s_{t}) \big) }{t}{0}\right|_{\sH}.
\end{equation}
If we note $\phi_{r'(t)}\big( \chi(s_{t}) \big)=\phi\big( r'(t),\chi(s_{f}) \big)$, the derivative in \eqref{eQbetadphiSigma} with respect to $t$ reads
\begin{equation}
\Dsdd{ \phi\big( r'(t),\chi(s_{0}) \big) }{t}{0}+\Dsdd{ \phi\big( r',\chi(s_{t}) \big) }{t}{0}.
\end{equation}
The second term of $\beta d\varphi\Sigma$ is
\begin{align*}
\left. \Dsdd{ \phi_{r'}(\chi(s_{0}))^{-1}\phi_{r'}(\chi(s_{t})) }{t}{0}\right|_{\sH}
		&=\left. \Dsdd{ \phi_{r'}\big( \chi(s_{0}^{-1}s_{t}) \big) }{t}{0}\right|_{\sH}\\
		&=\left. d\phi_{r'}d\chi(s_{0}^{-1}s'(0))\right|_{\sH}.
\end{align*}
From all that we want to define
\begin{equation}
   \alpha^{S}_{(r',s_{0})}\Sigma=\left.d\phi d\chi\big(s_{0}^{-1}s'(0)\big)\right|_{\sH}+\left.\phi_{r'}\big( \chi(s_{0}) \big)^{-1}\Dsdd{ \phi_{r'(t)}\chi(s_{0}) }{t}{0}\right|_{\sH},
\end{equation}
and we would not have $\alpha^{S}(\xi_{X})=0$.


\subsection{Horizontal lift}
%------------------------------

Since the spin component of the path of $\xi_{X}$ is constant, we have $\alpha(\xi_{X})=0$, so equation  \eqref{eq_xiXprojXsharp} says that
\begin{equation}
\overline{ X^{\sharp} }=\xi_{X}.
\end{equation}
Let us recall that an equivariant function (which defined a section of an associated bundle) is
\begin{equation}
\begin{aligned}
 \hat{\psi}\colon P&\to V \\ 
\hat{\psi}(\xi\cdot g)&= \rho(g^{-1})\hat{\psi}(\xi). 
\end{aligned}
\end{equation}
General definition of an equivariant derivative (theorem \ref{tho_dercovassoequiv}) leads to
\[ 
  \widehat{    \nabla_{X^{\sharp}}\psi    }=\overline{ X^{\sharp} }\cdot\hat{\psi}=\xi_{X} \cdot \hat{\psi}.
\]
In our setting, the equivariance of $\hat{\psi}$ reads, for all $a\in\Spin(2,1)$, 
\[ 
  \hat{\psi}\big( ([r],s)\cdot a \big)=\hat{\psi}\big( [r],sa \big)\stackrel{!}{=}\rho(a^{-1})\hat{\psi}\big( [r],s \big).
\]
We check the equivariance of $\widehat{\nabla_{X^{\sharp}}\psi}$ by the following computation :
\[ 
\begin{split}
\widehat{\nabla_{X^{\sharp}}\psi  }\big( ([r],s)\cdot a \big)&=\widehat{\nabla_{X^{\sharp}}\psi}( [r],sa )\\
		&=(\xi_{X}\cdot \hat{\psi})([r],sa)\\
		&=\Dsdd{ \hat{\psi}\big( [r e^{tX}],sa \big) }{t}{0}\\
		&=\Dsdd{ \rho(a^{-1})\hat{\psi}\big( [r e^{tX}],s \big) }{t}{0}\\
		&=\rho(a^{-1})(\xi_{X}\cdot \hat{\psi})\big( [r],s \big)\\
		&=\rho(a^{-1})\widehat{  \nabla_{X^{\sharp}}\psi  }\big( [r],s \big).
\end{split}  
\]

We define $\tilde{\psi}\colon AN/S\to \Lambda W$ by 
\[ 
  \tilde{\psi}([r])=\hat{\psi}( [r],e ),
\]
so that
\begin{equation}
\hat{\psi}([r],s)=\rho(s^{-1})\tilde{\psi}([r]).
\end{equation}
We can conclude
\[ 
\begin{split}
\widetilde{ \nabla_{X^{\sharp}}\psi  }([r])&=\widehat{\nabla_{X^{\sharp}}\psi}([r],e)\\
		&=\xi_{X}\hat{\psi}([r],e)\\
		&=\Dsdd{ \hat{\psi}\big( [r e^{tX}],e \big) }{t}{0}\\
		&=\Dsdd{ \tilde{\psi}\big( [r e^{tX}] \big) }{t}{0}\\
		&=\tilde X_{[r]}\tilde{\psi}([r]).		
\end{split}  
\]
So
\begin{equation}
\widetilde{\nabla_{X^{\sharp}}\psi}=\tilde X_{[r]}\tilde{\psi}.
\end{equation}

\subsection{Spin structure on \texorpdfstring{$AdS_3$}{AdS3} }
%+++++++++++++++++++++++++++++++++++++++++++++++++++++++++

\subsubsection{Spin structure on the whole \texorpdfstring{$AdS_3$}{AdS3} }

\begin{probleme}
	The following seems to contradict what I find in Michelson-Donaldson
\end{probleme}
The central fact is that
\[ 
  \Spin(2,1)\simeq\Delta\simeq\SL(2,\eR)
\]
where $\Delta=\{ (g,g)\tq g\in\SL(2,\eR) \}\subset G_0$. We take as notations: $G_{0}=\SL(2,\eR)$ and $\overline{G}=G_0\times G_{0}$. 

\begin{lemma}
We have the following homogeneous space isomorphism :
\[ 
  \overline{G}/\Delta\simeq\SL(2,\eR).
\]
\end{lemma}

\begin{proof} 

We have an action $\overline{G}\times AdS_3\to AdS_3$,
\begin{equation} \label{EqActghgxh}
  (g,h)x=gxh^{-1}
\end{equation}
where $x\in\SL(2,\eR)$ is seen as in $AdS_3$ by the usual isomorphism. Moreover we consider the isomorphism
\begin{align}
\overline{G}/\Delta&\simeq\SL(2,\eR)\\
[x_1,x_2]&\mapsto x_1x_2^{-1}
\end{align}
which is well defined because $[x_1g,x_2g]\mapsto x_1gg^{-1}x_2^{-1}=x_1x_2^{-1}$. In particular, $[g,g]\mapsto e\in\SL(2,\eR)$. So $\overline{G}$ acts on $\SL(2,\eR)$ and the elements which fix $e$ are the one of $\Delta$. It proves the lemma.
\end{proof}

We are going to take the following structure :
\begin{equation}  \label{EqScSpinAdS}
  \xymatrix{%
   \overline{G} \ar[rr]^{\displaystyle\varphi}\ar[dr]_{\displaystyle\pi}	&	&	\overline{G}/\eZ_{2}\ar[ld]\\
   						& M	
}
\end{equation}
where $M$ is $G_0$ seen as $M=\overline{G}/\Delta\simeq \SL(2,\eR)\simeq AdS_3$, and the projection $\pi\colon \overline{G}\to M$ is given by $\pi(g,h)=gh^{-1}$. The action of $\Delta\simeq \Spin(2,1)$ on $\overline{G}$ is given by formula $(xg,g)\cdot (a,a)=(xga,ga)$. First, let us prove the following.
\begin{proposition}
The frame bundle over $AdS_3$ can be seen as 
\[ 
  \SO(AdS_3)\simeq \overline{G}/\eZ_{2}
\]
where $\overline{G}=\SL(2,\eR)\times\SL(2,\eR)$.
\end{proposition}

\begin{proof}
In the fiber bundle $\pi\colon \overline{G}\to M$, the fibre over $x\in\SL(2,\eR)$ is the set of $(g,h)$ such that $gh^{-1}=x$, or
\[ 
  \overline{G}_{x}=\{ (xg,g) \}\subset \overline{G}.
\]
We will give a surjective map $\overline{G}_{x}\to\SO(M)_{x}$, the fibre of the frame bundle over $x\in AdS_3$. For this, we see a basis of $AdS_3$ as an isometric map $b\colon \sG_0\to T_{x}M$ where $\sG_{0}=\mathfrak{sl}(2,\eR)$, and we define
\begin{equation}
\begin{aligned}
 \psi_{x}\colon \overline{G}_{x}&\to \SO(M)_{x} \\ 
\psi_{x}(xg,g)(X)&=(dL_{x})_{e}\big( \Ad(g^{-1})X \big)
\end{aligned}
\end{equation}
for all $X\in\sG_0$. Let us study the kernel of this map, i.e. elements such that $\psi(xg_1,g_1)=\psi(xg_2,g_2)$. It needs, for all $X\in\mathfrak{sl}(2,\eR)$,
\[ 
  \Ad(g_1^{-1})X=\Ad(g_2^{-1})X,
\]
but we know that the requirement $\Ad(g)X=X$ is the fact the $g$ is in the center of the group. In our case, it results that $g_2^{-1}g_1=\pm\id$, so
\[ 
  \psi(xg_1,g_1)=\psi(\pm xg_1,\pm g_1)
\]
where the same $\pm$ has to be taken in both appearances of the right hand side. Now we put all the $\psi_{x}$ together to get $\psi\colon \overline{G}\to \SO(M)$. Once again we look in which cases $\psi(g_1,h_1)=\psi(g_2,h_2)$. We put this condition under the form
\[ 
  \psi(g_1h_1^{-1}h_1,h_1)=\psi(g_2h_2^{-1}h_2,h_2)
\]
which immediately gives $h_1=\pm h_2$. But on the other hand the base point of $\psi(g_{i}h_{i}^{-1},h_{i})$ is $g_{i}h_{i}^{-1}$, so that the condition also ask $g_1h_1^{-1}=g_2h_2^{-1}$ which in turn gives $g_1=\pm g_2$ with the same $\pm$ as in $h_1=\pm h_2$. We conclude that $\eZ_{2}$ is the problem for the inverse of $\psi$. This proves the proposition.
\end{proof}
We will usually use the same notation, $\psi$, to denote the map from $\overline{G}$ and the one from $\overline{G}/\eZ_{2}$. The following lemma will prove useful to study the actions of the structure groups in the picture \eqref{EqScSpinAdS}.

\begin{lemma}
The map
\begin{equation}
\begin{aligned}
 \SL(2,\eR)&\to \SO_{0}(1,2)\\
   g&\mapsto\Ad(g).
\end{aligned}
\end{equation}
is a double covering.
\end{lemma}
\begin{proof}
No proof.
\end{proof}
The action of $a\in\SO_{0}(1,2)$ on $(xg,g)\in\overline{G}/\eZ_{2}$ is defined by
\begin{equation}
\psi\big( (xg,g)\cdot a \big)=(dL_{x})_{e}\Ad(a^{-1}g^{-1}).
\end{equation}
On the other hand, let us see how does $(a,a)\in\Delta\simeq \Spin(2,1)$ acts on $\overline{G}$ and how does it reflects on the $\psi$ level. Since $(xg,g)\cdot (a,a)=(xga,ga)$, we have
\[ 
  \psi\big( [xg,g]\cdot a \big)=\psi\big( (xg,g)\cdot(a,a) \big),
\]
and then
\[ 
  \varphi\big( (xg,g)\cdot a \big)=(xg,g)\cdot(a,a).
\]
This proves that our structure is a spin structure.

\subsubsection{Reduction to one open orbit}
%----------------------------------------

We will use this isomorphism between $AdS_3$ and $\SL(2,\eR)$ :
\[ 
  \begin{pmatrix}
u\\t\\x\\y
\end{pmatrix}\mapsto
\begin{pmatrix}
u+x&y-t\\y+t&u-x
\end{pmatrix}.
\]
Then the famous point $[u]=\begin{pmatrix}
0&1\\-1&0
\end{pmatrix}\in AdS_3$ corresponds to the element $J:=\begin{pmatrix}
0&1\\-1&0
\end{pmatrix}\in \SL(2,\eR)$. This is our base point of the open orbit. We could also take
\[ 
  k_{0}=\frac{ \sqrt 2 }{ 2 }\begin{pmatrix}
1&1\\-1&1
\end{pmatrix}\in K_{0}
\]
where $K_{0}$ is the ``$K$'' of $\SL(2,\eR)$. 
\begin{probleme}
	I think that $J$ is also a complex structure. To be checked.
\end{probleme}
We have $J=k_{0}^{2}$ and following the action \eqref{EqActghgxh}, we have $J=(k_{0},k_{0}^{-1})e$. The subgroup $\overline{R}\subset\overline{G}$ acts on $AdS_3$, and we want to know the stabilizer of $J$. The condition is $(r,r')\cdot J=J$, or
\[ 
  r=\AD(J)r',
\]
but $\AD(J)=\theta$ (the Cartan involution). So an element $(r,r')\in\overline{R}$ stabilises $J$ if it is of the form $(r,\theta r)$, thus
\[ 
  \mfs=\text{Lie algebra of the stabiliser of $J$}=\{ (X,\theta X)\tq X\in\sR_{0} \}\cap\sR,
\]
where the intersection with $\sR$ is important because $\theta$ can send out of $\sR_{0}$. Note that when $X$ has a $\sN$ component, then $\theta X$ has a $\overline{ \sN }$ component, so $(X,\theta X)\in(\sA\oplus\sN,-\sA\oplus\overline{ \sN })$ where the minus sign comes from the fact that $\theta(\sA)=-\sA$. Then $X$ cannot have a $\sN$ component and finally,
\[ 
  \mfs=\eR (H,-H)\in\sQ.
\]
The group $R'$ is
\begin{equation}
R'= e^{\sR'}=\{ (an,an')\tq n,n'\in N_{0} \}
\end{equation}
because $\sR'$ is $\sR$ minus the stabiliser, i.e. $\sR'=\eR(H,H)\oplus\sN$. We have the identification $r'\mapsto r'\cdot J$ between $R'$ and the open orbit $\mU$. As usual, the action is $(g,h)\cdot x=gxh^{-1}$ if $r'=(g,h)$. Notice in particular that $R'\neq R_{0}'\times R_{0}'$.

Up to now we studied the fiber $\overline{G}\to M$; we are now able to restrict it to $\overline{G}|_{\mU}\to\mU$ and to establish an isomorphism with the trivial bundle $R'\times G_0\to R'$. The fiber over $x\in\mU$ is
\[ 
  \overline{G}_{x}=\{ (xg,g) \}.
\]
We define the isomorphism as follows:
\begin{equation}
\begin{aligned}
 \tau\colon R'\times G_0&\to \overline{G}|_{\mU} \\ 
(r',g)&\mapsto (r'\cdot Jg,g) 
\end{aligned}
\end{equation}
and we have the following picture: 
\[ 
  \xymatrix{%
   R'\times G_0 \ar[r]^-{\displaystyle\tau}\ar[d]	&	\overline{G}|_{\mU}\ar[d]^{\displaystyle\pi}\\
   R' \ar@{.>}[r]^{\displaystyle\tau}		&	\mU
}
\]
in which are defines by
\[ 
  \xymatrix{%
   (r',g) \ar[r]^-{\displaystyle\tau}\ar[d]	&	(r'\cdot Jg,g)\ar[d]^{\displaystyle\pi}\\
   r' \ar@{.>}[r]^{\displaystyle\tau}		&	r'\cdot J
}
\]
where the dotted line denotes the induced map from $\tau$, which is denoted by the same symbol. The map $\tau\colon R'\to \mU$ is just the restriction of the original $\tau$ to $g=e$. Notice that this $\tau$ provides a diffeomorphism of the basis spaces $R'$ and $\mU$.

\subsubsection{Spin connection}
%---------------------------

The spin connection on $\overline{G}|_{\mU}$ is given by
\begin{equation}    \label{EqDefConnAdS3}
  \alpha^{S}_{(g,h)}\Sigma=\left[ dL_{(g,h)^{-1}}\Sigma \right]_{\sH},
\end{equation}
or
\begin{equation}
\alpha^{S}_{(g,h)}=\pr_{\sH}\circ\big( dL_{(g,h)^{-1}} \big)_{(g,h)}.
\end{equation}
Notice that when we write $\sH$, we think about $\Delta$ : the group by which quotient  $\overline{G}$ in order to get $\SL(2,\eR)\simeq AdS_3$.
Our task now is to transfer this connection to $R'\times G_0$ by defining $\alpha'=\tau^*\alpha^{S}$. If $\Sigma\in T_{(r',g)}(R'\times G_0)$, we define
\begin{equation}
\alpha'_{(r',g)}\Sigma=\alpha^{S}(d\tau\Sigma).
\end{equation}
Let us take $X\in\sG_0$ and $0\in\sR'$ and let us compute $d\tau(0,X)$. More precisely, we consider
\begin{align*}
d\tau(0\oplus \tilde X_{g})_{(r',g)}&=d\tau\Dsdd{ r',g e^{tX} }{t}{0}\\
		&=\Dsdd{ r'\cdot Jg e^{tX},g e^{tX} }{t}{0}\\
		&=\big( \tilde X_{(r'\cdot Jg)},\tilde X_{g} \big).
\end{align*}
The next step is to compute $d\tau\Sigma$ in the case where $\Sigma=(\utilde Y\oplus -1)_{(r',g)}$ with $Y\in\sR'\subset\sR_{0}\oplus\sR_{0}$. We have
\begin{align}
d\tau\Sigma&=\Dsdd{ \tau\big(  e^{tY}r',g \big) }{t}{0}\\
	&=\Dsdd{ ( e^{tY}r'\cdot J)g,g }{t}{0}
\end{align}
where, if $r'=(r_{1},r_{2})$, we consider $Y=\big((\utilde Y_{1})_{r_{1}},(\utilde Y_{2})_{r_{2}}\big)  $. This appears to be difficult to be computed. This reflects the fact that the connection should be complicated in the trivial bundle $R'\times G_0$.

But there are no fate. We remember that $\tau$ furnish a diffeomorphism between the basis spaces, so one can consider the bundle 
\[ 
  \xymatrix{%
   \overline{G}|_{\mU} \ar[d]^{\displaystyle\tau^{-1}\circ\pi}\\		
   R'
}
\]
Vectors of $\sH$ are of the form $(X,X)$ with $X\in\mathfrak{sl}(2,\eR)$, thus $A\in T_{(xg,g)}\overline{G}|_{\mU}$ fulfils $\alpha^{S}(A)=0$ if and only if
\[ 
  dL_{(xg,g)^{-1}}(A)=(X,-X)
\]
for a certain $X\in \mathfrak{sl}(2,\eR)$. All this makes that the horizontal space over $(xg,g)$ is given by
\begin{equation}
\horsp(xg,g)=\big\{ (\tilde X_{xg},-\tilde X_{g})\tq X\in \sG_0=\mathfrak{sl}(2,\eR) \big\}.
\end{equation}
The strategy now is to project that on $R'$ and express Dirac operator in terms of the result. Let us make this simple computation :
\begin{align*}
d\pi(\tilde X_{xg},\tilde X_{g})&=\Dsdd{ \pi\big( xg e^{tX},g e^{-tX} \big) }{t}{0}\\
		&=\Dsdd{ xg e^{tX} e^{tX}g^{-1} }{t}{0}\\
		&=\Dsdd{ x e^{2t\Ad(g)X} }{t}{0}\\
		&=2(dL_{x})_{e}\Ad(g)X.
\end{align*}
This result has to be brought from $\mU$ to $R'$ by $\tau^{-1}$. Now we take a $\tilde Y\in\cvec(R')$ and we want to know which is the corresponding $X$, i.e. the $X\in\mathfrak{sl}(2,\eR)$ such that
\[ 
  d\tau^{-1}d\pi(\tilde X_{xg},-\tilde X_{g})=\tilde Y.
\]
From the previous computation, $\tilde Y=2d\tau^{-1}dL_{x}\Ad(g)X$, so
\begin{equation}  \label{EqXfracAdY}
  X=\frac{ 1 }{2}\Ad(g^{-1})dL_{x^{-1}}d\tau\tilde Y. 
\end{equation}
We now precise our idea: 
\begin{equation}   \label{EqtildeYrunrdeux}
  \tilde Y_{(r_{1},r_2)}=\big(    (\tilde Y_{1})_{r_1},(\tilde Y_{2})_{r_2}   \big)=\Dsdd{ r_1 e^{tY_{1}},r_2 e^{tY_{2}} }{t}{0}
\end{equation}
for $Y_{i}\in\sR'_{0}$ and $r_1$, $r_2\in R_{0}$. In this case, the ``$x$'' in equation \eqref{EqXfracAdY} is $(r'\cdot J)^{-1}$. Let us begin by taking $s'\in R'$ and compute $L_{(r'\cdot J)^{-1}}\tau(s')$. Remember that $r'\cdot J=r_1Jr_2^{-1}$ from the general action \eqref{EqActghgxh}, so if $r'=(r_1,r_2)$,
\begin{align*}
  dL_{(r'\cdot J)^{-1}}\tau(s')&=(r'\cdot J)^{-1}s_1 Js_2^{-1}\\
		&=(r_1Jr_2^{-1})^{-1}s_1Js_2^{-1}\\
		&=-r_2Jr_2^{-1}s_1Js_2^{-1}.
\end{align*}
Now, we apply that result on computation of \eqref{EqXfracAdY} with \eqref{EqtildeYrunrdeux} :
\begin{align*}
dL_{(r'\cdot J)^{-1}}d\tau\tilde Y&=\Dsdd{ -r_2Jr_1^{-1}r_1 e^{tY_{1}}J e^{-tY_{2}}r_2^{-1} }{t}{0}\\
		&=\Dsdd{ \AD(r_2)\big( -J e^{tY_{1}}J e^{-tY_{2}} \big) }{t}{0}\\
		&=\Dsdd{ \AD(r_2) e^{-tY_{2}} }{t}{0}+\Ad(r_2)\Ad(J)Y_{1}\\
		&=-\Ad(r_2)Y_{2}+\Ad(r_2)\theta(Y_{1}),
\end{align*}
and finally,
\begin{equation}  \label{EqValeurXAdtheta}
\begin{aligned}
X&=\frac{ 1 }{2}\Ad(g^{-1})dL_{(r'\cdot J)^{-1}}d\tau\tilde Y\\
		&=\frac{ 1 }{2}\Ad(g^{-1})\big[ \Ad(r_2)\theta(Y_{1})-\Ad(r_2)Y_{2} \big].
\end{aligned}
\end{equation}
For this $X$, the horizontal lift of $\tilde Y\in\cvec(R')$ is $(X,-X)\in T\overline{G}|_{\mU}$.

\subsection{Left invariance of Dirac}
%------------------------------------

Sections of the spin bundle over the open orbit $\mU$ are given by equivariant functions $\hat{\psi}\colon \overline{G}|_{\mU}\to \eR^{2}$. The action of $\Delta\simeq\Spin(2,1)$ on $\overline{G}$ is 
\[ 
  (g,h)\cdot(a,a)=(ga,ha).
\]
We define $\tilde{\psi}$ by
\begin{equation}
\tilde{\psi}(g)=\hat{\psi}(g,e)
\end{equation}
for $g\in\mU$. We get back the original $\hat{\psi}$ by formula
\begin{equation}
\hat{\psi}(g,h)=\rho(h,h)^{-1}\tilde{\psi}(gh^{-1}).
\end{equation}
Our intention is now to compute $\widehat{\nabla_{Z}\psi}(\xi)=\overline{ Z }_{\xi}(\hat{\psi})$ with $\xi=(xg,g)\in\overline{G}|_{\mU}$ (hence $x\in\mU$) and $Z\in\cvec(R')$. For instance we choose a left invariant $Z=\tilde Y=(\tilde Y_{1},\tilde Y_{2})$ for $Y_{1}$, $Y_2\in\sR_{0}'$. Recall that $\tilde Y$ is given by equation \eqref{EqtildeYrunrdeux}. From definition of the covariant derivative associated with the connection,
\[ 
  \widehat{\nabla_{\tilde Y}\psi}(\xi)=\overline{ \tilde Y }_{\xi}(\hat{\psi})=\overline{ \tilde Y }_{(xg,g)}(\hat{\psi})
\]
where $\overline{ \tilde Y }_{(xg,g)}$ is an horizontal vector at $(xg,g)$ whose projection is $\tilde Y$. From our previous work,
\[ 
  \overline{ \tilde Y }_{xg,g}=(\tilde X_{xg},-\tilde X_{g})
\]
with $X=\frac{ 1 }{2}\Ad(g^{-1})\big( \Ad(r_2)\theta Y_{1}-\Ad(r_2)Y_{2} \big)$. Let us understand the link between $(r_{1},r_2)$ and $g,x$. The vector $(\tilde X_{xg},\tilde X_{g})$ actually projects to a vector at $\tau^{-1}\circ\pi(xg,g)=\tau^{-1}(x)$. The fact that $x\in\mU$ guarantees existence and uniqueness of $(r_1,r_2)\in R'$ such that $r_1 Jr_2^{-1}=x$. We have
\[ 
\begin{split}
\protect\widetilde{\nabla_{\protect\tilde Y}\psi}(x)&=\widehat{\nabla_{\tilde{Y}}\psi}(x,e)\\
		&=(\tilde X_{x},-\tilde X_{e})\hat{\psi}\\
		&=\dsdd{ \hat{\psi}\Big( x e^{tX}, e^{-tX} \Big) }{t}{0}\\
		&=\dsdd{ \rho( e^{tX}, e^{tX})\tilde{\psi}(x e^{2tX}) }{t}{0}.
\end{split}  
\]
The first term of the derivation (the one with $t=0$ in the $\rho$) gives $2\tilde X_{x}\tilde{\psi}$. This is left invariant.
The second is 
\[
     \dsdd{ \rho( e^{tX}, e^{tX})\tilde{\psi}(x) }{t}{0}.
\]
We want to test the condition \eqref{EqDefLxinvarop} on this term. Let us pose
\[ 
  (E\tilde{\psi})(x)=(\tilde X_{x},\tilde X_{e})\hat{\psi}=\dsdd{ \rho( e^{tX}, e^{tX})\tilde{\psi}(x) }{t}{0}
\]
with $X$ given by equation \eqref{EqValeurXAdtheta}. On the one hand,
\begin{subequations}
\begin{equation}
L_{y}(E\tilde{\psi})(x)=(E\tilde{\psi})(yx)
		=\dsdd{ \rho( e^{tX_{a}}, e^{tX_{a}})\tilde{\psi}(yx) }{t}{0}
\end{equation}
with 
\begin{equation}
X_{a}=\frac{ 1 }{2}\Big( \Ad(r_2)\theta Y_{1}-\Ad(r_2) \Big)Y_{2}
\end{equation}
\end{subequations}
where $(r_1,r_2)$ is given by $yx$. On the other hand,
\begin{subequations}
\begin{equation}
E(L_{y}\tilde{\psi})(x)=\dsdd{ \rho( e^{tX_{b}}, e^{tX_{b}})\tilde{\psi}(yx) }{t}{0}
\end{equation}
with 
\begin{equation}
X_{b}=\frac{ 1 }{2}\big( \Ad(s_2)\theta Y_{1}-\Ad(s_2)Y_{2} \big)
\end{equation}
\end{subequations}
where $(s_1,s_2)$ is given by $x$.

The problem is that the choice of $y$ is arbitrary, so that $X_{a}$ and $X_{b}$ could be too different. Ok. That's the proof that Dirac is not invariant. Here is the proof that Dirac is invariant.

Following equation \eqref{EqDefConnAdS3}, the spin connection form is
\[ 
  \alpha^{S}_{(g,h)}\Sigma=\left( dL_{(g,h)^{-1}}\Sigma \right)_{\sH}.
\]
If $L_{(x,y)}$ is the left translation by $(x,y)$ we have
\[ 
  \big( L^{*}_{(x,y)}\alpha \big)_{(g,h)}\Sigma=\alpha_{(xg,yh)}\big( dL_{(x,y)}\Sigma \big)
		=\left( dL_{(g,h)^{-1}}\Sigma \right)_{\sH}.
\]
Thus we have $L^*_{(x,y)}\alpha^{S}=\alpha^{S}$. Now we consider the formula $\widehat{\nabla_{\tilde Y}\psi}(\xi)=(\tilde X_{xg},-\tilde X_{g})\hat{\psi}$, and we will check that 
\begin{equation}
 \big( L_{\eta}\widehat{\nabla_{Z}\psi} \big)(\xi)=\widehat{  \nabla_{Z}(L_{\eta}\psi)     }(\xi).
\end{equation}
with $\xi=(xg,g)$ and $\eta=(a,b)$. On the one hand,
\begin{align*}
  \big( L_{(a,b)}\widehat{\nabla_{Z}\psi} \big)(xg,g)&=\widehat{\nabla_{Z}\psi}(axg,bg)\\
		&=\widehat{\nabla_{Z}\psi}(axgg^{-1}b^{-1}bg,bg)\\
		&=\big( \tilde X_{(axb^{-1})bg},-\tilde X_{bg} \big)\hat{\psi}.
\end{align*}
On the other hand,
\begin{align*}
  \widehat{  \nabla_{Z}(L_{(a,b)}\psi)   }(xg,g)&=(\tilde X_{xg},-\tilde X_{g})\widehat{L_{(a,b)}\psi}\\
		&=\dsdd{ \widehat{L_{(a,b)}\psi}\big( xg e^{tX},g e^{-tX} \big) }{t}{0}\\
		&=\dsdd{ \hat{\psi}\big( axg e^{tX},bg e^{-tX} \big) }{t}{0}\\
		&= (\tilde X_{axg},-\tilde X_{bg}) \hat{\psi}\\
		&=( \tilde X_{(axb^{-1})bg},-\tilde X_{bg} )\hat{\psi}.	
\end{align*}



\section{Dirac operator on \texorpdfstring{$AdS_{4}$}{AdS4}}
%-----------------------------------------------------------

% Ok, il ne reste pas grand chose dans ce fichier, mais c'est appellé à grandir avec le temps.

\section{Dirac operator on \texorpdfstring{$AdS_{l}$}{AdSl}}
%+++++++++++++++++++++++++++++++++++++++++++++++++++++++++++

\subsection{Frame bundle}\index{frame!bundle!on $AdS_{l}$}
%------------------------

Construction of the frame bundle and the spin structure is a straightforward adaptation of theorem 2.2 (chapter ???) in \cite{AnnikFranc}, while Dirac operator and connection issues are adapted from proposition 1.3 (chapter III)

A \defe{basis}{basis} of a $m$ dimensional vector space $V$ is a free and generating part; it only has the structure of a set. A frame of the vector space $V$ is a nondegenerate map $b\colon \eR^{m}\to V$. Let us give an example in three dimensions the difference. If $\{ v_{1},v_{2},v_{3} \}$ is a basis of $V$, of course $\{ v_{2},v_{1},v_{3} \}$ is the same basis. Order has no importance. But if $\{ e_{1},e_{2},e_{3} \}$ is the canonical basis of $\eR^{3}$, the \emph{frames} $b(e_{1})=v_1$, $b(e_{2})=v_2$, $b(e_{3})=v_{3}$ and $c(e_{1})=v_2$, $c(e_{2})=v_1$, $c(e_{3})=v_{3}$ are not the same.

Now we consider $AdS_l=G/H=\SO(2,l-1)/\SO(1,l-1)$, the Lie algebra $\sG$ has a reductive homogeneous space decomposition $\sG=\sQ\oplus\sH$ and we consider the canonical projection $\pi\colon G\to AdS_l$.

Let the map (see relation \eqref{EqInclAdHSOq})
\begin{equation}
\begin{aligned}
 \alpha\colon H&\to \SO(\sQ) \\ 
h&\mapsto \Ad(h)|_{\sQ}.
\end{aligned}
\end{equation}
We consider, on $G\times\SO(\sQ)$, the equivalence relation $(g,A)\sim(g',A')$ if and only if there exists $h\in H$ such that $g'=gh$ and $A'=\alpha(h^{-1})A$. We denote by $G\times_{\alpha}\SO(\sQ)$ the set of equivalence classes. Now we have a principal bundle
\begin{equation}   \label{EqPrincPreB}
\xymatrix{%
   \SO(\sQ) \ar@{~>}[r]		&	G\times_{\alpha}\SO(\sQ)\ar[d]^{p}\\
   				&	   G/H
}
\end{equation}
where $p[g,A]=[g]$ and the action is given by $[g,A]\cdot B=[g,A]$. The fact that the projection fulfils $p\big( [g,A]\cdot B \big)=p[g,A]$ is evident, and the fact that the action is well defined is a simple computation :
if $[g',A']=[g,A]$, we have a $h\in H$ such that
\[ 
  [g',A']\cdot B=[g',A'B]
		=[gh,\alpha(h^{-1})AB]
		=[g,AB]
		=[g,A]\cdot B.
\]

\begin{proposition}
Let $\tau(g)\colon AdS_l\to AdS_l$ be the action of $g\in G$ on $AdS_l$ : $\tau(g)[g']=[gg']$, and $B$ be the frame bundle. We also consider the map $\sigma\colon \eR^{1,l-1}\to \sQ$ the isometry which sends the canonical basis of $\eR^{1,l-1}$ to the usual basis $\{ q_0,q_{1},\ldots,q_{l-1} \}$ of $\sQ$. The map
\begin{equation}
\begin{aligned}
 \beta\colon G\times_{\alpha} \SO(\sQ)&\to B \\ 
[g,A]&\mapsto d\tau(g)_{\mfo}A\circ\sigma 
\end{aligned}
\end{equation}
provides a principal bundle isomorphism between the principal bundle  \eqref{EqPrincPreB} and the frame bundle over $AdS_l$. 

\end{proposition}

By abuse of notation, we will not always write the $\sigma$.
\begin{proof}
We have to prove first that the map $\beta\colon G\times\SO(\sQ)\to B$ respects the classes. For that, consider $(g,A)\sim(g',A')$ and remark that
\[ 
\begin{split}
\beta(gh,\alpha(h^{-1}))&=d\tau(gh)_{\mfo}\alpha(h^{-1})A
		=d\tau(g)d\tau(h)d\pi \Ad(h^{-1})d\pi^{-1}A\\
		&=d\tau(g)d\tau(h)d\pi\Ad(h^{-1})d\pi^{-1} A
		=d\tau(g)d\pi dR_{h}d\pi^{-1}A\\
		&=d\tau(g)_{\mfo} A
		=\beta(g,A).
\end{split}  
\]
where we used equation \eqref{EqdpiAdpi} and the fact that $\pi\circ L_{g}=\tau(g)\circ\pi$. The frame bundle is
\begin{equation}   \label{EqPrincB}
\xymatrix{%
   \SO(1,l-1) \ar@{~>}[r]	&	B \ar[d]^{p}\\
   				&	   G/H
}
\end{equation}
where the fibre $B_{[g]}$ in $B$ over $[g]$ is the set of isometric maps $\eR^{1,l-1}\to T_{[g]}(AdS_l)$. So an element of $B$ is of the form $\big( [g],\tilde f\circ\sigma \big)$ where $g\in G$ and $\tilde f\colon \sQ\to T_{[g]}(AdS_l)$ contains the main information while $\sigma$ is the previously explained isometry. The action of $h\in\SO(1,l-1)$ on $\big( [g],\tilde f\circ\sigma \big)$ is defined by means of any fixed isomorphism $\varphi_{0}\colon \SO(1,l-1)\to \SO(\sQ)$ by
\begin{equation}
  \big( [g],\tilde f\circ\sigma \big)\cdot h=\big( [g],\tilde f\circ\varphi_{0}(h)\circ\sigma \big).
\end{equation}
The map $\beta$ is a morphism of principal bundle because
\[ 
\beta[g,A]\cdot\varphi_{0}^{-1}(B)	=\big( [g],d\tau(g)A\circ\sigma \big)\cdot \varphi_{0}^{-1}(B)
					=\big( [g],d\tau(g)A\circ B\circ \sigma \big)
					=\beta\big( [g,A]\cdot B \big).
\]
It remains to be proved that $\beta$ is a bijection. Surjectivity is natural: since $d\tau(g)$ is an isometry, $d\tau(g)A$ runs over the whole $\SO\big(T_{[g]}(AdS_l)\big)$ when $A$ runs over $\SO(\sQ)$. Injectivity is as follows; let's suppose $\beta[g,A]=\beta[g',A']$. It is immediate that in this case, $\exists h\in H$ such that $g'=gh$. Using the fact that $d\pi\circ dR_{h^{-1}}\circ d\pi^{-1}=\id$ and $d\tau(h)d\pi=d\pi dL_{h}$, we have
\[ 
d\tau(g)_{\mfo}A=d\tau(g)d\tau(h)A'
		=d\tau(g)d\tau(h)d\pi dR_{h^{-1}}d\pi^{-1}A'
		=d\pi\Ad(h)d\pi^{-1}A'
		=\alpha(h)A'.
\]

\end{proof}
From now on, we identify $G\times_{\alpha}\SO(\sQ)$ with the frame bundle over $AdS_l$. 

\subsection{Spin structure}
%--------------------------

We consider the principal bundle
\begin{equation}   
\xymatrix{%
   \Spin(1,l-1) \ar@{~>}[r]		&	G\times_{\tilde\alpha}\Spin(1,l-1)\ar[d]^{p}\\
   				&	   G/H
}
\end{equation}
where $\times_{\tilde\alpha}$ is the following equivalence relation on $G\times\Spin(1,l-1)$. We say that $(g,s)\sim(g',s')$ if and only if there exists a $h\in H$ such that
\begin{enumerate}
\item $g'=gh$,
\item $\chi(s')=\Ad(h^{-1})\chi(s)$.
\end{enumerate}
Notice that the second condition implies that $\Ad(h)\in \SO_0(\sQ)$. It is easy to prove that the given structure is well defined and is a principal bundle.  Now we consider the spin structure as follows:
\begin{equation}   
\xymatrix{%
   \Spin(1,l-1) \ar@{~>}[r]&G\times_{\tilde\alpha}\Spin(1,l-1)\ar[dr]\ar[rr]^-{\displaystyle\varphi}&&G\times_{\alpha} \SO(\sQ)\ar[dl]&\SO(\sQ)\ar@{~>}[l]  \\
   				&	   &G/H
}
\end{equation}
where $\varphi[g,s]=[g,\chi(s)]$. It is well defined since when $[g,s]=[g',s']$, there exists a $h\in H$ with $\chi(s')=\Ad(h^{-1})\chi(s)$ such that  $\varphi[g',s']=\varphi[gh,s']=[gh,\chi(s')]=[gh,\Ad(h^{-1})\chi(s)]=[g,\chi(s)]=\varphi[g,s]$.

\subsection{Reduction of the structural group}
%---------------------------------------------

The case of $AdS_l$ can be seen in the setting of subsection \ref{subsecCanConCovDer}. Let us show now that the bundle
\begin{equation}   \label{EqPrincHzGM}
\xymatrix{%
   H_0 \ar@{~>}[r]		&	G\ar[d]^{\pi}\\
   				&	G/H
 }
\end{equation}
is a reduction to $H_0$ (the identity component of $\SO(\sQ)$) of
\begin{equation}
\xymatrix{%
   G \ar@{~>}[r]		&	r(G)\ar[d]^{\pi}\\
   				&	G/H.
 }
\end{equation}
 Indeed, $u\colon G\to r(G)$ given by $u(g)=r(g)$ provides the reduction homomorphism: $r(gh)X=d\pi dL_{gh}X$ while $\big( r(g)\cdot h \big)X$ is the same.
 
\begin{lemma}
The tangent space $T(G/H)$ is an associated bundle of $r(G)$ trough the identification
\begin{equation}
\begin{aligned}
 \beta'\colon r(G)\times_{\rho}\sQ&\to T(G/H) \\ 
  [r(g),X]&\mapsto r(g)X 
\end{aligned}
\end{equation}
where $\rho(h)X=\Ad(h)X$, so that the quotient is given by $[g,X]=[gh,\Ad(h^{-1})X]$.
\end{lemma}
\begin{proof}
The proof is entirely similar to the one of lemma \ref{LemBazHGGH}.
\end{proof}


\chapter{Relativistic fields and group theory}
\section{Mathematical framework of field theory}
%++++++++++++++++++++++++++++++++++++++++++++++

    This is a short review; the aim is to see why the quantum theory of fields needs representations of the Poincaré group. It will be mostly physics oriented. References dealing with field theory including gauge theory and representations are \cite{Sternberg,Preparation,Naber,QMVirtanen,MQSenechal,GpAlgLie_Faraut,WormerAngular,Weinberg}.

\subsection{Axioms of the (quantum) relativistic field theory}
%-------------------------------------------------------------

The quantum mechanics is based on a few number of axioms:
\begin{enumerate}\label{pg:axiomes}
\item We have a Hilbert space $\pH$. A physical state is given by a \defe{ray}{ray} in $\pH$, i.e. a set 
       \[
           \rR=\{\xi\psi:|\xi|=1\}
       \]
for a certain $\psi\in\pH$ with $\scalh{\psi}{\psi}=1$. In other words, the set of physical sates is the quotient of the set of unital vectors in $\pH$ by the relation $\psi\sim\psi'$ if and only if $\psi=\xi\psi'$ for some unimodular complex number $\xi$. We denote by $\Ray\pH$\nomenclature{$\Ray\pH$}{Rays in a Hilbert space} the set of all rays in $\pH$.
       
\item\label{ax:vaps} The observables are represented by hermitian linear operators on $\pH$. A state $\rR$ has value $\alpha$ for the observable $A$ if $A\rR=\alpha\rR$, where the action of $A$ on the ray is obvious (and well defined because $A$ is linear).

\item If one has a system described by a state $\rR$, and if one want to measure if it is in one of the state $\rR_1,\ldots,\rR_n$ (orthogonal rays), the answer will be $\rR_i$ with probability 
        \[
	    P(\rR\to\rR_i)=|\scalh{\rR}{\rR_i}|^2.
        \]
If the $\rR_n$ form a complete system, one has a theorem which states that 
\[
   \sum_i P(\rR\to\rR_i)=1.
\]

\item\label{ax:reprez}  The rays of $\pH$ furnish a representation of the (identity component of) Poincaré group.
\end{enumerate}

This last point can look strange; we will see later (page \pageref{pg:poincare_act}) how it comes. It is the expression of a relativistic theory. That axiom is the reason why one make intensive use of representation theory in relativistic (quantum) field theory \ldots or maybe the intensive use of representation theory is the reason of that axiom. However, we will make an intensive use of representation theory developed in chapter \ref{ChapThoComsGroupes}.

\subsection{Symmetries and Wigner's theorem}
%-------------------------------------------

Consider the following situation: someone observes a system in a state $\rR$, and makes measures $P(\rR\to\rR_i)$. An other person observes the same system which is, for him, in a state $\rR'$ and observes $P(\rR'\to\rR'_i)$.

If two observers are related by a transformation of the Hilbert state which induces $\rR\to\rR'$ and $\rR_i\to\rR_i'$, there are said \defe{equivalent}{equivalence!of observer} if 
\begin{equation}\label{eq:sym_isom}
   P(\rR\to\rR_i)=P(\rR'\to\rR'_i).
\end{equation}

Let us say it more precisely from a mathematical point of view. A \defe{symmetry}{symmetry!of quantum system} is an invertible operator $\dpt{T}{\Ray\pH}{\Ray\pH}$ such that for any $\phi_i\in\rR_i$, $\phi_i'\in T\rR_i$ and $\phi''_i\in T^{-1}\rR_i$, 
\begin{equation}\label{eq:}
  |\scalh{\phi'_1}{\phi'_2}|^2=|\scalh{\phi_1}{\phi_2}|^2=|\scalh{\phi''_1}{\phi''_2}|^2
\end{equation}


\begin{remark}
Here, neither $\rR$ nor $\rR'$ are measurable: the $P$'s only are measurable.
\end{remark}

The following can be found in  \cite{Weinberg} p.91, \cite{Sternberg} p.354.
\begin{theorem}[Wigner]\label{tho:Wigner}
Any symmetry $T$ is induced by an operator $U$ on $\pH$ such that $\psi\in\rR$ implies $U\psi\in\rR'$. This operator is either unitary and linear, either anti-unitary and antilinear.
\end{theorem}

So, the symmetry operator must satisfy
\begin{subequations}
\begin{align}  
  \scalh{U\psi}{U\phi}&=\scalh{\psi}{\phi}\\
  U(\xi\psi+\eta\psi)&=\xi U\psi+\eta U\phi,
\end{align}  
\end{subequations}
or
\begin{subequations}
\begin{align}
  \scalh{U\psi}{U\phi}&=\scalh{\psi}{\phi}^*\\
  U(\xi\psi+\eta\psi)&=\xi^* U\psi+\eta^* U\phi.
\end{align}  
\end{subequations}
In the anti-linear case operator, we do not define $U^{\dag}$ by $\scalh{\phi}{U^{\dag}\psi}=\scalh{U\phi}{\psi}$ because the left-hand side should be anti-linear with respect to $\psi$ while the right-had should be linear. In place, for an antilinear operator $A$, we define $A^{\dag}$ by
\begin{equation} 
   \scalh{\phi}{A^{\dag}\psi}=\scalh{A\phi}{\psi}^*=\scalh{\psi}{A\phi}.
\end{equation}
In this way, the definitions of unitary and anti-unitary in term of dagger are the same: $U^{\dag}=U^{-1}$.

For any transformation $\dpt{T}{\Ray \pH}{\Ray \pH}$, the Wigner's theorem provides an operator $\dpt{U(T)}{\pH}{\pH}$ which induces $T$ on $\Ray$. If the operator $T$ depends on a parameter $\theta$, the operator $U(T(\theta))$ depends on $\theta$. If $T$ depends continuously on the parameter then the family $U(T(\theta))$ only contains unitary/linear operators or only antiunitary/antilinear operators.

In physical cases, $T(\theta)$ is mostly a Poincaré transformation: $\theta=(\Lambda,p)$. But $T(\mtu,0)$ is the identity which is represented by $U(\mtu,0)=\mtu$. Then all the (connected to identity) Poincaré transformations are represented by linear and unitary operators on $\pH$.

We will follow the proof given in \cite{Weinberg}. An other form of the proof can be found in \cite{Sternberg}. The latter use a slightly different formalism in the axioms of the quantum mechanics; this is explained in appendix \ref{app:Wigner}. It is now time to prove the theorem. 

\begin{proof}[Proof of Wigner's theorem]
We consider an orthonormal basis $\{\psi_k\}$ of $\pH$ with $\psi_k\in\rR_k$, and a choice of $\psi'_k\in T\rR_k$. From this and the assumptions, we have
\[
|\scalh{\psi'_k}{\psi'_l}|^2=|\scalh{\psi_k}{\psi_l}|^2=\delta_{kl}.
\]
Then $\scalh{\psi'_k}{\psi'_k}=0$ whenever $k\neq l$ and, since $\scalh{\psi'_k}{\psi'_k}$ is real and positive, $\scalh{\psi'_k}{\psi'_k}=1$. So $\scalh{\psi'_k}{\psi'_l}=\delta_{kl}$. 

The set $\psi_k'$ is also complete in $\pH$. Indeed suppose that we have a vector $\psi'\in\pH$ such that $\scalh{\psi'}{\psi'_k}=0$ for all $k$. If $\psi'\in\rR$, we consider a $\psi''\in T^{-1}\rR$ and we have
\[
   |\scalh{\psi''}{\psi_k}|^2=|\scalh{\psi'}{\psi'_k}|^2=0,
\]
which contradicts the fact that the $\psi_k$'s form a complete set. Now we have to fix a phase convention for the $\psi_k$. Since there are no canonical choice of phase, we fix with respect to an arbitrary one of the $\psi_k$, say $\psi_1$. We put
\begin{equation}
  \gamma_k=\us{\sqrt{2}}(\psi_1+\psi_k)\in\rC_k
\end{equation}
for $k\neq 1$. Any $\gamma'_k\in T\rC_k$ can be written in the basis $\{\psi'_k\}$:
\begin{equation}\label{eq:gamma_psi_k}
  \gamma'_k=\sum_l c_{kl}\psi'_l.
\end{equation}
From assumption \eqref{eq:sym_isom} and the fact that $|c_{kl}|^2=|\scalh{\gamma'_k}{\psi'_l}|^2$, we find, for $k,l\neq 1$
\[
  |c_{kl}|^2=\frac{1}{2}\delta_{kl}.
\]
We can choose the phase of $\gamma'_k$ and $\psi'_k$ in order to get $c_{kk}=c_{k1}=1/\sqrt{2}$. For this, we begin to fix $\gamma'_k$ in such a manner to get $c_{k1}=1/\sqrt{2}$ (from $|c_{k1}|=|\scalh{\gamma'_k}{\psi'_1}|$), and next we fix $\psi_k'$ for the $c_{kk}$. From now on, the so chosen $\gamma'_k$ and $\psi'_k$ are denoted by $U\gamma_k$ and $U\psi_k$. 

What we  did until now is to take a basis $\{\psi_k\}$ of $\pH$ and define $\gamma_k=1/\sqrt2(\psi_1+\psi_k)$. Next we had chosen the phases of $\psi'_k\in T\rR_k$ and $\gamma'_k\in T\rC_k$ in order to have
\begin{equation}\label{eq:c_kl}
\begin{aligned}
   c_{kk}=c_{k1}&=1/\sqrt 2&&\forall k,\\
   c_{kl}&=0 &&\text{if $l\neq k$ and $l\neq 1$}.
\end{aligned}   
\end{equation}
This allows us to check a certain linearity for the operator $U$:
\begin{equation}
\begin{aligned}
U\left( \us{\sqrt 2}(\psi_k+\psi1)  \right)
   &=U\gamma_k\\
   &=\gamma'_k\\
   &=\us{\sqrt 2}\psi'_1+\us{\sqrt 2}\psi'_k&&\text{from \eqref{eq:gamma_psi_k} and \eqref{eq:c_kl}}\\
   &=\us{\sqrt{2}}\big(  U\psi_1+U\psi_k   \big).
\end{aligned}
\end{equation}
Now we have to build $U$ on a general vector $\psi=\sum_k\psi_k\in\mR$. Any vector $\psi'\in T\mR$ can be decomposed with respect to the basis $\{\psi'_k=U\psi_k\}$: 
\begin{equation}\label{eq:dev_Upsi}
  \psi'=\sum_k C'_kU\psi_k.
\end{equation}
From the conservation of probability $|\scalh{\psi_k}{\psi}|^2=|\scalh{U\psi_k}{\psi'}|^2$ and $|\scalh{\gamma_k}{\psi}|^2=|\scalh{U\gamma_k}{\psi'}|^2$, we find
\begin{subequations}\label{eq:deux_C_k}
\begin{align}
          |C_k|^2&=|C'_k|^2, \label{eq:deux_C_k_a} \\
       |C_k+C_1|^2&=|C_k'+C'_1|^2.
\end{align}
\end{subequations}
If one writes $C_k=a_k+ib_k$, one finds $\Reel(C_k/C_1)=(a_ka_1+b_kb_1)/|C_1|^2$. By doing the same with $C'_k$ and using \eqref{eq:deux_C_k},
\begin{equation}\label{eq:C_k_C_1}
  \Reel(C_k/C_1)=\Reel(C'_k/C'_1).
\end{equation}
Equation \eqref{eq:deux_C_k_a} also imposes
\begin{equation}\label{eq:frac_C_k}
 |C_k/C_1|^2=|C'_k/C'_1|^2,
\end{equation}
while compatibility between \eqref{eq:frac_C_k} and \eqref{eq:C_k_C_1} requires
\begin{equation}\label{eq:C_k_C_1_im}
  \Imag(C_k/C_1)=\pm\Imag(C'_k/C'_1).
\end{equation}
Equations \eqref{eq:C_k_C_1} and \eqref{eq:C_k_C_1_im} show that $C_k$ and $C'_k$ must satisfy 
\begin{subequations}
\begin{align}
  C_k/C_1&=C'_k/C'_1\\
\intertext{xor}  
  C_k/C_1&=(C'_k/C'_1)^*.
\end{align}
\end{subequations}
For a given $\psi$ we have to show that the choice must be the same for all the $C_k$\footnote{We will show later that for a given $T$, the choice must be the same for all the $\psi$.}. Let $l\neq k$ and suppose that $C_k/C_1=C'_k/C'_1$ and $C_l/C_1=(C'_l/C'_1)^*$; we will show that in this case, one of the two ratios is real. So we can suppose $k\neq 1\neq l$. We consider the vector $\Phi=\us{\sqrt 3}(\psi_1+\psi_k+\psi_l)$, 
\[
  \Phi'=\frac{\alpha}{\sqrt 3}\big( U\psi_1+U\psi_k+U\psi_l  \big)
\]
where $\alpha\in\eC$ satisfies $|\alpha|=1$. The conservation of probability $|\scalh{\Phi}{\psi}|^2=|\scalh{\Phi'}{\psi'}|^2$ gives $|C_1+C_k+C_l|^2=|C'_1+C'_k+C'_l|^2$. Since $|C_1|^2=|C'_1|^2$, we can divide the left hand side by $|C_1|^2$ and the right one by $|C'_1|^2$. We find
\[
\left|1+\frac{C_k}{C_1}+\frac{C_l}{C_1}\right|^2=\left|1+\frac{C'_k}{C'_1}+\frac{C'_l}{C'_1}\right|^2.
\]
Using the assumption $C_k/C_1=C'_k/C'_1$ and $C_l/C_1=(C'_l/C'_1)^*$, we are in a case of an equation of the form $|u+v|^2=|u+v^*|^2$ with $u$, $v\in\eC$. If we write $u=a+bi$ and $v=x+iy$, we find $b+y=\pm(b-y)$, so that it leaves the choice $y=0$ or $b=0$ which corresponds to $(C_k/C_1)\in\eR$ or $(C_l/C_1)\in\eR$. So the coefficients $C'_k$ ($k\neq 1$) in the expansion \eqref{eq:dev_Upsi} must satisfy
\begin{subequations}
\begin{align}
  C_k/C_1&=C'_k/C'_1\quad\forall k \label{eq:rap_C_a} \\
\intertext{xor}  
  C_k/C_1&=(C'_k/C'_1)^*\quad\forall k  \label{eq:rap_C_b}.
\end{align}
\end{subequations}
Note that the phase of $C_1$ is not yet fixed. We naturally choice $C_1=C'_1$ or $C_1={C'_1}^*$ following the case. We define $\dpt{U}{\pH}{\pH}$ by
\begin{subequations}\label{eq:def_U}
\begin{align}
   U\left( \sum_k C_k\psi_k \right) &= \sum_kC_kU\psi_k&&\text{if \eqref{eq:rap_C_a}},
\label{eq:def_U_a}   
\intertext{xor}  
  U\left( \sum_k C_k\psi_k \right) &= \sum_kC_k^*U\psi_k&&\text{if \eqref{eq:rap_C_b}}.
\label{eq:def_U_b}  
\end{align}
\end{subequations}
One can explicitly check that it preserves the probability because $|\scalh{\psi}{\psi_k}|^2=|C_k|^2$ while $|\scalh{U\psi}{U\psi_k}|$ is equal to $|C_k|^2$ or $|C^*_k|^2$ (which are the same) following the case \eqref{eq:def_U_a} or \eqref{eq:def_U_b}.

Now we have to prove that the choice \eqref{eq:def_U_a} or \eqref{eq:def_U_b} is fixed by the data of $T$ and must be the same for all the $\psi\in\pH$. For, let us consider two vectors $\phi=\sum A_k\psi_k$, $\varphi=\sum B_k\psi_k$ and suppose that
\[
   U\phi=\sum_k A_k U\psi_k\text{ but } U\varphi=\sum_kB^*_k U\psi_k.
\]   
In order to see that it is impossible, looks at the conservation of probability $|\sum_k A_kB^*_k|^2=|\sum_k A_kB_k|^2$, then
\begin{equation}
\sum_{kl}\big(   B^*_lB_kA_lA^*_k-B^*_lB_kA^*_lA_k    \big)=\sum_{kl}B^*_lB_k\Imag(A_lA^*_k)=0.
\end{equation}
Since $A_lA^*l\in\eR$, we can regroup each term $(k,l)$ with the corresponding term $(l,k)$. We get
\begin{equation}\label{eq:imim_zero}
\begin{split}
  0=\sum_{kl}\Imag(A_lA^*_k)(B^*_lB_k-B^*_kB_l)
   =\sum_{kl}\Imag(A^*_kA_l)\Imag(B^*_kB_l).
\end{split}   
\end{equation}
We can find a vector $\sum_kC_k\psi_k$ such that 
\begin{subequations}\label{eq:choix_C}
\begin{align}
 \sum_{kl}\Imag(C^*_kC_l)\Imag(A^*_kA_l)&\neq 0  \label{eq:choix_C_a}\\
\intertext{and}
 \sum_{kl}\Imag(C^*_kC_l)\Imag(B^*_kB_l)&\neq 0.
\end{align}
\end{subequations}
In order to see how to find such a vector, let us show that there always exists a choice $(i,j)$ such that $B^*_iB_j$ is not real. Let us say $B_1=x+iy$ and $B_k=a_k+bi$. If $y\neq 0$, the condition $\Imag(B_1^*B_k)=0$ gives $B_k=\frac{b_k}{y}B_1$. It is always possible to find a sequence $(b_k)$ which gives $1$ as norm for $\sum B_k\psi_k$; the problem is not there. The problem is that $B_k/B_1\in\eR$, so that the choice \eqref{eq:def_U}  is not a true choice. For the same reason, all the $B^*_iB_k$ can't be pure imaginary.

Now we can find the vector which satisfy \eqref{eq:choix_C}. There are several cases. If there is a pair $(k,l)$ such that $A^*_kA_l$ and $B^*_kB_l$ are both complex, we can take all $C_i$'s zero for $k\neq i\neq l$ and choose $C_k$ and $C_l$ in such a way that $C^*_kC_l$ is not real. If there is a pair $(k,l)$ with $A^*_kA_l$ complex and $B^*_kB_l$ real, we consider a pair $(m,n)$ such that $B^*_mB_n$ is complex. If $A^*_mA_n$ is complex, we take all the $C_i$'s zero except $C_m$ and $C_n$ such that $\Imag(C^*_mC_n)\neq 0$. If $A^*_mA_n$ is real, we take all the $C_i$'s zero except $C_k,C_l,C_m,C_n$ which we choose in such a way that $\Imag(C^*_mC_n)\neq 0$ and $\Imag(C^*_kC_k)\neq 0$.

Equation \eqref{eq:choix_C_a} makes that the same choice must be made for $\sum A_k\psi_k$ and $\sum C_k\psi_k$ (if it was not the case, we would have an equation of the form of \eqref{eq:imim_zero}). For the same reason, the same choice must be made for $\sum B_k\psi_k$ and $\sum C_k\psi_k$. So we conclude that the data of $T$ fixes the choice between \eqref{eq:def_U_a} and \eqref{eq:def_U_b} and that this choice must be the same for all the vectors of $\pH$.

We have to show that the possibility \eqref{eq:def_U_a} makes $U$ linear and unitary while the possibility \eqref{eq:def_U_b} makes $U$ antilinear and antiunitary. For we consider $\psi=\sum_k A_k\psi_k$ and $\phi=\sum_k B_k\psi_k$. If \eqref{eq:def_U_a} works,
\begin{equation}
\begin{split}
  U(\alpha\psi+\beta\phi)&=U\Big(  \sum_k(\alpha A_k+\beta B_k)\psi_k     \Big)\\
                         &=\sum_k(\alpha A_k+\beta B_k)U\psi_k\\
			 &=\alpha U\psi+\beta U\phi,
\end{split}
\end{equation}
and
\begin{equation}
\scalh{U\psi}{U\phi}=\sum_{kl}A^*_kB_l\scalh{U\psi_k}{U\psi_l}=\sum_kA^*_kB_k,
\end{equation}
so that $\scalh{U\psi}{U\phi}=\scalh{\psi}{\phi}$. Thus in this case $U$ is linear and unitary. In the case where \eqref{eq:def_U_a} works, the computations are almost the same:
\begin{equation}
\begin{split}
  U(\alpha\psi+\beta\phi)&=U\Big(  \sum_k(\alpha A_k+\beta B_k)\psi_k     \Big)\\
                         &=\sum_k(\alpha^* A_k^*+\beta^* B^*_k)U\psi_k\\
			 &=\alpha^* U\psi+\beta^* U\phi,
\end{split}
\end{equation}
and
\begin{equation}
\scalh{U\psi}{U\phi}=\sum_{kl}A_kB^*_l\scalh{U\psi_k}{U\psi_l}=\sum_kA_kB^*_k,
\end{equation}
so that $\scalh{U\psi}{U\phi}=\scalh{\psi}{\phi}^*$. In this case, $U$ is antilinear and antiunitary.

\end{proof}

\subsection{Projective representations}
%-------------------------------------------------

If $T_1(\rR_n)=\rR_n'$ and $\psi_n\in\rR_n$, then $U(T_1)\psi_n\in\rR_n'$. If $T_2(\rR')=\rR''$, then $U(T_2)U(T_1)\psi_n\in\rR_n''$. But $U(T_2T_1)\psi_n$ also belongs to $\rR_n''$. Then there exists a $\phi_n(T_2,T_1)\in\eR$ such that 
\[
   U(T_2)U(T_1)\psi_n=e^{i\phi_n(T_2,T_1)}U(T_2T_1)\psi_n.
\]
Note that for any $\psi\in \pH$, there exists a $\lambda\in\eR$ such that $\|\lambda\psi\|=1$. Since a real can be sent out the $U(T)$'s, for \emph{any} $\psi\in \pH$, there exists a $\phi$ which only depends on $\psi/\|\psi\|$ such that
\begin{equation}
    U(T_2)U(T_1)\psi=e^{i\phi(T_2,T_1)}U(T_2T_1)\psi
\end{equation}

\begin{proposition}
The $\phi$ doesn't depend at all on the $\psi$:
\begin{equation}
  U(T_2)U(T_1)=e^{i\phi(T_2,T_1)}U(T_2T_1).
\end{equation}
\end{proposition}

\begin{proof}
Let us consider a $\psi_A$ and a $\psi_B$ which are not proportional each other. One has a $\phi_{AB}(T_2,T_1)$ such that
\begin{equation}
\begin{split}
e^{i\phi_{AB}(T_2,T_1)}U(T_2T_1)(\psi_A+\psi_B)&=U(T_2)U(T_1)(\psi_A+\psi_B)\\
                                               &=e^{i\phi_A(T_2,T_1)}U(T_2T_1)\psi_A\\
                                               &\quad +e^{i\phi_B(T_2,T_1)}U(T_2T_1)\psi_B.
\end{split}
\end{equation}
Now, we apply $U(T_2T_1)^{-1}$ to both sides. If it is unitary, the $e^{i\phi}$ get out without problems; else is get out as $e^{-i\phi}$:
\begin{equation}
  e^{\pm i\phi_{AB}}(\psi_A+\psi_B)=e^{\pm i\phi_A}\psi_A+e^{\pm i\phi_B}\psi_B.
\end{equation}
Since $\psi_A$ and $\psi_B$ are linearly independent, the only solution is $e^{i\phi_{AB}}=e^{i\phi_A}=e^{i\phi_B}$.

\end{proof}

Since the operators $U(T)$ must only fulfil
\begin{equation}\label{eq:projectif}
   U(T_2)U(T_1)=e^{i\phi(T_2,T_1)}U(T_2,T_1),
\end{equation}
these form a \defe{projective representation}{projective!representation} of the symmetry group on the physical Hilbert space $\pH$.

\begin{remark}
In order to have some physical relevance, this demonstration supposes that a state $\psi_A+\psi_B$ exists in nature. If one can divide the particles in several ``incompatibles''\ classes labeled by $a,b$ such that $\psi_a+\psi_b$ doesn't exist, then equation \eqref{eq:projectif} is false and one has to write 
\[
   U(T_2)U(T_1)\psi_a=e^{i\phi_a(T_2,T_1)}U(T_2T_1)\psi_a
\]
because we can't show that $\phi_a=\phi_b$ from the simple fact that $\psi_a+\psi_b$ doesn't exist !

For example, physicists think that there are no superposition of state of integer and semi-integer spin.
\end{remark}

\begin{remark}
If the group satisfies some requirements, one can choose $\phi=0$. From now we suppose that we are in this case: we work with ``true''\ representations.
\end{remark}


\subsection{Representations and power expansions}
%------------------------------------------------

Let $G$ be an arc connected Lie group whose elements are denoted by $T(\theta)$ with $\theta$, a continuous family of parameters (from a local chart). The multiplication law is given by a function $\dpt{f}{\eR^n\times\eR^n}{\eR^n}$:
\begin{equation}\label{eq:T_groupe}
T(\theta')T(\theta)=T\big(f(\theta',\theta)\big).
\end{equation}
If $ \theta=0$ is the coordinate of the identity, 
\begin{equation}\label{eq:f_0}
f(0,\theta)=f(\theta,0)=\theta.
\end{equation}

We suppose that $G$ acts on the rays of a Hilbert space $\pH$, so that there are represented on $\pH$ by unitary operators $U\big(T(\theta)\big)$. We denote by $W$ the group of transformations of $\pH$; roughly speaking, 
\[
     W=U(G).
\]
Now, we are going to cheat a little. We know that there exists a normal neighbourhood\index{normal!neighbourhood} of $e$ in $W$. In simple words, the map $\dpt{\exp}{\lW}{W}$ is a diffeomorphism between the elements of $\lW$ ``close''\ to $0$ and the ones of $W$ close to $e$. By \textit{close to}, we mean that the components of $\theta$ are small enough. If $\{it_a\}$ is a basis of $\lW$, we define
\begin{equation}\label{eq:U_expo}
   U(T(\theta))=e^{i\theta^at_a}.
\end{equation}
In other words, one considers the exponential map for a neighbourhood of identity.

The cheat is the fact that $U(T(\theta))$ is actually defined by Wigner's theorem from the data of the group $G$. So equation \eqref{eq:U_expo} should be seen as a requirement in the choice of the basis $\{t_a\}$.

\begin{remark}
The $i$ in the exponential in \eqref{eq:U_expo} and in the definition of the basis $\{it_a\}$ is a convention in order the $t_a$'s to be hermitian. Indeed, the Lie algebra of a group of unitary matrices is made of \emph{anti}hermitian matrices.
\end{remark}

With all that,
\begin{equation}\label{eq:dev_U}
   U(T(\theta))=\mtu+i\theta^at_a+\frac{1}{2}\theta^b\theta^ct_{bc}+\ldots
\end{equation}
where $t_{bc}$ is defined (among other requirements) to absorb the ``intuitive''\ minus sign in the third term.

Now we are going to explore some consequences of equation \eqref{eq:T_groupe}. Equation \eqref{eq:f_0} makes the expansion of $f$ as
\begin{equation}\label{eq:dev_f}
f^a(\theta',\theta)=\theta^a+{\theta'}^a+f^a_{bc}{\theta'}^b\theta^c+\ldots
\end{equation}
From expansions \eqref{eq:dev_f} and \eqref{eq:dev_U} of $f$ and $U(T(\theta))$, ``group structure''\ equation \eqref{eq:T_groupe} gives (at order two):
\begin{equation}\label{eq:t_ab}
   t_{bc}=-t_bt_c-if^a_{bc}t_a
\end{equation}
and nothing for the first order. Then, providing that one knows the group structure (the $f$), one knows the second order of the representation from the first one.
From equation \eqref{eq:U_expo}, one finds the value of $t_{ab}$:
\[
   e^{i\theta^at_a}=1+i\theta^at_a+\frac{1}{2}(i)^2(\theta^at_a)(\theta^bt_b),
\]
up to constant coefficients, one can choose $t_{ab}$ to be symmetric with respect to $a$ and $b$:
\[
   t_{ab}=\frac{1}{2}(t_at_b+t_bt_a).
\]
Taking this convention and computing $t_{bc}-t_{cb}$ from \eqref{eq:t_ab}, we find
\begin{equation}
  [t_a,t_b]=iC_{ab}^ct_c
\end{equation}
with $C_{ab}^c=f_{ab}^c-f_{ba}^c$.

On the other hand, one knows that if a group is abelian, its algebra is also abelian; we can see it here by considering that if $G$ is abelian, $f(\theta,\theta')=f(\theta',\theta)$, then $f_{ab}^c$ is symmetric and $[t_a,t_b]=0$. We can say more about $f$ Since the $t_a$ commute, equations \eqref{eq:T_groupe} and \eqref{eq:U_expo} make that
\begin{equation}
  e^{if(\theta,\theta')^at_a}=e^{i\theta^at_a}e^{i{\theta'}^bt_b}\\
                             =e^{i(\theta^a+{\theta'}^a)t_a},
\end{equation}
so that 
\[
   f(\theta,\theta')=\theta+\theta'.
\]

\section{The symmetry group of nature}
%-------------------------------------

\subsection{Spin and double covering}\label{subsec:sym_nature}
%-----------------------------------

Some of literature carry an ambiguity in the choice of the right space-time symmetry group in the quantum field theory. A very good and deep discussion about the choice of the space-time symmetry group of nature is given in the book \cite{Naber} which will be used here. An other enlightening review can be found in \cite{ModavePoincarre}.

From a relativistic point of view, the group is the Poincaré group of all the maps $\eR^4\to\eR^4$ which leaves invariant the quantity $s^2=-t^2+x^2+y^2+z^2$.  At this point we can already make an important remark: the so defined quantity $s$ is in fact \emph{not} a relativistic invariant. Indeed if I follow a (spatially) closed path, I will measure $\Delta t\neq 0$ and $\Delta x=\Delta y=\Delta z=0$ because in \emph{my} frame, my displacement is zero. A guy who keeps at my starting point will measure (between the beginning and the end of my travel) $\Delta 't\neq 0$ and also $\Delta x=\Delta y=\Delta z=0$. If $s=s'$, then $\Delta t=\Delta t'$.

So the relativistic invariance is only local: $ds^2=ds'^2$, and as far as relativity is concerned, one can work with infinitesimal transformations only. In this case, the distinction between the \emph{groups} $L_+^{\uparrow}$ and $\SLdc$ is no relevant. Intuitively, we choose $L_+^{\uparrow}$ to be the space-time symmetry group. As we will see the difference will reveal to be crucial in relativistic field theory because $L_+^{\uparrow}$ has no half-integer spin representations.

This group naturally splits into two parts: the translations and the rotations (and boost). As far as I know, the translation part makes no difficulties. For the other one, there are some difficulties to find the \emph{minimal} group of symmetry. First, one often want to separate the space-time inversions $P$ and $T$ from the remaining: the group then becomes the homogeneous orthochrone Lorentz group $L_+^{\uparrow}$. An other often presented group is \ldots $\SLdc$. This is our choice here. The physical reason of this choice is all but immediate. As we will see during the following pages, an elementary particle is an irreducible representation of the symmetry group.

For massive particles, the relevant subgroup of $\SLdc$ reveals to be $SU(2)$. If we had chosen the most intuitive $L_+^{\uparrow}$, we would have found $\SO(3)$. There is an important difference between $SU(2)$ and $\SO(3)=SU(2)/\eZ_2$: the first one admits representations of any integer and half-integer spin while the second only posses the integer spin representations (cf. page \pageref{pg:reprez_SO3}).

Let us now be more precise about the relation between $L_+^{\uparrow}$ and $\SLdc$. A know result is
\[
   L_+^{\uparrow}=\frac{\SLdc}{\eZ_2}.
\]
Let $\dpt{\mSpin}{\SLdc}{L_+^{\uparrow}}$ be the surjective homomorphism with kernel $\pm\mtu_{2\times 2}$ giving this relation.
We will not give a complete proof, but we will explain how $\SLdc$ acts by isometries on $\eR^4$. First, we remark that there exists a bijection between $\eR^4$ and the $2\times 2$ complex hermitian matrices:
\begin{equation}
  v=\begin{pmatrix} t+z & x-iy\\x+iy&t-z\end{pmatrix}=\begin{pmatrix}t\\x\\y\\z\end{pmatrix}.
\end{equation}
If $\lambda\in\SLdc$, the matrix $\lambda v\lambda^{\dag}$ is also hermitian and $\|v\|^2=\det v$. Thus 
\begin{equation}
\begin{aligned}
 \Lambda(\lambda)  \colon \eR^4&\to \eR^4 \\ 
   v&\mapsto \lambda v\lambda^{\dag} 
\end{aligned}
\end{equation}
is a Lorentz transformation if and only if $|\det\lambda|=1$. Moreover,
\[
     \Lambda(\lambda\lambda')=\Lambda(\lambda)\Lambda(\lambda').
\]
If $\lambda'=e^{\phi}\lambda$, then $\Lambda(\lambda')=\Lambda(\lambda)$, thus it is natural to impose $\det v=1$ and to consider $\SLdc$ instead of $L(2,\eC)$ to fit $L_+^{\uparrow}$. Now, $\Lambda(\lambda)=\Lambda(-\lambda)$, and we wish to consider $\SLdc/\eZ_2$.

I think the problem is the following: as far as the action of the ``nature group''{} on the space-time is concerned, it is sufficient to consider $L_+^{\uparrow}$. But the group which acts on the state space is wider: it must be $SL(2,\eC)$.

From now, when we say ``Poincaré group'', we mean $\SLdc\times\eR^4$ while  ``Lorentz''\ means $\SLdc$ acting on $\eR^4$ by $\Lambda(\lambda)v=\lambda v\lambda^{\dag}$.

Let us continue the discussion of page \pageref{pg:reprez_SOt}. A know result is the fact that the map $\mSpin$ restricts to a surjective homomorphism $\dpt{\mSpin}{SU(2)}{\SO(3)}$ with kernel $\pm\mtu$ giving the relation $\SO(3)=SU(2)/\eZ_2$. If one considers a representation $\dpt{\rho}{\SO(3)}{GL(V)}$, then $\tilde\rho=\rho\circ\mSpin$ is a representation of $SU(2)$ on $V$. So every representation of $\SO(3)$ comes from a representation of $SU(2)$.

As far as the transformation rule of a (quantum mechanical) wave function under a rotation $R\in \SO(3)$ is concerned, one can see (it is done in \cite{Naber}) that the try
\[
   \begin{pmatrix}
     \psi_1\\
     \psi_2
   \end{pmatrix}\to
T(R)
\begin{pmatrix}
     \psi_1\\
     \psi_2
   \end{pmatrix}
\]
doesn't works if $T(R)$ is a representation of $\SO(3)$ on $\eC^2$. If one allows $T$ to be a representation of $SU(2)$, then our choice ---for an electron--- should naturally be the spin one half representation $T=D^{(1/2)}$. Let us do it. The remaining problem is the following. Let's consider that in a certain frame, an electron is described by the wave function $\begin{pmatrix} \psi_1&\psi_2\end{pmatrix}$, the question is to know the wave function observed by a guy which use another frame linked to the first frame by $R\in \SO(3)$. We always have exactly two elements in $\SU(2)$ projected to $R$ by $\mSpin$; namely $\mSpin(\pm g)=R$; so how to choose between
\[
   D^{(1/2)}(g)
\begin{pmatrix}
     \psi_1\\
     \psi_2
   \end{pmatrix}
\quad\text{and}\quad
   D^{(1/2)}(-g)
\begin{pmatrix}
     \psi_1\\
     \psi_2
   \end{pmatrix}\; ?
\]
The trick is to remark that a change of frame is not the mathematical process described by a single element $R$ of $\SO(3)$, but a physical \emph{continuous} process which begins at the identity and stops at $R$. In other word, we have to ask ourself \emph{how to go from a frame to another} ? Taking as example the rotations around the $x$ axis, we can look at two different path in $\SO(3)$ from $\mtu$ to $\mtu$ given by the same expression
\[
  R_1(t)=R_2(t)=
\begin{pmatrix}
1&0&0\\
0&\cos t&\sin t\\
0&-\sin t&\cos t
\end{pmatrix},
\]
but considering $t\colon 0\to 2\pi$ for $R_1$ and $t\colon 0\to 4\pi$ for $R_2$. The covering map $\dpt{\mSpin}{SU(2)}{\SO(3)}$ allows us to lift any path in $\SO(3)$ to a path in $SU(2)$ in an unique way providing a starting point. In other words, if $\mSpin(g)=R$,
\[
\begin{split}
  \exists !\, \tilde{R}(t)\in SU(2)\text{ such that }\mSpin\circ\tilde{R}=R \text{ and } \tilde{R}(0)=\mtu,\\
  \exists !\, \tilde{R}(t)\in SU(2)\text{ such that }\mSpin\circ\tilde{R}=R \text{ and } \tilde{R}(0)=-\mtu.
\end{split}
\]
The question is now: how to choose the right path among these two ? The answer comes from the homotopy of $\SO(3)$: the path $R_1$ and $R_2$ belongs to two different classes.

Considering the ``change of frame''{} as a continuous process, the initial point is naturally chosen to be $\mtu$. With this choice, the lift of $R_1$ and $R_2$ are given by
\[
  g_1(t)=g_2(t)=
\begin{pmatrix}
\cos\frac{t}{2}&-i\sin\frac{t}{2}\\
-i\sin\frac{t}{2}&\cos \frac{t}{2}
\end{pmatrix}
\]
with $\dpt{t}{0}{2\pi}$ for $g_1$ and $\dpt{t}{0}{4\pi}$ for $g_2$. In $SU(2)$, the ending point of $g_1$ is $-\mtu$ while the one of $g_2$ is $\mtu$.

It is still possible to say a lot of interesting thinks about the space-time symmetry group of nature; let's just conclude saying that $SU(2)$ is more adapted to the rotations of non zero spin than $\SO(3)$. (it is  not intuitive !)
\subsection{How to implement the Poincaré group}
%----------------------------------------------
\label{pg:poincare_act}

We are not making physics here, but differential geometry and group theory; so we will not discuss the physical relevance of the Poincaré group from a ``speed of light''\ point of view. We consider the \defe{Poincaré group}{Poincaré!group} as the group of all the affine isometries of metric $\eta=diag(-1,1,1,1)$ and the \defe{Lorentz group}{Lorentz!group}\index{group!Lorentz} as the subgroup of rotations and boost.

A Poincaré transformation of $\eR^4$ is given by $(\Lambda,a)$ with $\Lambda$ a $4\times 4$ matrix and $a\in\eR^4$, a translation vector. The composition of $(\Lambda,a)$ with $(\Lambda',a')$ is given by $(\Lambda'\Lambda,\Lambda'a+a')$, the inverse is $(\Lambda^{-1},-\Lambda^{-1} a)$, the neutral is $(\mtu,0)$, and $(\det\Lambda)^2=1$.

The axiom \ref{ax:reprez} at page \pageref{pg:axiomes} gives us a group of transformation of the rays in $\pH$ parametrised by $(\Lambda,a)$ such that
\begin{equation}
  T(\Lambda',a')T(\Lambda,a)=T(\Lambda'\Lambda,\Lambda'a+a'),
\end{equation}
$\dpt{T(\Lambda,a)}{\Ray \pH}{\Ray \pH}$. Then Wigner's theorem defines a representation of the Poincaré group on $\pH$ by unitary matrices :
\[
\psi\to U(\Lambda,a)\psi.
\]

\begin{remark}
Wigner only ensure existence of \emph{projective} representations. Here we suppose that our symmetry group (maybe slightly different that Poincaré) is such that any projective representations can be turn into a classical representation. We will therefore use the composition law
\begin{equation}\label{eq:composition_U}
  U(\Lambda',a')U(\Lambda)=U(\Lambda'\Lambda,\Lambda'a+a')
\end{equation}
instead of $U(\Lambda',a')U(\Lambda,a')=e^{i\phi(\Lambda,a,\Lambda',a')}U(\Lambda'\Lambda,\Lambda'a+a')$.
\end{remark}

By axiom, the (connected) Poincaré group acts on rays of $\pH$, and we have the representation $U$ which form a group acting on $\pH$. The Lie algebra acts also :
\begin{equation}
u\psi=\Dsdd{U(t)}{t}{0}\psi:=\Dsdd{U(t)\psi}{t}{0}.
\end{equation}
This definition is natural because $\pH$ is a vector space: it can be identified with its tangent space: $U(t)\psi$ is a path in $\pH$ and its derivative at $t=0$ is still a well defined element in $\pH$. Now recall that the operators $U$ are unitary, so that the corresponding operators $u$ are hermitian (therefore diagonalisable).

Let us consider an abelian subgroup $A$ of Poincaré with Lie algebra $\lA$. One can find a basis of $\pH$ made of common eigenvectors of a basis of $\lA$. In other words, one can find a basis of $\pH$ which simultaneously diagonalises all $\lA$. If $\{a_i\}$ is a basis of $\lA$, one can find a basis $\{\ket{\psi_{\lambda}}\}$ (here $\lambda$ labels a basis of $\pH$ : it might take continuous values) such that
\begin{equation}
   a_i\ket{\psi_{\lambda}}=\lambda_i\ket{\psi_{\lambda}}.
\end{equation}

\subsection{Momentum operator}
%-----------------------------

Of course, there exists an abelian subgroup of Poincaré: the pure translations, $A=\{U(\mtu,a)\}$. A basis of the Lie algebra is given by four vectors labeled as $P^{\mu}$ and defined by 
\[
P^{\mu}=\Dsdd{U(\mtu,te^{\mu})}{t}{0}
\]
 where $e^{\mu}$ is the unit vector following the direction $\mu$ (for $\mu=0$, $e^0=(1,0,0,0)$). One can consider a basis which diagonalises the $P^{\mu}$'s: 
\begin{equation}
  P^{\mu}\ket{p,\sigma}=p\hmu\ket{p,\sigma}
\end{equation}
where by definition,
\begin{equation}\label{eq:def_P}
   P\hmu\ket{p,\sigma}=\Dsdd{U(\mtu,te\hmu)\ket{p,\sigma}}{t}{0}.
\end{equation}

\begin{remark}
Be careful on a point: we don't say anything about the symbol ``$p$''\ in the ket. The only property is that it labels a Hilbert space $\pH$. But nothing is already imposed to $\pH$ : it must just carry a representation of the Poincaré group on its rays. In particular, it is \emph{a priori} false to say that $p$ is a ``momentum $4$-vector'' and that $p\hmu$ is a component of $p$. Naturally, our notations are adapted to think that ! Maybe it is a pedagogical mistake; I don't know. 
\end{remark}

This remark can be disturbing: why is generally $\ket{p,\sigma}$ called ``a state of momentum $p$''\ ? Since $U(\mtu,a)$ is unitary, $P\hmu$ is hermitian; the $p\hmu$ are eigenvalues for an hermitian operator, so by axiom \ref{ax:vaps} (page \pageref{pg:axiomes}) they are candidate to be physical values. But equation \eqref{eq:def_P} shows that $P\hmu$ is what a physicist should call an ``infinitesimal translation'', so that Noether suggests us to interpret the eigenvalue as momentum. We are safe !

The parameters $\sigma$ are not yet defined neither. It will come later. For the moment, we include into the definition of a \defe{one particle state}{state!one particle} that $\sigma$ takes discrete values.
 
 Since $U(\mtu,a)=e^{a_{\mu} P\hmu}$,
\[
   U(\mtu,a)\ket{p,\sigma}=e^{ia_{\mu} p^{\mu}}\ket{p,\sigma}.
\]
Now we are interested in the determination of $U(\Lambda,a)\ket{p,\sigma}$.

\begin{proposition}
The operators $P^{\mu}$ are subject to the ``transformation law''
\begin{equation}
   U(\Lambda,a)P^{\mu} U(\Lambda,a)^{-1}=\Lambda^{\mu}_{\nu} P^{\nu}.
\end{equation}
\end{proposition}

\begin{proof}
Since operators $U(\Lambda,a)$ are linear, they can be putted in the derivative which defines $P\hmu$. Using the composition law  \eqref{eq:composition_U} we find :
\begin{equation}
\begin{split}
   U(\Lambda,a)P\hmu U(\Lambda,a)^{-1}&=\Dsdd{ U(\Lambda,a)U(\mtu,te\hmu)U(\Lambda,a)^{-1} }{t}{0}\\
                                     &=\Dsdd{U(\mtu,t\Lambda e\hmu)}{t}{0}.
\end{split}				     
\end{equation}
The $\Lambda$ can be putted out of derivative; let us see it for a sum of two terms (here it is four) :
\begin{equation}
\begin{split}
  \Dsdd{ U(\mtu,t(e\hmu+e\hnu)) }{t}{0}&=\Dsdd{U(\mtu,te\hmu)U(\mtu,te\hnu)}{t}{0}\\
                                      &=\Dsdd{U(\mtu,te\hmu)U(\mtu,0)}{t}{0}
				        +\Dsdd{U(\mtu,0)U(\mtu,te\hnu)}{t}{0}\\
				      &=P\hmu+P\hnu.
\end{split}
\end{equation}
Thus 
\begin{equation}
    \Dsdd{U(\mtu,\Lambda\hmu_{\nu} e\hnu)}{t}{0}=\Lambda\hmu_{\nu}\Dsdd{U(\mtu,te\hnu)}{t}{0}
                                              =\Lambda\hmu_{\nu} P\hnu.
\end{equation}
\end{proof}

\subsection{Pure Lorentz transformation}
%--------------------------------------

Now we consider a pure Lorentz transformation $U(\Lambda)\equiv U(\Lambda,0)$, and we want to look at $U(\Lambda)\ket{p,\sigma}$. In order to see its decomposition into others $\ket{k,\sigma'}$, we apply a $P\hmu$ :
\begin{equation}
\begin{split}  
  P\hmu U(\Lambda)\ket{p,\sigma}&=U(\Lambda)\Big( U(\Lambda)^{-1} P\hmu U(\Lambda)  \Big)\ket{p,\sigma}\\
                                &=U(\Lambda) {(\Lambda^{-1})}\hmu_{\nu} P\hnu \ket{p,\sigma}\\
				&=(\Lambda^{-1})\hmu_{\nu} p\hnu U(\Lambda)\ket{p,\sigma}.
\end{split}				
\end{equation}
Thus the vector $U(\Lambda)\ket{p,\sigma}\in\pH$ has $(\Lambda^{-1})\hmu_{\nu} p\hnu$ as eigenvalue for $P\hmu$. If the $p\hmu$'s are seen as components of a $4$-vector $p$, one can write
\[
    P\hmu U(\Lambda)\ket{p,\sigma}=(\Lambda p)\hmu U(\Lambda)\ket{p,\sigma};
\]    
     thus we naturally write  
\begin{equation}     
  U(\Lambda)\ket{p,\sigma}=\sum_{\sigma'}C_{\sigma'\sigma}(\Lambda,p)\ket{\Lambda p,\sigma'}.
\end{equation}
       

Note that we had not yet given anything about the nature of the $p$ in the ket $\ket{p,\sigma}$ so we can \emph{define} the product $\Lambda p$ by the fact that the ket $\ket{\Lambda p,\sigma}$ has eigenvalue $(\Lambda^{-1})\hmu_{\nu} p\hnu$ for the operator $P\hmu$. So it is one of the $\ket{p',\sigma'}$.

\subsection{Rebuilding of a basis for \texorpdfstring{$\pH$}{H}}
%----------------------------------------

From general considerations about the Lorentz group (many physicists had written very better books than me about) anyone knows that the only functions of the $p\hmu$'s  which are invariant under all the Lorentz transformations are $p^2=\eta_{\mu}nu p\hmu p\hnu$ and the sign of $p^0$ when $p^2< 0$.

For any value of $p^2$ and sign of $p^0$, one consider a ``standard vector''\ $k$. For example :
\begin{subequations}
\begin{align}
   k&=(1,0,0,1)&&\text{for $p^2=0$},\\
   k&=(1,0,0,0)&&\text{for $p^2<0$, $p^0>0$},\\
   k&=(-1,0,0,0)&&\text{for $p^2<0$, $p^0<0$}.
\end{align}   
\end{subequations}
With this convention, $p$ can be written as $p=L(p)k$ for a suitable Lorentz transformation $L(p)$. The vector $U(L(p))\ket{k,\sigma}$ has eigenvalue $L(p)k$ for the operator $P$, thus it is a linear combination of some $\ket{p,\sigma'}$.

Now we will cheat and redefine our basis of the Hilbert space $\pH$. First, we consider a fixed $k$; in other words, we build the state space for a given particle which has given momentum $p$.  The basis vectors must be eigenvectors for the fours operators $P\hmu$. As far as we say no more, any eigenvalue is possible. Thus our basis must be labelled by at least an element $p$ of $\eR^4$ with only one constraint: the value of $p^2$ (plus eventually the sign of $p^0$). So we define the $\ket{k,\sigma}$ to be such that
\[
   P\hmu\ket{k,\sigma}=k\hmu\ket{k,\sigma}.
\]
Since we know that with this definition of $\ket{k,\sigma}$, the eigenvalue of $U(L(p))\ket{k,\sigma}$ for $P\hmu$ is $p\hmu$, we \emph{define} $\ket{p,\sigma}$ as
\begin{equation}
  \ket{p,\sigma}=N(p)U(L(p))\ket{k,\sigma}.
\end{equation}
where $N(p)$ is a normalization to be discussed later. With this construction, we have an eigenvector for any possible eigenvalue for $P\hmu$. We have to show that these vectors are linearly independent.

The set of the $\ket{p,\sigma}$ with different $p$ is free in $\pH$ because they are eigenvectors for different eigenvalue of an hermitian operator\quexto{I did not checked that it is sufficient}. There are no reason to think that the set of operators $P\hmu$ is complete; in other words, it remains not clear that there exist only one way to diagonalise the all the $P\hmu$. The function of the extra label $\sigma$ is to label different linearly independent vectors with same eigenvalue for $P$.

From now, we are interested in $\ket{k,\sigma}$ and $N(p)$.

\subsection{Little group}
%-----------------------

We have :
\begin{equation}
\begin{split}
  U(\Lambda)\ket{p,\sigma}&=N(p)U(\Lambda L(p))\ket{k,\sigma}\\
                          &=N(p)U(L(\Lambda p))U\big(  L(\Lambda p)^{-1}\Lambda L(p) \big)\ket{k,\sigma},
\end{split}
\end{equation}
So we will try to understand the operation $L(\Lambda p)^{-1}\Lambda L(p)$. First remark that
\[
   U(L(\Lambda p)^{-1})\ket{\Lambda p,\sigma}=N(\Lambda p)\ket{k,\sigma},
\]
and then compute :
\begin{equation}
\begin{split}
  U(L(\Lambda p)^{-1}\Lambda L(p) )N(p)\ket{k,0}&=U(L(\Lambda p)^{-1}\Lambda)\ket{p,\sigma}\\
                                               &=U(L(\Lambda p)^{-1})\sum_{\sigma'}
					         C_{\sigma'\sigma}(\Lambda,p)\ket{\Lambda p,\sigma'}\\
					       &=\sum_{\sigma'}C_{\sigma'\sigma}(\Lambda,p)
					         N(\Lambda p)\ket{k,\sigma'}.
\end{split}
\end{equation}

The \defe{little group}{little group} is the subgroup of the Lorentz transformations which leaves the chosen standard vector $k$ invariant: $Wk=k$. For any $W$ in the little group,
\[
   U(W)\ket{k,\sigma}=\sum_{\sigma'}D_{\sigma'\sigma}(W)\ket{k,\sigma'}
\]
With this definition, the $D$'s form a representation of the little group. Indeed for any $V,W$ in the little group,
\begin{equation}
\begin{split}
  \sum_{\sigma'}D_{\sigma'\sigma}(VW)\ket{k,\sigma'}&=U(VW)\ket{k,\sigma}\\
                                    &=U(V)\sum_{\sigma''}D_{\sigma''\sigma}(W)\ket{k,\sigma''}\\
				    &=\sum_{\sigma'\sigma''}D_{\sigma'\sigma''}(V)D_{\sigma''\sigma}(W)
				      \ket{k,\sigma'}.
\end{split}
\end{equation}
Since we want the $\ket{p,\sigma}$ with different $p$ and $\sigma$ to form a basis of $\pH$, they are linearly independent, then we can get rid of the sum over the $\sigma'$ and keep the equation
\[
D_{\sigma'\sigma}(VW)=\sum_{\sigma''}D_{\sigma'\sigma''}(VW)D_{\sigma''\sigma}(VW);
\]
if we adopt a more ``matricial''\ notation,
\begin{equation}
D(VW)=D(V)D(W).
\end{equation}

We are now able to perform a step in the study of the vector $U(\Lambda)\ket{p,\sigma}$. We naturally define $W(\Lambda,p)=L(\Lambda p)^{-1}\Lambda L(p)$. This belongs to the little group\footnote{Pay attention that $L(p)$ depends implicitly on the choice of $k$.}. Then,
\begin{equation}
\begin{split}
  U(\Lambda)\ket{p,\sigma}&=N(p)U(L(\Lambda p))U(W(\Lambda,p))\ket{k,\sigma}\\
                          &=N(p)\sum_{\sigma'}D_{\sigma'\sigma}(W)U(L(\Lambda p))\ket{k,\sigma'}\\
			  &=\frac{N(p)}{N(\Lambda p)}\sum_{\sigma'}D_{\sigma'\sigma}(W(\Lambda,p))
			      \ket{\Lambda p,\sigma'}.
\end{split}
\end{equation}

But we have no constraint on the $D$'s: it must just form a representation of the little group. Consequently, we are at a point in which our axioms are no more sufficient to continue the building of quantum field theory: we will get as many theories as representations of the little group.

The physical interpretation is the following\label{pg:phyz_reprez} : each type of particle has its own representation. When we consider a Hilbert space on which $U(\Lambda)$ acts via one given representation of the little group, we consider the Hilbert space which describes the corresponding particle. Note that the little group depends on the choice of $k$, and therefore depends on the particle which is studied (massive or not).

In this sense, a particle is a representation of the Poincaré group\quexto{I think that the irreducibility of a representation is related to \emph{elementary} particles.}. In particular, the nature of the index $\sigma$ can change from the one representation to the other.

\begin{remark}
As far as normalization is concerned, we will pose
\[
  N(p)=\sqrt{k^0/p^0}.
\]
There are some good reasons to take it; but it is irrelevant from our group point of view of the theory.
\end{remark}

\subsection{Positive mass}
%------------------------

This is the easy case. The choice of standard momentum is $k=\begin{pmatrix}1&0&0&0\end{pmatrix}$. One could believe that the little group is $\SO(3)$. It would be the case if we had chosen $L_+^{\uparrow}$ instead of $\SLdc$ --see point \ref{subsec:sym_nature}. In our hermitian representation of $\eR^4$, $k=\mtu$. Then a matrix of $\SLdc$ which leaves it invariant fulfills
\[
   \lambda k\lambda^{\dag}=\lambda\lambda^{\dag}=\mtu,
\]
this is $\lambda\in SU(2)$. By the way, note that $\SO(3)=SU(2)/\eZ_2$.

The celebrated ``law of transformation''\ of a massive particle of spin $j$ (integer or half integer) under the Lorentz transformation $\Lambda$ is 
\begin{equation}
  U(\Lambda)\ket{p,\sigma}=\sqrt{\frac{ (\Lambda p)^0 }{p^0}}
       \sum_{\sigma'}D^{(j)}_{\sigma'\sigma}(W(\Lambda,p))\ket{\Lambda p,\sigma'}
\end{equation}
where $\sigma$ runs from $-j$ to $j$ by step of $1$.

\subsection{Null mass}
%---------------------

In the case of a null mass, the standard vector is $k=\qvect{1}{0}{0}{1}$ and an element of the little group fulfils $Wk=k$. As the little group is part of the Lorentz group, this is an isometry, so
\begin{subequations}
\begin{align}
  \scalh{Wt}{Wk}&=\scalh{t}{k}\\
  \scalh{Wt}{Wt}&=\scalh{t}{t},
\end{align}  
\end{subequations}
for any $t\in\eR^4$. Taking in particular $t=\qvect{1}{0}{0}{0}$,
\begin{subequations}
\begin{align}
  (Wt)^{\mu}k_{\mu}&=t\hmu k_{\mu}=-1\\
  (Wt)^{\mu}(Wt)_{\mu}&=t\hmu t_{\mu}=-1.  
\end{align}  
\end{subequations}
If we write $Wt=({a},{b},{c},{d})$,  the first relation gives $d=a-1$, so that  $Wt=({1+\xi},{\alpha},{\beta},{\xi})$, while the second one gives $\xi=(\alpha^2+\beta^2)/2$. The conclusion is that $W$ acts on $t$ as a certain Lorentz transformation $S(\alpha,\beta)$ :
\begin{equation}
Wt=\begin{pmatrix}
     1+\xi\\
     \alpha\\
     \beta\\
     \xi
   \end{pmatrix}=
   \begin{pmatrix}
     1+\xi  & -\xi    & \alpha & \beta\\
     \alpha & -\alpha &   1    &   0\\
     \beta  & -\beta  &   0    &   1\\
     \xi    & (1+\xi) & \alpha & \beta.
   \end{pmatrix}
   \begin{pmatrix}
     1\\
     0\\
     0\\
     0
   \end{pmatrix}.
\end{equation}
Be careful: it doesn't means that $W=S$, but $Wt=St$. However it is an information: $S(\alpha,\beta)^{-1} W$ is a Lorentz transformation which leaves $t$ invariant. Then it is a spatial rotation. More precisely, since $W$ and $S$ conserve $\qvect{1}{0}{0}{1}$, it is a rotation around the $z$ axis: $S(\alpha,\beta)^{-1} W=R(\theta)$, and
\begin{equation}
  W(\theta,\alpha,\beta)=S(\alpha,\beta)R(\theta)
\end{equation}
is the most general element of the non massive little group.


\chapter{Relativistic fields and fibre bundle formalism}
% This is part of (almost) Everything I know in mathematics
% Copyright (c) 2013-2014
%   Laurent Claessens
% See the file fdl-1.3.txt for copying conditions.

This chapter actually don't deal with \emph{quantum} field theory in the sense that our wave functions aren't operators which acting on a Fock space. So this is just relativistic field theory. 

\subsection{A digression:  \texorpdfstring{$T_Y\yG$}{TYG} and \texorpdfstring{$\yG$}{G}}\label{subsec:digress}
%+++++++++++++++++++++++++++++++++++++++++++

We define two product: $G\times\yG\to TG$ and $\yG\times\yG\to\yG$. If $g\in G$ and $X\in\yG$, we put
\begin{subequations}
\begin{equation} \label{eq_gXdefa}
   gX=\dsdd{ge^{tX}}{t}{0},
\end{equation}
and if $X$, $Y\in\yG$,
\begin{equation}\label{eq:yGyGb}
   XY=\DDsdd{e^{tX}e^{uY}}{t}{0}{u}{0}.
\end{equation}
\end{subequations}
We naturally define the product of a $\yG$-valued $1$-form $A$ by an element $g\in G$ by $(gA)v=gA(v)$.

 Note that $gX$ does not belong to $\yG$ but to $T_{g}G$. Fortunately, in the expressions which we will meet, there will  always be a $g^{-1}$ to save the situation.

Let us give some precisions about derivatives as \eqref{eq:yGyGb}. We consider the expression
\[
  \frac{d}{du}\left( \left.\frac{d}{dt} c_u(t)\right|_{t=0}\right)_{u=0},
\]
which will be more simply written as :
\begin{equation}\label{eq:2307e1}
\DDsdd{ c_u(t) }{u}{0}{t}{0}
\end{equation}
with $c_u(t)\in G$ for all $u,t$; $c_u(0)=e$ for all $u$ and $c_0'(0)=Y\in\yG$ where the prime stands for the derivative with respect of $t$. So $\dsdd{c_u(t)}{t}{0}\in\yG$ for each $u$ and \eqref{eq:2307e1} belongs to $T_Y\yG$. But we know that $\yG$ is a vector space, then $T_Y\yG\simeq\yG$, the isomorphism being given by the following idea: if $\{\partial_i\}$ is a basis of $\yG$ and $\{e_i\}$ the corresponding basis of $T_Y\yG$, we define the action of $A^ie_i\in T_Y\yG$ on $\dpt{f}{G}{\eR}$ by $(A^ie_i)f:=A^i\partial_if$.

\begin{lemma}
Seen as an equality in $\yG$, for $\dpt{f}{G}{\eR}$ we have :
\begin{equation}
   \DDsdd{c_u(t)}{u}{0}{t}{0}f=\DDsdd{f(c_u(t))}{u}{0}{t}{0}.
\end{equation}
\end{lemma}

\begin{proof}
Let us consider $X_u=X_u^i\partial_i=c_u'(0)$ and $X_0=Y$. We naturally have
\begin{align}
   X_uf&=\dsdd{f(c_u(t))}{t}{0},&\text{ and } &&\dsdd{X_u}{u}{0}\in T_Y\yG.
\end{align}
Now, we consider a function $\dpt{h}{\yG}{\eR}$ and compute :
\[
  \Dsdd{X_u}{u}{0}h=\Dsdd{h(X_u)}{u}{0}
                   =\dsdd{ h(\Dsdd{c_u(t)}{t}{0}) }{u}{0}.
\]		   
If $\{\partial_i\}$ is a basis of $\yG$ and $\{e_i\}$, the corresponding one of $T_Y\yG$, thus
\begin{equation}
                   \Dsdd{X_u}{u}{0}h=\left.\dsd{h}{e_i}\right|_Y\DDsdd{c^i_u(t)}{u}{0}{t}{0}.
\end{equation}
So, we can write
\[
   \Dsdd{X_u}{u}{0}=\DDsdd{c^i_u(t)}{u}{0}{t}{0}\left.\dsd{}{e_i}\right|_Y,
\]
and, as element of $\yG$, we consider
\[
  \Dsdd{X_u}{u}{0}=\DDsdd{ c^i_u(t) }{u}{0}{t}{0}\partial_i|_e.
\]
Now, we can compute the action of $\dsdd{X_u}{u}{0}$ on a function $\dpt{f}{G}{\eR}$ as
\begin{equation}
\begin{split}
\Dsdd{X_u}{u}{0}f&=\DDsdd{c^i_u(t)}{u}{0}{t}{0}\left.\dsd{f}{x^i}\right|_e\\
                 &=\Dsdd{ \left.\dsd{f}{x^i}\right|_e\dsdd{c^i_u(t)}{t}{0}  }{u}{0}\\
		 &=\Dsdd{ \dsdd{f(c_u(t))}{t}{0} }{u}{0}.
\end{split}
\end{equation}
\begin{probleme}
Je ne sais pas pourquoi tout d'un coup la dernière équation était commentée, et donc la phrase n'était pas finie.
\end{probleme}

\end{proof}

Let us now see a great consequence of the definition \eqref{eq:yGyGb}
\begin{proposition} \label{prop:XY_YX}
The formula
\begin{equation}
   XY-YX=[X,Y].
\end{equation}
links the formal product inside the Lie algebra and the Lie bracket.
\end{proposition}

\begin{proof}
From this, we can precise our definition of $XY$ when $X$, $Y\in\yG$. The product $XY$ acts on $\dpt{f}{G}{\eR}$ by
\[
  (XY)f=\DDsdd{f(e^{tX}e^{uY})}{t}{0}{u}{0}.
\]
We can get a more geometric interpretation of this. We know how to build a left invariant vector field $\tilde Y$ from any $Y\in\yG$ : for each $g\in G$ we just have to define
\[
  \tY_g(f)=\Dsdd{f(gY(s))}{s}{0}.
\]
%
First remark: $\tY_g$ is precisely the object that previously named ``$gY$''. In order to construct the basis blocks of the formula $XY-YX=[X,Y]$, we turn our attention to $\tX_e\tY$. It is clear that $\tY(f)$ is a function from $G$ to $\eR$, so we can apply $\tX_e$ on it. If $X_t$ is a path which gives the vector $\tX_e$ (for example: $X_t=e^{tX}$), we have
\begin{equation}
  \tX_e(\tY(f))=\Dsdd{\tY(f)_{X_t}}{t}{0}\\
               =\DDsdd{f(X_tY(u))}{u}{0}{t}{0}\\
	       =\DDsdd{f(e^{tX}e^{uY})}{u}{0}{t}{0}.
\end{equation}
Thus we have: $XY=\tX_e\tY$, but it is clear that $[\tX,\tY]_e=\tX_e\tY-\tY_e\tX$. The claim reads now: $[\tX,\tY]_e=[X,Y]$. We can actually take it as de \emph{definition} of $[X,Y]$. It is done in \cite{Helgason}. The link with the definition in terms of successive derivations of $\AD_g(x)=gxg^{-1}$ is not trivial but it can be done. 
\end{proof}



Now, we can give a powerful definition of the wedge for two $\yG$-valued $1$-forms. If $A$, $B\in\Omega^1(M,\yG)$ and $v$, $w\in\cvec(M)$, we define
\begin{equation}
  (A\wedge B)(v,w)=A(v)B(w)-A(w)B(v).
\end{equation}
For $A^2$, we find back the usual definition :
\[
  (A\wedge A)(v,w)=A(v)A(w)-A(w)A(v)=[A(v),A(w)].
\]
%
One can see that for any section $\dpt{\salpha}{\mU_{\alpha}}{P}$, we have
\begin{equation}\label{eq:1907r2}
   \salpha^*(A\wedge A)=(\salpha^*A)\wedge(\salpha^*A).
\end{equation}


\begin{lemma}
If $A$ and $B$ are two $\yG$-valued $1$-forms, one can make  ``simplifications'' as
\begin{equation}
 (Ag)\wedge(g^{-1} B)=A\wedge B.
\end{equation}
\label{lem:simplif}
\end{lemma}

\begin{proof}
We just have to prove that for $A$, $B\in\yG$, $(Ag)(g^{-1} B)=AB$ with definitions \eqref{eq_gXdefa} and \eqref{eq:yGyGb}. Remark that $Ag=\Dsdd{e^{sA}g}{s}{0}$, so
\[
  e^{tAg}=\exp(t\dsdd{e^{sA}g}{s}{0})=\exp(\dsdd{e^{stA}g}{s}{0})=e^{tA}g.
\]
Therefore
\[
  (Ag)(g^{-1} B)=\DDsdd{  e^{tAg}e^{ug^{-1} B}  }{t}{0}{u}{0}=\DDsdd{  e^{tA}gg^{-1} e^{uB}  }{t}{0}{u}{0}=AB.
\]
\end{proof}

\begin{lemma}
\begin{equation}
    F_{\beta}=dA_{\beta}+A_{\beta}^2.
\end{equation}
\end{lemma}

\begin{proof}
This is  a direct consequence of \eqref{eq:1907r2} and $[\sbeta^*,d]=0$.
\end{proof}


\section{Connections}
%++++++++++++++++++++

\subsection{Gauge potentials}
%----------------------------

Let us consider a \defe{section}{section} $\salpha$ of $P$ over $\mU_{\alpha}$. It is a map $\dpt{\salpha}{\mU_{\alpha}}{P}$ such that $\pi\circ\salpha=\id$. A \defe{connection}{connection} on $P$ is a $1$-form $\dpt{\omega}{T_pP}{\yG}\in\Omega^1(P)$ which satisfies the following two conditions:
\begin{subequations}
\begin{align}
   \omega_p(Y^*_p)&=Y,   \label{conn_1}\\
   \omega(dR_g\xi)&=g^{-1}\omega(\xi)g.\label{conn_2}
\end{align}
\end{subequations}
The \defe{gauge potential}{gauge!principal potential} of $\omega$ with respect of the local section\label{PgLocSecConn} $\salpha$ is  the $1$-form on $\mU_{\alpha}$ given by
\begin{equation}
          A_{\alpha}(x)(v)=(\salpha^*\omega)_x(v).
\end{equation}
We will not always explicitly write the dependence of $A_{\alpha}$ in $x$.\nomenclature{$A_{\alpha}$}{Gauge potentials} Now we consider another section $\dpt{\sbeta}{\mU_{\beta}}{P}$ which is related on $\mU_{\alpha}\cap\mU_{\beta}$ to $\salpha$ by $\sbeta(x)=\salpha(x)\cdot g_{\alpha\beta}(x)$ for a well defined map $\dpt{g_{\alpha\beta}}{\mU_{\alpha}\cap\mU_{\beta}}{G}$.

\begin{proposition}
The gauge potentials $A_{\alpha}$ and $A_{\beta}$ are related by
\begin{equation}\label{trans_A}
                A_{\beta}=g^{-1} A_{\alpha} g-g^{-1} dg.
\end{equation}
\label{prop:trans_A}
\end{proposition}

\begin{proof}
By definition, for $v\in T_x\mU_{\alpha}$,
\[
   A_{\beta}(v)=(\sbeta^*\omega)_x(v)=
         \omega_{\salpha(x)\cdot g_{\alpha\beta}(x)}\big((d\sbeta)_x(v)\big).
\]
We begin by computing $d\sbeta(v)$. Let us take a path $v(t)$ in $\mU_{\alpha}$ such that $v(0)=x$ and $v'(0)=v$. We have :
\begin{equation}\label{eq:1407r1}
\begin{split}
   (d\sbeta)_x(v)&=\dsdd{\sbeta(v(t))}{t}{0}\\
                 &=\dsdd{\salpha(v(t))\cdot\gab(v(t))}{t}{0}\\
		 &=\Dsdd{\salpha(v(t))\cdot\gab(x)}{t}{0}
		    +\Dsdd{\salpha(x)\cdot\gab(v(t))}{t}{0}\\
		 &=dR_{\gab(x)}d\salpha(v)+\Dsdd{\salpha(x)\cdot\gab(x)e^{-ts}}{t}{0}\\
		 &=dR_{\gab(x)}d\salpha(v)+s^*_{\salpha(x)\cdot\gab(x)}
\end{split}
\end{equation}
where $s$ is defined by the requirement\label{pg:justif_s} that $\gab(x)^{-1}\gab(v(t))$ can be replaced in the derivative by $e^{-ts}$, so that we can replace $\gab(v(t))$ by $\gab(x)e^{-ts}$. As far as the derivatives are concerned, $e^{-ts}=\gab(x)^{-1}\gab(v(t))$, then
\[
     s=-\dsdd{\gab(x)^{-1}\gab(v(t))}{t}{0}=-\gab(x)^{-1} d\gab(v),
\]
this equality being a notation. Now, properties \eqref{conn_1} and \eqref{conn_2} make that
\[
   A_{\beta}(v)=\gab(x)^{-1}\omega_{\salpha(x)}(d\salpha(v))\gab(x)+s.
\]
The thesis is just the same, with ``reduced'' notations (see section \ref{subsec:digress}).
.
\end{proof}
An explicit form for this transformation law is :
\begin{equation}
    A_{\beta}(v)=\Dsdd{g^{-1} e^{tA_{\alpha}(v)}g}{t}{0}-\Dsdd{g^{-1}\gab(v(t))}{t}{0},
\end{equation}
where $g:=\gab(x)$.

\subsection{Covariant derivative}
%-------------------------------

When we have a connection on a principal bundle, we can define a covariant derivative\index{covariant!derivative} on any associated bundle. Let us quickly review it. An associated bundle is the semi-product $E=P\times_{\rho} V$ where $V$ is a vector space on which acts the representation $\rho$ of $G$. We denote the canonical projection by $\dpt{\pi_p}{E}{M}$. The classes are taken with respect to the equivalence relation $(p,v)\sim(p\cdot g,\rho(g^{-1})v)$.

A \defe{section}{section!of associated bundle} of $E$ is a map $\dpt{\psi}{M}{E}$ such that $\pi\circ\psi=\id$. We denote by $\Gamma(E)$ the set of all the sections of $E$. A section of $E$ defines (and is defined by) an equivariant function\index{equivariant!function} $\dpt{\hpsi}{P}{V}$ such that
\begin{subequations}
\begin{align}
  \psi(\pi(\xi))&=[\xi,\hpsi(\xi)],\\
  \hpsi(\xi\cdot g)&=\rho(g^{-1})\hpsi(\xi).
\end{align}
\end{subequations}
For a section $\psi\in\Gamma(E)$, we define $\dpt{\psi\bsa}{\mU_{\alpha}}{V}$ by
 \[
 \psi\bsa(x)=\hpsi(\sigma(x)).
 \]
We saw in \eqref{eq:nabla_coord} that a covariant derivative on $E$ is given by
\begin{equation}\label{3008e1}
  (D_X\psi)\bsa(x)=X_x\psi\bsa-\rho_*\Big((\salpha^*\omega)_x(X_x)\Big)\psi\bsa(x).
\end{equation}
Since $(d\psi)(X)=X(\psi)$, we can rewrite this formula in a simpler manner by forgetting the index $\alpha$ and the mention of $X$ :
\[
    D\psi=d\psi-(\rho_*A_{\alpha})\psi.
\]
If we note $(\rho_*A_{\alpha})\psi$ by $A_{\alpha}\psi$, we have
\begin{equation}
        D\psi=d\psi-A\psi.
\end{equation}
One has to understand that equation as a ``notational trick''\ for \eqref{3008e1}. By the way, remark that $(\rho_*A_{\alpha})$ is the only ``reasonable'' way for $A$ to act on $\psi$.

\section{Gauge transformation}
%++++++++++++++++++++++++++++

A \defe{gauge transformation}{gauge!transformation} of a $G$-principal bundle is a diffeomorphism $\dpt{\varphi}{P}{P}$ which satisfies
\begin{subequations}
\begin{align}
   \pi\circ\varphi&=\pi,\\
   \varphi(\xi\cdot g)&=\varphi(\xi)\cdot g.
\end{align}
\end{subequations}
In local coordinates, it can be expressed in terms of a function $\dpt{\tilde{\varphi}_{\alpha}}{\mU_{\alpha}}{G}$ by the requirement that
\begin{equation}\label{def_vpt}
    \varphi(\salpha(x))=\salpha(x)\cdot\tilde{\varphi}_{\alpha}(x).
\end{equation}

We have shown in proposition \ref{prop:vp_conn} that, if $\omega$ is a connection $1$-form on $P$, the form $\varphi\cdot\omega:=\varphi^*\omega$ is still a connection $1$-form on $P$. Of course, with the same section $\salpha$ than before, we can define a gauge potential $(\varphi\cdot A)_{\alpha}$ for this new connection. We will see how it is related to $A_{\alpha}$. The reader can guess the result (it will be the same as proposition \ref{prop:trans_A}). We show it.

\begin{proposition}		\label{Proptr_de_A}
\begin{equation}\label{tr_de_A}
     (\varphi\cdot A)=\tilde{\varphi}^{-1} A\tilde{\varphi}-\tilde{\varphi}^{-1} d\tilde{\varphi}.
\end{equation}
\end{proposition}
\begin{proof}
Let us consider $x\in\mU_{\alpha}$, and $v\in T_x\mU_{\alpha}$, the vector which is tangent to the curve $v(t)\in\mU_{\alpha}$. We compute
\[
    \salpha^*(\varphi^*\omega)_x(v)=\omega_{(\varphi\circ\salpha)(x)}((d\varphi\circ d\salpha)(v)),
\]
but equation \eqref{def_vpt} makes
\begin{equation}
\begin{split}
   (d\varphi\circ d\salpha)(v)&=\dsdd{\varphi(\salpha(v(t)))}{t}{0}\\
                          &=\dsdd{\salpha(v(t))\cdot\tilde{\varphi}_{\alpha}(v(t))}{t}{0}.
\end{split}
\end{equation}
Now, we are in the same situation as in equation \eqref{eq:1407r1}.
\end{proof}

If $\dpt{\psi}{M}{E}$ is a section of $E$, the gauge transformation $\dpt{\varphi}{P}{P}$ acts on $\psi$ by
\begin{equation}
   \widehat{\varphi\cdot\psi}(\xi)=\hpsi(\varphi^{-1}(\xi)).
\end{equation}
On the other hand, $\varphi$ acts on the covariant derivative (and the potential) :
$\varphi\cdot D$ is the covariant derivative  of the connection $\varphi\cdot\omega$. Of course, we define
\begin{equation}
    (\varphi\cdot D)\psi=d\psi-(\varphi\cdot A)\psi.
\end{equation}

\begin{lemma}
If $\dpt{\varphi}{P}{P}$ is a gauge transformation, then
\begin{enumerate}
\item $\varphi^{-1}$ is also a gauge transformation and
             $(\widetilde{\varphi^{-1}})_{\alpha}(x)=\tilde{\varphi}_{\alpha}(x)^{-1}$, \label{lem:i}
\item $(\varphi\cdot\psi)\bsa(x)=\rho(\tilde{\varphi}_x^{-1})\psi\bsa(x)$.\label{lem:ii}
\end{enumerate}
\label{lem:prop_gauge}
\end{lemma}

\begin{proof}
The first part is clear while the second is a computation :
\begin{equation}
    (\varphi\cdot\psi)\bsa=\widehat{\varphi\cdot\psi}(\salpha(x))
                         =\hpsi(\varphi^{-1}(\salpha(x)))
			 =\hpsi(\salpha(x)\cdot\tilde{\varphi}_{\alpha}(x)^{-1})
			 =\rho(\tilde{\varphi}_{\alpha}(x))\psi\bsa(x).
\end{equation}
\end{proof}

Now, we will proof the main theorem: the one which explains why the covariant derivative is ``covariant''.

\begin{theorem}
The covariant derivative $D$ fulfils a ``covariant'' transformation rule under gauge transformations:
\begin{equation}\label{eq:covariance_math}
      (\varphi\cdot D)(\varphi^{-1}\cdot \psi)=\varphi^{-1}(D\psi).
\end{equation}
\label{th:covariance}
\end{theorem}
\begin{remark}
Let us use more intuitive notations: we write \eqref{tr_de_A} under the form $A'=g^{-1} Ag-g^{-1} dg$. If we have two sections  $\psi$ and $\psi'$, they are necessarily related by a gauge transformation: $\psi'=g^{-1}\psi$. Then, the theorem tells us that the equation $D\psi=d\psi-A\psi$ becomes $D'\psi'=g^{-1} D\psi$ ``under a gauge transformation''. This is: $D\psi$ transforms under a gauge transformation as $d\psi$ transforms under a constant linear transformation. This is the reason why $D$ is a \emph{covariant} derivative. The physicist way to write \eqref{eq:covariance_math} is
\begin{equation}\label{eq:covariance_phys}
    D'\psi'=g^{-1} D\psi
\end{equation}
\label{rem:intuitif}
\end{remark}

\begin{proof}[Proof of theorem \ref{th:covariance}]
First, we look at $(\varphi\cdot A)\psi_{\alpha}$. Using all the notational tricks used to give a sens to $A\psi$, we write :
\[
   [(\varphi\cdot A)_X\psi]\bsa(x)=(\varphi\cdot A)_X\psi\bsa(x)=\rho_*(\varphi\cdot A(X))\psi\bsa(x).
\]
But we know that $\varphi\cdot A=\tilde{\varphi}^{-1} A\tilde{\varphi}-\tilde{\varphi}^{-1} d\tilde{\varphi}$, then
\begin{equation}\label{eq:en_deux}
\begin{split}
  (\varphi\cdot A)_X\psi\bsa(x)
  &=\rho_*(\tilde{\varphi}^{-1} A(X)\tilde{\varphi})\psi\bsa(x)\\
  &\quad-\rho_*(\tilde{\varphi}^{-1} d\tilde{\varphi}(X))\psi\bsa(x)\\
  &=\Dsdd{ \rho(\tilde{\varphi}^{-1} e^{tA(X)}\tilde{\varphi})\psi\bsa(x)}{t}{0}\\
 &\quad-\Dsdd{ \rho(\tilde{\varphi}^{-1}\tilde{\varphi}(X_t))\psi\bsa(x) }{t}{0}
\end{split}
\end{equation}
Now, we have to write this equation with $\varphi^{-1}\cdot\psi$ instead of $\psi$. Using lemma \ref{lem:prop_gauge}, we find :
\begin{equation}
\begin{split}
   (\varphi\cdot A)_X(\varphi^{-1}\cdot\psi)\bsa(x)
   &=\Dsdd{ \rho(\tilde{\varphi}^{-1} e^{tA(X)}\tilde{\varphi}\tilde{\varphi}^{-1})\psi\bsa(x)}{t}{0}\\
   &\quad-\Dsdd{ \rho(\tilde{\varphi}^{-1}\tilde{\varphi}(X_t)\tilde{\varphi}^{-1})\psi\bsa(x) }{t}{0}
\end{split}
\end{equation}
After simplification, the first term is a term of the thesis: $\tilde{\varphi}(x)^{-1}(A\psi)_{\alpha}(x)$ and we let the second one as it is. Now, we turn our attention to the second term of \eqref{eq:covariance_math}; the same argument gives:
\begin{equation}
\begin{split}
  d(\varphi^{-1}\psi\bsa)_xX
  &=\Dsdd{(\varphi^{-1}\cdot\psi)\bsa(X_t)}{t}{0}\\
  &=\Dsdd{\rho(\tilde{\varphi}(X_t)^{-1})\psi\bsa(X_t)}{t}{0}\\
  &=\Dsdd{\rho(\tilde{\varphi}(X_t)^{-1})\psi\bsa(x)}{t}{0}
  +\Dsdd{\rho(\tilde{\varphi}^{-1})\psi\bsa(X_t)}{t}{0}.
\end{split}
\end{equation}
The second term is $\tilde{\varphi}^{-1} d\psi_{\alpha}(X)$. In definitive, we need to prove that the two exceeding terms cancel each other:
\begin{equation}\label{eq:le_zero}
  \Dsdd{\rho(\tilde{\varphi}^{-1}\tilde{\varphi}(X_t)\tilde{\varphi}^{-1})\psi\bsa(x)}{t}{0}
  +\Dsdd{\rho(\tilde{\varphi}(X_t)^{-1})\psi\bsa(x)}{t}{0}
\end{equation}
must be zero.

One can find a $g(t)\in G$ such that $\tilde{\varphi}(X_t)=\tilde{\varphi} g(t)$, $g(0)=e$. Then, what we have in the $\rho$ of these two terms is respectively $g(t)\tilde{\varphi}^{-1}$ and $g(t)^{-1}\tilde{\varphi}^{-1}$. As far as the derivative are concerned, $g(t)$ can be written as $e^{tZ}$ for a certain $Z\in\yG$. So we see that $g(t)^{-1}=e^{-tZ}$ and the derivative will come with the right sign to makes the sum zero.
\end{proof}

\begin{remark}
If we naively make the computation with the notations of remark~\ref{rem:intuitif}, we replace $\psi'=g^{-1}\psi$ and $A'=g^{-1} Ag-g^{-1} dg$ in
\[
  D'\psi'=d\psi'-A'\psi',
\]
using some intuitive ``Leibnitz formulas'', we find :
$D'\psi'=dg^{-1}\psi+g^{-1} d\psi+g^{-1} A\psi+g^{-1} dg g^{-1}\psi$. It is exactly $g^{-1} d\psi+g^{-1} A\psi$ with two additional terms: $dg^{-1}\psi$ and $g^{-1} dg g^{-1}\psi$. One sees that these are precisely the two terms of the expression \eqref{eq:le_zero}. We will give a sens to this ``naive''\ computation in section~\ref{subsec:digress}.
\end{remark}
\section{A bite of physics}
%++++++++++++++++++++++++++

\subsection{Example: electromagnetism}\index{electromagnetism}
%-------------------------------------

Let us consider the electromagnetism as the simplest example of a gauge invariant physical theory. We first discuss the theory of free electromagnetic field (this is: without taking into account the interactions with particles) from Maxwell's\index{Maxwell} equations, see \cite{Schomblond_em,llf}. The electric field $\bE$ and the magnetic field $\bB$ are subject to following relations:
\begin{subequations}
\begin{align}
\nabla\cdot{\bE}&=\rho,    \label{M1}\\
\nabla\cdot{\bB}&=0,                     \label{M2}  \\
\nabla\times{\bE}+\partial_t{\bB}&=0,         \label{M3}\\
\nabla\times{\bB}-\partial_t{\bE}&=\overline{ j }. \label{M4}
\end{align}
\end{subequations}
Comparing \eqref{M1} and \eqref{M2}, we see that Maxwell's theory does not incorporate magnetic monopoles.
Suppose that we can use the Poincaré lemma. Equation \eqref{M2} gives a vector field $\bA$ such that $\bB=\nabla\times\bA$, so that \eqref{M3} becomes $\nabla\times(\bE+\partial_t\bA)=0$ which gives a scalar field $\phi$ such that $-\nabla\cdot\phi=\bE+\partial_t\bA$.

Now the equations \eqref{M1}--\eqref{M4} are equations for the potentials $\bA$ and $\phi$, and we find back the ``physical''\ field by
\begin{subequations}\label{BE_de_A}
\begin{align}
    \bB&=\nabla\times\bA,\\
    \bE&=-\nabla\phi-\partial_t\bA.
\end{align}
\end{subequations}

One can easily see that there are several choice of potentials\index{potential} which describe the same electromagnetic field. Indeed, if
\begin{subequations}\label{jauge_A}
\begin{align}
    \bA'&=\bA+\nabla\lambda,\\
    \phi'&=\phi-\partial_t\lambda,
\end{align}
\end{subequations}
the electromagnetic field given (via \eqref{BE_de_A}) by $\{\phi',\bA'\}$ is the same as the one given by $\{\phi,\bA\}$

The Maxwell's equations can be written in a more ``covariant''\ way by defining
\begin{equation}\label{def_F}
F=\begin{pmatrix}
0 & -E_x/c & -E_y/c & -E_z/c \\
\cdot & 0 & -B_z & \cdot \\
\cdot & \cdot & 0 & -B_x \\
\cdot & -B_y & \cdot & 0
\end{pmatrix},
\end{equation} $F^{\mu\nu}=-F^{\nu\mu}$ and\nomenclature{$F_{\mu\nu}$}{Electromagnetic field strength}
\[
 J=\begin{pmatrix}
c\rho & j_x & j_y & j_z
\end{pmatrix}.
\]\nomenclature{$J_{\mu}$}{Electromagnetic $4$-current}
We also define $\star F^{\alpha\beta}=\frac{1}{2} e^{\alpha\beta\lambda\mu}F_{\lambda\mu}$. With all that, Maxwell's equations read:
\begin{equation}
\begin{split}
\partial_{\mu} F^{\mu\nu}&=\mu_0J^{\nu},\\
\partial_{\alpha}\star F^{\alpha\beta}&=0.
\end{split}
\end{equation}
If we define
\begin{equation}\label{def_A}
  A=\begin{pmatrix}
\frac{\displaystyle\phi}{\displaystyle c} & -A_x & -A_y & -A_z
\end{pmatrix},
\end{equation}
%
the physical fields are given by
\[
  F_{\mu\nu}=\partial_{\mu}A_{\nu}-\partial_{\nu}A_{\mu}.
\]
The \defe{gauge invariance}{gauge!invariance!of electromagnetism} of this theory is the fact that
\begin{subequations}
\begin{equation}
  F'_{\mu\nu}=\partial_{\mu}A'_{\nu}-\partial_{\nu}A'_{\mu}=F_{\mu\nu}
\end{equation}
when
\begin{equation}
A'_{\mu}(x)=A_{\mu}(x)+\partial_{\mu}f(x)
\end{equation}
\end{subequations}
for any scalar \emph{function} $f$ (to be compared with \eqref{jauge_A}).

This is: in the picture of the world in which we see the $A$ as fundamental field of physics, several (as much as you have functions in $C^{\infty}(\eR^4)$) fields $A$, $A'$,\ldots\ describe the \emph{same} physical situation because the fields $\bE$ and $\bB$ which acts on the particle are the same for $A$ and $A'$.

Now, we turn our attentions to the interacting field theory of electromagnetism. As far as we know, the electron makes interactions with the electromagnetic field via a term $\overline{ \psi } A_{\mu}\psi$ in the Lagrangian. The free Lagrangian for an electron is
%
%\begin{center}
%\begin{fmffile}{monfey}
%\begin{fmfgraph*}(40,25)
%\fmfleft{i1}
%\fmfright{o1,o2}
%\fmf{fermion}{i1,v1}
%\fmf{photon}{v1,o1}
%%\fmf{fermion}{v1,o2}
%\fmflabel{$e^-$}{i1}
%\fmflabel{$e^-$}{o2}
%\fmflabel{$A_{\mu}$}{o1}
%\end{fmfgraph*}
%\end{fmffile}
%\end{center}
%
\begin{equation}\label{eq:freeL}
\mL=\overline{ \psi }(\gamma^{\mu}\partial_{\mu}+m)\psi.
\end{equation}
%
 The easiest way to include a $\overline{ \psi } A\psi$ term is to change $\partial_{\mu}$ to $\partial_{\mu}+A_{\mu}$. But we want to preserve the powerful gauge invariance of classical electrodynamics, then we want the new Lagrangian to keep unchanged if we do
 \begin{equation} \label{eq_jaugeA_em}
 A_{\mu}\rightarrow A_{\mu}'=A_{\mu}-i\partial_{\mu}\phi.
 \end{equation}
  In order to achieve it, we remark that the $\psi$ must be transformed \textit{simultaneously} into
\begin{equation}\label{eq:jaugepsi_em}
 \psi'(x)=e^{i\phi(x)}\psi(x).
\end{equation}

The conclusion is that if one want to write down a Lagrangian for QED\index{QED}, one must find a Lagrangian which remains unchanged under certain transformation $A\rightarrow A'$ and $\psi\rightarrow\psi'$. In other words the set $\{\psi,A\}$ of fields which describe the world of an electron in an electromagnetic field is not well defined from data of the physical situation alone: it is defined up to a certain invariance which is naturally called a \defe{gauge invariance}{gauge!invariance}.

\begin{remark}
In the physics books, the matter is presented in a slightly different way. We observe that the Lagrangian \eqref{eq:freeL} is invariant under
\begin{equation}\label{eq:globale}
\psi(x)\rightarrow\psi'(x)=e^{i\alpha}\psi(x)
\end{equation}
%
for any \emph{constant} $\alpha$. One can see that the associated conserved current (Noether) is closely related to the electric current. The idea (of Yang-Mills) is to develop this symmetry. Since the symmetry \eqref{eq:globale} depends only on a constant, we say it a \defe{global}{global symmetry} symmetry; we will simultaneously add a new field $A_{\mu}$ and upgrade \eqref{eq:globale} to a \defe{local}{local symmetry (physics)} symmetry:
 \begin{equation}\label{eq:locate}
\psi(x)\rightarrow\psi'(x)=e^{i\phi(x)}\psi(x).
\end{equation}
Then, we deduce the transformation law of $A_{\mu}$.
 \end{remark}

Because of the form of \eqref{eq:jaugepsi_em}, we say that the electromagnetism is a $U(1)$-gauge theory. The fact that this is an abelian group have a deep physical meaning and many consequences.

\subsection{Little more general, slightly more formal}
%--------------------------------------------

The aim of this text is to interpret the field $A$ as a gauge potential for a connection. But equation \eqref{eq_jaugeA_em} is not exactly the expected one which is \eqref{tr_de_A}. The point is that equation \eqref{eq_jaugeA_em} concerns a theory in which the gauge transformation of the field was a simple multiplication by a scalar field, so that simplifications as $e^{-i\phi(x)}A_{\mu}(x)e^{i\phi(x)}=A_{\mu}(x)$ are allowed.

Now, we consider a vector space $V$, a manifold $M$ and a function $\dpt{\psi}{M}{V}$ which ``equation of motion''\ is
\[
   L^i(\partial_i+m_i)\psi=0
\]
Where we imply an unit matrix behind $\partial$ and $m$; the indices $i,j$ are the (local) coordinates in $M$ and $a,b$, the coordinates in $V$. Let $G$ be a matrix group which acts on $V$. If $\psi$ is a solution, $\Lambda^{-1}\psi$ is also a solution as far as $\Lambda$ is a constant --does not depend on $x\in M$-- matrix of $G$. In other words, $L^i(\partial_i+m_i)\psi_a=0$ for all $a$ implies $L^i(\partial_i+m_i)((\Lambda^{-1})^b_a\psi_b)=0$.

The function, $\psi'(x)=\Lambda(x)^{-1}\psi(x)$ is no more a solution. If we want it to be solution of the same equation as $\psi$, we have to change the equation and consider
\[
   L^i(\partial_i+A_i+m_i)\psi=0.
\]
This equation is preserved under the \emph{simultaneous} change
\begin{equation}\label{eq:jaugeG}
    \left\{\begin{aligned}
	   \psi_a'&=(\Lambda^{-1})^b_a\psi_b\\
           (A_i')^a_b&=(\Lambda^{-1})^c_b (A_i)^d_c(\Lambda^a_d)-(\partial_i\Lambda^{-1})^d_b\Lambda^a_d.
          \end{aligned}\right.
\end{equation}
The second line show that the formalism in which $A$ is a connection is the good one to write down covariant equations. This has to be compared with \eqref{trans_A}. Logically, a theory which includes an invariance under transformations as \eqref{eq:jaugeG} is called a $G$-gauge theory.

\subsection{A ``final'' formalism}
%---------------------------------------

Now, we work with fields which are sections of some fiber bundle build over $M$, the physical space. More precisely, let $G$ be a matrix group. 

\begin{probleme}
	For sure, it also works for a much lager class of groups. Which one ?
\end{probleme}


We search for a theory which is ``locally invariant under $G$''. In order to achieve it, we consider a $G$-principal bundle $P$ over $M$ and the associated bundle $E=P\times_{\rho}V$ for a certain vector space $V$, and a representation $\rho$ of $G$ on $V$. Typically, $V$ is $\eC$ or the vector space on which the spinor representation acts.

The physical fields are sections $\dpt{\psi}{M}{E}$. If we choose some reference sections $\dpt{\sigma_{\alpha}}{M}{P}$, they can be expressed by $\psi\bsa(x)=\hpsi(\salpha(x))$. We translate the idea of a local invariance under $G$ by requiring an invariance under
\[
     \psi'\bsa(x)=\rho(g(x))\psi\bsa(x)
\]
for every $\dpt{g}{M}{G}$. By \ref{lem:ii} of lemma \ref{lem:prop_gauge}, we see that $\psi'\bsa(x)=(\varphi^{-1}\cdot\psi)\bsa(x)$, where $\dpt{\varphi}{P}{P}$ is the gauge transformation given by
\[
   \varphi(\salpha(x))=\salpha(x)\cdot g(x).
\]

We want $\psi$ and $\psi'$ to ``describe the same physics''. From a mathematical point of view, we want $\psi$ and $\psi'$ to \emph{satisfy the same equation}. It is clear that equation $d\psi=0$ will not work.

The trick is to consider any connection $\omega$ on $P$ and the gauge potential $A$ of $\omega$. In this case the equation
\begin{equation}\label{eq:Dpsi}
    (d-A)\psi=0\qquad\textrm{ or }\qquad D\psi=0
\end{equation}
is preserved under 
\[ 
 \begin{split}
A&\rightarrow\varphi\cdot A,\\
 \psi&\rightarrow\varphi^{-1}\cdot\psi.
\end{split} 
\]
 Theorem \ref{th:covariance} powa !
        
In this sense, we say that equation \eqref{eq:Dpsi} is gauge invariant, and is thus taken by physicists to build some theories when they need a ``local $G$-covariance''. This gives rise to the famous Yang-Mills theories.

In this picture the matter field $\psi$ and the bosonic field $A$ are both defined from a $U(1)$-principal bundle. When physicists say 
\begin{quote}
	$\psi$ transforms as ``blahblah'' under a $U(1)$ transformation,
\end{quote}
they mean that $\psi$ is a section of an $U(1)$-associated bundle; 
when they say
\begin{quote}
	$A$ transforms as ``blahblah'' under a $U(1)$ transformation,
\end{quote}
they mean that $A$ is the gauge potential of a connection on a $U(1)$-principal bundle. In each case, the ``blahblah'' denotes an irreducible\footnote{Irreducibility is for elementary particles} representation of $U(1)$.

\begin{remark}
The mathematics of equation \eqref{eq:Dpsi} only requires a $\yG$-valued connection on $P$. There are several physical constraints on the choice of the connection. These give rise to interaction terms between the gauge bosons. We will not discuss it at all. This a matter of books about quantum field theories.

The most used Yang-Mills groups in physics are $U(1)$ for the QED, $SU(2)$ for the weak interactions and $SU(3)$ for chromodynamic.
\end{remark}
\section{Curvature}
%+++++++++++++

\subsection{Intuitive setting}
%-------------------------------

From the $\yG$-valued connection $1$-form $\omega$ on $P$, we may define its \defe{curvature $2$-form}{curvature} :
\begin{equation}
     \Omega=d\omega+\omega\wedge\omega.
\end{equation}
As before, we can see $\Omega$ as a $2$-form on $M$ instead of $P$. For this, we just need some sections $\dpt{\salpha}{\mU_{\alpha}}{P}$ and define
\begin{equation}
        F_{\alpha}=\salpha^*\Omega.
\end{equation}
This $F$ is called the \defe{Yang-Mills field strength}{Yang-Mills!field strength}. The question is now to see how does it transform under a change of chart ? What is $F_{\beta}=\sbeta^*\Omega$ in terms of~$F_{\alpha}$ ?

First, note that we can't try to find a relation like $d(g\omega)=dg\wedge\omega+g\,d\omega$. Pose $A_x=g(x)\omega_x$:
\[
  A_x(v)=\dsdd{g(x)e^{t\omega_x(v)}}{t}{0}.
\]
Using
\[
   (d\alpha)(v,w)=v(\alpha(w))-w(\alpha(v))-\alpha([v,w]),
\]
we are led to write
\begin{equation}
     w(A(v))=d(A(v))w
            =\dsdd{A_{w_u}(v)}{u}{0}
	    =\dsdd{ \Dsdd{g(w_u)e^{t\omega_{w_u}(v)}}{t}{0} }{u}{0}.
\end{equation}
But at $t=u=0$, the expression in the bracket is $g(x)$, and not $e$. Then the whole expression is not an element of $\yG$. In other words, the problem is that for $\dpt{g}{M}{G}$, we have $\dpt{dg_x}{T_xM}{T_{g(x)}G\neq T_eG}$.

Now, remark that in our matter, the problem will not arise because in the expressions $A_{\beta}=g^{-1} A_{\alpha} g+g^{-1} dg$, each term has a $g$ and a $g^{-1}$.

\begin{lemma}
\begin{equation}
   d(g^{-1})_x(v)=-g(x)^{-1} dg(v)g(x)^{-1}.
\end{equation}
\label{lem:dgemu}
\end{lemma}

\begin{proof}
Let $v_t$ be a path which defines the vector $v$, and define $Y\in\yG$ such that as far as the derivative are concerned, we have $g(v_t)=g(x)e^{tY}$. Then,
\[
      g(g^{-1})(v)=\Dsdd{g(v_t)^{-1}}{t}{0}=\Dsdd{e^{-tY}g(x)^{-1}}{t}{0}.
\]
But on the other hand,
\[
  g^{-1} dg(v)g^{-1}=\Dsdd{g(x)^{-1} g(v_t)g(x)^{-1}}{t}{0}=\Dsdd{e^{tY}g(x)^{-1}}{t}{0},
\]
thus $d(g^{-1})_x(v)=-g(x)^{-1} dg(v)g(x)^{-1}$, as we want.
\end{proof}

\begin{theorem}
\begin{equation}
     F_{\beta}=g^{-1} F_{\alpha} g.
\end{equation}
\label{tho:trans_F}
\end{theorem}

\begin{proof}[Naive proof]
Let us accept $F_{\beta}=dA_{\beta}+A_{\beta}\wedge A_{\beta}$. With proposition \ref{prop:trans_A}, we can perform a simple computation with all the intuitive ``Leibnitz rules''\ :
\[
   dA_{\beta}=-g^{-1} dg\, g^{-1}\wedge A_{\alpha} g+g^{-1} dA_{\alpha} g+g^{-1} A_{\alpha}\wedge dg-g^{-1} dg\,g^{-1}\wedge dg,
\]
and
\[
  A_{\beta}\wedge A_{\beta}=g^{-1} A_{\alpha} g\wedge g^{-1} A_{\alpha} g+g^{-1} A_{\alpha} g\wedge g^{-1} dg+g^{-1} dg\wedge g^{-1} A_{\alpha} g+g^{-1} dg\wedge g^{-1} dg.
\]
The sum is obviously the announced result.
\end{proof}
This proof seems too beautiful to be false\footnote{More precisely, it is as beautiful as we want it to be true.}. Here is the full proof.

\begin{proof}[Ultimate proof of theorem \ref{tho:trans_F}]
First we compute $d(g^{-1} A_{\alpha} g)$. In order to do this, remark that the $1$-form $g^{-1} A_{\alpha} g$ is explicitly given on $v\in\cvec(M)$ by
\[
   (g^{-1} A_{\alpha} g)(v)_x=\Dsdd{g(x)^{-1} e^{tA(v)_x}g(x)}{t}{0}.
\]
For all $x\in M$, this expression is an element of $\yG$; then we can say that $(g^{-1} A_{\alpha} g)(v)$ is a map $\dpt{(g^{-1} A_{\alpha} g)(v)}{M}{\yG}$. So it is unambiguous to write $w((g^{-1} A_{\alpha} g)(v))\in\yG$ for $w\in T_xM$.

We will use the formula
\[
   d(g^{-1} A_{\alpha} g)(v,w)=v(g^{-1} A_{\alpha} g)(w)-w(g^{-1} A_{\alpha} g)(v)-(g^{-1} A_{\alpha} g)([v,w]).
\]
As $w((g^{-1} A_{\alpha} g)(v))=d((g^{-1} A_{\alpha} g)(v))w$, we have
\begin{equation}
\begin{split}
    w((g^{-1} A_{\alpha} g)(v))&=\dsdd{(g^{-1} A_{\alpha} g)(v)_{w_u}}{u}{0}\\
                &=\dsdd{ \Dsdd{g(w_u)^{-1} e^{tA(v)_{w_u}}g(w_u)}{t}{0}  }{u}{0}\\
		&=\dsdd{  \Dsdd{g(w_u)^{-1}}{u}{0} e^{tA(v)_x}g(x)  }{t}{0}\\
		&\quad+\dsdd{ g(x)^{-1} \Dsdd{e^{tA(v)_{w_u}}}{u}{0} g(x)  }{t}{0}\\
		&\quad+\dsdd{  g(x)^{-1} e^{tA(v)_x} \Dsdd{g(w_u)}{u}{0}  }{t}{0}\\
		&=d(g^{-1})(w)A(v)_xg(x)\\
		&\quad+g(x)^{-1} w_x(A(v))g(x)\\
		&\quad+g(x)^{-1} A(v)_x dg(w).
\end{split}
\end{equation}
On the other hand, one easily finds that
\[
     (g^{-1} A_{\alpha} g)([v,w])=g(x)^{-1} A([v,w])g(x).
\]
 Using lemma \ref{lem:dgemu}, we have
\begin{equation}
\begin{split}
   d(g^{-1} A_{\alpha} g)_x(v,w)&=-g(x)^{-1} dg(v)g(x)^{-1} A(w)_xg(x)+g(x)^{-1} v(A(w))g(x)\\&\quad+g(x)^{-1} A(w)_xdg(v)_x\\
                   &\quad+g(x)^{-1} dg(w)_xg(x)^{-1} A(v)_xg(x)-g(x)^{-1} w(A(v))g(x)\\&\quad-g(x)^{-1} A(v)_xdg(w)\\
		   &\quad-g(x)^{-1} A([v,w])g(x).
\end{split}
\end{equation}
We can regroup the terms two by two in order to form $dA_{\alpha}$ and some wedge; with simpler notations, we can write :
\begin{equation}\label{eq:dA_1}
  d(g^{-1} A_{\alpha} g)=-(g^{-1} dg\,g)\wedge(A_{\alpha} g)-(g^{-1} A)\wedge dg+(g^{-1} dA g).
\end{equation}
We compute $d(g^{-1} dg)$ in the same way; the result is
\[
   (g^{-1} dg)(v)_x=\Dsdd{g(x)^{-1} g(v_t)}{t}{0}\in\yG.
\]
For $v$, $w\in\cvec(M)$, we have :
\begin{equation}
\begin{split}
   w\big((g^{-1} dg)(v)\big)&=\dsdd{ (g^{-1} dg)(v)_{w_u} }{u}{0}\\
                   &=\DDsdd{  g(w_u)^{-1} g(v_{w_u}(t))  }{u}{0}{t}{0}\\
		   &=\DDsdd{  g(w_u)^{-1} g(v_t)  }{t}{0}{u}{0}
		      +\DDsdd{  g(x)^{-1} g(w_u(t))  }{t}{0}{u}{0}\\
		   &=d(g^{-1})(w)dg(v)+\Dsdd{ g(x)^{-1} dg(v_{w_u}) }{u}{0}\\
		   &=-g^{-1} dg(w)g^{-1} dg(v)+g(x)^{-1} w(dg(v))
\end{split}
\end{equation}
where $w_u$ is a path such that $w'_0=w_x$ and $v_{w_u}(t)$ is, with respect of $t$, a path which gives the vector $v_{w_u}$. On the another hand, we have
\[
   (g^{-1} dg)([v,w])=g^{-1} dg([v,w]).
\]

Remark that the term $g(x)^{-1} w(dg(v))$ of  $w((g^{-1} dg)(v))$ together with the same with $v\leftrightarrow w$ and $(g^{-1} dg)([v,w])$ which comes from  $(g^{-1} dg)([v,w])$ will give $g(x)^{-1}(d^2g)(v,w)=0$ when we will compute $d(g^{-1} dg)$.
Finally,
\begin{equation}\label{eq:dA_2}
   d(g^{-1} dg)=-(g^{-1} dg\,g^{-1}\wedge dg).
\end{equation}
The equations \eqref{eq:dA_1} and \eqref{eq:dA_2} allow us to write :
\begin{equation}\label{eq:dA}
\begin{split}
    (dA_{\beta})&=d(g^{-1} A_{\alpha} g)+d(g^{-1} dg)\\
              &=-(g^{-1} dg\,g^{-1}) \wedge(A_{\alpha} g)-(g^{-1} A_{\alpha})\wedge dg\\
	      &\quad+(g^{-1} dA_{\alpha} g)-(g^{-1} dg\,g^{-1})\wedge dg.
\end{split}
\end{equation}
Notice that the term $(g^{-1} dA_{\alpha} g)$ corresponds to the first one in $F_{\beta}=g^{-1}(dA_{\beta}+A_{\beta}\wedge A_{\beta})g$.

For anyone who had understood the whole computations up to here, it is clear that
\begin{equation}
\begin{split}
     [A_{\beta}(v),A_{\beta}(w)]&=\DDsdd{ e^{tA_{\beta}(v)}e^{tA_{\beta}(w)} }{t}{0}{u}{0}\\
                            &\quad-\DDsdd{ e^{tA_{\beta}(w)}e^{tA_{\beta}(v)} }{t}{0}{u}{0}\,,
\end{split}
\end{equation}
so that
\begin{equation}
\begin{split}
  A_{\beta}\wedge A_{\beta}&=g^{-1} A_{\alpha} g\wedge g^{-1} A_{\alpha} g
                         +g^{-1} A_{\alpha} g\wedge g^{-1} dg\\
		       &\quad+g^{-1} dg\wedge g^{-1} A_{\alpha} g
		       +g^{-1} dg\wedge g^{-1} dg.
\end{split}
\end{equation}
Lemma \ref{lem:simplif} allows us to write it under the form
\begin{equation}\label{eq:AA}
\begin{split}
  A_{\beta}\wedge A_{\beta}&=g^{-1} A_{\alpha} g\wedge g^{-1} A_{\alpha} g
                         +g^{-1} A_{\alpha} g\wedge g^{-1} dg\\
		       &\quad+g^{-1} dg\wedge g^{-1} A_{\alpha} g
		       +g^{-1} dg\wedge g^{-1} dg.
\end{split}
\end{equation}
Here the term $(g^{-1} A_{\alpha}\wedge A_{\alpha} g)$ corresponds to the second one in $F_{\beta}=g^{-1}(dA_{\beta}+A_{\beta}\wedge A_{\beta})g$. The sum of \eqref{eq:dA} and \eqref{eq:AA} is
\[
    F_{\beta}=g^{-1} F_{\alpha} g.
\]
\end{proof}




\subsection{The electromagnetic field \texorpdfstring{$F$}{F}}
%--------------------------------------
Now, we are able to interpret the field $F$ introduced in equation \eqref{def_F}.  We follow \cite{Preparation}. From now, we use the usual Minkowski metric $g=diag(-,+,+,+)$.
From the vector given by \eqref{def_A}, we define a (local) potential $1$-form
\[
   A=A_{\mu}dx^{\mu}=-\phi dt+A_x dx+A_y dy+A_z dz.
\]
The field strength\index{Yang-Mills!field strength} is $F=dA$. We easily find that
\begin{equation}
\begin{split}
   F&=(dt\wedge dx)(\partial_x\phi+\partial_t A_x)+\ldots\\
    &\quad+(dx\wedge dy)(-\partial_z A_x+\partial_x A_y)+\ldots
\end{split}
\end{equation}
But the fields $\bB$ and $\bE$ are defined from $\bA$ and $\phi$ by \eqref{BE_de_A}, so
\begin{equation}
\begin{split}
  F&=-E_x(dt\wedge dx)-E_y(dt\wedge dy)-E_z(dt\wedge dz)\\
   &\quad+B_x(dy\wedge dz)+B_y(dz\wedge dx)+B_z(dx\wedge dy).
\end{split}
\end{equation}

We naturally have $dF=d^2A=0$. But conversely, $dF=0$ ensures the existence of a $1$-form $A$ such that $F=dA$. If we define\footnote{\emph{i.e.} we consider $F$ as the main physical field while $\bE$ and $\bB$ are ``derived'' fields.} $\bB=\nabla\times\bA$ and $\bE=-\nabla\phi-\partial_t\bA$, equations \eqref{M2} and \eqref{M3} are obviously satisfied. So in the connection formalism, the equations ``without sources''\ are written by
\begin{equation}\label{M23}
 dF=0.
\end{equation}
In order to write the two others, we introduce the current $1$-form\index{current!$1$-form} :
\[
  j=j_{\mu}dx^{\mu}=-\rho dt+j_x dx+j_y dy+j_z dz.
\]
%
One sees that
\begin{equation}
\begin{split}
  \delta F:=\star d\star F&=-dt(\nabla\cdot\bE)\\
                &\quad+ dx(-\partial_t \bE_x+ (\nabla\times\bB)_x )\\
		&\quad+ dy(-\partial_t \bE_y+ (\nabla\times\bB)_y )\\
		&\quad+ dz(-\partial_t \bE_z+ (\nabla\times\bB)_z ),
\end{split}
\end{equation}
so that equation $\delta F=j$ gives equations \eqref{M1} and \eqref{M4}. Now, the complete set of Maxwell's\index{Maxwell} equations is :
\begin{subequations}
\begin{align}
   d F&= 0\label{SM1}\\
   \delta F &=j\label{SM2}
\end{align}
\end{subequations}
with
\begin{subequations}
\begin{align}
j&=-\rho dt+j_x dx+j_y dy+j_z dz,\\
  \bB&=\nabla\times\bA\\
  \bE&=-\nabla\phi-\partial_t\bA
\end{align}
\end{subequations}
where $A$ is a $1$-form such that $F=dA$ whose existence is given by \eqref{SM1}.
\section{Inclusion of the Lorentz group}\label{subsec:incl_Lorentz}
%+++++++++++++++++++++++++++++++++++++++

Up to now we had seen how to express the \emph{gauge} invariance of a physical theory. In particle physics, a really funny field theory must be invariant under the Lorentz group; it is rather clear that, from the bundle point of view, this feature will be implemented by a Lorentz-principal bundle and some associated bundles. A spinor will be a section of an associated bundle for spin one half representation of the Lorentz group on $\eC^4$. In order to describe non-zero spin particle interacting with an electromagnetic field (represented by a connection on a $U(1)$-principal bundle), we have to build a correct $\SLdc\times U(1)$-principal bundle. We are going to use the ideas of \ref{subsec:sym_nature}.

A \defe{space-time}{space-time} is a differentiable \defe{pseudo-Riemannian}{pseudo-Riemannian} $4$-dimensional manifold. The pseudo-Riemannian structure is a $2$-form $g\in\Omega^2(M)$ for which we can find at each point $x\in M$ a basis $b=(\overline{ b }_0,\ldots,\overline{ b }_3)$ which fulfils
\[
  g_x(\overline{ b }_i,\overline{ b }_j)=\eta_{ij}.
\]
When we use an adapted coordinates, the metric reads $g=\eta_{ij}dx^i\otimes dx^j$.

One says that $M$ is \defe{time orientable}{time!orientable} if one can find a vector field $T\in\cvec(M)$ such that $g_x(T_x,T_x)>0$ for all $x\in M$. A \defe{time orientation}{time!orientation} is a choice of such a vector field. A vector $v\in T_xM$ is \defe{future directed}{future!directed vector} if $g_x(T_x,v)>0$.

The Lorentz group $L$ acts on the orthogonal basis of each $T_xM$, but you may note that $L$ don't act on $M$; it's just when the metric is flat that one can identify the whole manifold with a tangent space and consider that $L$ is the space-times isometry group. In the case of a curved metric, the Lorentz group have to be introduced pointwise and the building of a frame bundle is natural.

Now, we are mainly interested in the frame related each other by a transformation of $L_+^{\uparrow}$. An arising question is to know if one can make a choice of some basis of each $T_xM$ in such a manner that 

\begin{enumerate}
\item pointwise, the chosen frames are related by a transformation of $L_+^{\uparrow}$,
\item the choice is globally well defined.
\end{enumerate}
The first point is trivial to fulfil from the definition of a space-time. For the second, it turns out that a good choice can be performed if and only if there exists a vector field $V\in\cvec(M)$ such that $g_x(V_x,V_x)>0$ for all $x\in M$. We suppose that it is the case\footnote{That condition is rather restrictive because we cannot, for example, find an everywhere non zero vector field on the sphere $S^n$ with $n$ even.}.

So our first principal bundle attempt to describe the space-time symmetry is the $L_+^{\uparrow}$-principal bundle of orthonormal oriented frame on $M$ :
 \begin{equation}\label{bun:LpF}
\xymatrix{
    L_+^{\uparrow}  \ar@{~>}[r] & L(M) \ar[d]^{p_L} \\
    &M
  }
\end{equation}
The notion of ``\defe{relativistic invariance}{relativistic invariance}'' has to be understood in the sense of associated bundle to this one. The next step is to recall ourself (see subsection \ref{subsec:sym_nature}) that the physical fields doesn't transform under representation of the group $L_+^{\uparrow}$ but rather under representations of $\SLdc$. So we build a $\SLdc$-principal bundle
 \[
\xymatrix{
    \SLdc  \ar@{~>}[r] & S(M) \ar[d]^{p_S} \\
    &M
  }
\]
In order this bundle to ``fit''{} as close as possible the bundle \eqref{bun:LpF}, we impose the existence of a map $\dpt{\lambda}{S(M)}{L(M)}$ such that

\begin{enumerate}
\item $p_B(\lambda(\xi))=p_S(\xi)$ for all $\xi\in S(M)$ and
\item $\lambda(\xi\cdot g)=\lambda(\xi)\cdot\mSpin(g)$ for all $g\in\SLdc$.
\end{enumerate}
You can recognize the definition of a \defe{spin structure}{spin!structure}\label{pg_spinenphyz}. Notice that the existence of a spin structure on a given manifold is a non trivial issue.

Now a physical field is given by a section of the associated bundle $E=S(M)\times_{\rho} V$ where $\rho$ is a representations of $\SLdc$ on $V$. For an electron, it is $V=\eC^4$ and $\rho=D^{(1/2,0)}\oplus D^{(0,1/2)}$. That describes a \emph{free} electron is the sense that it doesn't interacts with a gauge field. So in order to write down the formalism in which lives a non zero spin particle, we have to build a $U(1)\times\SLdc$-principal bundle. For this, we follow the procedure given in section \ref{sec:produit_bundle}

\section{Interactions}
%++++++++++++++++++

\subsection{Spin zero}
%--------------------

The general framework is the following :
\[
 \xymatrix{
    U(1)  \ar@{~>}[r] & P \ar[d]_{\displaystyle \pi} && E=P\times_{\rho} V \\
                      & M&\mU_{\alpha}\ar@{^{(}->}[l]  \ar[ur]_{\displaystyle\phi} \ar[lu]_{\displaystyle\sigma_{\alpha}}
  }
\]
a $U(1)$-principal bundle over a manifold $M$ (as far as topological subtleties are concerned, we suppose $M=\eR^4$) and a section $\phi$ of an associated bundle for a representation $\rho$ of $U(1)$ on $V$. We consider $M$ with the Lorentzian metric but, since we are intended to treat with scalar (spin zero) fields, we still don't include the Lorentz (or $\SLdc$) group in the picture. We also consider local sections $\dpt{\sigma_{\alpha}}{\mU_{\alpha}}{P}$, a connection $\omega$ on $P$ and $\Omega$ its curvature. We define $A_{\alpha}=\sigma_{\alpha}^*\omega$.

Now we particularize ourself to the target space $V=\eC$ on which we put the scalar product
\begin{equation}		\label{EqProdScalVeCU}
	\scal{z_1}{z_2}=\frac{1}{2}( z_1\overline{ z }_2+z_2\overline{ z }_1  ),
\end{equation}
and the representation $\dpt{\rho_n}{U(1)}{GL(\eC)}$,
\[
  \rho_n(g)z=g\cdot z=g^nz
\]
where we identify $U(1)$ to the unit circle in $\eC$ in order to compute the product. A property of the product \eqref{EqProdScalVeCU} is to make $\rho_n$ an isometry: for all $g\in U(1)$, $z_1,z_2\in\eC$,
\[
  \scal{\rho_n(g)z_1}{\rho_n(g)z_2}=\scal{z_1}{z_2}.
\]
Our first aim is to write the covariant derivative of $\phi$ with respect to the connection $\omega$. For this we work on the section $\phi$ under the form $\dpt{\phi_{(\alpha)}}{M}{V}$ and we use formula \eqref{3008e1} :
\begin{equation}
  (D_X\phi)\bsa(x)=X_x\phi\bsa-\rho_*\big( (\sigma^*_{\alpha}\omega)_xX_x \big)\phi\bsa(x).
\end{equation}
Let us study this formula. We know that $(\sigma_{\alpha}^*\omega)_x=A_{\alpha}(x)\,:\,T_x\mU_{\alpha}\stackrel{\sigma}{\to}T_{\sigma_{\alpha}(x)}P\stackrel{\omega}{\to}u(1)$. Thus $A_{\alpha}(x)X_x$ is given by a path in $U(1)$; it is this path which is taken by $\rho_*$. Therefore (we forget some dependences in $x$)
\begin{equation}
\begin{split}
  \rho_*\big( A_{\alpha}(x)X_x \big)\phi\bsa(x)&=\Dsdd{ \rho_n\big( (A_{\alpha} X)(t)  \big)\phi\bsa(x) }{t}{0}\\
                                                &=\Dsdd{ (A_{\alpha} X)(t)^n }{t}{0}\phi\bsa(x)\\
						&=n\Dsdd{(A_{\alpha} X)(t)}{t}{0}\phi\bsa(x)\\
						&=nA_{\alpha}(X)\phi\bsa(x).
\end{split}						
\end{equation}
Thus the covariant derivative is given by
\begin{equation}
  (D_X\phi)\bsa(x)=X_x\phi\bsa-nA_{\alpha}(x)(X_x)\phi\bsa(x).
\end{equation}

One can guess an electromagnetic coupling for a particle of electric charge~$n$. If this reveals to be physically relevant, it shows that the ``electromagnetic identity card'' of a particle is given by a representation of $U(1)$. This has to be seen in relation to the discussion on page \pageref{pg:phyz_reprez} where the ``type of particle'' was closely related to representations of the Lorentz group. It is a remarkable piece of quantum field theory: the properties of a particle are encoded in representations of some symmetry groups.

Now we are going to prove that $\|D\phi\|^2$ is a gauche invariant quantity. The first step is to give a sense to this norm. We consider $X_i$ ($i=0,1,2,3$), an orthonormal basis of $T_xM$ and we naturally denote $D_i=D_{X_i}$, $\partial_i=X_i$ and $A_{\alpha i}=A_{\alpha}(\partial_i)$. Remark that 
\begin{equation}
   A_{\alpha}(x)X_x=(\sigma_{\alpha}^*)_xX_x
                 =\omega( d\sigma_{\alpha} X_x )
                 =\omega\Dsdd{ \sigma_{\alpha}(X(t)) }{t}{0}\in u(1),
\end{equation}
so this is given by a path in $U(1)$ which can be taken by $\rho$. Let $c(t)$ be this path, then
\[
   A_{\alpha}\phi\bsa(x)=\Dsdd{ e^{ic(t)}\phi\bsa(x) }{t}{0},
\]
so that under the conjugation, $\overline{A_{\alpha}\phi\bsa(x)}=-A_{\alpha}\overline{ \phi }\bsa(x)$. Now our definition of $\|D\phi\|^2$ is a composition of the norm on $V$ and the one on $T_xM$ :
\begin{equation}
  \|D\phi\|^2=\eta^{ij}\scal{ D_i\phi\bsa }{D_j\phi\bsa}
\end{equation}
Using the notation in which the upper indices are contractions with $\eta^{ij}$, we have
\[
\|D\phi\|^2=\Big(  (\partial_i\phi\bsa)(x)-nA_{\alpha i}\phi\bsa(x)   \Big)
            \Big(  (\partial^i\overline{ \phi }\bsa)(x)+nA_{\alpha}^i\overline{ \phi }\bsa(x)   \Big).
\]

\subsubsection{Gauge transformation law}
%////////////////////////////////////////

A gauge transformation $\varphi$ is given by an equivariant function $\dpt{\tilde{\varphi}_{\alpha}}{\mU_{\alpha}}{U(1)}$ which can be written under the form
\[
   \tilde{\varphi}_{\alpha}(x)=e^{i\Lambda(x)}
\]
for a certain function $\dpt{\Lambda}{\mU_{\alpha}}{\eR}$. From the general formula \ref{lem:ii} of lemma \ref{lem:prop_gauge},
\begin{equation}
  (\varphi\cdot\phi)\bsa(x)=\rho_n(e^{-i\Lambda(x)})\phi\bsa(x)
                          =e^{-ni\Lambda(x)}\phi\bsa(x).
\end{equation}
The transformation of the gauche field $A$ is given by equation \eqref{tr_de_A}. Let us see the meaning of the term $d\tilde{\varphi}$. For $v\in T_x\mU_{\alpha}$,
\begin{equation}
  (d\tilde{\varphi}_{\alpha})_xv=\Dsdd{\tilde{\varphi}_{\alpha}(v(t))}{t}{0}
                   =\Dsdd{ e^{i\Lambda(v(t))} }{t}{0}
		   =i\Dsdd{\Lambda(v(t))}{t}{0}e^{i\Lambda(v(0))}
		   =i(d\Lambda)_xve^{i\Lambda(x)}.
\end{equation}
Thus $\tilde{\varphi}_{\alpha}^{-1}(x)(d\tilde{\varphi}_{\alpha})_x=i(d\Lambda)_x$. Since $U(1)$ is abelian, $\tilde{\varphi}^{-1} A\tilde{\varphi}=A$. Finally,
\begin{equation}
  (\varphi\cdot A)_{\alpha}(x)=A_{\alpha}(x)+i(d\Lambda)_x.
\end{equation}
Now we are able to prove the invariance of $\|D\phi\|^2$. First,
\begin{equation}
\begin{split}
  (\varphi\cdot A)_{i\alpha}(x)=(\varphi\cdot A)_{\alpha}(\partial_i)
                           =A_{i_{\alpha}}(x)+i(\partial_i\Lambda)(x);
\end{split}
\end{equation}
second,
\begin{equation}
\partial_i\left(  e^{-ni\Lambda(x)}\phi\bsa(x)  \right)=-ni(\partial_i\Lambda)(x)\phi\bsa(x)+e^{-in\Lambda(x)}(\partial_i\phi\bsa)(x).
\end{equation}
With these two results, 
\begin{equation}
\partial_i(\varphi\cdot\phi)\bsa(x)+n(\varphi\cdot A)_{\alpha i}(\varphi\cdot\phi)\bsa(x)=e^{-in\Lambda(x)}( nA_{\alpha i}(x)+\partial_i\phi\bsa(x) ).
\end{equation}

The Yang-Mills \defe{field strength}{field!strength} is given by $F\bsa=\sigma^*_{\alpha}\Omega$ (cf. page \pageref{pg:curv_princ}). Since $U(1)$ is abelian, $dF\bsa=0$, so that the second pair of Maxwell's equations is complete without any Lagrangian assumptions.

 The full Yang-Mills action is written as
\[
  S(\omega,\phi)=\int_M\left[  -\frac{1}{4}F\bsa_{ij}F\bsa^{ij}+\frac{1}{2}\|D\phi\|^2+\frac{1}{2} m\phi\bsa\overline{\phi\bsa}  \right].
\]
The Euler-Lagrange equations are
\begin{subequations}
\begin{align}
(\partial_i-inA_{\alpha i})(\partial^i-inA_{\alpha}^i)\phi_{\alpha}+m^2\phi_{\alpha}&=0\\
        \partial_iF\bsa^{ij}&=0.
\end{align}
\end{subequations}
So the Yang-Mills Lagrangian only gives the first pair of Maxwell's equations while the second one is given by the geometric nature of fields.

As explained in \cite{AlexModaveII}, the topology of the physical space has deep implications on the physics of Yang-Mills equations. The absence of magnetic monopoles for example is ultimately linked to the (simple) connectedness of $\eR^4$. When one consider the $U(1)$ Yang-Mills on a sphere, some topological charges appear and magnetic monopoles naturally arise.

\subsection{Non zero spin formalism}
%++++++++++++++++++++++++++++++++

The formalism for a non zero spin particle in an electromagnetic field is described in section \ref{sec:produit_bundle}. We consider the spinor bundle
\[
\xymatrix{
    \SLdc \ar@{~>}[r] & S(M)\ar[d]^{p_S}\\ &M 
  }
\]
with the spinor connection on $S(M)$, and $\rho_1$, a representation of $\SLdc$ on $V$. For an electron, it is $V=\eC^4$ and $\rho_1=D^{(1/2,0)}\oplus D^{(0,1/2)}$, so for $g_1\in\SLdc$,
\begin{equation}
  \rho_1(g_1)\begin{pmatrix}
  z_1\\\vdots\\z_4
             \end{pmatrix}
=
\begin{pmatrix}
  g_1\\\\
&(\overline{g_1}^t)^{-1}
\end{pmatrix}
\begin{pmatrix}
  z_1\\\vdots\\z_4
             \end{pmatrix}.
\end{equation}
On the other hand, we consider the principal bundle
\[
\xymatrix{
    U(1) \ar@{~>}[r] & P\ar[d]^{p_U}\\ &M 
  }
\]
with a connection $\omega_2$ which describes the electromagnetic field. As representation $\dpt{\rho_2}{U(1)}{GL(\eC^4)}$ we choose the multiplication coordinate by coordinate :
\begin{equation}
\rho_2(g_2)\begin{pmatrix}
  z_1\\\vdots\\z_4
             \end{pmatrix}
=
\begin{pmatrix}
  g_2z_1\\\vdots\\g_2z_4
             \end{pmatrix}.
\end{equation}
The physical picture of the electron is now the principal bundle
\[
\xymatrix{
    \SLdc\times U(1) \ar@{~>}[r] & S(M)\circ P\ar[d]^{p}\\ &M,
  }
\]
and the field is a section of the associated bundle $(S(M)\circ P)\times_{\rho}\eC^4$.





\appendix
\setcounter{chapter}{0}
\setcounter{section}{0}
\renewcommand{\theequation}{\Alph{chapter}.\arabic{equation}}
\renewcommand{\thenumtho}{\Alph{chapter}.\arabic{numtho}}
\renewcommand{\thechapter}{\Alph{chapter}}




\chapter{Complements}
\section{Alternative formalism for the quantum mechanics}\label{app:Wigner}
%++++++++++++++++++++++++++++++++++++++++++++++++++++++++

We can a little reformulate the axioms of the quantum mechanics. Since we are in a Hilbert space $\hH$ we can speak about orthogonal projections; if $\phi\in\mR$, we can consider the projection on the space spanned by $\phi$:
\[
   P_{\phi}e_k=\frac{\scal{\phi}{e_k}}{\|\phi\|}\phi
\]
where $\{e_i\}$ is a basis of $\hH$. It is pretty clear that 
\begin{equation}
\tr(P_{\phi} P_{\psi})=\frac{| \scal{\psi}{\phi} |^2}{\|\psi\|\|\phi\|}.
\end{equation}
If $\psi\in\mR$ and $\phi\in\mR'$ are unimodular, then
\begin{equation}
P(\mR\to\mR')=\tr(P_{\phi}P_{\psi}),
\end{equation}
so we can express the axioms in terms of projections instead of rays. For notational convenience, we put
\begin{equation}
\hH_1=\{\psi\in\hH\tq\|\psi\|=1\}.
\end{equation}

We denote by $\hS$ the space of the projections into one dimensional subspaces of $\hH$ (in other words $\hS$ is the space of physical states) and for $P$, $Q\in\hS$, the transition probability is $P\cdot Q=\tr(PQ)$.
Now a \defe{quantum symmetry}{symmetry!of quantum system} is a map $\dpt{T}{\hS}{\hS'}$ such that $(TP)\cdot(TQ)=P\cdot Q$.

One can prove the following :

\begin{theorem}
If $\dpt{T}{\hS}{\hS'}$ is a quantum symmetry, then there exists an operator $\dpt{U}{\hH}{\hH'}$ such that
\begin{enumerate}
\item $P_{U\phi}=TP_{\phi}$,
\item $U(\xi+\eta)=U(\xi)+U(\eta)$,
\item $\scal{U\xi}{U\eta}=\kappa(\scal{\xi}{\eta})$ \label{item:cond_3}
\end{enumerate}
where $P_{\psi}$ is the projection onto the one dimensional space spanned by $\psi$ and $\dpt{\kappa}{\eC}{\eC}$ fulfils $\kappa(\lambda)=\lambda$ or $\kappa(\lambda)=\overline{\lambda}$ and 

\begin{enumerate}
\item $U(\lambda\xi)=\kappa(\lambda)\xi$.
\end{enumerate}
\label{tho:pre_Wigner}
\end{theorem}
Here is why this implies Wigner's theorem as given by theorem \ref{tho:Wigner}. Let us consider some $\varphi_i\in\hH$ such that $\|\varphi_i\|=1$ and $P_{\varphi_i}$, the corresponding projections. Let
\[
   \Delta(P_1,P_2,P_3)=\scal{\varphi_1}{\varphi_2}\scal{\varphi_2}{\varphi_3}\scal{\varphi_3}{\varphi_1}.
\]
It is clear that this expression doesn't depend on the choice of $\varphi_i$ in its ray. We have
\begin{equation}\label{eq:Delta_T}
\begin{split}
\Delta(TP_1;TP_2,TP_3)&=\Delta(P_{U\varphi_1},P_{U\varphi_2},P_{U\varphi_3})\\
                      &=\scal{U\varphi_1}{U\varphi_2}\scal{U\varphi_2}{U\varphi_3}\scal{U\varphi_3}{U\varphi_1}\\
		      &=\kappa(\scal{U\varphi_1}{U\varphi_2})\kappa(\scal{U\varphi_2}{U\varphi_3})\kappa(\scal{U\varphi_3}{U\varphi_1})\\
		      &=\kappa\big(\Delta(P_1,P_2,P_3)\big).
\end{split}
\end{equation}
We can see from this that the choice of $\kappa(\lambda)=\lambda$ or $\kappa(\lambda)=\overline{\lambda}$ is determined by the data of $T$ if $\dim\hH\geq 2$. In the case where $\dim\hH=1$, $\Delta$ is always equals to $1$ and the equality \eqref{eq:Delta_T} don't give any informations. In the case $\dim\hH\geq 2$, we can choice $\varphi_1$ and $\varphi_2$ such that $\scal{\varphi_1}{\varphi_2}$ takes any value $z\in\eC$ with $\|z\|\leq 1$. Taking $\varphi_3=\varphi_1+\varphi_2$, we find
\[
  \Delta(P_1,P_2,P_3)=z(1+\overline{z})^2.
\]
which is easily non real for a suitable choice of $z\in\eC$. Let us suppose that we have an operator $U$ which satisfies the theorem \ref{tho:pre_Wigner}. If $\kappa(\lambda)=\lambda$, then
\begin{equation}
U(z\psi+z'\phi)=U(z\psi)+U(z'\phi)
               =zU(\psi)+zU(\phi)
\end{equation}
and 
\begin{equation}
\scal{U\psi}{U\phi}=\kappa(\scal{\psi}{\phi})
	=\scal{\psi}{\phi},
\end{equation}
so that $U$ is linear. If $\kappa(\lambda)=\overline{z}$, then 
\begin{equation}
 U(z\psi)=\overline{z}U\psi
\end{equation}
and
\begin{equation}
\scal{U\xi}{U\eta}=\kappa(\scal{\xi}{\eta})
                  =\overline{\scal{\xi}{\eta}}.
\end{equation}


% \begin{proof}[Proof of theorem \ref{tho:pre_Wigner}]
% We begin with $\dim\hH=1$. We have $\dpt{T}{\hS}{\hS'}$; for a given $\varphi\in\hH_1$, we can finf $\varphi'\in\hH_1'$ such that $P_{\varphi'}=TP_{\varphi}$ and define $U\varphi=\varphi'$. We define $U$ on the other vectors of $\hH$ by $U(z\varphi)=\kappa(z)U\varphi$ and $U(\xi+\eta)=U(\xi)+U(\eta)$. In order to see that $\scal{U\xi}{U\eta}$, remark that there exists $z,z'\in\eC$ such that $\eta=z\xi$ and $\xi=z'\varphi$. Then
% 
% \begin{equation}
% \begin{split}
% \scal{U\xi}{U\eta}&=\scal{U\xi}{U\xi}\kappa(z)\\
%                   &=\overline{\kappa(z')}\scal{U\varphi}{U\varphi}\kappa(z)\kappa(z')\\
% 		  &=\overline{\kappa(z')}\kappa(z)\kappa(z'),
% \end{split}
% \end{equation}
% but 
% \begin{equation}
% \begin{split}
% \scal{\xi}{eta}&=\scal{z'\varphi}{zz'\varphi}\\
%                &=\overline{z'}zz'.
% \end{split}
% \end{equation}
% 
% Now we turn our attention to the case $\dim\hH\geq2$. The condition $P_{U\varphi}=TP_{\varphi}$ fix $U$ up to a complex number: if $U(\xi)$ works, then $U'(\xi)=z(\xi)U(\xi)$ equaly works. The condiction \ref{item:cond_3} gives
% 
% \begin{equation}
% \begin{split}
%   \kappa(\scal{\xi}{\eta})&=\scal{U'\xi}{U'\eta}\\
%                           &=\overline{z(\xi)}z(\eta)\kappa(\scal{\xi}{\eta})
% \end{split}
% \end{equation}
% for any $\eta,\xi\in\hH$. Then $z$ is a constant of module $1$. Let $\xi'$ be a candidate to be $U\xi$. We can define $\varphi$ by $\xi=\|\xi\|\varphi$ (then $\|\varphi\|=1$). Since $P_{U\varphi}=TP_{\varphi}$ and $U\varphi=\xi'$ up to a multiplicative constant, then $\xi'$ is an arbitrary vector of nomr $\|\xi\|$ in the space given by $TP_{\varphi}$. Then
% 
% \[
%   \|\scal{\xi'}{\eta'}\|=\|\scal{\xi}{\eta}\|
% \]
% beause $\kappa$ preserves the norms. If $\{\varphi_1\ldots,\varphi_m\}$ is an orthonormal set of vectors in $\hH$, then $\{\varphi_1,\ldots,\varphi_m\}$ is a set of orthonormal vectors in $\hH'$ because $\|\scal{\varphi_i'}{\varphi'_j}\|=\|\scal{\varphi_i}{\varphi_j}|=\delta_{ij}$. Note that $\scal{\varphi_i}{\varphi_i}=1$ because is must be real and positive.
% 
% \end{proof}

% This is part of (almost) Everything I know in mathematics
% Copyright (c) 2013-2014
%   Laurent Claessens
% See the file fdl-1.3.txt for copying conditions.

\section{Statement of some results}
%+++++++++++++++++++++++++++++++++

We provide here some results that are at undergrad level. Some of them are contained in the part «Pour l'agrégation» in French.

\begin{lemma}[Schur's lemma] \label{lem:Schur}
If $\dpt{\phi}{\lG}{\gl(V)}$ is irreducible, then the only endomorphism of $V$ which commutes with all $\phi(\lG)$ are multiples of identity.
\end{lemma}
%TODO : look the link with the Schur'lemma given in Agregation.

\begin{theorem}[Cayley-Hamilton]\index{Cayley-Hamilton theorem} \label{ThoCayleyHamilton}
    A square matrix on \( \eR\) or \( \eC\) satisfies its own characteristic equation. If \( A\in\eM_n(\eK)\), we consider the polynomial \( p(\lambda)=\det(A-\lambda\mtu)\) in \( \lambda\). Thus \( p(A)=0\).
\end{theorem}
For a proof see \wikipedia{en}{Cayley–Hamilton_theorem}{Wikipedia}.

\begin{proposition}[\cite{SerreMatrices}]     \label{PropMtrDiagablaUnit}
    If \( M\) is a complex \( n\times n\) matrix, then there exists an unitary matrix \( U\) such that \( U^*MU\) is upper triangular.
\end{proposition}

\begin{definition}
    A \defe{positive defined}{positive!defined matrix} matrix is a matrix $B$ such that
    \begin{equation}
        \sum_{ij}B_{ij}\overline{ x_i }x_j
    \end{equation}
    is real and positive for every complex vector $x$.
\end{definition}

\begin{proposition}
    A positive defined matrix is Hermitian.
\end{proposition}

\begin{proof}
    We define the Hermitian matrices $M=(B+B^*)/2$ and $N=(B-B^*)/2i$, so $B=M+iN$ and
    \begin{equation}
        \bar x Bx=\bar x M x+i\bar x Nx.
    \end{equation}
    The matrices $M$ and $B$ being Hermitian, the numbers $\bar xMx$ and $\bar xNx$ are real. If $\bar xBx$ has to be real, we need $\bar xNx=0$ for every $x$. This shows that $N=0$, so that $B=N$.
\end{proof}

\begin{theorem}
Let $G$ be a Lie group and $H$ a subgroup (with no special other structures) of $G$. If $H$ is a closed subset of $G$ then there exists an unique analytic structure on $H$ such that $H$ is a topological Lie subgroup of $G$.
\label{Helgason2.3}
\end{theorem}
This comes from \cite{Helgason}, chapter 2, theorem 2.3.

\begin{lemma}
Let $G$ be a connected Lie group with Lie algebra $\mG$ and let $\varphi$ be an analytic homomorphism of $G$ into a Lie group $X$ with Lie algebra $\mathcal{X}$. Then

\begin{enumerate}
\item The kernel $\varphi^{-1}(e)$ is a topological Lie subgroup of $G$. Its Lie algebra is the kernel of $d\varphi_e$.
\item The image $\varphi(G)$ is a Lie subgroup of $\mathcal{X}$ with Lie algebra $d\varphi(\mG)\subset\mathcal{X}$.
\end{enumerate}
\label{Helgason5.1}
\end{lemma}
This comes from \cite{Helgason}, chapter 2, lemma 5.1.

\begin{lemma}
Let $G$ and $H$ be two Lie group, whose Lie algebra are $\mG$ and $\mH$. If $\dpt{\theta}{G}{H}$ is a surjective map, then we have $\mH\simeq\mG/Ker\,d\theta_e$.
\label{1203r1}
\end{lemma}

\begin{theorem} \label{1503t1}
Let us consider $\dpt{Ad}{SU(2)}{GL(3)}$, $Ad(U)X=UXU^{-1}$. We have the following properties:

\begin{enumerate}
\item $Ad$ is a linear homomorphism,
\item it takes his values in $\SO(3)$; then we can write $\dpt{Ad}{SU(2)}{\SO(3)}$,
\item it is surjective,
\item $Ker\,Ad=\eZ_2$,
\item all these properties show that \[\SO(3)=\frac{SU(2)}{\eZ_2}.\]
\end{enumerate}
\end{theorem}

\begin{definition}
If $(a_k)$ is a sequence in $\eR$, its \defe{upper limit}{upper limit} is the real number
\[
  \lim\sup_{n\to\infty}a_n=\lim_{l\to\infty}\sup\{a_k:k\geq l\}.
\]
\end{definition}

\begin{lemma}
If $\omega$ is a $k$-form (not specially a symplectic one), and $\nabla$ a torsion free connection, one has
\begin{equation}\label{eq:d_omega_nabla}
  (d\omega)(X_0,\ldots,X_k)=\sum_{i=0}^k (-1)^i(\nabla_{X_i}\omega)(X_0,\ldots,\hX_i,\ldots X_k).
\end{equation}
\end{lemma}

\begin{remark}
The link between $d$ and $\nabla$ comes from the fact that in the left hand side of \eqref{eq:d_omega_nabla} appears some commutators $[X_i,X_j]$, but since the connection is torsion-free,
\[
  [X_i,X_j]=\nabla_{X_i}X_j-\nabla_{X_j}X_i
\]
\end{remark}
The main consequence of this lemma is that $\nabla\omega=0$ implies $d\omega=0$. 

\begin{proposition} \label{prop:fdefint}
    Consider a function $\dpt{f}{X\times E}{\overline{\eR}}$ and $z_0\in E$ such that
    \begin{itemize}
        \item for all $z\in E$, the function $x\to(x,z)$ is integrable,
        \item for (almost) all $x\in X$, the function $z\to f(x,z)$ is continuous at $z_0$,
        \item there exists a function $g\geq 0$ such that for all $z\in E$, $| f(x,z) |\leq g(x)$ almost everywhere in $X$.
    \end{itemize}
    Then the function $\dpt{h}{E}{\eR}$ defined by $h(z)=\int_Xf(x,z)$ is continuous at $z_0$.
\end{proposition}
This proposition is (up to ``almost'') the theorem \ref{ThoKnuSNd}.

% \begin{definition}

% If $(a_k)$ is a sequence in $\eR$, its \defe{upper limit}{upper limit} is the real number
% \[
%   \lim\sup_{n\to\infty}a_n=\lim_{l\to\infty}\sup\{a_k:k\geq l\}.
% \]
% \end{definition}
% 
% \begin{proposition}\label{prop:cv_lim_sup}
% If $(a_k)$ is a sequence in $\eR$ such that there exists a $a\in\eR$ for which for any $k\in\eN$,
% \[
% \lim\sup_{n\to\infty}a_n\leq a\leq a_k
% \]
% then $(a_k)$ admits a limit and $\lim_{n\to\infty}a_k=a$.
% \end{proposition}
% 
% \begin{lemma}
% If $\omega$ is a $k$-form (not specially a symplectic one), and $\nabla$ a torsion free connection, one has
% 
% \begin{equation}\label{eq:d_omega_nabla}
%   (d\omega)(X_0,\ldots,X_k)=\sum_{i=0}^k (-1)^i(\nabla_{X_i}\omega)(X_0,\ldots,\hX_i,\ldots X_k).
% \end{equation}
% \end{lemma}
% 
% \begin{remark}
% The link between $d$ and $\nabla$ comes from the fact that in the left hand side of \eqref{eq:d_omega_nabla} appears some commutators $[X_i,X_j]$, but since the connection is torsion-free,
% \[
%   [X_i,X_j]=\nabla_{X_i}X_j-\nabla_{X_j}X_i
% \]
% \end{remark}
% The main consequence of this lemma is that $\nabla\omega=0$ implies $d\omega=0$. 


\ifthenelse{\value{VerifForwardRef}=0}{
    \bibliographystyle{unsrt}           % unsrt fait que la biblio arrive dans l'ordre de citation au lieu de l'ordre alphabétique.
    \bibliography{mazhe}
}{}


\addcontentsline{toc}{chapter}{\nomname}
        \printnomenclature
% C'est le paquet tocbibind qui fait que l'index vient tout seul dans la table des matières.
        \printindex




\ifthenelse{\value{VerifForwardRef}=0}{
    \bibliographystyle{unsrt}           % unsrt fait que la biblio arrive dans l'ordre de citation au lieu de l'ordre alphabétique.
    \bibliography{mazhe}
}{}


\addcontentsline{toc}{chapter}{\nomname}
        \printnomenclature
% C'est le paquet tocbibind qui fait que l'index vient tout seul dans la table des matières.
        \printindex
\end{document}
